\documentclass{article} 
%\documentclass[draft]{article} 

%\sloppy
%\usepackage{tikz}
%\usetikzlibrary{shapes, arrows}
%\usetikzlibrary{positioning}
%\usepackage{amsmath,amssymb}
%%\usepackage{fullpage}
%\usepackage{comment}
%\usepackage{graphicx}
%%\pagestyle{empty}
%\usepackage{mathtools}
%\usepackage{algpseudocode}
%\usepackage{graphicx}
%%\usepackage{geometry,amssymb,latexsym,amsmath,graphicx}
%\usepackage{stmaryrd}%======================================================================
%\usepackage{bm,xcolor}
%\usepackage{ifdraft}




%code julia makematrix
%# expansion robuste
%expand(x::Int) = [x]
%expand(x::UnitRange{Int}) = collect(x)
%expand(x::AbstractVector) = isempty(x) ? Int[] : reduce(vcat, expand.(x))
%
%
%
%function makematrix(A, B)
%
%    A = expand.(A)
%
%    B = expand.(B)
%
%
%
%    # largeur max réelle
%
%    allcols = vcat(vcat(A...), vcat(B...))
%
%    maxcol = isempty(allcols) ? 0 : maximum(allcols)
%
%
%
%    # nombre de lignes
%
%    R = 2max(length(A), length(B)) + 1
%
%    C = maxcol
%
%
%
%    # matrice vide
%
%    M = fill(" ", R, C)
%
%
%
%    # remplir les tirets
%
%    for (i,l) in enumerate(A), j in l
%
%        M[2i-1, j] = "-"
%
%    end
%
%
%
%    # remplir les barres
%
%    for (i,l) in enumerate(B), j in l
%
%        M[2i, j] = "|"
%
%    end
%
%
%
%    # sortie LaTeX
%
%    "\\begin{matrix}\n" *
%
%    join([join(M[i,:]," & ") * " \\\\" for i=1:R], "\n") *
%
%    "\n\\end{matrix}"
%
%end

\usepackage[utf8]{inputenc}
\usepackage[LAE,T1]{fontenc}  % Ajoute un encodage compatible avec l'arabe
\usepackage[arabic, english,french]{babel}
\usepackage{CJKutf8} 

%======\usepackage{CJKutf8} 
%========================
\usepackage{tikz}
\usepackage{tikz-cd}
\usetikzlibrary{positioning, shapes, arrows.meta}

\usepackage{amsmath, amssymb, mathtools, bm, xcolor}
\usepackage{algpseudocode}
\usepackage{stmaryrd}
\usepackage{ifdraft}
\usepackage{graphicx}  % Pour includegraphics
\usepackage{caption} 
\usepackage{mdframed}  % Pour créer le cadre
%\usetikzlibrary{arrows.meta,positioning}
\usepackage{ytableau}
\usepackage{tikz-cd}
%\usepackage{circletikz} % Ensure circletikz is loaded

\newlength{\wider}
\setlength{\wider}{1.6cm} %{1.6cm}%2.8
\addtolength{\textwidth}{2\wider}
\addtolength{\oddsidemargin}{-\wider}
\addtolength{\evensidemargin}{-\wider}
\addtolength{\textheight}{2.5cm}
\addtolength{\topmargin}{-1.5cm}
%\newcommand{\Sucs}{%  \tikz[baseline=-0.5ex]{\draw[->, line width=1mm] (0,0) -- (0,-0.8);}}
 \newcommand{\Suct}{%
  \tikz[baseline=(current bounding box.center)]{
    \draw[->, line width=0.4mm] (0,0) -- (0,-1em);}}
    
     \newcommand{\Sucs}{%
  \tikz[baseline=(current bounding box.center)]{
    \draw[->, line width=0.4mm] (0,0) -- (1em,0);}}

 \newcommand{\Dect}{%
  \tikz[baseline=(current bounding box.center)]{
    \draw[<-, line width=0.4mm] (0,0) -- (0,-1em);}}
    
     \newcommand{\Decs}{%
  \tikz[baseline=(current bounding box.center)]{
    \draw[<-, line width=0.4mm] (0,0) -- (1em,0);}}
%======================================================================

\newcommand{\ket}[1]{\ensuremath{|{#1}\rangle}}
\newcommand{\bra}[1]{\ensuremath{\langle{#1}}|}
\newcommand{\inner}[2]{\ensuremath{\langle{#1}|{#2}\rangle}}
\newcommand{\zo}{\ensuremath{\{0,1\}}}
\newcommand{\level}{\ensuremath{{diag}}}
\renewcommand{\mid}{\ensuremath{\,:\,}}

\newcommand{\gftwo}{{\mathcal{F}_2}}
\newcommand{\quadratic}{q}
\newcommand{\qoracle}{\mathsf{Q}}
\newcommand{\adjoracle}{\tilde{\mathsf{Q}}}
\newcommand{\adj}{\texttt{A}}
\newcommand{\sigmax}{\mathsf{X}}
\newcommand{\sigmaz}{\mathsf{Z}}
\newcommand{\deltaloc}{\delta_{\textup{loc}}}
%\newcommand{\zo}{\ensuremath{\{0,1\}}} 
\newcommand{\dom}{\{0,1,-1\}}

\newcommand{\one}{1_{\textup{lim}}}
 
\newcommand{\op}[1]{\mathsf{#1}} 
\newcommand{\meas}[1]{\op{#1}} 
\newcommand{\symmdiff}{\Delta} 

\newcommand{\timestikz}{\tikz[baseline=(X.base)] \node[inner sep=0pt] (X) {\textbf{\Large x}};}

\newcommand{ \smp }{:_{\mathfrak s}}
\newcommand{\cre}{:_{\mathfrak t}}
\newcommand{\icp}{{\mathfrak i}_{\textup{cp}}}
\newcommand{\va}{\mathfrak v}

\newcommand{\Wl}{\rule{2pt}{5pt}}
\newcommand{\Wlt}[1]{\rule{2pt}{5pt}^{(#1)}}

\newcommand{\Comp}{\textup{Comp}} 
\newcommand{\Verif}{\textup{Verif}} 
\newcommand{\Prod}{\textup{Prod}} 
\newcommand{\Seq}{\textup{Seq}} 
\newcommand{\Seqd}[2]{\textup{Seq}_{#1_{|#2}}} 
\newcommand{\Expl}{\textup{Expl}} 
\newcommand{\Grad}{\textup{Grad}} 
\newcommand{\Next}{\textup{Next}} 
\newcommand{\Proc}{\mathcal{P}} 
\newcommand{\Bdf}{\textup{Bout}} 



\newcommand{\Infl}{\textup{Infl}} 
\newcommand{\Liane}{\textup{Liane}} 
\newcommand{\Line}{\textup{Cstr}} 
\newcommand{\mot}{\textup{mot}} 
\newcommand{\Supseq}{\textup{Supseq}} 
\newcommand{\Defn}{\textup{Def}} 
\newcommand{\Cpt}{\textup{Concept}} 
\newcommand{\Ppt}{\textup{Percept}} 
\newcommand{\Find}{\textup{Find}} 
\newcommand{\Read}{\textup{Read}} 
\newcommand{\Imcp}{\textup{Imcp}} 


\newcommand{\odd}{\textup{Odd}} 
\newcommand{\even}{\textup{Even}} 
\newcommand{\cutrank}{\textup{Cutrk}} 
\newcommand{\sgn}{\textup{sgn}} 
\newcommand{\tres}{\textup{ti-res}} 
\newcommand{\gtres}{\textup{gti-res}} 
\newcommand{\ts}{\textup{ts}} 
\newcommand{\jar}{\textup{Jard}} 
\newcommand{\tid}{\textup{Tid}} 
\newcommand{\Tid}{\mathcal{T}\textup{id}} 
\newcommand{\Btid}{\mathcal{B}\textup{Tid}} 

\newcommand{\Univ}{\textup{Univ}} 
\newcommand{\Expand}{\textup{Expand}}
\newcommand{\Create}{\textup{Create}}

\newcommand{\St}{\textup{SUniv}}  
\newcommand{\Gap}{\textup{Gap}} 

\newcommand{\FPV}{\textup{FPV}} 
\newcommand{\AtC}{\textup{AncToCell}} 

\newcommand{\gap}{\textup{Gap}} 
 \newcommand{\tU}{\textup{tU}} 

\newcommand{\JS}{\textup{JS}} 
\newcommand{\NJS}{\textup{NJS}} 
\newcommand{\CJS}{\textup{CJS}} 
\newcommand{\BJS}{\textup{BJS}} 

\newcommand{\canJ}[1]{\textup{canJ}_{#1}} 

\newcommand{\tg}{\mathfrak t _g} 
\newcommand{\desc}{\textup{desc}} 

\newcommand{\tim}[2]{\mathfrak t _{#1} (#2)}
\newcommand{\ti}[1]{\mathfrak t _{#1}}

\newcommand{\degexp}{\textup{deg}} 

%\newcommand{\Val}{\mathfrak v}
%\newcommand{\Tmd}{\textup{Time}}
%\newcommand{\Spc}{\iota}
%\newcommand{\Spcd}{\textup{Space}}
%\newcommand{\SpcTm}{\text{@}}
%\newcommand{\SpcTmd}{\text{SpaceTime}}
%\newcommand{\Hum}{\mathfrak h}
%\newcommand{\Humed}{\textup{Hum}}
%
%
%\newcommand{\Prc}{\textup{Process}}
%\newcommand{\HEnc}{\textup{HomoEncode}} 
%\newcommand{\IEnc}{\textup{ImEncode}} 
%\newcommand{\FEnc}{\textup{FullEncode}} 
%\newcommand{\HDec}{\textup{HomoDecode}} 
%\newcommand{\IDec}{\textup{ImDecode}} 
%\newcommand{\FDec}{\textup{FullDecode}} 


\input{Custom}


%\newcommand{\ngx}[2]{\mathfrak n_{\textup{#1}_{#2}}} 
%\newcommand{\ngx}[2]{\mathfrak n_{\textup{#1}_{#2}}} 


\newcommand{\exlp}[1]{\textup {hExt}_{\textup{lp}_{#1}}}
\newcommand{\exrp}[1]{\textup {hExt}_{\textup{rp}_{#1}}}
\newcommand{\exla}[1]{\textup {hExt}_{\textup{la}_{#1}}}
\newcommand{\exra}[1]{\textup {hExt}_{\textup{ra}_{#1}}}

\newcommand{\pexl}{\textup {ex}_l}
\newcommand{\pexr}{\textup {ex}_r}

\newcommand{\sepimd}{\brown{|}}
\newcommand{\sepffd}{\red{|_{FF}}}
\newcommand{\sepld}{\purple{|_{PL}}}
\newcommand{\sepesp}{\blue{|}}

\newcommand{\vgx}[2]{\mathfrak v_{\textup{#1}_{#2}}} 
\newcommand{\igx}[2]{\iota_{\textup{#1}_{#2}}} 
\newcommand{\ngx}[2]{\mathfrak n_{\textup{#1}_{#2}}} 

\newcommand{\Enc}{\textup{Enc}}
\newcommand{\Dec}{\textup{Dec}}


\newcommand{\vg}{\mathfrak v_{0\textup{-}Gap_1}} 
\newcommand{\ig}{\mathfrak n_{0\textup{-}Gap_1}} 
\newcommand{\iogi}{\iota_{0\textup{-}Gap_1}} 

\newcommand{\vog}{\mathfrak v_{O0\textup{-}Gap_1}} 
\newcommand{\iog}{\mathfrak n_{O0\textup{-}Gap_1}} 


\newcommand{\zg}{(0\textup{-}Gap_1)} 
\newcommand{\zbg}[1]{(#1\textup{-}Gap_1)} 
\newcommand{\oog}{(O0\textup{-}Gap_1)} 
\newcommand{\obg}[1]{(O#1\textup{-}Gap_1)} 


\newcommand{\vbg}[1]{(\mathfrak v_{#1\textup{-}Gap_1})} 
\newcommand{\vobg}[1]{(\mathfrak v_{O#1\textup{-}Gap_1})} 

\newcommand{\ibg}[1]{\mathfrak  n_{#1\textup{-}Gap_1}} 
\newcommand{\iobg}[1]{\mathfrak  n_{O#1\textup{-}Gap_1}} 

\newcommand{\idg}{\mathfrak  n_{Gap_1}} 
\newcommand{\iodg}{\mathfrak  n_{OGap_1}} 

\newcommand{\ibgi}[1]{\iota_{#1\textup{-}Gap_1}} 
\newcommand{\iobgi}[1]{\iota_{O#1\textup{-}Gap_1}} 

\newcommand{\idgi}{\iota_{Gap_1}} 
\newcommand{\iodgi}{\iota_{OGap_1}} 

\newcommand{\extenc}{\textup{ExtEnc}} 
\newcommand{\extdec}{\textup{ExtDec}} 


\newcommand{\Jz}{J_\emptyset}

\newcommand{\patc}[2]{\textbf{at}_{\frac {#1} {#2}}}
\newcommand{\jat}[2]{#1\textbf{at}_{#2}}
\newcommand{\Jat}[3]{\textbf{at}_{#1} \frac {#2} {#3}}


\newcommand{\co}{\mathfrak c}


\newcommand{\z}{\mathfrak z}
 \newcommand{\ii}{\mathfrak i} 
\newcommand{\p}{\mathfrak p}

 \newcommand{\ir}{\purple{\mathfrak i_r}} 
 \newcommand{\il}{\red{\mathfrak i_l}} 
 \newcommand{\iim}{\brown{\mathfrak i_g}} 
 \newcommand{\iimi}[1] {\brown{\mathfrak i_{#1}}} 

\newcommand{\Fa} {\mathfrak F}
\newcommand{\Spin} {\{\uparrow,\downarrow\}}
\newcommand{\inv} {inv}

%\newcommand{\jard}{\mathcal{B}\textup{T}} 
%\newcommand{\Jard}{\mathcal{B}\textup{Tid}} 
 

\newcommand{\bounce}{\textup{bounce}} 

\newcommand{\blue}[1]{{\color{blue} #1}}
\newcommand{\red}[1]{{\color{red} #1}}
\newcommand{\purple}[1]{{\color{purple} #1}}
\newcommand{\green}[1]{{\color{green} #1}}
\newcommand{\brown}[1]{{\color{brown} #1}}
\newcommand{\lightblue}[1]{{\color{cyan!50} #1}} 
\newcommand{\darkblue}[1]{{\color{blue!70!black} #1}}

\newcommand{\Z}{\mathcal{Z}}
\newcommand{\D}{\mathcal{D}}
\newcommand{\C}{{\mathbb{C}}}
\newcommand{\oN}{\overset{\circ}{\mathbb N}}
\newcommand{\mab}[1]{\bm{(\mkern-5mu|}#1\bm{|\mkern-5mu)}}

 \newcommand{\interp}[1]{\left\llbracket #1 \right\rrbracket}



%\newcommand{\smd}{\mathrel{\ooalign{\hidewidth$\rightarrow$\hidewidth\cr$+$}}}
\newcommand{\smd}{\overset{\rightarrow} +}
\newcommand{\prd}{\overset{\rightarrow}\Pi}

\newcommand{\bJ}{\overset{\leftrightarrow} J}
\newcommand{\lJ}{\overset{\rightarrow} J}
\newcommand{\rJ}{\overset{\leftarrow} J}

%\newcommand{\sepffd}{\red{|_{FF}}}

%{\mathrel{\ooalign{\hidewidth${\overset{ \, }{\rightarrow}$\hidewidth\cr$\Pi$}}}

\newcommand{\pld}{\mathrel{\ooalign{\hidewidth$\rightarrow$\hidewidth\cr$+$}}}
\newcommand{\pllr}{\mathrel{\ooalign{\hidewidth$\leftrightarrow$\hidewidth\cr$+$}}}
\newcommand{\plr}{\mathrel{\ooalign{\hidewidth$\leftarrow$\hidewidth\cr$+$}}}
\newcommand{\pldc}{\mathrel{\ooalign{\hidewidth$\rightarrow$\hidewidth\cr$+\mkern-11mu+$}}}
\newcommand{\plrc}{\mathrel{\ooalign{\hidewidth$\leftarrow$\hidewidth\cr$+\mkern-11mu+$}}}
\newcommand{\pllrc}{\mathrel{\ooalign{\hidewidth$\leftrightarrow$\hidewidth\cr$+\mkern-11mu+$}}}
\newcommand{\upl}{\mathrel{\ooalign{\hidewidth$\-$\hidewidth\cr$+\mkern-11mu+$}}}
\newcommand{\uplu}{\mathrel{\ooalign{\hidewidth$\Uparrow$\hidewidth\cr$+\mkern-10mu+$}}}
\newcommand{\tdf}{\mathrel{\ooalign{\hidewidth$=$\hidewidth\cr$\timestikz$}}}

%\newcommand{\rs}[1]{_{|_^{\mathfrak s} #1}}%\mathrel{\ooalign{\hidewidth$=$\hidewidth\cr$\times$}}}

\newcommand{\rs}[1]{_{|^\mathfrak s}} 
\newcommand{\rr}[2]{_{|_{#1} #2}}%\mathrel{\ooalign{\hidewidth$=$\hidewidth\cr$\times$}}}

\newcommand{\re}[2]{_{{|_{#1}}  #2}}%\mathrel{\ooalign{\hidewidth$=$\hidewidth\cr$\times$}}}
\newcommand{\res}[3]{_{|_{#1}}^{\mathfrak s_{#2}}  #3} 

\newcommand{\li}{{\mathfrak{l}}}
\newcommand{\cJ}[1]{can\textup{J}_{#1}} 
\newcommand{\cJD}[1]{oldcan\textup{JD}_{#1}} 

%==================================================================
\newcommand{\dmat}{\scriptscriptstyle\begin{matrix}
|&   \\
  & - \\
 | &   \\
\end{matrix}} 
\newcommand{\pmat}{\scriptscriptstyle\begin{matrix}
&  & \\
 - &&  -\\
  &|&  \\
\end{matrix}}
\newcommand{\qmat}{\scriptscriptstyle\begin{matrix}
& |&  \\
  -&&-  \\
  &&   \\
\end{matrix}}
 \newcommand{\bmat}{\scriptscriptstyle\begin{matrix}
& |  \\
 - &  \\
  & |  \\
\end{matrix}}

%====================================================================== 
 
\newcommand{\scom}{\color{blue}} 
\newcommand{\cz}[2]{\Delta_{#1}(\textsc{Z}_{#2})} 
 

\newcommand{\tensorprod}{\bigotimes}
 



\def\squareforqed{\hbox{\rlap{$\sqcap$}$\sqcup$}}
\def\qed{\ifmmode\squareforqed\else{\unskip\nobreak\hfil
\penalty50\hskip1em\null\nobreak\hfil\squareforqed
\parfillskip=0pt\finalhyphendemerits=0\endgraf}\fi}
\newtheorem{theorem}{Theorem}
\newtheorem{lemma}[theorem]{Lemma}
\newtheorem{corollary}[theorem]{Corollary}
\newtheorem{proposition}[theorem]{Proposition}
\newtheorem{definition}[theorem]{Definition}
\newtheorem{conjecture}[theorem]{Conjecture}
\newtheorem{fact}[theorem]{Fact}
\newtheorem{claim}[theorem]{Claim}
\newtheorem{informaltheorem}{Théorème}
\renewcommand{\theinformaltheorem}{}
\newenvironment{proof}{\begin{trivlist}\item[]{\flushleft\bf Proof }}
{\qed\end{trivlist}}
\newenvironment{proofof}[1]{\begin{trivlist}\item[]{\flushleft\bf 
Proof of~#1 }}
{\qed\end{trivlist}}
% MM 210906
\DeclarePairedDelimiter\ceil{\lceil}{\rceil}
\DeclarePairedDelimiter\floor{\lfloor}{\rfloor}
\title{Atelier nu de réflexions: Que pourraient être les  théories de l'émergence:  pour un  essai sur le doute à l'usage de ceux qui  paraissent  savoir %qui  savent ou peut etre pas
\\ \emph{Understanding understanding}
\\ Attention: objet vivant!}
%\title{ Essai Entre Philosopher les Mathématiques   et Mathématiser la  Philosophie [Ceci Formalise   Ma Compréhension]$_{[\textup{Grenoble} \textup{ Mehdi}_J \, 2025]}$ \emph{A Proposal for  Mathematical Humanities}
%}%(from sets)}
\author{M.M}
\begin{document}
\maketitle
 
$$\begin{matrix}
\text{Essayons }&\leftarrow&-&-&                              -&\text{quand il est}&-&   &\text{explicable} &\rightarrow \\
\uparrow &		&	&	&	&	&	&		&|\\        
|&               &                &\text{de comprendre}&      -&\text{comprendre}&&&| \\
\\          &	               &                   &  |&-&|&\\
\\&|&-&-&-&-&-&-\\
|&\text{ quand d'explication}& \text{enseignable} \\
-&-&\downarrow &\\
\end{matrix}
$$

\section{Résumé de la thèse émergeante}

Dans cet essai, ce sont  nos failles, nos déchirures, qui font crier en nous ce besoin d'existence, à la vérité vraie ou à l'indépendance, à ces constructions explicatives qui semblent plâtrer les fissures. Ce ne serait pas la cohérence qui produit la signification, mais    la blessure qui la rend nécessaire localement.

  A chercher des structures qui supportent  l'interruption c'est elle qui semble  émerger comme mère universelle  de toutes les structures 
 dont je ne peux nier la localisabilité et qui paraissent  transmissibles    puis  de celles localisées qu'on dira constructibles.
 
 Ici, le possible n’existera pas  sans un minimum de temps, et le temps n’existera pas sans l'interruption.
 
 L'imperfection sera partout présente: quand un fracture ouvre devant moi un hier et un demain,  je n'ai pas la garantie que les mots 
que je peux utiliser pour expliquer mon lien avec hier sont assez riches pour couvrir les possibilités de demain mais je peux essayer en répétant et interagissant avec ce que je ne sais pas reproduire.
 %Quelque part, c'est elle qui coupe produisant un avant et un après, qui suspend produisant l'intervalle et puis qui reprend produisant le rythme.
\footnote{ [a chercher plus tard entree math] : On voudra probablement construire des homologies du sens ou les cycles sont des stabilisations locales, les bords sont des interruptions, les classes d’équivalence sont des interprétations, et ou le symbole ||  sera comme une barrière de support compact,
qui permet de définir une cohomologie locale.}

\section{ Le trait ses formes et ses manières}
Le trait pour nous aura par défaut deux formes $-$ et $|$, un sens "normal" dont le degré de normalité dépendra des cultures, du temps, de l'espace ou d'autre choses. Pour le narrateur ici, ce sera $\rightarrow$ et $\downarrow$ comme une manière implicite de produire la forme qui sera la première nommée quand on discutera des séquences de  manières  de faire que j'imaginerai plus tard  car première a avoir émergé comme idée dans la formation de ma pensée.
  
 De loin le trait ressemble à un point qui pourrait tout aussi bien être droit ou couché. De près en n'importe quel endroit  en fonction de comment on regarde [zoome],  on verra de l'espace et plusieurs  traits décrivant celui qui semblait un, continu. Regarder finement interrompt la simplicité assumée par naturalité et force à reconcevoir l'unité. Une question ici quelque-part sera d'essayer d'explorer tout les manières de raisonner avant d'oublier l'épaisseur du trait.
  
   \subsection{Demi corps ouvrant  "["  et demi corps fermant "])}
 
 Dans cet essai je veux tenter l'expérience de ne présupposer vrais ou admis, les pré-requis  mathématiques construits sur ensembles, catégories,  ou toute autre axiomatisation formalisable, qu'entre  les parenthèses arrondies "(" et ")" [ces formes aux courbures étranges  dont la continuité semble résister à la résolution de mon affichage]. J'utiliserai donc pour les aparté ces formes de [au moins]  trois traits faites, comme substitution qui me semble  plus accessible  en imagination de gestes "[", "]" qu'on appellera les demi-corps.
 
 Un demi-corps a en haut deux bras dont je ne vois qu'un "-" , au milieu   une âme ou une tête qui cache un coeur  "|" et en bas deux pieds dont je vois qu'un "-".
 A deux ils enserrent ce dont la disparition pourrait passer inaperçue quand on regarde de loin.
 
 Nous voyons déja ce choix fort de multi-ancrage en équipant le symbole abstrait de geste, de manière de faire, d'images et en imaginant la séquence de liens associés pour tous ceux qui le voient et s'imaginent le tracer et aussi tous ceux que je peux envisager comme ayant pu le tracer.
 
 Ainsi quand je dis trois traits je pourrais plus tard dire que c'était un trait se déplacant ou  un trait dans un mouvement et un autre se découpant ... Je m'autoriserais ainsi [ou m'imposerais à savoir gérer] de découvrir plus tard que ce je voyais tantôt comme pied m'apparait, plutôt comme bras, au vu des échanges depuis. Une notion pour moi ne vit dans que dans une temporalité relative liée à une émergence estimée. 
  Cela ne m'empêche pas de tenter ici d'expliquer ma vision du monde ou plus précisément, ce dont je ne peux pas au vu de mon vécu, distinguer [imaginant oublier toute tentative de pensée qui voudrait émerger après le moment [évalué] ou j'ai écrit le "Cela" commençant cette phrase],  de ce qui aurait pu être candidat à faire partie d'une objectivité. 
 
 On verra plus tard, qu'un corps séparant ce qui m'apparait comme deux définitions indépendantes doit garder en lui le potentiel de retrouver un lien caché qui ne m'apparaitra que plus tard.
  
  \subsection{Enseigner vu comme se transcender.}% pour faire émerger une existence ressentie pairée d'indépendence}
  Apprendre à quelqu'un à  reproduire un geste c'est  gagner pour un instant le don de la prophétie, c'est contraindre cette imprévisibilité de la mort en ce disant 
  d'un coup "j'existe, je sais je transmet". Et ce,  comme un bout de causalité certifiable par traits  puisque je peux dire regardez! je peux faire faire la même chose à tous ceux de ma communauté.  J'ai donc atteint  ainsi ce qui me dépasse, me sépare et me justifie  [pour ceux qui agréent] ce que je vais clamer, mais aussi du coup, va me rassurer en me faisant sentir que j'ai insufflé, dans une rupture créatrice  d'avant et après [l'apprentissage], une famille qui pousse dont j'aimerais me clamer l'auteur. Et ce,  même si par les faits  ce que je clame en fait c'est la croissance de ma croyance en le fait que de mon existence la technique s'est détachée.
   
   On va voir émerger dans l'humanité ces enseignements de techniques de traits comme quête d'expressivité de ce ressenti de séparabilité qu'on appelle parfois objectivité.
   
   On va à un moment se concentrer sur toutes ces techniques de réécritures  rudimentaries qui se revendiquent  chacune de ce même nom noble ou effrayant 
   pour les gens c'est selon : les "mathématiques".  Je dirais que le premier point commun est une quête de contractibilité dans l'expressibilité de la justificatiion, en réduisant la forme et limitant l'expression. Cherchant à faire  exister une justification qu'on devra chercher parmi des règles simples afin de valider l'émergence de tel ou tel motif dans un ouvrage donné. Gonflant ainsi l'égo de ceux qui les pratiquent puisque   sous réserve de taille [nombre de symboles] tout apôtre a de fait  un pouvoir constant d'oracle.
   
%   \subsection{Vouloir ne pas détruire}
%Dans tout cet essai et probablement au delà je m'opposerais

\section{Intuition}
%Ce texte explore comment peuvent émerger suctuitioncessivement :
%’espace, le nombre, la vérification, le calcul et la causalité,
%à partir de simples gestes et traces.
Nous allons tenter ici, de décrire les naissances   du sens, de l’espace, du nombre, de la vérification, du calcul et du quantique.
C’est une proposition de  cosmogenèse cognitive sous forme d'un laboratoire de réflexions tentant une dérive d'explicitations  des implicites dans  les transformations de la pensée en gestes transmissibles .

%Une théorie du sens comme dynamique transmissible imparfaite
%où les mathématiques, la logique, le calcul, et la physique
%ne sont que des cristallisations locales.

%\subsection {Intuition}
Pour une intuition rapide, dans cet essai,  réussir  un enseignement  c'est   établir un langage vu comme un système de critères de validation de la répétabilité  potentiellement approximatifs et fabriquer des processus   dont on ne peut borner le dernier succès tant qu'on garde ce système. Ces processus vont être anthropomorphisés comme être interactifs imaginaires bi-localisables qui peuvent essayer de  confirmer ou d'infirmer quand on le leur demande avoir utilisé un langage qu'on leur désigne par une trace en parlant d'une \textbf{autre}  trace qu'on leur impute.\footnote{
on voit déja émerger quelque chose qui semble être d'un niveau inférieur mais qui est vu dans une dynamique, ce qui est  primoridal puisque les testabilité devra être vue dans un essai de raffinement qui se répète.}

On cherchera  ici a construire progressivement un moyen de donner une forme avec une représentation rigide  à ce qui n'est pas encore, mais qui pourrait être, 
en s'appuyant sur la restriction qu'il soit lié, soit à une trace du passé comme reproduction d'un apprentissage reçu ou une découverte soit une trace du futur comme intention de transmission par la trace de ce qu'on pense pouvoir être une invention ou de ce qu'on a soi-même appris. On va essayer d'approcher l'infini par le non-fini, en acceptant de  l'accepter, en fonction du sujet tentant de l'appréhender et  tel qu'il est dans son imperfection.  On va proposer ensuite un modèle de collapses  pour reproduire les langages usuels. On  devrait voir   apparaitre  en relachant l'universalité de la norme pour des constructions locales et en entrelaçant localement  la composition algébrique et la décomposition analytique  maintenant un lien  à la trace locale sans avoir besoin de savoir si elle vient du passé ou du futur.  \footnote{l'idée est de figer les suffixes va faire tendre vers les *algebres approchéees et la communtation comme localité et figer les préfixes va faire tendre vers le produit tensoriel des espaces comme localité avec l'espoir d'être dans le triangle formé audessu par les jeux de pseudotélépathie quantiques au dessus de RE et co-RE}   

Je veux  tenter un méta-langage pour décrire toutes les manières que j'ai d'imaginer, sans lynchage ni sacralisation, un imparfait qui ne cesse de se reconstruire lui même en injectant dans les décontructions qui ont des traces et en projetant à partir des constructions qui peuvent en créer mais en alternant les tentatives de rejoindre l'un des ces sens traditionnels par l'une ou l'autre méthode.

Le problème de la dimentionnalité : le trait peut imaginer un point d'un coté ou de l'autre potentiellement avec un espace qui les sépare  ou qui le coupe tout ca en maintenant sa direction la changeant ou donnant deux directions. Le point s'il pouvait exister tout seul serait problématique, c'est la raison du rajout primitif de la pré-cognition de cette notion de transmissibilité ou d'enseignement et à partir du moment ou on a le questionnement on aura cette structure sur les traits, ne serait ce que comme temps d'attente avant réponse ou changement de sujet. Un point ni dessinable ni dessiné est hors sujet car trop parfait pour la tentative de compréhension dans le mouvement de l'imparfait.

Je conjecture qu'on va probablement  voir apparaitre des structures de communautés bi-typables par deux systèmes de types dynamiques qu'on peut voir comme indices et exposants avec des mécanismes d'alternance localement définis et des systèmes de traductions qui ressemblent à s'y méprendre à des équivalences. Par contre les définitions simples apparaissent quand on borne un coté ou l'autre par des processus interactifs d'échange qui deviennent sources de définitions formelles typées usuelles.\footnote{Par désir de sincérité ceta atelier va tenter d'essayer de rester au plus proche de la forme que j'imagine à la pensée: dense, riche, souvent difficile à déchiffrer, pouvant même paraitre chaotique ou désordonné, avec une lecture-écriture sans cesse revisitée, mais tentant de respecter les faits qui paraissent essayant de  subsumer sans écraser, requérant des traces justifiant l'effacement permettant d'imaginer la reconstruction en remontant le temps.}

La notion d'interruption semble émerger comme l'acte  minimal permettant la  co-construction de référence, d'espace, et de causalité avant même le que symbole ou le modèle   formel ne puisse exister: l'interruption marque la  référence, la référence génère l' espace, l' espace  génère l'impression locale de  compréhension. 

Les deux mécanismes fondateurs appelés partage  - et composition | expriment le lien entre l'interruption et l'explication : comment voir que l'interruption est elle -meme un explication et que l'explication est elle même une interruption : ce sont deux faces d'une pièce impalpable que nous allons théâtraliser pour lui donner de la consistance et de la transmissibilité en assumant l'implicite comme source d'information et un autre implicite comme source de formulation.

Pour tenter de capturer, ne serait-ce qu'en essence [symboliquement] toutes les significations, dans leurs sens [ultimement tracables] on devra débrider notre imagination.
% \emph{ A on devra débrider et le risque de noyade on devra assumer mais la beauté intrigue, je me laisse donc tenter.}
 \emph{ A se laisser tant errer, le risque de noyade devient réalité, mais devant 
 %la beauté,
 ce cri si vrai,  du désir ressenti de ces mots contractés voulant s'étirer. En recherche de cause je me laisserai tenter avec le risque de dérive, comme posture assumée, devenant le cheval d'une expérience de pensée, qui a un peu l'air d'une folle chevauchée.}% on devra débrider et le risque de noyade on devra assumer mais la beauté intrigue, je me laisse donc tenter.

\section{Exercice de style ou idée de fond ? Que devient comprendre la vie avec des mots  quand on accepte la vie des mots dans une autre temporalité}

Si vous entrez dans cet essai et que vous avez envie de vous laissez prendre au jeu acceptez l'espace de cette lecture [ou refusez, je n'imposerai jamais rien à personne ou en tout cas je ferais de mon mieux pour éviter d'avoir à le faire] que nous sommes le temps des idées qui nous traversent.

Nous allons tenter de voir l'émergence de la compréhension transmissible présignifiée comme consensus intime balloté par des interactions

Nous nous intéresserons aux mathématiques, à l'informatique, à la physique quantique, à la physique,à la causalité et aux langues naturelles pour les voir comme processus similaires à travers deux prismes:
\begin{itemize}
\item celui de la \textbf {composition} ou |:
L'amour [de soi ?], la schizophrénie,   la réflexion la copie la reproduction le multiple, l'être qui se [com]plait dans sa répétition, ou la simplicité de répéter et attendre, quand on croit sincèrement savoir  ce qu'on dit, c'est cette équivalence avec les autres qui me semble être ce qu'on ressent au fond de nous, c'est la tendresse de la dilatation. C'est le trait, rayon de lumière  vu comme un mur sur le quel rebondit le trait.[Penser féminin]
\item celui du  \textbf{partage} ou -:
La décomposition , la paranoia , le réductionnisme , le sacrifice ou la création d'espace, le désir de posséder, le fractionnement. C'est le trait vu comme un rayon de lumière comme un point qui nous suggère une simplification réductrice, et qui fait émerger ce qui est plus petit que ce qu'on croyait l'étre.[Penser masculin]
\end{itemize}

L'idée directrice de la méthode est de se dire que je n'ai pas encore dit ce que je suis en train de dire ou de penser et donc il manque toujours intrinsèquement quelque chose à mon expression dans l'appréhension des potentialités de son expressibilité pour quelqu'un d'autre ou moi même pour le quel il ne serait pas parfaitement reproductible et donc similaire à une manière rigide de pensée. 

Conscient de l'imperfection de la description de ce à quoi on admet une liberté qui nous dépasse en prédictibilité ou en connaissance de l'histoire ou de l'intérieur on va anthropomorphismer ce qui n'est pas parfait et en ce qui me concerne dans cet essai je prends   la décision, vous l'imposant du reste et je m'en excuse, d'exclure le champ %eschatologique 
du divin ou de l'induction sans modèle de l'objet de ce que j'espère comprendre ou appréhender. Dans le sens compositionnel interruptible, le parfait n'existera que comme croyances et ne prendra sens que  dans des cadres quand et tant que leurs comportement sera consistant avec une dynamique ultimement descriptible et  je ne peux m'empêcher d'ordonner mes plus grandes certitudes au fur et à mesure qu'elles émergent dans mes pensées.\footnote{ Je tiens à noter que cette exclusion me chagrine, et n'est nullement un choix arbitraire mais l'expression d'un ressenti d' impuissance à avancer clairement qui me rend le sacrifice inévitable pour la survie de mon questionnement permanent même si je me dis que ceux qu'il n'y a  que ça que ca intéresse gagneront du temps en évitant de me lire.}

\emph{Cet atelier peut être lu comme une invitation à revoir sa pensée quand on peut, pour voir si celle-ci ce matin avoit desclose.}

\subsection{Pourquoi "que pourraient être" comme contrainte imposée}

L'existence est bien trop dure pour moi à clamer qui suis-je pour aux gens que penser ! Choisir c'est en partie  projeter, et la copie locale que je garde ne saurait qu'imiter. Le sens existentiel pour moi nécessairement disparait et se fond pas essence dans la diversité. Le chemin unique  ou la négation du possible ne peut être  que limite, du moment que d'abord j'insiste à rappeler que je ne peux qu'essayer de limiter l'effet de mon imperfection. 

 Je ne sais voir l'intention dans  la trace que par imagination et l'indépendance tout court ne saurait remplacer cette humanité  qui fait questionner.    
Peut -etre est  du fait que je ne saurais lister explicitement ce qui me parait sincèrement pour l'instant "n'avoir rien à voir" , peut être est-ce la réticence de se soumettre par essence à ceux qui posent ces barrières qui voudraient bloquer le pouvoir savoir, peut-être...

Entre un coeur qui bat et un cerveau qui impulse, je me sens balancer entre deux  rythmes à fréquences pourtant si distanciées mais si présents partout tellement entremêlés. J'imagine l'un dans l'autre, dans un sens et puis l'autre, ancrant mes pensées  comme des interruptions comme une séquence de traits qui ne sait ou elle est et qui pourtant avance sans discontinuer en imaginant l'avant mais aussi l'après.

Si l'échelle est fixée pour un clic de pensée,  plus ces deux traits sont loin plus ma confiance est grande et je suis rassuré quand ils sont hors de portée. Mais la ou je me vois, pour ce temps de pensée, c'est juste avant deux traits qui veulent représenter le milieu du chemin stoppé pour une idée que j'ancre en faisant lien avec celles que je refais avec ma mémoire fragile en forçant à coller les traits que j'imagine avec ceux en qui je croyait. 

Les libertés proposées aux autres [donc à tous comme diversité] sont des contraintes imposées à chacun  [comme communauté] ou à tous ceux qui ont envie de jouer à ce mode de pensée. J'estime en effet à ma charge  d'apprendre à communiquer avec celui qui prend comme base pour compter, son nombre d'années.\footnote{J'imagine ou plutot j'espere que  les limites tronquées avec une norme imposée permettront de retrouver le deuxième niveau  universel des langages formels $\Pi_2$ }
 
 La posture adoptée ici inverse l’idée de synchronie par défaut. Au lieu de se dire : je n’ai pas besoin de préciser que la multiplication est prioritaire sur l’addition, tout le monde dans ma communauté le sait, je poserai au contraire que si je ne le dis pas explicitement, alors tous les parenthésages implicites sont permis, et que je me dois d’assumer les explications qu’ils rendent nécessaires.


 
\subsection{Et la simplicité ?}
Plutot que de rendre premier la simplicité de lecture, l'effort est d'abord mis ici sur l'authenticité ressentie, ou comment coller au mieux ce que je pense et ce que je dis. C'est d'ailleurs, je le pense, pour qui veut comprendre, d'abord l'effort de penser [soi-même] qui doit être rendu simple, avant pendant et après, que l'on retranscrive des pensées décédées. Je me dis même  parfois que j'ai du mal à résoudre par la paire [penseur|dactylographe] celui qui ne peut s'expliquer sa pensée sans avoir recours à celle des autres. Bref, un équilibre semble à trouver dans la valorisation entre l'effort et la simplicité...

En tout cas , c'est la certitude auto-suffisante de la constante indépendentte avec la tranquilité qu'elle amène que je n'arrive pas à appréhender comme une unité et que je lie donc aux battements de mon coeur pour marquer  les instants d'occurence de pensée et déduire d'entre eux cette notion de durée.
Quelque part je dois rajouter quelque chose pour faire quelque chose de cette arrogance de celui qui clame "ça y est j'ai compris, je n'ai plus besoin d'écouter", 
car je ne peux donner un sens à une partie que si elle résiste à l'extensionnalité : si on vérifie comment l'idée va se comporter en rajoutant du possible au milieu ou on la conçoit et en acceptant l'existence même infime de la possibilité de se reconcevoir. Je ne prendrais donc pas les chemins les plus courts, mais si vous êtes ici, je me dis que nous partageons peut-être cet amour des détours.

\subsection{Pourquoi la poésie?}
Si le vide, le rien ou le point me rappellent l'imperfection de mon contrôle d'une notion de stabilité ou d'espace d'origine, si la mort ou l'infiniment grand me rappelle mon incapacité en général a parfaitement prédire le futur ou l'indépendance ou la variabilité parfaite ou de temps final, la poésie s'impose à moi pour me rappeler qu'écrire c'est échanger avec les mots et qu'ils ont leur mot à dire dans l'habillement de similaires, mots, entre les quels je m'amuse à les placer pour étayer un argument. La poésie est la pour me rappeler la vie des mots et éviter l'arrogance de les croire  simple outil rigide dont je contrôlerais le sens, ils peuvent se jouer de moi comme je joue avec eux et c'est dans cette interaction que nait mon raisonnement comme un ordre imposé rompant la superposition. A me laisser porter par leurs arrangements permet de rappeler qu'ils ne sont pas soumis mais forment des vagues ou on peut naviguer.  

\subsection{ -Partir de l'imparfait subissant la localibilisation}
 Imaginez  XX au huits bords pointés est-ce un fond est ce une forme ou une femme incarnée?
 En tout cas ce que je vois c'est ma perplexité sur comment recoller ces 4 bords libérés, par le désir de faire un de ces deux moitiés.
 
Pour retirer le moins projetons sur le plus et faisons le de surcroit comme exercice matinal afin de rythmer le sens de cette cohérence née de ce processus 
d'autorégulation par questionnement et échange, avant  réflexion.

Si la forme et le fond je ne peux distinguer si le tout et le rien semblent une vague idée pourquoi ne pas partir de l'imparfait ou je peux bien me voir même me transformer de l'idée de moi -même comme étant quelque trait à ce trait du locus né de mon interpellation je veux associer un conte de la signification.  
 
 Et j'espacerais donc en faisant ce dessin relégant la responasibilité d'interprétation du plein et du vide à ceux qui sentent pouvoir et vouloir le faire.
$
{\scriptscriptstyle\begin{matrix}
- &   & - &   & - &   & - \\
  & | &   &   &   &  | &   \\
- &   & - &   & - &   & - \\
\end{matrix}}$
 
XX   :  Quand  on voudra ici oser une histoire on dirait qu'au début était la chaine parfaite, cette séquence infinie de couples entrelacés mais dont la perfection faisait rejeter sans même les entendre  tous ces prétendants,  ces définitions vouant  à localisation. Un jour un accident fit tomber un bras et tentant de reproduire sa propre création ce bras libéré causa séparations, décrochant ça et là, créant une dynamique,  une danse incessante de restructurations qui ressemble bien fort à des ondulations.  Ce ne fut plus une mais un jardin de significations qui ainsi naquit de ce qui fit l'évènement.\footnote{ On peut imaginer les approximation locales de la structure par blocs de 4 en partant de la moitié : XXXX, $\begin{array}{c} XXXX\\  XXXX \\XXXX \\ XXXX\end{array}$.}
   
\section{Méthodologie:}

Quoi de mieux que l'humain pour représenter ce qu'on ressent mais ce qu'on n'arrive à dire, pour incarner la multitude et le changement de sens comme signification, cette main qui nous livre le livre et dont on voudrait parler va être le narrateur de l'histoire qu'elle apporte. La main est déja trop riche nous épurerons encore plus plus tard appelant anthropomorphisme ce qui pourrait sembler  humain dans un cadre comme un personnage de pièce de théatre.\footnote{anachroniquement ce nous ca sera le narrateur, moi qui écrit et Dolorès le rôle que je joue par défaut dans l'histoire que je tente d'écrire.} 
 
\subsection{ Le conte théatralisé comme support de pourquoi les quatre coins de cet impossible qui nous définit } 

La méthodologie présentée ici ne voit  les objets finis, que comme troncation contextualisée rigidifée non compositionnelle. Ils ne prennent vraiment sens que dans des environments dynamiques vu comme stabilisation de gestes interprétables sous contrainte de transmissibilité imparfaite. La cloture de toutes  les manières de comprendre n'a pas besoin d'être atteignable  mais la diversité avec la non exclusion à priori  de toute nouvelle méthode qui émerge est un pré-requis de la formulation "que pourraient être" comme mode de questionnement. Si vous n'aimez pas la gymnastique de l'esprit mais que vous avez eu le courage de lire jusqu'ici  je suggère [en pensant à ce dont je remet l'essentiel à plus tard] un GOTO section \ref{sec:simple} qui contiendra au fur et à mesure [je le pense sincèrement] que ça pousse la partie nettoyée pour penseurs n'aimant pas marcher sur l'eau.

Une communauté partage une forme d'expression d'échanges et une impression d'intercompréhension. Nous allons ici voir une communauté comme en même temps deux types:
\begin{itemize}
\item ATM: un anthropomorphisme ou plus précisément une séquence  indexable d'athropomorphismes qui la composent. L'indexation est d'abord une date de naissance suivi des chemins imaginaires de futures descendances. 
\item RTM une forme rigide de pensée comme la constitution ou séquence indexée de règles écrites convenues pas ses membres pour vivre ensemble et les histoires des récits historiques ou fictifs de ses membres. Leur indexation est d'borde une date de décès suivi de chemins discutables pour trouver l'anthropomorphisme originel de la méthode.
\end{itemize}

L'indexation est une forme rigide de pensée qui finit ultimement sur des traits unaires séparables et montrables dans une interaction assurant la satisfaction des membres de leur cohérence dans leurs croyances par des suites de questions et de copies d'encodages partiels locaux.
 

Toute communauté définie comme l'un peut être traduite en l'autre et vice vera.

Dans une communauté comme dans une pièce de théatre on peut placer des personnages, rôles d’interprétation qui peuvent être occupés successivement par un même observateur. Ces noms offrent un ancrage  comme des adresses  qui permettront de cristalliser si on s'arrête 4 communautés compatibles avec le sens transmissible elle symbolisent la multitude des interprétations possibles du lecteur de l'intention qu'on avec le quel le narrateur doit se dépatouiller en imaginant imparfaitement toutes les lectures que le spectateur pourrait faire de la pièce qu'il est en train de voir.
.
\begin{itemize}
\item Dolores : L'ame du narrateur dans le personnage principal c'est l'anthropomorphisme ou cet humain esquissé dessiné par des mots partiellement connus du lecteur mais interagissant comme un lien avec les autres personnages. C'est la douleur du personnage principal du drame théatral c'est le présent dans sa confusion dans le mélange des émotions c'est ce doute qui nous remplit, qui nousf fait ce qu'on est:  fragiles et imparfaits. Quand je pense à elle : c'est le moteur de la passion humaine qui nous pousse à continuer c'est le $\le$ de gauche à droite qui amène le plus petit à toucher ce qui le dépasse. C'est la quête incessante d'un sens inachevé. C'est dans la gestuelle de la forme le haut de cet arrondi du coté nord ouest du cercle qui se cache dans le 'd'  c'est $\downarrow \rightarrow$ ou simplement leur trace pour qui n'a pas vu le geste $|^-$. C'est le début haut gauche de cette parenthèse que j'imagine mais que mes capacités de représentation réellle réduira aux trois traits  chaleureux de l'accolade. La manière de lui associer une qualité et une qualité est par essence contextuelle car elle sera peut être la $k$ ième de $n$ parenthèse ouvrantes si je rigidifie l'associativité et force une interprétation parfaite complètement explicitée dans une réécriture en terme d'accolades ou terme d'accolades enrichi de symboles transportés que je dois amener à des des amis sachant leur donner du sens.  Comme méthode de lecture c'est se mettre juste en dessous ce qu'on estime et fair un pas en avant. Elle est le temps de celui qui cherche dans sa mémoire avant de parler et qui commence don par remonter le temps de l'histoire qu'elle va raconter son sens par défaut est l'opposé de celui des aiguilles d'une montre. Son coté naturel est la gauche. Son sommet naturel est le haut à gauche. Son type est RAR comme un anthropomorphisme connectant deux manières rigides de penser.



\item Qudsia : l’origine, la constance, la répétition dans sa complétion.  Quand je pense à elle : c'est la lettre 'q' avec le geste appris  $\leftarrow\uparrow$ ou $|_-$ mais c'est aussi  $<$ C'est le concept, l'origine de l'espace le début de la lecture.  La ou  l'homme ou plus généralement l'anthropomorphisme  commence ça lecture c'est les constantes les plus fiables. C'est cet infini bosonique qui permet d'utiliser  l'unité comme étalon c'est l'origine des quantités. Comme méthode de lecture c'est de partir du plus petit qu'on connait et fair un pas en avant. C'est la mémoire imaginaire cumulative de Dolorès qui est explorée de la plus récente à  la plus vieille son sens par défaut est l'inverse de celui de l'horloge. Son sommet naturel est le bas gauche. Son coté naturel est le bas. Son type est RA comme une ressource en train d'être épuisée par les anthropomorphismes.

\item Perséphone : C'est celle qui traverse le Styx c'est l'image de la mort, la dissolution, l’indifférencié, l’incertitude controlée parce que bornée par l'imaginaire de Dolorès.  Quand je pense à elle : c'est la lettre 'p'  écrit $ \leftarrow$ ou $ .^-|$. Mais c'est aussi $\le$ mais lu de droite a gauche et renoté $\ge$ pour synchroniser le sens effectif de lecture.
C'est le percept ou le phénomène c'est la direction de demain c'est la ou est la variable c'est une vue de ce   que l'on va dire  ou tout est rangé et qui peut fournir de l'aléa  qui peu sembler parfait au gens limités par leur principes ou des variables pures qui bougent tellement vite que  personne à mon échelle ne peut donner de leur avenir qu'one confiance en sa prediction. Elle est le temps de celui qui écoute l'histoire et son temps par défaut respecte le sens de l'horloge.
Comme méthode de lecture c'est se mettre juste au dessus de ce qu'on espère et faire un pas en arrière. Son coté naturel est le haut. Son sommet naturel est le haut droit.  Son type est ARA comme un canal de transmission entre deux anthropomorphismes.

\item Bénédicte : la limite, l’ultime cohérence, l’inaccessible. Quand je pense à elle : C'est la lettre b dans la gestuelle $\downarrow $  de trace $_-|$, c'est le $<$ lu de droite à gauche et renoté >  pour maintenir le sens de lecture de Dolorès. C'est ce qui pourrait lier la dernière fin au premier début  en tant que parenthèse.
C'est le pourquoi de l'impossibilité d'explication unique et la nécessité de la diversité. 
 C'est ce qui dépasse plus grand des plus grands  et qui crée le plus petits des plus petits dont ce qui m'est accessible par le symbole ,ou   ce dont je peux imaginer avoir le droit de parler en ayant des limites pour ma stabilité qui me permet de rester sur la même idée et à ma vitesse de distinction qui me permet d'envisager de répondre à la question la plus élémentaire en coupant ce que je sais en deux parties presque égales fabriquant un chemin d'identification efficace . C'est ce percept ultime de changement  fermionique dont la variabilité dépasse l'entendement qui donne naissance au concept premier. 
  C'est elle qui voit le miracle du poids des mots découpant de nouvelle pensées.  La bénédiction vue par un être synchronisable, c'est le pouvoir de réécriture : celui de se tromper de tomber et de pouvoir se relever pour reconstruire avec les traces d'hier dans le corps. C'est la ou vivent les métaphores et ou naissent les concepts. Comme méthode de lecture c'est partir du tout et revenir. Son sens naturel est de remonter le temps. Son coté naturel c'est la droite. Son sommet naturel c'est le bas à droite. Son type est AR c'est l'anthropomorphisme sans passé.

\end{itemize}

A elles quatre, elles forment un espace dynamique qui me semble pouvoir contenir chaque jour toute mes compréhensions  de la veille ou j'ai  pensé  déployer  l'émergence entre deux magnifiques constructions imaginaires de i'humanité bien ancrées  dans l'imaginaire commun :  le point et le trait cherchant à faire le cercle.  Cet imaginaire   est tellement ancré en nous que cela peut sembler difficile a certains de chercher à comprendrece que sont les objets finis de ces dessins.  On verra ici qu'on peut toujours essayer de sortir de la soumission si on le désire vraiment et prendre le risque de  décomposer le point ou compléter le cercle quand on a l'imagination et le désir de le faire.

Vous noterez que les deux premières ont été sommet avant de se voir cotés contrairement aux deux dernières.


%}\footnote{ce paragraphe explicatif est un commentaire suggéré par la machine après la lecture du draft mais que je trouve adapté et utile et n'entravant pas trop la liberté d'imaginer du lecteur.}

Dolores illustrera aussi le concept de la constante ou le nombre représenté[encodé]  comme quantité fiable  c'est l'optimisme c'est les nombres qui quand on leur demande se disent positifs et Qudsia la limite de son origine. Ma représentation de moi-meme précède ma représentation de la mort ou du bout du chemin.

Persèphone formera le percept  ou le nombre représenté [encodé] comme qualité  variable différentiante c'est le pessimisme la représentation de la mort précède la représentation de soi et Bénédicte sera sa limite.


 Les actions opératoire sur les verticales peuvent être associées à des mouvement d'adresse et celles horizontales  à des mouvement de valeur encodée.




%La fracture du double ou de l'étalage 

Un trait se balance dans un mouvement incessant  entre un presque rien :  point intérieur qui ne lui apparait que comme trace et un presque cercle ce tout extérieur qui le ramène à lui même.

L'axe Dolores Qudsia est celui de l'ellipse de la contraction de l'espace jusqu' à son invariant minimal pour rejoindre l'axe Dolores Persephone à l'inverse des aiguilles d'une montre. 

L'axe Dolores Perséphone est celui de la périphrase avec une extension de base et suit le sens des horloges. Mais n'oubliez pas qu'on peut ressentir  la mort comme  une interruption  de quelque chose  avant de trouver la vague de  mots qui tentera d'expliquer ce qu'elle coupe 

Ils  se rencontrent dans la métaphore sur la diagonale pour faire émerger un espace nouveau.
Un peu comme un bout local [un quart] d'une contraction de méduse.  Les reste du produit fibré est le prémisse d'un nouvel espace de pensée.

\footnote{On aimerait symétriser parfaitement sur place avec  ces 4 diagonlaes d'échange de bouts qui se repoussent mais souvent  Bénédicte n'est accessible que dans la métaphore et c'est vers elle que nait le nouveau bout d'espace nécessitant de tout reconfigurer par ce que par miroir le dernier émergeant étant le moins connu et donc le plus variable il sera central pour l'incapacité de perfection de Dolorès ou ce coin haut gauche inaccessible dans la boucle du d.}
\footnote{  Ajouter un et recommencer permet de récupérer l'arithmétique par répéter répéter qui fait naitre localement les opération d'ordre supérieur ou  hyperopérations [somme, produit, puissance, tour de puissance et ainsi de suite]. Bénédicte à la chance d'être après la fin des calculs elle n'a pas à calculer juste à expliquer le pourquoi du chemin qui peut faire ce qu'on voit elle défait les construction avec l'antisomme - et la division | qui sont une manière de  comprendre la suffisance des deux traits pour capturer l'expressible.} 


\subsection{Règle opératoire du théatre des imaginations}

Tout Si est le Quand d'un personnage [anthropomorphisme].

La dérive sera de faire un pas  comme une mise à jour de jardin suspendu [rajouter un fruit  a un jardin ou un trait a une liane identifier le passe et imaginer les futurs associes] d'une interruption :
Voir  les  horizontales de décomposition qui découpent et diffractent [dp] puis  recolleront a linfini  toutes les découpes imaginables [bq]
 et les verticale de composition qui detendent [bq] et contractent [qd] comme un "au suivant" de quelque chose que l'on trace. 
 
 Comme les mots qui brillent réflexions de miroirs sur égos explosant de multiplicité ou les mots qui brulent, qui déchirent, qui explosent , qui nous découpent en morceaux,  avec des bribes d'égo qu'on ramasse à la pelle.
 

Pour une pièce donnée, on fait état ou objet par défaut  ce dont la nécessité de décrire les états d'âme semble superflue dans une  lecture du passé tel que reconstitué, on jette donc par dessus bord ce qui n'intéresse pas le narrateur pour aller droit au but mais n'oubliez pas le préfixe pour une pièce donnée.
 
%Un état est un changement que l'on fige par absence de preuve de  la nécessité d'interaction future 
 J'explicite ici la morale cachée telle que je l'imagine ou ose l'espérer : je pense sincèrement que  la soumission à une norme commune n'est pas un prérequis pour une inter-opérabilité fonctionnelle localement satisfaisante.

\subsection{Méthode de respiration transcendentale}
%\emph{[attention melange ubiquitaire le discours intermediaire et form de l'article a evoquer ousupprimer]}
Cet essai tente une construction pour faire émerger le sens dans le parcours dans toute sa richesse, en naviguant dans l'ambiguité comme on l'est dans nos sens mais  proposera 
une voie construite plus facile pour l'utilisation en boite noire de ce qui veut expliquer les boites noires mais qui devrait je le pense être utile quand même.
 

 Le projet est de proposer un discours de la méthode sur les méthodes de discours qui acceptent l'interruptibilité sous réserve d'impression locale
 approximativement vérifiable par échange d'intercompréhension ou dit autrement une tentative de donner un sens raisonnable et raffinable de tous les  sens imaginables du presque-mot ambigu en train de sortir de la bouche. On voudrait donc à l'image de ce qu'on veut représenter autoriser le lecteur à prendre le train en marche ou il veut et à interrompre quand il veut en maintenant une possibilité de normalisabilité locale par extension de bouts imaginaires de complètion, ou par coupures de bouts passés inachevés  de simplification, et d'avancer sans tomber dans la trivialité tautologique ce qui permet de  ne jamais avoir l'impression de recommencer des origines qui du doute seront probablement  indéfinissable pour l'imparfait.
  
  On osera  donc un pré-langage graphique à trait pour exprimer   les langages encours  en essayant d'être au plus fidèle de ce qu'ils pourraient etre et ou on pourrait s'être trompé et auquel cas savoir  comment corriger en gardant juste à l'esprit les meilleures adéquations des gestes qu'on peut imaginer avec les discours qu'on voit. On le collapsera ensuite à sa forme optimiste assumant le succès de la méthode de la construction et la finalité de l'argument perdant la compositionnalité pour la simplicité avec une représentation imagée comme des carrés de cotés un dans un plan.
  
Il est rarement facile de se détacher ce qu'on connait ou croit connaitre pour arriver à minimiser l'effet de l'implicite et des présupposés sur notre manière 
d'appréhender la signification. Ce qui est proposé ici en première lecture n'est pas une linéarisation prémachée pour transmission réfléchissante mais une mise en raisonnance par couches. Pour comprendre: "il faut lire" et noter, lire et noter, lire et noter et puis respirer et prendre le temps de se relire faire les liens et noter les incohérences différentiantes de possibles pour fabriquer des paires de cellule. Et puis si on veut [GOTO  "il faut lire" ].

Ce n'est pas grave de ne pas comprendre, la compréhension est une dérive sous forme de ballade avec d'incessants allers-retours imaginaires  dans une mémoire fragile.   Une incompréhension est toujours susceptible d'être une contradiction nécessitant de mieux localiser le propos, de réorganiser, ou de raffiner  la ou on voulait l'ancrer pour séparer ce qui est compatible et ce qui ne l'est pas avec ce qu'on veut y mettre. N'oubliez pas que nous autres, avons la chance que même si l'on ne peut parfois la mesurer on peut sentir que toute interaction requiert un temps qui permet d'imaginer un système d'indices à tout ce qu'on imagine plausible si on l'équipe d'une manière  de la raconter qui ne peut se permettre le luxe de l'oubli que si quelqu'un connait l'histoire. 
 
 On n'a pas besoin de terminer de répondre ou de prétendre savoir, 
 
 Le trois permettant de mieux maitriser le un en le renforçant deux fois en l'infirmant par deux fois ce qui l'oppose ou en augment l'incertiude en arguant une fois pour chaque.

L'idée est que si on accepte d'aller assez profondément avant de bâtir ses certitudes on arrivera a voir des traces primitives d'interprétation élémentaires 
générant pour l'ensemble un jardin suspendu qui semble s'accrocher au ciel.
 
 Comme dans un film construit par bouts épars anachroniques l'effort est requis pour atteindre la stabilisation du sens, mais la joie sans effort l'est-elle vraiment?
  
\section{De l'émergence de l'espace ou de la définition du "trait qui imagine le point " qui pourrait bien être le "point qui imagine le trait".   }
L'ambiguité dans le titre est la clé de la méthode: du moment que je sais que je peux me tromper je peux estimer ma confiance dans le sens que je clame malgré moi [car forcé de tracer], et dans la possibilité du contraire quand je vois [ce que j'imagine voir]. Je me sais, deplus, incapable de distinguer ce percept intentionnalisé de ce que je fais sans sa projection [accompagné de reproductions imaginaires] dans  mon estimation de la séquence de manières dont cela pourrait arriver.


 
Nous allons 
ancrer des postures co-nommables  incluant une part d'interprétation inspécifiable [entre  deux choix le losange ou .  et le rectangle ou | \footnote{qui seront peut-être enrichis plus tard par  O ou le Oxygene-Hydrogène, out, oui ou impair strictement positif et le N pour nitrogène in , ou pair strictement positif?}].

"Ceci est un point est vu comme : Le carré est un losange" représenté par  les milieux des arètes du presque carré imaginaire liant les 4 bouts d'illusions diagonales notées X ou les sommets du presque losange liant les bouts de + qui marquent tous deux un locus localisable car protégé par ses bords imaginaires représenté par ce qui donne l'impression du point précis trace de la surimpression due  au recouvrement des traits,  coffre à trésor ou on voudrait mettre la lettre et ce qu'elle contient.
 
$$\begin{matrix}
& p \text{ comme Perséphone}&\\
d \text{ comme Dolores}&&b\text{ comme Bénédicte}\\
&q\text{ comme Qudsia}&
\end{matrix}$$

  "Ceci est un trait est vu comme : le carré est un rectangle " représenté par les sommet du presque carré imaginaire liant les 4 bouts de X ou les mileux des cotés du presque losange liant les 4 bouts de  +
 $$\begin{matrix}
& d\text{ comme Dolores}&&p\text{ comme Perséphone}\\
&&\\
&q\text{ comme Qudsia}&&b\text{ comme Bénédicte}\\
\end{matrix}$$ 


sont deux représentations indistinguables d'un même monde que l'un voit par des équations [complétables en structure régulières qui vont ressembler a des parfois/par défaut  à des successions de groupes simples de répétitions 
%différentiellement alternées produites par ppcm [je vois ici l'accélération, l'inverse de la température, la verticale]] 
et l'autre des distinctions.
% [qui peuvent etre vue comme des symmetries simples amplifiables par deformation et copie, systemes evolutif de differences produites par pgcd[je vois ici le temps , a distance l'horizontale]].


Nous allons raffiner  en :
Voir Dolores comme un sommet en haut à gauche et se demander si. 

\begin{itemize}
\item  le carré est un rectangle plat ou -
 $\begin{matrix}
& d&&p\\
&&.&b_p\\
&q&b_q\\
\end{matrix}$ 
\item ou  le carré est un rectangle plat droit |
 $\begin{matrix}
& d&&p\\
&&.&b_p\\
&&\\
&q&b_q&\\
\end{matrix}$
\end{itemize}

Pour plus tard, n'oubliez pas [ou essayez] que dans un rectangle allongé  pointer  sur le long coté assez loin du centre marque un carré [quitte a faire glisser les noms].  

Pour le losange, quand on voit Dolorès comme une arète de bord, on ne sait pas vraiment si 

\begin{itemize}
\item  l'ancètre de de Dolores est juste dessous $|$
$$\begin{matrix}
& p &\\
d &.&b\\
&q&\\
b_q&\\
\end{matrix}$$
\item ou il était la un peu avant comme $-$
$$\begin{matrix}
&& p &\\
&d &.&b\\
&&q&\\
b_p&&\\
\end{matrix}$$
\end{itemize}

et j'espère que vous avez assez zoomé pour voir que cette spécification à fait naitre une graine de possible ce petit point que ls les nouvelles barrières de lettres protègent. Vous avez assisté à une naissance ou comment du raffinement émerge l'espace nouveau.



On entrevoit déja   la diversité  intrinsèque aux communautés en se disant " si on cherche, un carré offre toujours des choix non triviaux."\footnote{Ce choix important de la formulation  toujours conditionnel par une intention anthropomorphisable va nous placer dans des famille de langages qui essayent d'être [je l'espère encore pour l'instant]  ise dirigenat vers $\Pi_2$ comme les jeux de pseudotélépathie quantiques c'est à dire ne peuvent être contenu dans du décidable ou du non décidable. Nous manipulemeront ceopendant essentiellement les fragments avec au moins une dimension finie non triviale au niveau de l'ancrage central indéfinissable. }  

A noter que si on voit le chemin dans l'amour de composer ensemble [vertical] comme l'espace des symétries ou des théories équationnelles, et de le raffinement par méta-projections comme l'espace des différences ou des antisymétrie le point ne nait ni dans l'un ou dans l'autre et il va falloir marcher pour lui donner un nom comme une équation caractéristique vue comme une adresse constante dont il est la régularité naturelle et une qualité ou une valeur différentiante comme cette antisymétrie qu'il génère entre le plus petit avant lui et le candidat a être nouveau plus petit par l'acte de sa niassance.

L'abstraction ici [remplacement de Dolorès par d est un contraction fidèle : si un dans une expression dont on ne voit pas le sens sera toujours dolores par défaut et n'hésitez pas à vous demander ce qu'elle fait là. Ainsi se cachait dans le dessin précédant un indice la bénédicte d'avant aurait pu venir par le chemin cohomologique des constantes $b_q$ ou peut être que dolorès est née d'elle en passant par persephone par le chemin homologique des variables $b_p$. Mais seules les Dolorès à grande mémoire s'en souviennent, les autres voient juste une superposition  comme un espace isotropique de chemins.
 
Dans cette quête de ses ancètres en fonction de jusqu'ou on pense aller on voit de chemin qui arrivent par le  bout et qui pour ceux dont la memoire est plus courte que ce trait selon indistinguables
$\begin{matrix} 
&-&-&-_d\\
|&&&\\
 |&&&\\
|&&&
\end{matrix}$
 et 
$\begin{matrix} 
&&&-_d\\
&&|&\\
 &-&&\\
|&&&
\end{matrix} $ Notez que si l'on préfixe on retirant une ligne les fins deviendraient distinguables. %de même en allant vers 


Un choix ou une intention ou une restriction peut se voir comme un marquage d'une, deux ou trois personnes et peut n'émerger qu'à partir d'un collage spécifique d'une description préalable sous réserve qu'elle est une description communautariste acceptant la diversité,
 qu'on va  définir localement avec une vue par cadran avec une direction de référence qui parait fixe et une direction alternante. 
Par exemple, [ce qui sera le choix par défaut ici]  si on dit qu'on est dolores et qu'on ne veut pas douter de l'origine de l'espace [référence  losange]: nous pouvons avoir comme choix  descendre  des losanges et avancer  $p$ comme losange ou rectangle.   

Ce qu'on retire du cadre de cette méthode pour cause de perfection que je ne saurais pas gérer car l'absolutisme est incompatible avec l'imperfection inhérente au geste c'est l'absence de pièce  
$\begin{matrix} 
&\ldots&\\
\vdots&&\vdots\\
 &\ldots&\\
\end{matrix}$ et les avant premières sans spectateurs 
$\begin{matrix} 
&-&\\
|&&|\\
 &-&\\
\end{matrix}$.

Une chose à se rappeler est que la vision la plus proche qu'on peut avoir de soi c'est entre deux formes qu'on sent ne pas être mais qu'on peine à séparer avec des mots ce qui fait qu'on ne cesse de rebondir, dérivant entre deux bords an morceaux élastiques et dynamiques qui peuvent donner l'illusion du continu et que des traces de ces rebonds on peut imaginer émerger la cognition. Pour l'humain, on peut imaginer ces bords  comme le coeur qui bat et le cerveau qui impulse.


\subsection{ Qui est le message et qu'est le messager?}

Si je vous parle c'est que j' imagine un lecteur et que je m'imagine l'image qu'à le lecteur de moi pour adapter les mots comme dans une pièce de théatre ou mon moi construit parlerait au spectateur que j 'imagine avec au moins cet implicite assomption que le spectateur n'est pas aveugle ou qu'il ne soit pas sourd.
On retrouve l'impuissance du metteur en scène d'il y a deux cents ans face à un public sourd et aveugle et encore cette acceptation d'imperfection qui s'impose à nous jusqu'à ce que les technologies permettent d'accéder à d'autre sen ou qu'en sais enfin un je ne sais pas qu'il  ne peut traiter avec son imaginaire de l'époque.

La compréhension transmissible suppose un sujet qui parle [que j'appellerais Dolorès] et un sujet qui écoute que j'appellerais [Perséphone] dans un échange fini dans l'expression.

Voir ce dont on  est  \textbf { en train} de parler sommet  c'est un peu comme un objet catégorique tentant de se construire une identité ou de se rajouter des qualités. Il faut au moins un bout de flèche sortant et un bout de flèche sortant dont je ne puis être confiant de la connection dans une illusion de figer le temps
que par des processus de vérification qui me rassurent ou augment mon doute cet Acknowledgement que je n'arrête de relire et dont je discute par bribes par ci par là en me disant que la constance des réponses des autres me donne une rassurance pour quantifier sa stabilité pour moi.

Le voir comme une arête ou un foncteur  c'est prendre un sens pour le lire par petit bouts si besoin avec un coté de sommet objectififié  d'un coté et un
coté d'un sommet obejctifié de l'autre.   

On voit la la justification de l'hubris du fini  qui voudrait  décrire ce qui bouge. mais dont la compositionnalité raisonnable requiert la remise en cause et aussi la justification du 4 comme manières de rentrer dans l'ouverture du sommet pour s'y reposer un moment donnant l'impression à qui regarde d'une sorte de stabilité justifiant la notion de locus comme cadre de localisation.

L'unifcation des nombres comme implicitable sans confusion sur une réalité induisant un identité implicite me semble cacher un principe fondamental d'unité implicte rigide à jamais et consultable par tous pour se synchroniser. 
Relacher la perfection permet de donner sens à trouver la joie dans les échanges de ceux qui essayer de s'en approcher avec des prérequis moins exigeants tenant compte des attentes de chacun et respectant ses règles plutôt que d'imposer la normalisation par l'autoritarisme de l'unification simpliste avec un 0 commun et des ressources infinies préexistantes d'aléa parfait qui est pourtant un cadre très intéressant mais n'en est pas moins un cadre restricteur dont le penseur en quête de liberté peut rêver le dehors avec des mots échangeables.

Cette liberté n'est pas toujours un choix elle est des fois forcée dans des communautés pour maintenir la cohérence des propos comme en informatique quantique ou la paire [et un peu plus] ]précède l'unité. \footnote{Une mesure a besoin d'un espace et d'une base}.

\subsection{Des limites}
Maintenant qu'on a parlé de ce qu'on croit pouvoir exprimer de manière assez partagée pour la discussion aie la possibilité   de donner l'impression de la potentialité de faire avancer  [corriger, raffiner, étendre, lier] les constructions conceptuelle

Le changement d'échelle nous force à   imaginer la transmission intergénérationnelle:  ce que j'imagine que Perséphone transmettra à ses enfants 
et eux mêmes à leurs enfants, cet être limite que j'imagine au bout de la chaine de transmissibilité sans jamais pouvoir l'atteindre parceque le repoussant chaque fois que je m'avance vers lui, cet être limite des chaines de décomposition sera appelé Bénédicte l'inatteignable je lui associerai l'espace parfaitement symétrique ou parfaitement isotropique. 
   
Similairement,  mais différemment ma conception de moi-même dans ce que ressens être cette interaction me lance dans une introspection mémorielle 
jusqu'aux origines cette quête de reconstruction de l'Histoire aussi fidèlement que ma mémoire me le permet avant la transformation pour mettre le costume du Dolorès je l'appelerais Qudsia Je lui associerais la constance ou ce point imaginaire dont beaucoup [et ils en ont le tout à fait le droit] postulent l'existence.

\subsection{Des mythes originels pour ancrer un imaginaire communautariste}
Les mythes [que j'imagine] disent que c'est le sacrifice de Bénédicte qui a fabriqué ce point originel qui lierait le temps à l'espace et qui pour moi
sera ces traits concrets d'énergie observable qui me permettent de pousser plus loin les limites de l'infiniment grand vers Bénédicte et de l'infiniment petit vers Qudsia. 

Ce mythe séparateur c'est celui de la naissance de la paire indissociable : ces deux axes qui me permettent d'imaginer [des variables raffinables d'espace que j'appelerais locus ou voisinages qui supportent sous réserve de multiplicité d'êtres traduites en chaines de restrictions pairées indissociablement avec des constantes extensibles de temps qui supportent l'union sous réserve de  non trivialisation].     

Une interaction transformante met donc en scène deux personnages et deux représentations des fois on a Dolorès qui repose sur le poids de l'histoire peut utiliser les constantes [de moins en moins constantes plus elles sont vieilles] héritées de Qudsia pour indexer les variables [de moins en moins variables plus elles sont proches] héritées de Bénédicte.

\subsection{ Repésentation de l'impossibilité de savoir}

La pensée bien éduquée se heurte rapidement à la pression castratrice de la perfection des nombres [ou des définitions dans certains cadres] qui ont l'audace de se présenter comme s'ils étaient pareils et inéluctables. Je ne voulais pour autant pas me résoudre à des extrêmes  je voulais juste les libérer des lois des algèbres les laissant casser ce plâtre comme le quatrième mur de la scène théatrale. Je n'ai toujours pas compris pourquoi au commencement était la commutativité ou l'associativité  ou tel ou tel cadre qui serait unique cadre pour ceux avec qui je peux jouer.

Un nombre libéré à le droit de choisir sa base et il a même le droit d'en changer au fur et à mesure que je le lis puisque sa lecture n'est pas instantanée. Un nombre équipé d'un processus de lecture qui pourra jouer le rôle de  vérifieur ou de prouveur dans les interactions. 

A l'origine  soit dans la fin soit au début Dolorès  doit trouver une trace pour pouvoir vérifier ou prouver en sacrifiant l'intimité de son locus pour en montrer les qualités au spectateur.

Cette trace dans plus simple appareil est le trait avec les deux gestes qui pourraient enseigner à tenter de le reproduire. Je vais donc revenir au trait et  dans un effort anthropologique penser à son apprentissage imaginant un geste assumant l'imperfection de la prédiction de par le respect pour le droit des autres à ne pas respecter mes règles.

Comme il y a le trait que je vois et celui que j'imagine, je vais toujours écrire sur une ligne horizontale ou verticale  ce qui me  permet de distinguer le trait monstrateur  de qualités qui traverse ma ligne et dont j'entr'aperçois ou j'imagine la trace  et le trait mesureur de quantités que je suis sur mon chemin jusqu'à ce qu'un bord me superpose dans des choix que je ne veux pas prendre tant que je ne m'y sens pas contraint.\footnote{Une des question qui date le long de ce raisonnement est celle de la possibilité de simplifier les écritures mathématiques pour avoir un cohomotopeur et un coégalisateur naturel  et des systèmes biunivoques permettant d'unir comme un cet assemblage immense.} \footnote{Je me rends compte que le toujours remplit cette section, pensez bien à un ou presque qui lui est inéluctable donc que je n'ai pas explicitement dit maintenant voici l'accolade qui délimite le cadre du sens de cette phrase.} 


Si le paragprahe précédent était un choix volontaire voici son pendant insoumission nécessaire: plus j'y pense plus je me vois me résoudre à [en?] cette idée :  l'existence sans contexte me dépasse des 4 cotés : celui du sens ou idées ressenties, celui du rien  ou de l'avant,  celui du tout ou de l'après, et même ce quatrième coté qui me sert à faire naitre un semblant de continuité et que je ne peux m'empêcher de percer en levant le bras pour former une pensée marquant ainsi l'interruption du suivant. Est-ce juste de l'auto-confirmation ? Peut-être, mais peut-être pas et on va tenter ici l'expérience de cette folle idée de faire naitre des formes qui donnent l'illusion de prendre vie, et de voir en miroir  cette perte de contrôle dans l'impossibilité de rigidifier l'origine sans apauvrir à outrance a force de soumissions et chercher la limite de cette séparation.   

\section{Parler avec deux traits : Et si notre cerveau était une fabrique d'illusions presque parfaites }

\subsection{Diagonale ? vous avez dit diagonale?}
Comme si le cerveau était un correcteur presque parfait pour nous convaincre de voir  ce qu'il construit par collages corrigés.
Chacun a en soit une image du tout qui se reproduit et s'altère mais distord [corrige] ce qu'on sent avec ses deux mécanismes pour faire naitre un semblant de stabilité. 

\subsection{Des formes conventionnelles [canonisées] des objets, presque parfaits du savoir représentable, version décompositionnelle}
\footnote{la canonisation est un peu comme le rôles une de ces conventions générant un implicite pour ceux qui les connaissent et obfusquant le sens pour ceux qui les apprennent sans se demander d'ou elles viennent. Je dirais juste ici qu'on a pris le losange en premier parce qu'on a vu naitre l'espace d'abord et par ce qu'il me semble naturel de penser qu'on soit tous similaires avant de commencer à vouloir séparer le rien du tout}
\footnote{Dans une version précédente la version à 3 trait était la complémentaire de celle a un trait  mais j'ai prefere revenir sur le mouvement de balance préférant pour l'instant ce semblant de continuité de mouvement  o l'identification dans le nom de la boite et du contenant}

\begin{itemize}
\item Dolores cherchant ses repères en montrant plus ou moins d'ou elle vient: 
\subitem $d_\downarrow$ ou $b_{q_|}\le b_{q_{||}}$: 
$.,
\begin{matrix}
 && \\
 |&&\\
&&
\end{matrix},
\begin{matrix}
 &-& \\
 |&&\\
&&
\end{matrix},
\begin{matrix}
 &-& \\
 |&&\\
&-&
\end{matrix}.$ ou $\downarrow, \uparrow \rightarrow,\leftarrow ._\downarrow  \rightarrow$
\item Perséphone s'en allant après avoir écouté:   
\subitem  $p_\leftarrow$ou  $b_{q_{||}}\le b_{q_{|}}$: 
$.,
\begin{matrix}
 &-& \\
 &&\\
&&
\end{matrix},
\begin{matrix}
 &-& \\
 &&|\\
&&
\end{matrix},
\begin{matrix}
 &-& \\
 |&&|\\
&&
\end{matrix}.$  ou $\leftarrow, \rightarrow \downarrow, \uparrow .^\rightarrow \downarrow$
\item Qudsia ramenant les croyances: 
\subitem $q_\rightarrow$ ou  $b_{q_{|}}< b_{q_{||}}$:
 $.,
\begin{matrix}
 && \\
 &&\\
&-&
\end{matrix},
\begin{matrix}
 && \\
 |&&\\
&-&
\end{matrix},
\begin{matrix}
 && \\
 |&&|\\
&-&
\end{matrix}$.  ou $\rightarrow,  \leftarrow \uparrow ,\downarrow ._\rightarrow \uparrow$
\item Bénédicte fabriquant les variables 
\subitem $b_\uparrow$ ou  $b_{q_{||}} < b_{q_{|}}$: 
$.,
\begin{matrix}
 && \\
 &&|\\
&&
\end{matrix},
\begin{matrix}
 && \\
 &&|\\
&-&
\end{matrix},
\begin{matrix}
 &-& \\
 &&|\\
&-&
\end{matrix}$. ou $\uparrow, \downarrow \leftarrow,.\rightarrow .^\downarrow \leftarrow$

\end{itemize}

On a associé à chaque personnage un questionnement par défaut sur la taille des variables les grands infinis pouvant se permettre le luxe d'être stricts, et ceux qui se dépatouillent avec des dimensions qu'il ne peuvent augmenter qu'au fur et à mesure qu'il les comptent de se voir obligé de ranger des mots dans les mêmes tiroirs. 

\subsection{Exemple ondulatoire : traits qui se retrouvent et se repoussent}
Si on voit les fléches comme sens de déplacementon peut imaginer  les cotés , bouts décollables,  se heurter sur la diagonale et revenir.
$
\begin{matrix}
._d&\leftarrow\uparrow &|\cdots |\\
\uparrow\leftarrow \\
\begin{matrix}
-\\ \vdots \\ -
\end{matrix}
\end{matrix}
$


\subsection{Digression autour de la fracture centrale : de la politique ou de la répartition des richesses }
Les gens de droite qui semblent avoir pris les rênes du pouvoir presque partout pensent qu'ils faut donner au plus riches il veulent insérer à droite voient $b_p$ et se sentent pas très loin de l'infini regardant les autres du haut de leur tour cela génère l'aléa le plus complet dans mes tentatives, ridicules mais conscientes de l'être , de prédiction de l'avenir du monde. L'exemple illustratif est ces droites américaines ou israéliennes qui enivrées par leur puissance ont oublié ce qui est écrit dans leurs livres et qu'en regardant bien on verrait qu'un indicateur monstrable d'humanité est écrit au delà des bornes de leurs tribu  fondateur d'un équivalence qu'ils se complaisent à oublier par convenance.

Les gens de centre droit pensent qu'il faut donner aux moins riches mais parmi les riches, parceque quand même ils l'ont mérité, mais bon trouver le moins riche parmi les plus riches n'est pas un protocole aisé.
D'autant que les plus riches à commencer par le deuxième ont peur de se retrouver plus proche de la faille de la moitié séparatrice. La somme alternée de ceux qui sont des uns qui se voudraient puissances [et qui le deviennent de fait en imposant leur imaginaire aux autres] est  ce qui gouverne la force  qui poussée par des idéaux de droite agit à gauche [mais pas trop] à réduire la fracture sociale.
 %des plus roches et peut pousser  à donner de l'autre coté pour augmenter la chance de maintenir leur  statut.

Les humanistes pensent qu'il suffit d'être pour mériter d'avoir et veulent simplement donner aux plus pauvres mais malheureusement, je crois que  c'est  une espèce en voie de disparition on ne prend même plus la peine d'arguer le manque de ressources tellement ces clameurs sont loin de nos oreilles. 

Mais revenons au centre, mais de l'autre cotés vers ceux fractionnés qui se partagent les restes. De manière similaire mais différente, 
%On peut voir la gauche perdue par le fait que 
personne n'aime à se voir comme le plus riche parmi les pauvres parce que l'action de le devenir fait cesser le droit de représenter ceux dont on est fier d'être le premier. Le un du bas peut même s'offrir le luxe d'ignorer despotiquement le processus de la somme  globale de moitiés des moitiés [Passez  ou  GOTO 4 mots en arrière de la drôle de parenthèse tant que vous en avez la patience] et tout seul de manière presque paradoxale creuser le fossé social dont il se plaint, en faisant passer la barre pour maintenir son statut du un guidant  des puissances d'un f[r]action. 

On comprend donc dans ce jeu une possible confusion apparente vers le centre ou les idéaux des uns peuvent faire appliquer la politique des autres avec cette construction de désirs moteurs d'impulsion ou chacun ne pense qu'en couches communautaristes centrées autour de lui même.

Mais  à imaginer cette classe moyenne [décalée sur soi ou ailleurs] en  excluant les bords  on risque de ne pas prendre conscience que l'exclusion communautariste par le bas c'est dit Dolorès  et c'est ma conviction la plus pronfonde l'antipode du bon sens que je ressens ou qui s'impose selon moi à celui qui prend le temps d'écouter son coeur.  Je ne peux cependant m'empêcher de voir   que fermer les yeux sur les souffrances des autres pour sa convenance même si je ne saurais lui donner  du sens ne cesse de me noyer   par sa dominance et seule une lutte incessante [contre soi -même] peut laisser entrevoir une lueur d'espoir.

\subsection{La normalisation :  Des liens comme formes presque parfaites ces canaux sans erreurs mais avec un curseur sur l'imprévu qui pourrait devenir mais qui n'est toujours pas}\footnote{Je n'aime pas ce mot car il pourrait faire penser que la norme existerait avant qu'on ne la définisse mais la normalisation,  les blecs diagonaux, et les formes normales ont été martelées tellement fort en moi qu'il serait malhonnête de le nommer autrement. Je m'y soumettrai donc avec juste une note de bas de page pour manifester une protestation.}
$(n=1,i=1,k=n)$ j'ai lu $n$ trait sur mon premier coté je dois tourner ma langue dans ma bouche $i$ fois . Le = apparait la comme valeur par défaut à un nom bien localisé indistinguable de tout autre dans la section avec un lien qui les lie [en dupliquant d'autres noms].

un ou deux intervalles [matrices circulante après le passage d une gomme sur une bande c'est l extension]
comme représentant [une fois décodés) les trois nombres de la quantité \emph{doit a-on attendre la naissance des nombres pour en parler?}

\subsection{Représentation interne d'un coté de soi-même}
Mais avant de représenter quoi que ce soit comment ne pas aborder les représentations de soi-même vue comme une mémoire locale?

Sur un coté ca fait miroir :
${\scriptscriptstyle
\begin{array}{cc} 
d \cdots \\
|& \\
|& \\
|& \\
|& \\
\vdots &\\
q&\\
\vdots&
\end{array} 
}$
 devient ${\scriptscriptstyle
\begin{array}{ccc} 
d \cdots \\
|& \\
|& -_1\\
=&|&|\\
|& -_2&|\\
|& -_3&-_4\\
\vdots &\\
q&\\
\vdots&
\end{array} 
}$ 

Si on voit les traits horizontaux de la deuxième colonne comme un doigt qui montre et qu'on lit le long du coté attribué [proche] à d
on verra le 1 puis le trait 6 si on continue.


A regarder sur un losange ca fait rotation ou shift rescale


\footnote{to clean  la dérive sera un exercice de simple répétition/réflexion sur/entre les 4 cotés de soi-même avec copie locale en miroir/transmisison fidèle. 
to do $n|k$ coverings avec  une fonction finie de POP repetable sur [ (k) copies dun graphe fini]  pour couvrir  n copies d'un grpahe infini decouvrable par POP de lui avec  un taux.\emph{ je ne sais pas genre ? $(s_b x n^{s_q 1/y}$ }}


Quand dolores croit voir une trace elle ne sait pas si c'est une trace du passé ou un début de futur mais ce qu'elle sait c'est qu'elle induit une distinction entre deux directions différentes sur son chemin dont une seule peut rebondir et revenir en un temps fini pour devenir similaire à l'autre si je suis sur une régularité [hors de l'hyperplan d'équivalences que j'imagine] et qui ne se rejoindront jamais si je suis sur une singularité [j'étais sur le point critique de l'équivalence] 

Définir c'est déchirer l'espace temps, qui nous permet d'échanger dans nos  communautés  pour localiser, pour faire place à ce choix qu'on impose aux autres, c'est cette faille qui sépare le trait au milieu du coté du losange mais qui  est sur le sommet du rectangle. 

C'est introduire comme des des murs induisant le rebond la ou on veut pointer le locus ou a l'impression de voir.

Un rebond est un changement de direction qui se note.  

Dès qu'un élément nait candidat à la compréhension, son image le suit de l'autre coté du mur liée par l'unisson qui  aimerait se  voir un mais qui construit son chemin liant deux murs allonges par un chemin tout droit ou bien deux murs tout droits par un chemin allongé.

\section{Interlude : assemblages}



 
 
 \subsection{ six citations :}
 L'absence de citation ce n'est pas faire comme si  cette pensée était mienne est pas aux autres, comme si, comme c'est de nos jours courant, il était parfaitement naturel de posséder des idées de les attribuer aux gens et de soumettre ceux qui s'y refusent à l'autorité des amendes. 
 Je ne suis que les pensées qui n'ont pas fait que me traverser mais qui ont raisonné assez en moi pour que je me reconnaisse en elles et m'enchevêtre à elles en les redisant comme je m'en souviens ou plus précisément comme je les ressens..
 Je suis évidemment tout ce que j'ai lu tout ce que j'ai vu et tout ce que j'imagine pouvoir être par cet acte de le penser liquéfié dans le verbe.
 Dolorès ne sait pas si elle a écrit avec ses larmes ou avec son sang mais elle écrit parce qu'elle ressent devoir le faire ou est ce vouloir elle ne le sait déja plus très bien. Il me semble pourtant que la protection des idées contre les détournements n'est ni possible, ni légitime, ni souhaitable.
  
  \begin{otherlanguage}{arabic}
 حين اليقين عليك ينادي هل العين ترى أم القلب يشادي
\end{otherlanguage} 
 
\emph{quand on sait en son for intérieur est ce la raison ou bien le coeur.} 

La première réponse du soi au moi lors de la rédaction de cet essai  était une citation :"Cachez moi ce sein que je ne saurais voir"
Je me permet donc cet assemblage  qui représente allégoriquement ce travail.

\begin{CJK}{UTF8}{min} 
%\textbf{正} [せい / しょう) 
猿
も
木
か
ら
落
ち
る
\end{CJK}
 Saru mo ki kara ochiru.\emph{le singe aussi  peut tomber de l'arbre}
  
                                             \emph{innama ennasou soutouron kutibet laken bime} Les gens cont des lignes écrites, mais avec de l'eau
 \begin{center}

                                                             \begin{otherlanguage}{arabic}
إنّما النّاس سطور كتبت لكن بماء 
                                    \end{otherlanguage}             
\end{center}
                        
												\emph{translated access:}  Par phénomène j’entends tous les sens de phénomène  mais avec une attitude différente  					                        
                        														
                             \emph{ ma assaba an tahwa emraatan laisa laha inwan} Qu'il est difficle d'aimer une femme, ô mon fils, qui n'a pas d'adresse
                              
                      \begin{center}

                                                             \begin{otherlanguage}{arabic}

          ما أصعب أن تهوى إمرأة يا ولدي ليس لها عنوان                                 
               
                                    \end{otherlanguage}             

\end{center}
                        Fossoyeur il est beau de contempler les ruines des cités; mais, il est  plus beau de contempler les ruines des humains
                                          
                                                                                                 \emph{Per causa sui intelligo id cujus essentia involvit existensiam}
                                                                                               Par cause de soi j'entends ce dont l'essence induit/implique/contient/involve l'existence   
               
    
 \subsection{ Les éléments}
 \[
 \begin{matrix}
 \text{Dolores écrit avec}   &  \vdots&     \text{ Persephone lui montre sa mort comme } \\
 \text{avec ses larmes ou son  sang} &\vdots&\text{  la fin des temps ancrée solidement}\\
 &&\text{dans un nuage gazeux} \\
 \cdots & &\cdots \\
  \text{ Qudsia a laissé les traces dans} &\vdots & \text{ Bénédicte pourrait savoir }\\
  \text{ la roche de l'origine de l'espace} &\vdots&     \\
  \text{ Qu'on voit dans le brouillard} &\vdots&     \\
    \end{matrix} \]
 
 
 
 La folie est de dépasser Perséphone car la on n'a plus aucune équivalence noyé dans les percepts tout semble différents on n'arrive pas à oublier que le locus ne cesse de changer et que demain est d'abord différent d'aujourd'hui. C'est l'espace isotropique ou toutes les positions de départ sont pareilles et toutes les directions interchangeables. C'est pour  ça qu'elle est la source de l'incertitude parfaite de ces variables dont on sait que déja  avec notre manière de voir ells ont eu au moins deux locus dont une trace témoigne la distingabilité. Le fou est celui dont le raffinement me dépasse.
 
 Bénédicte fait le lien entre  l'infini donné par Perséphone hier qui donne par son inverse un nouveau candidat pour l'unité de demain. Elle supporte les deux topologies résiste à l'intersection infinie d'un coté et à l'union infinie de l'autre, derrière elle est ce que  celui qui a conscience de sa finitude ne peut espérer décrire. 
 
 Qudsia montre le trait qui est la plus petite unité qui est capable de faire sens de tous le reste des encoches  par les multiples et le rescaling c'est la théorie équationnelle qui se construit sur les répétition du commun multiple. Nous mortels qui avons besoin de dimension finie pour agir on ne pourra rien voir que ce qu'on exprime par multiple [du moins dans cette dérive, la prochaine évolution permettra peut être d'aller plus vite et de voir émerger la vie la ou ne voyait rien. Le Nihilisme se cache  avant Qudisa car si tout était pareil on ne pourrait rien distinguer c'est pour ça qu'ellle est l'origine des constantes, ces principes qu'on veut mettre pour batir le sens constructible ou cohomologique. Le Nihiliste est celui dont le refus de séparation m'atterre.
 
 Dans les vagues des émotions ballotés par les incertitudes avec les larmes dans les yeux quand elle regarde l'Hisitoire et ce sentiment d'inconsistance quand elle essaye de voir au de là sa mort Dolores s'adresse par défaut à Bénédicte la meilleure représentation qu'elle s'était fait avant de commencer à parler de celle qui pourrait savoir, elle espère que Perséphone construira le savoir pour entrevoir le chemin de la fin car malgré tout on ne peut pas être vraiment vraiment vraiment vraiment sur de l'impossibilité de finir. 

\section{Sous-texte : de la personnalité des personnages pour l'interopérabilité}

\begin{itemize}
\item  Dolores  et Qudsia sont optimistes et disent oui par défaut
\item Perséphone et Bénédicte elles sont pessimistes et disent non par défaut
\item Dolores et Perséphone ont un timeout connu leur forfait de parole coupe après un temps connus de communication ou de silence: elle ont un budget d'occupation de la ligne conversationnelle
\item Qudsia et Bénédicte ont le luxe de répéter tant qu'ells ont l'impression qu'elle ne sont pas comprise et sont du coup souvent interrompue par celui qui échange avec elles.
\end{itemize}

Quand elles ont deux arugments si elles étaient une proposition: elles seraient 
\begin{itemize}
\item Dolores : $._| \le ._{||} $
\item Qudsia : $._| < ._{||} $
\item Perséphone : $._| \ge  ._{||} $
\item Bénédicte : $._| > ._{||} $
\end{itemize}

 Ici l'indice est une première position certifiée qui précède une seconde: Si j'identifie deux locus et que j'arrive à étiqueter le premier que j'ai imaginer avec un "|" et le deuxième comme "||" .\footnote{Même si on interprète || comme une fracture,  celui qui la conçoit peut la concevoir dans un environnement ou des compositions l'ont précédée.}
\section{Nombres | : Que pourrait-on dire quand on apprend à  compter mais qu'on se trompe encore parfois: L'espoir, La rassurance et l'incertitude}

Pour comprendre un nombre il faut imaginer un protocole ou je vais chercher un nombre, puis un autre en fonction de ce que j'ai décodé/échangé et de ma mémoire et ainsi de suite jusqu'à arriver à des bords sur les les quels je dois compter des encoche unaire. En refaisant le chemin inverse je vais produire ma conviction dans le oui le non ou le je ne sais pas.
Si sur le bord la base est unaire [en encoches] sur les couche supérieures on peut ganer des dimensions : un role est un locus partiellement spécifié et on 
voudrait quelque part avoir des termes $(\tilde b^{\lambda \tilde m}_{\tilde v})$ as $(\forall \lambda \ge 1)$ comme des rayons qui devraient être superposés selon leur décomposition.

\subsection{ Imaginons deux  erreurs possibles pour trois indicateurs de comptage}
  
  \section{[Section clé] Propriétés des points de vues entre message et messager}
Si un anthropomorphisme peut prendre 4 rôles, il porte [ou  je peux du moins toujours imaginer que celui qui le reçoit lui impute] toujours un message: ou plus précisément l'information qui le décrit lui prête toujours une intention  par exemple d'ou il vient et ce qu'il veut.


\begin{itemize}
\item la multiplicité de l'incertitude  est une partition de trace  O de causes: 
\subitem $O_b$:  Obfuscation de valeur  comme confusion  morale
\subitem $O_q$:  Obfuscation d'adresse comme moteur de l'oppression par contrôle de l'information sortante faire croire que ce qu'on n'arrive pas à voir ne peut pas exister
\item $n_h$ : L'espoir que mon comptage en partant des constantes soit bon 
\item $n_v$ : La rassurance dans le changement suffisant des variables
\end{itemize}


%println(makematrix([[2:5],[1:3],[],[3:5],[2]],[[2],[1:2,4,6],[1,3,5,6],[3]]))
${\scriptscriptstyle\begin{matrix}
  & - & - & - & - &O_b   \\
  & | &   &   &   &   \\
- & - & - &   &   &   \\
| n_h&  &   & n_v| &   & n_{FAIL}| \\
  &   &   &   &   &   \\
&   &  &|   & & | \\
  &   &  & - & - &   \\
  &   &  &   |&   &   \\
  O_q& - &   &   &   &   \\
  &   &   &   &   &   \\
  &   &   &   &   &   \\
\end{matrix}}$



%\sction{De la forme parfaite}

%suel la femme parfaite  qui ne peut etre localisée saurait recolles les bouts exterieurs et comme le fond et la forme je ne sais séparer 
%on e va utiliser que celle que vois mais que je pourrais compléter. 

\section{En essayant de le dire  avec des traits}

\subsection{Commandes anachroniques pour ceux qui savent compter }
notation:
\begin{itemize}
\item  sequence des  séquence de - puis
\item  sequence des séquences de | 
\item   : pour sequence croissante de 1  
\item  coeff pour multiplicite
\item [] pour e ligne d'espaces
\end{itemize}

Attention.  ce code n'a de sens que comme deuxième lecture d'une copie en mémoire ayant été validée  après un signal recu :
 si on tracait au vol il faiudrait alterner une lecture de l'un puis une de l'autre pour voir synchroniser l'émergence d'un sens des deux matrices de traits par leur représentation.
 
%println(makematrix([[2:4],[2:4],[2:4]],[[1,5],[1,5]])) %2fois cote a cote
code  draw [2[1,3,5,7]] [[2,6]]:
La cellule espacée pour être manipulable  $
{\scriptscriptstyle\begin{matrix}
- &   & - &   & - &   & - \\
  & | &   &   &   &  | &   \\
- &   & - &   & - &   & - \\
\end{matrix}}$

Le Losange  vu par diagonales repoussées à l'extérieur [2draw [3[2:4]] [2[1,5]]]%][draw [3[2:4]] [2[1,5]]]
%
%${\scriptscriptstyle\begin{matrix}
%  & - & - & - &   \\
%| &   &  ._d &   & | \\
%  & - & - & - &   \\
%| &   &  ._q &   & | \\
%  & - & - & - &   \\
%  &   &   &   &   \\
%  &   &   &   &   \\
%\end{matrix}
%\begin{matrix}
%  & - & - & - &   \\
%| &   &  ._p &   & | \\
%  & - & - & - &   \\
%| &   &  ._b &   & | \\
%  & - & - & - &   \\
%  &   &   &   &   \\
%  &   &   &   &   \\
%\end{matrix}}$
%ou avec ses cotés distingués 

${\scriptscriptstyle\begin{matrix}
  & - & - & - &   \\
 \blue{|}&   &  ._d &   & | \\
  & - & - & - &   \\
| &   &  ._q &   & | \\
  &\blue{-}  &\blue{-}  &\blue{-}  &   \\
  &   &   &   &   \\
  &   &   &   &   \\
\end{matrix}
\begin{matrix}
  &\red{-}  &\red{-}  &\red{-} &   \\
| &   &  ._p &   & | \\
  & - & - & - &   \\
| &   &  ._b &   &\red {|}  \\
  & - & - & - &   \\
  &   &   &   &   \\
  &   &   &   &   \\
\end{matrix}}$

Le rectangle vu par coins repliés à l'intérieur.

${\scriptscriptstyle\begin{matrix}
  &  & | & ._p| &   \\
 & -  &   &   & - \\
|._d &    &  &  &  | \\
| &   &   &   & ._b| \\
  & - &  & - &   \\
  &   & |._q  &|   &   \\
  &   &   &   &   \\
\end{matrix}}$

La diffraction explicitée

${\scriptscriptstyle\begin{matrix}
  & - & - & - & - &   \\
| &   &   &   &   & | \\
  &   &   &   &   &   \\
| &   |&   &   & | & | \\
  &   &   &   &   & - \\
| &  | &   &   & | &  \\
- &   &   &   &   &   \\
  & | &   &   &  | &   \\
  &   &   &   &   &   \\
\end{matrix}}
$



Entre miroirs bien polis qui savent réfléchir c'est un  trait insoumis que l'on voit surgir. C'est cet interruption ou bien cet accident 
qui est le point d'ancrage de cette histoire que je ressens. Le sens aura pour cadre cet espace qui nait par les transformations de ce jaillissement.
Ce miroir polissson n'est pas dur a polir : tournez donc la tete vous saurez ce que je veux dire. 

 
A entendre cette histoire, c'est à se demander,  si  définir la vie ne se  pourrait pas  voir,  comme une manière.
C'est en  plongeant dans les mots que l'on tisse les liens et cet acte de tisser est le plus proche de ce qu'on est quand on se voit mouvement.

De la folie des autres: Certains voient la folie dans vouloir espérer  répéter la même chose et voir émerger, quelque chose de nouveau dans la dérivée.
Mais chez  les être imparfaits  pour qui  la mesure ne peut être qu'estimée, comment bien savoir si en des lieux, ou temps  séparés c'est bien la même chose qui s'est déroulée. Si le poids des actes créent sous moi un sillon, si l'espace se déforme à chaque pas que je prends, j'ai beau ne pas le voir, je sens bien arriver rien que par répéter, au moins un  renforcement, de  ma rassurance en le fait, que je ne me suis pas  trompé, encore... 


\footnote{A TROUVER  : dessin pour la dilatation contraction la barrière qui recule sur les horizontales je la double pour marquer le temps de la rencontre 
- -- | -  - -  |-  -| devient - -- | -  - -  |-  -|| - |  et - -- | -  - -  | marqueur ? }

 

Tentative de définition de 5 ou l'intervalle $[5]$ ou plutot sa moitié sans copie locale avec ces 5 préfixes comme nombres premier test de définitions locales dans le contexte de l'essai valides par défaut jusqu'a réinterprétation explicite. La redéfinition du gros 1 induit la redéfinition par paires en commencant par le bas, indice puis explosants.
cette définition ne peut pas avoir de sens si onn'attends pas 7 clics d'horloge donc elle nait au plus tot le jour local $(t_1 > 5)$ à justifier plus si nécessaire en fonction de la précision souhaitée. 

$ \left({ \scriptscriptstyle\begin{matrix}

  &  | &  &   &   &   \\
  1_{t_1}^{4^1}:& - &   &   &   &   \\
  &   & | &   &   &   \\
  1_{t_1}^{4^2}:& - & - &   &   &   \\
  &   &   & | &   &   \\
  1_{t_1}^{4^3}:& - & - & - &   &   \\
  &   &   &   & | &   \\
  1_{t_1}^{4^4}:& - & - & - & - &   \\
  &   &   &   &   & | \\
  1_{t_1}^{4^5}:& - & - & - & - & - \\
  &   &   &   &   &   \\
  &   &   &   &   &   \\
\end{matrix}}   \right)$

En payant le cout de pousser la barre à chaque temps d'horloge et en cherchant un comportement reproductible on va vouloir mettre des repoussoirs au bords qui profitent de son énergie pour faire rebondir le trait 
en imaginant un zone centrale repoussante $|_{._d}$ quand on arrive et $|_{._p}$ quand on repart et la zone de bord ou on part et celle ou on rebondit $|_{._q}$ et $|_{._b}$
On peut donc imaginer  5 comme un coté de rectangle assumant la représentation interne d'un coté de soi-même.


A VERIFIER :Les quatre roles  en tant que méthode  de lecture se réinterprèteraient donc avec un jour ultérieur à la définition interruptible  de 5 avec copie locale et || comme zone protégeant le milieu espére.% comme les préfixes de  |||||- pré-construit de gauche à droite en 7 temps.

 \begin{itemize}
 \item $d$   comme    $._5-._3-._1 ||._2-._4$ 
 \item $p$ comme     $._4-._2||._1-._3-._5$ 
 \item $q$ comme     $|._1-._2-._3-._4-._5$ 
 \item $b$ comme     $._5-._4-._3-._2-._1|$ 
 \end{itemize} 
 
 %A TESTER impulsion electriques du cerveau vs battement du coeur en terme d'attente inconnue avant la premiere pulse  puis comptage du nombre de positions juste avant ou juste après l'écoute du mouvement de reference.
  
\subsection{Demi-diffraction et demi-rebond en deux temps trois mouvements l'arrivée le contact le départ}  
  
  on peut imaginer
  \begin{itemize}
\item     la décomposition
  $\bmat$ comme un trait arrivant $-$ un trait rencontré $|$ et un trait séparé $\begin{matrix} | \\ | \end{matrix}$  
  \item le rebond $- |$  comme un trait arrivant $-$ un trait rencontré $|$ et un trait rebondissant  $- - |$
  \end{itemize}
  
 \subsection{L'interruption dure le temps d'une rencontre qui peut être parfois  éclatante}
 Il y a des rencontrent qui marquent non seulement les personnages, qui repartent en trois morceaux chacun, mais aussi le lieu de la rencontre.
 
  Imanginez $]||[$
 les deux traits centraux représentent le passé en tant que deux traits arrivant mais aussi le lieu estimé central du futur ou chaque trait reparti est une moitié d'accolade. 
 
 Les petits bouts [horizontaux]  prolongés [verticalement] ressemblent à s'y méprendre au deux premier traits  || juste en deux séparés, si on occulte ou oublie l'histoire [pourtant centrale] de la rencontre.
 
  ${\scriptscriptstyle
  \begin{matrix}
  
  |&&\\
  &-&\\
  &&|\\
  &&\\
  &&|\\
   &-&\\
    |&&\\
\end{matrix}
  \begin{matrix}
  &&|\\
  &-&\\
  |&& \\
  &&\\
  |&&\\
   &-&\\
    &&|\\
\end{matrix}}$
 
 
\subsection{Expression d'un envie de dénomination [locale]}
Ces histoires de rencontres  donnent envie d'équiper les espaces de nom autour  d'un miroir central, et d'imaginer un feuilleton avec des espaces d'expressions.
 $._{d_\le}]._{b_>}||._{q_<}[._{p_\ge}$
  
\section{|Version contractée compositionnelle des interprétations transmissibles souffrant l'interruption, vue d'après pointant vers l'origine les partants ou mâles}
Avec la douceur de la lèvre qu'on notera "| "et la rigueur du doigt qu'on notera "-" nous allons tenter de signifier ce qui peut prétendre à une compréhension transmissible pour un être discret.
Le mécanisme central est l'interruption de la bouche représentée par une lèvre et sa jumelle "||" avec cette élasticité de cet espace entre elles qui permettra d' y  insérer des symboles qui trainent avant de leur accorder une signification vue ici comme un percept exprimable ou un nom vu comme adresse ou procédure visant à caractérisation.

Les demi-corps explicités  seront ces symboles qui voudront expliquer ce qui est implicite dans le trait son origine ou son devenir ils vivent sur quatre  symboles "$d_m$ " ou "$\leftarrow$" ou 
${\scriptscriptstyle\begin{matrix}
|&   \\
  & - \\
 | &   \\
\end{matrix}}$, 
"$p_m$" ou "$\uparrow$ "ou ${\scriptscriptstyle\begin{matrix}
&  & \\
 - &&  -\\
  &|&  \\
\end{matrix}}$,
" $q_m$" ou "$\downarrow$" ou  ${\scriptscriptstyle\begin{matrix}
& |&  \\
  -&&-  \\
  &&   \\
\end{matrix}}$, et "$b_m$" ou "$\rightarrow$" ou $\bmat$.

%
%${\scriptscriptstyle\begin{matrix}
%& |  \\
% - &  \\
%  & |  \\
%\end{matrix}}$,

Les paires naturelles ne voulant pas voir leur fin la cacheront au signe arguant que par leur taille, les gens de leurs famille saurant les distinguer
 on dira donc aussi 
 \begin{itemize}
\item $m_d$ pour $d_mp_m$ ou  $\leftarrow \uparrow$ ou $\dmat \pmat$ ou $q_md_mp_m$ vu comme $\begin{matrix} d&p \\ q & \end{matrix}$ ou 
$\begin{array}{cc} \dmat & \pmat \\ \qmat& \end{array}$
\item $m_p$ pour $p_mb_m$ ou  $\uparrow\rightarrow $ ou $\pmat \bmat$ ou $d_mp_mb_m$ vu comme $\begin{matrix} d&p \\  & b\end{matrix}$ ou
 $\begin{array}{cc} \dmat & \pmat \\ &\bmat \end{array}$
\item $m_q$ pour $q_md_m$ ou  $\downarrow \leftarrow $ ou $\qmat \dmat$ ou $b_mq_md_m$ vu comme $\begin{matrix} d& \\ q & b \end{matrix}$ ou 
$\begin{array}{cc} \dmat &\\  \qmat & \bmat \end{array}$
\item $m_b$ pour $b_mp_m$ ou  $\rightarrow \downarrow$ ou $\bmat \pmat$ ou $q_mb_mp_m$ vu comme $\begin{matrix} &p \\ q &b \end{matrix}$ ou
 $\begin{array}{cc} &\pmat \\ \qmat&  \bmat \end{array}$
\end{itemize}
 
 Quand imagine  les quatre on aurait donc  comme un trois recouvrement de l'esprit du  carre  $\begin{matrix} d&p \\ q &b \end{matrix}$
  $\begin{array}{cc} \dmat &\pmat \\ \qmat&  \bmat \end{array}$
 
On verra   dp la découpe du multiple, pb la dilatation bq le recollement et qd la contraction comme quatre mouvements dans la sensation du un comme clic qui nous pousse à réessayer.
 
 On montrera du doigt en les superposant toutes ces belles  insoumissions suspectées, estimées ou envisageables qui enrichissent le possible qu'on rigidifie par le mot de tout un champs de possibles nouvelles ambiguités. Ce sont ces comment, ces alternatives  ou même si la trace est indistinguable d'ou on est  la manière de les faire aurait pu être  distingable si on pouvait non pas voir le dessin mais voir dessiner.
 C'est cet espace complet d'autres  manières de faire dont l'acte de tracer  force a collapser, et dont la vie imaginaire pour celui qui raconte  enrichit le possible de nouvelles ambiguites : $ \uparrow \leftarrow$ ou $ \pmat \dmat$ pourraient peut etre utiliser l'indexation pour encoder le pointage?
 
 Les deux traits de l'interruption || peuvent se lire comme un désir de faire tout ce qui est possible pour ne pas détruire  ce que j'ai l'impression de ne pas comprendre parfaitement, ici  cette durée cachée entre les deux séparés qui vu son ambiguïté est plus vue  comme espace que comme largeur de trait.
  
 Notes\footnote{ La où les espaces de probabilités écrasent le temps avec des équivalences irrévocables. La démarche ici sera  de vouloir ne pas tuer ce que je ne sais reproduire et de n'identifier les faits que par nécessité,.
 [PROPOSITION] Pour un élément canonique il suffit d'indexer. Ainsi $d_-$ pour $[dp]_-$   ou  $ \uparrow \leftarrow$ ou $ \pmat \dmat$ 

 [Attention] ne pas oublier qu'un trait peut exploser en d'autres.
 
[Rq] j'ai l'impression que pour la bouche je prend le comme trait pouvoir le mur gauche et le p comme flech pour garder le mur droit  comme si je prenais son image miroir par la droite de 45 degre qui est la diagonale qu'il est cense representer....
% $\begin{matrix} q&d \\ p & \end{matrix}$ ou 
%$\begin{array}{cc} \dmat & \pmat \\ \qmat& \end{array}$

 [TODO Proposition, ] dp a la pgcd mais version a partir de shift du multiple a partior de  $1\mapsto 2$, $ 1_2 \mapsto 3$, $ 1_3 \mapsto 4$ coupée a $k$ copies de 1 et répétée  
 pb : en dilatation pour simuler le produit par 1+ r ou r est le reste de la division euclidienne du trancon haut
 
 bq: recollement en shift vers le maps gauche a droite a vec les meme param ?
 
 qd : contraction, par le meme quotient 
 puis rajouter 1 le tout étant l'intéreur d'un ++?
 }
 
 \section{ De grandir}
 Un personnage se construit quand je lui associe plus de mots : le fait d'avoir une dérive au lieu de cycles permet quelque part 
 à une machine de se redéfinir au fur et à mesure  des expériences. 
 
 L'idée d'avoir juste le même format de séquence de traits espacés permet de revoir à une machine de Turing de revoir son automate et de ne plus avoir la même réponse en fonction de quand et ou on demande donnant une impression de "vie" parcequ'elle agit sur les règles même de sa dynamique.
 
 \section{ Du Contrôle}
On peut chercher  deux phases dans un mouvement, une phase libre et une phase controlée ou conditionnelle, celle ou le pré-requis conditionne le propos.
 On ne peut certifier l'infidélité que sous réserve qu'il n'y a pas erreur sur la personne.
   
\section{ [Section cle]Le problème central auquel je me ramène : VERIF est-ce que ce que je regarde est  une constante ou une variable ?}
Protocole unaire de vérification
$
\begin{matrix}
\text{local view}&w_Hww_F&&&&&&\\
\text{Age en clic}&|&|&|&|&|&|&|&|\\
\text{Change}&|&|&|&&&&|\\
\text {Meme}&|&&&|&&|&\\
\text{FAIL}&&&&&|&&&|\\
\end{matrix}
 $
 
 Le FAIL est une absence de réponse.\footnote{Ceci laisse entrevoir la possibilité de réintégrer les objets parfait paires avec l'incertitude en leur perfection}
 
 Le Change me dit que ce n'est pas une constante sur l'intervalle autour, Le Meme que la vitesse de variation si elle existe elle est plus lente que ce qu'il m a fallu pour parcourir l'intervalle ou alors que je mon temps de lecture est sur un multiple de la vitesse de changement.
 
 Je peux fabriquer une garantie  partielle en mesurant des temps perturbés  [ consécutifs dans une direction]  ou en vérifiant  le vérificateur.
 
 \subsection{VERIF VERIF}
  Je peux verifier à partir d'une copie locale sur une mémoire Verif si le verif  est constant sur ce dont je suis plus sur de la constance puis que garanit par un Verif.
 
 Je peux enchainer mes vérifications en choisissant ou vérifier si j'ai une variable aléatoire : on voit donc que la paire de mon estimation du type de mon adresse va me permettre de lancer deux chaines en parallèles dont l'idée est d'amplifier la confiance  en un intervalle de stabilité ou de variabilité.
 
 \section{Dynamiques de traits comme dérives opérationelles}
\subsection{l'helice}
%$${\scriptscriptstyle \begin{array}{ccccccccccccccc}
%  &   &   &   &   & - & - & \red{-} &\red{ -} &\red{ \rightarrow} &   &   &   &   &   \\
%  &   &   &   & | &   &   &   &   &   & | &   &   &   &   \\
%  &   &   &   &   &   & - & - & - &   &   &   &   &   &   \\
%  &   &   &   &| & | &   &   &   & | & | &   &   &   &   \\
%  &   &   &   &   &   &   & - &   &   &   &   &   &   &   \\
%  &   &   &   & | & | & | &   & | & | & | &   &   &   &   \\
%  &   &   & - &   &   &   &   &   &   &   & - &   &   &   \\
%\blue{\uparrow} &   &   &   &   & | & | &   &   & | & | &   &   &   & | \\
%  &   &   & - &-  &   &   &   &   &   & - & - &   &   &   \\
%\blue{|} & | &   &   &   &   & | &   & | &   &   &   &   & | & | \\
%  &   &   & - & - &-  &  &   &   & -& - & - &   &   &   \\
%\blue{|} & | & | &   &   &   &   &   &   &   &   &   & | & | & \red{|} \\
%  &   &   & - &  -& - &   &   &   &  -& - & - &   &   &   \\
%| & | &   &   &   &   & | &   & | &   &   &   &   & | & \red{|} \\
%  &   &   & - & - &   &   &   &   &   & - & - &   &   &   \\
%| &   &   &   &   & | & | &   & | & | & | &   &   &   & \red{\downarrow} \\
%  &   &   & - &   &   &   &   &   &   &   & - &   &   &   \\
%  &   &   &   & | & | & | &   & | & | & | &   &   &   &   \\
%  &   &   &   &   &   &   & - &   &   &   &   &   &   &   \\
%  &   &   &   & | & | &   &   &   & | & | &   &   &   &   \\
%  &   &   &   &   &   & - & - & - &   &   &   &   &   &   \\
%  &   &   &   & | &   &   &   &   &   & | &   &   &   &   \\
%  &   &   &   &   & \blue{\leftarrow} & \blue{-} & \blue{-} & - & - &   &   &   &   &   \\
%  &   &   &   &   &   &   &   &   &   &   &   &   &   &   \\
%  &   &   &   &   &   &   &   &   &   &   &   &   &   &   \\
%\end{array}}$$
avec la version losange?
$${\scriptscriptstyle \begin{array}{ccccccccccccccc}
  &   &   &   &   & &  &  && &| & ._p  &|      &   &  \\
  &   &   &   &   & - & - & \red{-} &\red{ -} &\red{ \rightarrow} &  &-   &   &   &  \\
  &   &   &   & | &   &   &   &   &   & | &   &   &   &   \\
  &   &   &   &   &   & - & - & - &   &   &   &   &   &   \\
  &   &   &   &| & | &   &   &   & | & | &   &   &   &   \\
  -&   &   &   &   &   &   & - &   &   &   &   &   &   &   \\
  ._d&|   &   &   &  \brown{|} &  \brown{|}  &  \brown{|}  &   &  \brown{|}  &  \brown{|} &  \brown{|}  &   &   &  &   \\
  -&   &   & \brown{-} &   &   &   &   &   &   &   & \brown{-}&   &   &   \\
\blue{\uparrow} &   &   &   &   &  \brown{|} &  \brown{|}  &   &     \brown{|}  &  \brown{|}   & &   &   &   & | \\
  &   &   & \brown{-} &\brown{-}  &   &   &   &   &   & \brown{-} & \brown{-} &   &   &   \\
\blue{|} & | &   &   &   &   &  \brown{|} &   &  \brown{|}  &   &   &   &   & | & | \\
  &   &   & \brown{-} & \brown{-} &\brown{-}  &  &   &   & \brown{-}& \brown{-} & \brown{-} &   &   &   \\
\blue{|} &  |  &  |  &   &   &   &   &   &   &   &   &   & | & | & \red{|} \\
  &   &   & \brown{-} &  \brown{-}& \brown{-} &   &   &   &  \brown{-}& \brown{-} & \brown{-} &   &   &   \\
| & | &   &   &   &   &  \brown{|} &   &  \brown{|} &   &   &   &   & | & \red{|} \\
  &   &   & \brown{-} & \brown{-} &   &   &   &   &   & \brown{-} & \brown{-} &   &   &   \\
| &   &   &   &   &  \brown{|} &  \brown{|}  &   &     \brown{|}  &  \brown{|}   &  &   &   &   & \red{\downarrow} \\
  &   &   & \brown{-} &   &   &   &   &   &   &   & \brown{-} &   &   &  - \\
  &   &   &   & \brown{|} &  \brown{|}  &  \brown{|}  &   &  \brown{|}  &  \brown{|} &  \brown{|} &   &   &  | &._b   \\
  &   &   &   &   &   &   & - &   &   &   &   &   &   &  - \\
  &   &   &   & | & | &   &   &   & | & | &   &   &   &   \\
  &   &   &   &   &   & - & - & - &   &   &   &   &   &   \\
  &   &   &   & | &   &   &   &   &   & | &   &   &   &   \\
  &   &   &  - &   & \blue{\leftarrow} & \blue{-} & \blue{-} & - & - &   &   &   &   &   \\
  &   &   |&  ._q &|   &   &   &   &   &   &   &   &   &   &   \\
  &   &   &   &   &   &   &   &   &   &   &   &   &   &   \\
\end{array}}$$


et tentative de la version rectangle
 $$\begin{array}{ccccccc} &\dmat & \dmat & \begin{matrix}._p|\\ \,|\end{matrix}&  \pmat & \pmat &\\
   &\dmat & \dmat &||&  \pmat & \pmat &\begin{matrix}-\\._b\end{matrix}\\
  -&-&-&&-&-&-\\
  \begin{matrix}._d\\-\end{matrix}& \qmat&  \qmat&||&  \bmat& \bmat&\\ 
  & \qmat&  \qmat& \begin{matrix} |\,\\ |._q \end{matrix}&  \bmat& \bmat &  
      \end{array}$$

ou de bouts constructifs

\begin{tabular} {|c|c|c|c|}
\hline
$m_b$ as $\begin{matrix} &|\\ - &._{f_{dp}m_{pd}}\end{matrix}$ & 
 $m_p$ as $\begin{matrix} - &._{f_{bq}m_{qb}}\\&| \end{matrix}$& 
 $m_d$ as $\begin{matrix}  ._{f_{bq}m_{qb}} &-\\ | & \end{matrix}$ & 
 $m_q$ as $\begin{matrix} | & \\ ._{f_{pd}m_{dp}}&- \end{matrix}$\\ 
\hline
$f_d$ as $\begin{matrix} -&-\\ |& ._{m_{q}}\end{matrix}$ &
$f_q$ as $\begin{matrix} |& ._{m_{d}} \\ -&- \end{matrix}$ &
 $f_b$ as $\begin{matrix}  ._{m_{p}} & | \\ -&-  \end{matrix}$ &
 $f_p$ as $\begin{matrix} -&-\\  ._{m_{b}}&|\end{matrix}$  \\ 
\hline
$\begin{matrix} \downarrow &\\ &\rightarrow \end{matrix}$&
$\begin{matrix}  &\rightarrow \\\uparrow& \end{matrix}$&
$\begin{matrix} \leftarrow &\\ &\uparrow \end{matrix}$&
$\begin{matrix} &\downarrow \\ \leftarrow & \end{matrix}$\\
\hline
\end{tabular}
Les flèches indiquent le chemin appris vers le centre pour le point de vue f  , le point est indice par le type de motif supporté.


\subsection{Notation Boite ou espace}
On verra que dans la représentation jardin on a dessiné  une version  boite 
$"d_f"$ comme  $\begin{matrix}
  &-&&-&\\
  &|&&&\\
  &-&&-&\\
  \end{matrix}
  $   , $"p_f"$ comme 
 $\begin{matrix}
  &-&&-&\\
  &|&&|&\\
  &&&&\\
 \end{matrix}$, $"b_f"$ comme  $\begin{matrix}
  &-&&-&\\
  &&&|&\\
  &-&&-&\\
  \end{matrix}
  $   , $"q_f"$ comme 
 $\begin{matrix}
  &&&&\\
  &|&&|&\\
  &-&&-&\\
 \end{matrix}$

et la liane de haut en bas qui commence par $d$ et continue si étendue par  des $qbpd$ coupables sera naturellement $f_d$.
NOTE : je suspecte $m$ d'etre $q$ et $f$ d'etre $b$ mais dans le doute commencant par separer le cote maniere ou dolors des roies pour clarifier , ma sexuation est un outil temporaire de clarification visant a être retiré
 \subsection{ rebondir entre  $p_f$ et $q_f$ découpe en potentiels (binaire) ou compteur de niveau d'eau (unaire)}
 Un passage horizontal entre un $p_f$ allongé et un $q_f$ allongé serait comme:
  
 $\begin{matrix}
  &-&&-&\\
  &|&&|&\\
  &|&&|&\\
  &|&&|&\\
  &|&&|&\\
  &|&&|&\\
  &&&&\\
 \end{matrix}$ puis 
 $\begin{matrix}
  &&&&\\
 &|&&|&\\
  &|&&|&\\
  &|&&|&\\
  &|&&|&\\
  &|&&|&\\
  &-&&-&\\
 \end{matrix}$ 
  En normalisant en intervalles de meme taille pointer les interstices du bord droit  par des flèches va être comme certifier  en unaire le nombre d'échelons franchis 
  (ou le prochain à atteindre) 
  Interpréter les marquages tourné comme des mots sur l'alphabet $-|$ retrouve l'encodage binaire (avec les deux chois partir de la puissance d'avant ou de la suivante)
   
   
 Plus tard il y aura un histoire et le nombre d'exploration dépendra du choix d'ouverture à chaque tour  ce qui limitera le potentiel par rapport à ce qu'on pourra voir  par limite d'espace.

L'épaisseur du doigt qui montre le niveau d'eau, l'unité de mouvement vertical minimale avec la quelle je l'imagine se mouvoir  et l'épaisseur des interstices pour une de coupe uniforme du coté sont des éléments de détail de la description.

On peut imaginer, l'age du narrateur comme  lui donnant  accès à des unités plus fines car il faut forcément des temps de pensées pour les construire avant de les utiliser comme longueur de trait. 

\section{ Les simples dérives}
La dérive simple peut être vu comme élement primitif ou ontologie minimale du trait n'ayant jamais souffert une décomposition.

C'est l'histoire de Dolores, vue comme séquence de traits apparaissant consécutivement sur 4 axes.Un bout dhistoire élémentaires pour dolores ou de dérive est donc   soit le meme nombre de traits mais dessinés sur les 4 axes, vu comme un dessin fait en
\begin{itemize}
\item  trouver l'axe en cours dans l'exploration d (horizontale gauche) puis  q (bas) puis  b (droite) puis  p (haut)  
 \item  identifier les traits avec les traits d'avant
\item rajouter un  trait 
 \end{itemize}

famille simple avec identification sans croisement vu comme déplacement  d'une sequence d'horizonles vue comme matching

\subsection{un mouvement en cache [parfois-souvent] un autre}
aller-retour rigide en haut mais  ondulation  en bas (surqd)  séparé par fracture 

$
\begin{matrix}
\text {ancre:} &  =  &&\\
\\
&-\\
&&\uparrow&\downarrow \\ 
&-\\
&&\updownarrow\\
&=
\\
&&\updownarrow\\
&-
\end{matrix}$

\begin{tabular}{cllll}
&Vers Haut et Extension&Vers Haut et Contraction&Vers Bas et Contraction&Vers Bas Et extension\\
avant& 
$\begin{matrix}||&&|&&||&&|\end{matrix}$ &
$\begin{matrix}||&&|&&||&&|\end{matrix}$ &
$\begin{matrix}||&|&&||&&|\end{matrix}$&
$\begin{matrix}||&|&&||&&|\end{matrix}$& \\
après& $\begin{matrix}||&|&&&||&&&|\end{matrix}$ & 
$\begin{matrix}||&|&||&|\end{matrix}$ & %$\begin{matrix}||&&|&&||&&|\end{matrix}$ &
 $\begin{matrix}||&&|&||&|\end{matrix}$&
 $\begin{matrix}||&&|&&&||&&&|\end{matrix}$\\
\end{tabular}

\subsection{topologie relative}
le statut de  liane  ouverte ou fermée  dépend de qui regarde : c'est ouvert si je n'ai pas encore de pointeur qui permet de synchroniser les  mouvements translation rigide  et d'extension contraction %au dessus de la fracture d'attache et 

\subsection{ une relativité bien relative}

\subsection{des espérences ou des formes grandissantes de moyens de  moyenner}
La fracture centrale  est établie avec une séquence d'estimateurs qui grandit, moyenne , médiane, moyenne géométrique, etc comme famille de nomres à raffiner  entre norme 1 et norme infinie
\subsection{rigidificatio des ancres}
Un jardin une sequence de blocs  ancres  , cahque bloca est un bout de jardin (qdpb) (potentiellement interrompu)

A voir comment gérer les failles au niveau des blocs.

Une idée c'est que toutes les ancres peuvent êtres vues comme des traces d'un meme axe on a donc du coup naturellemet deux jardins
celui qu'on voit et celui dans notre mémoire, ou celui qu'on voit et celui qu'on nous demande de prédire

  \section{Les Jardins Suspendus}
 
 L'unité sémantique  aurait pu être le carquois qui revendique des flèches qui pourraient être des carquois elles mêmes et dont 
on n'est pas vraiment sur qu'il ne soit pas juste  une flèche d'un carquois rêvant d'être un carquois.
Mais comme c'est dur de voir quelle flèche est la plus vieille dans un carquois, on va utiliser  les  jardins suspendu : des lianes élastiques qui supportent les torsions, des ancres et des pommes ou autres fruits. 
 
 On peut nommer placer deformer le jardin et on peut même dessiner des jardins dans les fruits imaginant ce qui poussera de leurs graines.
 
 C'est une trace dépliée de dérive qui explique la décision locale dans un cadre comme une adresse mémoire les plus basses étant les plus anciennes et donc les plus soumises à la chute et l'oubli. 
 
 Les marqueur pointage horizontaux montrent des cadres ou des liens. 
 
 \subsection{  Liane et jardin suspendus }

 la Liane est ce qui va jouer le role du groupe si on la rigidifie [c'est le trancon qd] porteur des constantes
  
  la pomme primordiale est la marque a plus basse et a certaines position les fruits sont des reperes sur un contenaire d'eau 
%  
%  Une Liane :
%  $\begin{matrix}
%  &-&\\
%  &&|\\
%  &|&\\
%  &-&\\
%  |&&\\
%  &-&
%  \end{matrix}
%  $



  Une Liane :
  $\begin{matrix}
  & &||& &\\
  &-&&-&\\
  &|&._d&&\\
  &-&&-&\\
  &&|&&\\
  &&&&\\
  &|&._q&|&\\
  &-&&-&\\
  &&|&&\\
&-&&-&\\
 &|&._p&|&\\
  &&&&\\
  &&|&&\\
  &-&&-&\\
    &&._b&|&\\
  &-&&-&\\
  &&|&&\\
  &-&&-&\\
  &|&._d&&\\
  &-&&-&\\
  \end{matrix}
  $  
  un cellule de le liane peut être vue comme pouvant contenir un angle de torsion ou dans version contractee pour la trace le nombre de rotation a gauche- le nombre de rotations a gauches exprimes avec l'untie du locus
  
  le lcm harmonise la liane  pour rendre les nombres lisibles avec l'unite de la liane 
  la séquence de pop ayant généré le nombre permet le retire en arrière pour voir plus finement l'imprecision d'un tour précédent
  
  sa représentation cristallisée sera l'interpretation comme des séquence de slots  avec deux marqueurs  sur ce qu'ils pensent être un slot
mais plusieurs slots 

unaire : sequence de traits  | espacés par -   [image du liquide qu'on peut remplir ou faire secouler]


la ligne  [dp] est 
binaire :  des pointeurs sur des lianes siginifiantes d on sait deja nb lianes potentielles , nb markeurs,  nb de pretravail de reptition  [ngn,nr], 


 
  Pour faire sens il faut un anthropomorphisme vu comme système d'alternance  encodage du pouvoir de choisir qui s'impose au lecteur de traces 
  [mais la lecture etant partielle par un choix qui ne depend pas de l'encodeur cette commande n'explique pas seule: pour un sens communitautariste il faudra justifier la disparition du spectateur dans le propos  ]
   et deux lianes
% 
%\subsection{Graines de torsion dans la liane, cette communauté qui pourrait devenir  un jour un groupe}
%
%pousser le plus petit de sa taille d'origine a deux formes [composition/decomposition ou vertticale/horizontale)
%
%gcd style  s'il devient trop dangereusement pres de depasser   amputez  le plus grand des deux  et echangez les roles comme si les poids lourds descendent et les poids legers montent
%  
%pcm style sil depasse trop pousser le plus grand et recommencer
% s'il depasse pousser   
%
%\textbf{ a retourner}
%
%$
%\begin{matrix}
%| \longrightarrow&&\downarrow&&\\
%                  &-&-&-&- \\
%                  &&&&\\
%                   &-&-&-&- \\
%| \rightarrow&&&\uparrow&&\\                
%\end{matrix}
%$
%donne 
%$
%\begin{matrix}
%| \longrightarrow&&&&\\
%                  &-&-&.^{\longrightarrow} \downarrow &- \\
%                  & & &=& \\
%                   &-&-&.^\rightarrow \uparrow            &- \\
%| \rightarrow& & &&&\\                
%\end{matrix}
%$ 
%
%Marqueur d'état dans le process d'harmonisation 
%
%La liane est un processus avec marquer unaire un peu comme un accumulateur bosonique

  
\section{Du danger des extrêmes pour ceux qui savent [pensent savoir] compter}
J'adopte ici l'écriture standard par défaut des nombres.\footnote{ Je  les construire avec le comptage par paquets de 10 et la retranscription sur 10 symboles appris en commencant par le 1 en ayant passé le test de reconnaissance visuelle de l'école et les écrivant comme on me l'a appris de gauche à droite : l'interruption donne les préfixes. L'ordre de construction est le même que l'ordre d'apprentissage c'est l'ordre usuel des entiers, la aussi pas de rebellion locale mais si j'assume par défaut que les gens font comme moi sauf mention explicite du contraire je n'ai pas quantifé mon assurance pour ce fait et elle loin d'être la plus grande.}

Si j'arrive à sécuriser [ à me rassurer sur la stabilité pendant mon expérience avec une rassurance qui me satisfait] deux adresses mémoires $q_u$ et $q_{u'}$
avec des valeurs que je vois comme des nombres naturels $b^{v},b^{v'}$.

Les questions que je sépare sur la ligne 
$[Hist \begin{matrix} \downarrow \\ . \end{matrix} \rightarrow [v<v'] |[ v\le v'] |[ v'\le v] |[ v'<v] \begin{matrix} \ | \\ | \\|  \end{matrix} ]$  que je lie un par un dans l'ordre à  $w_1,w_2,w_3,w_4$ : 
devraient avoir exactement deux sorties vraies  les deux du mileu si égalité et un coté ou l'autre si j'arrive à séparer pour que ce soit compatible avec la logique usuelle que j'adopte par défaut cela dit en passant.
 
J'ai donc la paire qu'on me dit bonne et l'autre : deux paires d'adresses ou de liens. La première est  celle ou je fais confiance, la deuxième le sous-mot manquant en préservant l'ordre et que j'écris en assujettissant ma variable la plus soumise qui j'espère ne se rebellera pas avant la fin du paragraphe
que je noterai $i$ ici et celle ou j'ai juste un tout petit moins confiance mais bon qu'y puis-je co-nommons en enchainant le $j$  et espérant que cela tienne 
[en rajoutant 1 à l'incertitude vu mes capacité à ma nier deux traits à la fois]. Nous disions donc $[w_{i_1}|w_{i_2})-[w_{j_1}|w_{j_2}))$.

Si j'en ai trois une interprétation satisfait les trois et une interprétation manque. J'ai donc une paire d'adresse [la première est l'intersection c'est celle ou je fais confiance.

Si j'en ai une j'obtiens mon premier mot miroir en rajoutant à la suite celle des trois qui manquent et qui nous l'avons vu me remettra en troisième.


Si j'en ai aucune ou 4 j'avoue mon désarroi, je rajoute 1 à mon incertitude et je passe à la suite. Je repense le titre en mon incapacité à comprendre les extrêmes qui est la cause du danger du désarroi de ne savoir que faire qui me force à ignorer ce que je vois mettant juste un point dans le visible puisque j'ai gagné un endroit que les scribes qui transcriront pourront ignorer en racontant l'histoire résumée à ce qui a compté.

On voit ici l'esprit de la méthode, j'ai amélioré ma confiance dans mes certitudes en me laissant la possibilité de continuer.
  
\section{Extension et nouvelle liane}

Il ne peut y avoir qu'une liane par temps quitte a revoir en raffinant ce qu'est le temps [une peu comme les fermions dans leurs slots]

\section{Premiers pas cherchant les rigidifications vers les mathématiques}
\subsection{Un  peu d'imagination pour libérer les nombres du joug de la perfection: nombre enseignables}

%Un nombre est enseignable s'il est de la forme $n=_{{\mathbb N_{>0}}} 2^{[-1]^{._{s_{dp}} ._d}}2^{[-1]^{._{s_{dp}}._{p'[}}}[ [2._b+1]] -[2._q]]$ 


Un nombre enseignable  est de la forme $(2(1+._{t_v}))^{d-p}(2(b-q+(1+._{t_h}))+1)$ et un paire $(2(1+._{t_v}))^{(d-p)+1}(2(b-q+(1+._{t_h)})+1)$ ou  quand je mets des parenthèses standard c'est que j'utilise les représentations standard complétées opératoirement comme vous l'avez appris ou que vous voulez.\footnote{javant c'eatit $X$, la j'ai mis $1+X$ pour voir d'abord le nombre comme successeur et si j'efface modulo les $+$ qui trainent on retrouverait une valeur raisonable,   mais ca aurait aussi  pu etre $1\pm X^{\pm 1}$ quelqeu part ce choix dit qu'on a pense l'horizontale avant la verticale (la premiere interruption est || elle coupe l'horizontale d'ou le - qui est moins long. On va essayer de garder des représentation justifiables pour se souvenir  des orientations par défaut}



l'obfuscation de valeur échange l'horizontale [puissances de (2)] et la verticale [puissance de (1/2)], c' est un changement de sigme de l'amplitude

l'obfuscation de phase change le calcul de phase de retirer le pair  $(2q)$ à l'impair $(2b+1)$ à retirer l'impair $(2b+1)$ au pair $(2q)$.

Le pair est le a la bosonique, l'impair à la fermionique. 

Le +1 c'est la séparation centrale pour maintenir la symétrie pour grader des copies locales.

Les constructions pas à pas c'est dire si on part d'une valeur $b$ sous sa forme enseignable et qu'on veut aller a $(b+b')$ la séquence 
communautaire de  (+1)  à faire qu'on enseigne à l'école : le mécanisme opératoire que les gens bien éduqués sont supposés respecter.

 \emph{
 $1._t$ est mon candidat a être la plus petite distance mémorisée que j'utilise pour imaginer mettre des trais de part et d'autre de mon nouveau coté pour voir s'il see rejoignent au centre.  4 manières de faires des paquets comme on avait 4 manieres de faire des fractions continues c sont nos 4 saisons qui sont capable de s'enchainer. $|$ est peut être  la plus petite distance entre desux traits - sur la verticale que j'essaye d'utiliser comme marquer sur l'horizontale suivante coupée en ce que j'espère être sa moitié.} 
 
 \subsection{La Fissurabilité la plus fine $._{t_v}=1$ dans la communauté}

Dans la paire  enseignable $(2(._{t_v}))^{(d-p)+1}(2(b-q+(1+._{t_h}))+1)$ on voit que $._{t_v}$ l'unité d'espace vertical fixe l'échelle de la diférentitation alors que  $1+._{t_h}$ fixe transitivement l'intégartion. 

 
$ \begin{matrix}
 (\text{terme})& \text{Commande} & \text{flèche}\\
( X|2)&X_q++&\downarrow\\
 (2X)&X_p++&\rightarrow\\
(X+1)_2&X_d++&\leftarrow\\
(X-1)_2&X_p++&\uparrow
\end{matrix}$
C'est la deuxieme étape de l'addition rendue opératoire car il faut d'abord trouver la puissnace avant d'incrémenter : le sens de la position de la fléche de val

On ne peut faire mieux qu'uniformiser l'infini  bornable et en plus ce faisant il devient contractible.
Ecriture tassée :  au lieu  d'avoir lignes de - ou espaces et lignes de | ou espace qui alternent  on projette en macro lignes 


 \subsection{ Nombres monolithiques tronqué delayés}
 On va revenir à une moitié de soi même 
 et voir ici les 4 roles comme on pourrait les imaginer s'ils étaient parfaits et seules au monde, prendre les 4 diagonales 
  simplement itérées et imaginer les modes de lecture itérés comme ces lianes contenant des cellules toutes dessinées comme trois traits avec une ouverture synchronisée.
 
 Ecritures sortantes tronquées monolythiques || indique le mileu estimé des adresses pointées:
 \begin{itemize}
 \item $d$   ou.     $._5-._3-._1 ||._2-._4$ 
 \item $p$ ou.     $._4-._2||._1-._3-._5$ 
 \item $1-2 q$ ou.     $|._1-._2-._3-._4-._5$ 
 \item $1-2 b$ ou.     $._5-._4-._3-._2-._1|$ 
 \end{itemize} 

en termes de séqeucnes de positions pointées\footnote{le pointage est le doigt qui montre, la main du noyé ou le sexe masculin en érection. Je me permets l'indécence car je  ne veux pas me priver du droit à l'image, puisque c'est la métaphore qui me semble le plus naturellement cacher la création. Il me semble en effet implicitemnt sous-entendre  une  adhèration  de fait, comme si c'était vérité, qu'on ne peut discuter entre communautés. Cependant  sincèrement je pense que jamais je n'oserais offenser autre chose  qu'une idée que je construis en moi, d'une action , d'un être ou  d'un objet,  et qui de fait ne peut y être identifié car  je le sens ainsi et donc je le redis c'est pur hubris que de croire que l'on peut parfaitement savoir, et j'ai peu de phrases que je mets à droite de cette dernière.} 
 
\subsubsection{draft notes propositions pour plus tard} 
quelque part on aurait 
\begin{itemize}
\item $bPOP$ comme donner le Block de position droite en écriture miroir ( poids fort) puis faire un $bPOP$
\item $qPOP$ comme donner le Block de position gauche en écriture miroir  (poids  faible) puis faire un $qPOP$
\item $dPOP$ comme donner le Block de position juste en dessous ou égal à  l'apriori [moyenne localement certifiée]  puis faire un $pPOP$ 
\item $pPOP$ comme donner le Block de position  juste au dessus ou égal à  l'apriori [moyenne localement certifiée] t puis faire un $dPOP$
\end{itemize} 

on pourait tenter 

$.^w-.^k| .^l$  dans le sens  attendre (descendre dans la liane $w$ Blocks avec point de retour si erreur la barrière juste avant) ensuite utiliser la bariiere en cours et celle $l$ plus bas comme barriere rebondissante et compter jusqu'à $k$ dans un mouvemet osciallant puis répéter la deuxième parte : remonter au block après l'attente  et recompter jusqu'à $k$ . Quelque part l'ancre de la liane est un locus et on peut raffiner en cherchant  des jardins dans les cellules

si on supposait les quantités bien calculées $|^{m-m'}$, $-_{n|n'}$ pour encodages unaires ou niveau d'eau (normalement c'est m parmi m' mais si on dépasse on aura une technique de gestion localement définie?
 et aussi  utiliser le $+^l_{j|n}$ pour dire qu'on a un encodage avec au plus $n$ pointeurs que c'est le $j$ème élément sachant qu'on lit avce la méthode $l$ le potentiels $2^n$ positions

 \section{De la logique honnête ? }
 Ayons l'honnêteté de voir la décomposition pour ce qu'elle est : un aveu d'impuissance à comprendre les choses dans leur ensemble 
 qui est aussi, quand on le prend par recouvrement de la multitude de manière de faire, au lieu d'en prendre une, notre, seule manière d'appréhender no sens en y collant des mots. Quelque part on pourrait tout résumer en "$\neg \exists ! \quad := \quad -$" lu de gauche à droite  comme un principe : quand je n'arrive pas à définir de manière unique je ne peux que me résigner à décomposer et imaginer des qualités qui scindent le perceptible, c'est la que le choix opère, le chemin de retour n'est possible que si on gardé trace de la multitude de ce qu'on aurait pu imaginer d'autre.
\section{Exemples de compréhensions transformables concrétisables dans la diversité}

\section{De la décence ou  du faire semblant pour le dernier message, Jardins présentables avec nombres et paires en uniformes}
Soi on se restreint  aux discours non interrompus pour retrouver le sens traditionnel, on va pouvoir fermer les carres et les changer en paire de nombres portant  un poids [qui est un nombre] et les lianes comme temps d'attente ou longueur en simple nombre, ....

Les fruits sont prêts a être cueillis je n'ai plus besoin de savoir que leur fraction d'existence ou l'impossibilité [fraction rationnelle positive/negative]
La liane contient l'information des fruits qu'elle contient et le jardin se décrit en parcourant ses lianes.
 
 la comparaison des bords revus par renormalisation pour arriver a la trace finale update l'information sur la valeur et la fiabilité  de ce qu'on construit en partant du comptage de traits par cote acknowledge par test de stabilité de la copie locale.
 
Les  discours non interrompus vont avoir ce semblant de perfection puisqu'il s'ancrent entre un début et une fin ou une source et un puits.
Il vont donc gagner la simplicité de la forme parfaite comme correction du discours, à l'échelle locale et globale.

les paires corrigées $\begin{matrix} ._1&-&._2\\ t_2&&\\&&t_3\\&&t_1\end{matrix} $

on va maintenir que les lianes s'accordent avec leur cellules par ppcm les cellules apres messages echanges sont parfaites les liens qui lient les cekllules dans les lianes  sont les espaces qui les séparent

on peut comparer un message vu comme a un nombre a ce que perséphone en dit 
et sur  $n$ , $Get_p n$  et puis   $Get_p  Get_p  n$ on peut essayer de voir la valeur comme égale a son estimation supérieure 
un protocole qui en decoule pourra identifier les trait a des construction a la fraction continue ou le - est la différence et le | le trait de fraction
avec entiers comme  $n\mapsto n++-|$ et sinon | | LOC INV REM


on va presque retrouver les mathématiques standard mais on aura l'antisomme au lien de la différence
histoire de garder le processus operatoire du calcul par somme de la difference en fonction de l'intention

etant donne $(x-x')$ et$ (z-z')$  
l antisomme va decomposer ce qui donnera ultimement la valeur $(x+z) -( (x'+z')$ mais avec un ordre opératoire explicité par défaut en fonction du locus et altérable.

On a  et doit garder constemment à l'esprit quatre endroit d'écriture / lecture potentielle ..

si on veut mimiquer la decomposition il va falloir utiliser lespérence d'être au milieu retirer 1 comme d et retirer 2 pour q  rajouter un a la valeur initialement copiee  pour p rajouter 2 pour b


Les nombre positifs ou optimistes  pensent avoir   dans  leur tête  leur futur séparé comme préfixe. \emph{proposition temporaire } ils sont  de la forme  
, $\begin{matrix} | &- &-& | \\ {b_q}_| &b_{||}&-&{b_q}_| \end{matrix}$ 


\footnote{ plus tard on verra
Les approchants régularisés : méthode de lecture nous montre le point d'ancrage comme monome portant un poids
Un polynome  est relatif au  monome qui le regarde.  $[k|P]$     peut laisser une trace un espace pour écrure ou la $k$ème infirmation comme update de l'information locale. c'est la mise en packet des écriture usuelles lecture imposée par la puissance des plus forts qui choisissent la taille du bus le bit, l'ictet le x32 ou le x64.}


\section{Les mathématiques}
\begin{itemize}
\item Le losange est le point de vue fonctorialisé les rôles sont des foncteurs entres des objets qui sont des catégories ce sont d'abord des anthropomorphismes \textbf{ ATM} les objets sont les objets de fracture ou de dispute ce sont les objets vus par leur valeur  comme  manières rigides de penser \textbf{RTM}
les fractures sont les sommets, les équivalences les long des cotés.
\item Le rectangle est le point de vue objectifié les rôles sont d'abord des manières rigides de penser  identifiés par leur adresse comme une catégorie de mots qui suivent les ':' dans une definition, les fracture sont au milieu des arètes. Les équivalences sont les diagonales.
\item Compose | :  est le copier localement   pour distinguer c''est la \textbf {division } génératrice de groupe
\item Decompose -:est les grouper pour séparer c'est l' \textbf{anti-somme} qui fait des paquets de traits
\end {itemize}

\subsection{l'anti-somme ou pousser vers le centre }
Si ou plutot quand , car ici tout si est un quand,  
on a deux espace extensibles ou ou peut rajouter des traits $|._q ._d | | ._p    ._b|$

ajouter a gauche c'est augmenter a jouter à droite c'est diminuer la quantité étant l'unité locale d'action en distance de pas ou nombre de traits.

Si on imagine que chacun a un passé imaginaire  en fonction de la direction d'ou il vient. Vouloir écrire au de l'autre coté de son camp  tronque et fait f'ait rebondir de la différence. Cette action est localement indexée par le temps local en nombre de clics. 

Pour avoir la compositionnalité maintenant on doit gérer  deux paires  $|._q -._d | | ._p  -  ._b|$ avec une histoire qui connait le nombre de pas horizontaux et le nombre de pas verticaux  et les positions [temps local [en clics] de rebonds]

\subsection{graine de  pgcd par coupe et insertion de liane }


\begin{itemize}
\item chaque liane à la taille de la dernière poussée, 
\item on fait grandir la plus petite des deux pointées 
\item quand   ca semble pareil en taille on attend
\item quand ca dépasse on coupe le bout qui dépasse comme nouvelle liane qu'on insèrera
\end{itemize}


\subsection{La division ou graine de ppcm par extension de lianes graine de torsions de ce qui deviendra un jour un groupe}

\begin{itemize}
\item chaque liane à la taille de la dernière poussée, 
\item on fait grandir la plus petite des deux pointées 
\item quand   ca semble pareil en taille on attend
\item quand ca dépasse on indexe par un encodage du temps et on recommence la plus petite étant devenue l'autre
\end{itemize}




\textbf{ a retourner}

$
\begin{matrix}
| \longrightarrow&&\downarrow&&\\
                  &-&-&-&- \\
                  &&&&\\
                   &-&-&-&- \\
| \rightarrow&&&\uparrow&&\\                
\end{matrix}
$
donne 
$
\begin{matrix}
| \longrightarrow&&&&\\
                  &-&-&.^{\longrightarrow} \downarrow &- \\
                  & & &=& \\
                   &-&-&.^\rightarrow \uparrow            &- \\
| \rightarrow& & &&&\\                
\end{matrix}
$ 

Marqueur d'état dans le process d'harmonisation 
%\emph{draft temporaire}
%on peut  maintenir par des insertions quelque chose comme 
%$$(-1)^{O_{b_-}} (b_{d_{t-}}+ |\frac{b_d- b_p} {1+b_q}|^{(-1)^{O_{b_|}}})$$ 

%a lire comme% $\text{SWAP}_-^{O_{b_-}}  (b_{d_{t-}}+ \text{SWAP}^{O_{b_|}}_| |\frac{b_d- b_p} {1+b_q}|}})$ ou

% $\text{SWAP}_-^{O_{-}}    (b_{d_{t-}}+\text{SWAP}_-^{O_{|}}| |\frac{b_d- b_p} {Next_{b_q}}|})$

\emph{Quand la communication est assumée parfaite ? illus[rer les 4 manieres des 4 personnage a faire des fraction continues comme exemple concret?}

 


\subsection{Des traits aux nombres libérés et des nombres aux dessins avec un pixel en cours}
\subsection{Le pgcd et le ppcm }
\subsection{Les corrections d'erreurs par échanges de messages/ renormalisation}
\subsection{Discussions entre 4 arithmétiques}

\subsection{nombres et ideaux}
j'aimerais que chaque nombre soit premier dans son idéal et que qund j'injecte un idéal dans un autre je dois generer comme ilaginaire la dynamique de ce que ce dont je n etiens pas compte.

$(x|y-z)$ comme par $y$ pour $z$  pourait peut etre se faire en injectant dabord les deux dans l'ideal de z ?

\subsection{Partir de la bi-directionnalité c'est peut être commencer par la multiplication}
Et si le (+) etait une des 4 partitions du fois celui de la decomposition avec les permutations construite le jour t agrandissant sa liane ?, quand on pousse les doigts qui pointent sur la ligne on obtient des (+ -) comme mots binaires entre deux barres|+--+--| par exemple le +1 a la puissance de deux nait de répéter passer au dessus et pousser en bas les marqueurs un étage au-dessus.
Puis sections binaires avec contre poids? $|w||w'|$ ou $w$ et $w'$ sont des mots sur $+, -$ avec  $w, w' $ de longueur compatible avec le temps local et dont la longueur diffère d'au plus un. 

Dans ce cadre, l’addition ne serait pas considérée comme une opération primitive, mais comme l’une des quatre partitions directionnelles issues de la dynamique multiplicative. Elle apparaît comme une projection locale d’une trajectoire de croissance globale, encodée par une liane dont la structure se raffine à chaque étape.

\subsection{Liane fermée et Liane ouverte}

Si on dit qu'un liane fermée a un "canal" ou une connection qui mene a une de ses cellules ca fait un trait particulier qui la coupe en deux on a deux types de dynamiques : une symetrique avec des va et vient de l'ancre au trait et une tournante dans la partie du dessous si on dit qu'un changement de direction fait redessiner la liane on a juste le point d'entrée et le nombre de pas pour avoir une dynamique interruptible

une liane ouverte c'est une chemin descendant ou on n'a toujours pas pointé sur quelque chose indistinguable du passé

\subsection{des compositions linéaires et dans le plan de motifs}
soit $m$ un motif $\begin{matrix} m_{t_6}|_{t_2}m_{t_4}|_{t_1}m_{t_5}|_{t_3} \\ .  \end{matrix} \begin{matrix} m&.\\ \rightarrow&\uparrow \end{matrix}$
et  $$\begin{matrix} m^d_{t_5} &|_{t_4}& m^p_{t_7}\\-_{t_3}&&-_{t_2}\\m^q_{t_6}&|_{t_1}&._m \end{matrix}$$

\section{Elasticité et translation pour représenter les impressions de puissances}
Ayant bien pris le temps de penser  oin va tenter  de figer un peu plus les idées pour retrouver les mathématiques des autres.

Tout d'abord essayons une mécanique opératoire pour représenter ces écritures avec une base donc chaque puissance transporte une information.

Imaginons une structure de controle comme un demi-corps ouvert dans le bras allongé pointe du doigt (ou plutot de deux doigts pour maintenir le coté contrpoids des nombres binaires).

 $\begin{matrix} 
 - & -&-&-&&\\
 |& & &  |&\\
  & &- &  - &-&\\
 |& &\downarrow &&\downarrow\\
 -&&&\\
 \end{matrix}$
 
 
 Maintenant imaginons ses pieds solidaire à  un demi-corps ouvert  qui peut  etre vue comme un élastique oscillant (ATM?) ou une structure rigide en translation et rebonds (RTM). et qui peut etre lui meme vu comme  deux  demi corps  ouverts avec 4 choix de pointage
  $\begin{matrix} 
 - & \\
 |&  \\
 -&\\
 &\\
 -&\\
 |&\\
 -&\\
 \end{matrix}$
    
   On imagine le motif sur mur gauche  
    $\begin{matrix}
    -\\
    w_3\\
w_2 \\
w_1\\
-\\
-\\
\tilde{w_1}\\
\tilde{w_2}\\
\tilde{w_3}\\
-\\
\end{matrix}$ lu par des bras oscillant comme apte  vérifier l'égalité du motif et de son ombre  et à agir en fonction sur le bars de controle sur l'autre versant
et  $\begin{matrix}
-\\
    w_3\\
w_2 \\
w_1\\
-\\
\tilde{w_3}\\
\tilde{w_2}\\
\tilde{w_1}\\
-\\
\end{matrix}$ lu par des bras rigides (montant pour q ou descendants pour b   comme apte  vérifier l'égalité du motif et de sa translation   et à agir en fonction sur le bras de controle sur l'autre versant
  
 
\section{ [Section Cle] Informatique }
suppose that we have  POP actions from a list that try to get  $A -  B$  on local memory   and $REM$ is the list that lost $A$ and $B$ reconnected if disconnected by the action to have a coherent meaning as cell

and  rescale to make comparable using multiple $ [A-B]^{REM} |  [A]^{REM}-B$ 

\begin{itemize}
\item |:  rescale to make comparable using multiple $ [AB]^{REM} |  [A]^{REM}B$  with free associativity as  
\subitem POW POP ADJ POP REM
\item - :    $[A \text{ explains } B]^{REM} -A^{REM} \text{ explains } B$  which can be written with free associativity as 
\subitem ADJ POW REM POP POP
\end{itemize}
\section{Le Graphe comme process libérable}
Bout d'arêtes à connecter et ordre total sur la découvertes des ports autour d'un sommet. chaque sommer a un coté d'entrée horizontal ou vertical 
et l'écart entre le nombre de traits avant le mileu et ceux après le mileu ne peut être que d'au plus un on a donc un matching parfait qui aura le sens entrées sortie avec éventuelleement un trait tout seul pour le dessin ou l'attente dans la zone en cours.
Un croisement d'arete peut etre d'un type parfaitement distinguable d'un sommet mais il a besoin d'un locus déterminable qui le distingue par ou il est des sommets qui se veulent comme tel [ou on s'attarde plus longtemps].
  
\section{La complexité revisitée par des processus concrétisables dans la diversité}
On pourrait imaginer une structure dont les fragments finis seraient des courbes fermées dont la splittabilité uniforme en fragments dépend de l'age du narrateur.

On positionne une faille comme un vérifieur et deux bords d'un coté et de l'autre comme sa zone d'accès dans l'histoire.

La zone centrale balancée est la ou on couvre les langages décidable ou il couvre par "oui" [dans [partie gauche $\ge$ partie droite] mais au plus de 1 d'écart] et "non", défini dans partie droite plus grande ou égale mais pas plus de 1 d'écart. 
Et la recherche de préfixes , suffixes en expansion des bords éloignés permettrait de couvrir par prefixes et suffixes RE et co-RE 

Le prouveur ou le reste de la bande lui decide des inputs [valeurs fixes non encore lues] et des variables espace adresable non encore réclamées.

Les 4 manières de parcourir essayent de placer  un candidat-mot dans  4 familles de langages
\section{Réflexion du lundi a reclasser}
Clamer l'existence est nécessairement accompagné par la limitation induites par cette formulation sur les libertés perdues de l'environnement de cette définition. 
Nous devons donc accompagner le dessin en imaginant toutes les existences qui auraient pu avoir été illustrées. 
Néanmoins pour qu'un mot soit libérable il doit d'abord être sorti de son contexte pour retrouver toutes les libertés qu'il auraient pu avoir ou qu'il aie pu survenir, et c'est dans cet espace qu'on doit chercher la limite locale par le contexte de son écriture comme une projection puis la forme identifiée comme une par lecture de la lettre. 

Essai: on pourrait put etre dire [avec une formalisation a revoir dans un collapse partiel en cardinalités certifiées] $1_{6_{2|4}}$ si la filtration se fait par un element à 6 forme qui vu le passé ne peut plus avoir que 4 et  c'est la deuxième qui a été identifiée par lecture mais encore une fois le mélange nombre cardinal et nombre potentiel rend la forme illisible (en plus le nombre de dimension d'ouverture peut dépendre de la construction d'histoire de façon non monotone). 
 
 Un mot est comme une tentative de répétition d'un bloc radical de la liane de ses définitions localement accessibles.
 
 La terre est une classe d'équivalence pour tout ce que je peux imaginer comme ayant pu rêver y séjourner.
 
 Les motifs concrets sont mon moyen de connecter ce que je pressens être  avec les différences que je pourrais imaginer observer s'il était formalisé.
 
 La position de la faille est celle de mes apriori.
 
 Les changements prennent sens quand on les cumule et qu'on les ancre.
 
 La différentiation prend sens quand sait à peu près de quoi on parle : quand on a des locus relativement stables ou on peut aller pour comparer, certifier les différences.
  
  En partant de la multiplication des corps on n'a pas une algèbre sous-jacente mais une superposition d'algèbres imaginable d'avant la rupture qui fait l'existences pairées par un motif faisant son possible avec les algèbres qu'il peut imaginer pour l'après. Les moyens de pairer ainsi que ceux qu'il s'imagine qu'i aurait pu faire s'il s'imaginait être quelqu'un d'autres sont peut être  ce qu'une liane tente de représenter.
  
  Les premières libertés d'expressions vont arriver avec les premières limites qu'on peut mettre à la composition (multiplication,ppcm opératoire,composition) , décomposition (pgcd opératoire décomposition) décrites avec un point d'ancrage (point central du dessin dont le cadre d'interprétation (bord externe)  ne sera jamais trop difforme) et qu'on va faire dans une deuxième projection pour avoir un élément de base comme une paire , un  bord interne (aussi pas trop difforme) mais qui contient par sa trace des idéaux expliquant toutes les trajectoires autorisées. La restriction de trace certifiable par traits sur bord interne nous force à relacher la proportion et justifie qu'on n'ait accès qu'à des idéaux et non à des nombres: si pas assez d'expressivité imaginer refaire le tout en dilatant les quatre cotés.
  
  
 \subsection{humanités}
 Si l'historien honnête cherche des éléments pour revoir, agrandir des conditions et certifier et croiser les éléments pour faire émerger une Histoire [en gardant trace de ce qu'il présente comme contre-vérité en justifiant des intentions falsificatrices e leurs auteurs mais en laissant l'accès libre au futurs penseurs aux traces de son raisonnement], le politicien d'aujourd'hui veut faire croire que la mécanique qu'il propose est celle qu'il faut choisir pour une causalité dont il prétend [souvent] maitriser  ce qu'elle impliquerait. Son honnêteté tient à sa croyance en ce qu'il dit mais  son intégrité [aujourd'hui presque disparue] serait de dire lui-même les éléments [qu'ils connait bien souvent [ au moins en partie]] et qui sont les omissions qui font les exception nécessaires pour formaliser sa théorie sociale.
 
 Un journaliste doit-il restreindre ce qu'il dit en fonction de ce qu'il pense que la manière de présenter les choses entrainera comme réaction?  De ce fait soumettant le lecteur à la restriction de visibilité induite par la réaction anticipée par  le modèle que le narrateur a de ses lecteurs. 
 Peut-être parfois dans certains contextes, mais on doit toujours selon moi essayer de tendre vers le modèle du respect du lecteur et de sa responsabilisation plutôt que de son infantilisation car les dérives des modèles de soumission par  hiérarchisation imposée me semblent effrayantes au vu des modèles associés dans l'Histoire au-quels naturellement je l'associe.
 
 
 Avoir un formalisme qui cherche a exprimer les intentionnalités et les potentialités dans leur diversité est un peu essayer de ne pas fermer les yeux 
 sur ce que la tranquillité de l'esprit nous pousserait à rejeter car [prétendument évidemment]  "n'ayant rien à voir".
 Au lieu de juste choisir s'imposer un programme pour investiguer les "si" s'inscrivant dans un imaginaire qui permettrait en théorie de ne rien omettre.
 
\section{L'informatique quantique}
\subsection{Pourquoi sucesseur et lien conditionnelsuffisent}
Ontolgie minimaliste du jour 2 clics observables [2 paires intriquées quand je veux en parler d'une]:

$
\begin{matrix}
d:\rfloor &\ket{00}^{\mathcal H} |&p:\lfloor \ket{01}^{\mathcal H}&||&d:\rfloor \ket{00}^{\mathcal E}&| \,\,p: \lfloor \ket{01}^{\mathcal E}\\
&-&-&&-&-\\
q:\rceil &\ket{01}^{\mathcal H} |&b:\lceil \ket{11}^{\mathcal H}&||&q:\rceil \ket{10}^{\mathcal E}&| \, \, b:\lceil \ket{11}^{\mathcal E}\\
&-&-&&-&-\\
\end{matrix}
$

Du lien du jour 2 click

$H\otimes H :  \begin{matrix}
  | &|  &-&-&|&| \\
  |&-|&-&-&|&-|
\end{matrix}$ est la moitié d'une  unitaire opérante)


De la mixité du jour 5 : Pour un non clifford simulable avec Addone


${W_5}_U=
\begin{matrix}
4 &5 \\
1&4
\end{matrix}
-
 \begin{matrix}
1 &1 \\
5&1
\end{matrix}
$ vu comme matrices d'entiers générant par différence et rescale une unitaire non Clifford.

Une unitaire n'existerait que comme produit tensoriel de 4 unitaires localement équivalentes permettant d'encoder les effets de $U$, de $XU$, de $ZU$ et de $YU$ %sorte de $\sigma_i  \ket i P_i U \ket 
\emph{a creuser}
Du transport de lien comme émergence dans l'espace 

CNOT vue de dessus $ \begin{matrix}
  |  \end{matrix}$  pioche dans la colonne initiale et donne $ \begin{matrix}
& -\\
  | & |
   \end{matrix}$ et  $ \begin{matrix}
  |  \end{matrix} $ va aussi donner  $ \begin{matrix}
 &-&-&\\
  | & |&|
   \end{matrix}$ un élément de la classe de Toffoli  surspecifie puisque lien accroché au premier
 
 \subsection{ Qu'est ce qu'un changement d'état quantique [attempt]}
 \emph{proosition à étudier}
 Pour décrire un état quantique compatible il faut que $2^{a}$ ou $a$ est l'amplitude divise la dimension $d$ de l'espace qu'on décrit afin d'interpréter les choses comme multiplicités.
 
 L'espace des percept p a une dimension de plus que les pace mesurable d pour encoder l'inconnue, les espace bosonique et fermioniques associés ont co-dimension la dimension de ces systèmes. 
 
 $a {\ket{ \psi_{p_t}}} {\ket {\psi_{b_t}^{\mathcal F}}} {\ket {\psi_{q_t}^\mathcal{B}}} \ket{\psi_{d_t}} \ket{ \psi_{p_{t+1}}}$  
 
 ou $t\mapsto t+1$ comme nombre enseignable devrait  préserver $2^d|a$ si je peux faire une mesure a $d$ résultats possibles.
 
\subsection{$GHZ (abc|xyz)$  comme Bénédicte}


${\scriptscriptstyle
\setlength{\arraycolsep}{1pt}
\renewcommand{\arraystretch}{0.45}
\begin{array}{@{}*{12}{c}@{}}
  & - & - & - & - & - & - & - & - & - &  - \\
  & | &   &   &   & | y&   &   &   &&b | &   \\
- &   &   & - & - & - & - & - &   & &- &   \\
| &   &   & | &   &   &   &   & | &   & | &\\
  &   &   &   &   &   &   &   &   &   &  & \\
| &   &   &   &   &   &   &   &   &   & | &\\
  & - &   &   &   &   &   &   &  & -  &  & \\
| & |a &   &   &   &   &   &   &   && z| &  \\
  &   &   &   &   &   &   &   &   & &- &   \\
| & | &   &   &   &   &   &   &   & &|    &\\
  &   &   &   &   &   &   &   &   &   &   &\\
| & | &   &   &   &   &   &   &   & &c|   & \\
- &   &   &   &   &   &   &   & & -  &  & \\
  & |x &   &   &   &   &   &   &   & &|   & \\
  & - &   &   &   &   &   &   &   &   &   &\\
  &   &   &   &   &   &   &   &   & &|    &\\
  &   &   &   & - &   &   &   &   &   &   &\\
  &   &   &   &   & | &   &   &   & &|    &\\
  &   &   &   &   & - & - & - & - & -  &  & \\
  &   &   &   &   & | &   & | &   &   &   &\\
  &   &   &   & - &   &   &   &   &   &  & \\
  &   &   &   &   &   &   & | &   &   &   &\\
  &   &   &   &   &   &   & - & - & - & - &\\
  &   &   &   &   &   &   &   &   &   &   &\\
  &   &   &   &   &   &   &   &   &   &   &\\
\end{array}
}$

Une bénédicte sans erreur similaire à qu'on trouve dans le calcul local  utilise $--$ ou $-|$ comme préfixe autorisés [à compléter par symétrie pour les autres pour avoir un sens communautaire au dla du sens dans a plus petite communauté] pour nous pour que cela soit acceptable
et propose un chemin  qui peut faire des allez retours mais seule la chaine initiale minimale d'attente (répétition de -- n'est pas un espace perceptible. et n'est pas acceptées comme mot central dans une induction de la forme $|x-y|_{x|y}$ Try Get  $[a|x]_q- [b|y]_d- [c|z]_{p}$ Under allowed lengths.
et le splitteur c'est les deux sorties en bas en fonction si on accepte ou pas l instanciation comme  correlation proposee .

\subsection{Les stabiliseurs et contextuallité}
\subsection{ Rebond entre parois élastiques balancés par des chocs simulant un trou noir}
\subsection{tentative pour sequences de construction}
pour des nombre vus comme precision requise pour l'identite $n_I$ pour $X$ stabilité $n_X$ et pour $Z$ point fixe $n_Z$
on aurait pour tout nombre $N$ un idéal $\lambda w_N N$  pour $\lambda$ entier et $w_N$ pas trop grand par rapport à $N$
ou on aura pour tout $n$ de l'idéal $H_{x,n}\ket{ \psi}^{\otimes n}=_{n_X} \ket \psi^{\otimes n}$, $H_{z,n}\ket{ \psi}^{\otimes n}=_{n_z} \ket \psi ^{\otimes n}$ et $Tr(\ket \psi \bra \psi)^n=_{n_I}1$
 
 
\section{notes to fix }
Brider l'universel par l'améliorabilité: à partir d'imperfection inévitable mais toujours améliorable

Process: Les termes seront de 
\begin{itemize}
\item Confiance estimable , et
\item  d'estimation estimable
\item [ et d' estimation CLONE ligne précédente ]  Répéter ou sortir.									
\end{itemize}
en pouvoir reconstruire un identifiable de stabilité suffisante à partir de parties locales.

maintenir une sequence de definitions avec pointeur vers premiere connus et pointeur vers precedente?

On aimerait que si la représentation est restreinte a trois sur 4 consecutivement, on ait un schema de reconstuction dépendant de cette filtartion qui 
nous reproduise une expression commune comme une classe d'équivalence. Quelque part on affecterait des roles verifieur, prouveur a poyvroi borne et on cetiferait unlangage.

On pourraitt garder des "batch" relies devenenment 2 a deux distinguables en mainenant un espace comme une matrice carrée  $n^2$ avec une indexation
des evebnements distingues , puis des que evenement collusion reconstruire une nouvelle matrice en liant le un de la nouvelle case a sa case de collision estimee a l'epoque   et en mettant une taille de case supportant les permutations d'ordre local , on a une vision par colonne par lognes et par matrices , avec des sequence de matrices qui donc grandissent et dont les position ont un ordre et la taille des cases peut encider au moins le plys grandindice.



Reasonable history given context :

Two non communicating ATM,  A asked $x\in (0,1)$  answering $a\in (0,1)$ then B  asked $y\in (0,1)$ answering $b\in (0,1)$ repeatedly assuming the delay is reasonably short enough to made us believe that  they cannot transmit information in between the answers.

$$
\begin{matrix}
ax|yb&\begin{matrix} y&b\\ 0&0\end{matrix}&\begin{matrix} y&b\\ 0&1\end{matrix}&\begin{matrix} y&b\\ 1&0\end{matrix}&\begin{matrix} y&b\\ 1&1\end{matrix}&\\
\begin{matrix} a& x\\1 &1\end{matrix}& \leftarrow_{13} &\leftarrow_{6} & \rightarrow_{8} &\rightarrow_{16}\\
\begin{matrix} a& x\\1 &0\end{matrix}& \downarrow_{11} &\downarrow_{9} & \uparrow_{2} &\uparrow_{4}\\
\begin{matrix} a& x\\0 &1\end{matrix}& \leftarrow_{15} &\leftarrow_{7} & \rightarrow_{5} &\rightarrow_{14}\\
\begin{matrix} a& x\\0 &0\end{matrix}& \downarrow_{10} &\downarrow_{12} & \uparrow_{3} &\uparrow_{1}\\
\end{matrix}
$$

is read as I fixed $x$ to 0  then I fixed $y$ to 1 then I checked if  the answer was 01, then if it was 10, then if it was 00 then if it was 11
if everything is  with o error only one answer is registred per question so we should draw an up then a right then  a dawn then a left arrow.
It is a fully linearized setting that could fail or be interrupted at any step

or
 $$
\begin{matrix}
ax|yb&\begin{matrix} y&b\\ 0&0\end{matrix}&\begin{matrix} y&b\\ 0&1\end{matrix}&\begin{matrix} y&b\\ 1&0\end{matrix}&\begin{matrix} y&b\\ 1&1\end{matrix}&\\
\begin{matrix} a& x\\1 &1\end{matrix}& \leftarrow_{2} &\leftarrow_{14} & \rightarrow_{16} &\rightarrow_{8}\\
\begin{matrix} a& x\\1 &0\end{matrix}& \downarrow_{11} &\downarrow_{3} & \uparrow_{5} &\uparrow_{13}\\
\begin{matrix} a& x\\0 &1\end{matrix}& \leftarrow_{10} &\leftarrow_{6} & \rightarrow_{4} &\rightarrow_{12}\\
\begin{matrix} a& x\\0 &0\end{matrix}& \downarrow_{7} &\downarrow_{15} & \uparrow_{9} &\uparrow_{1}\\
\end{matrix}
$$
$e_{i,j}$ where $i$ is $(a|x)$ and $j$ is $(b|y)$


is read as I fixed $x$ to 0  then I fixed $y$ to 1 then I checked if  the answer was 01, 
then  I fixed $x$ to 1  then I fixed $y$ to 0 then I checked if  the answer was 10, 
then  I fixed $x$ to 0  then I fixed $y$ to 0 then I checked if  the answer was 11, 
then  I fixed $x$ to 1  then I fixed $y$ to 1 then I checked if  the answer was 00, 
etc

if everything is with no error every 4 clock we have a normal configuration.


A maximal (Full) Batch has 4 elements in the first count and 16 in the second.

when asking the same question to both and getting the same answer we want to write in the same spot, so we create a new Batch Matrix each has between $4$ and $4.2^{16}$ bits ?

A plausible cumulative  event is an even that has a past and a future that are normal.

Imagine obfuscation such that $O_B^0 (ab|xy)=(ab|xy)$ and $O_B^1 (ab|xy)=(ab|x(1-y))$and $O_A^0 (ab|xy)=(ab|xy)$ and $O_A^1 (ab|xy)=(ab|(1-x)y)$

Where normality means that for all $s,a,x$  $\Sigma_{b=0:1} N(O_B^s (ab|xy))=\Sigma_{b=0:1} N(ab|xy)$ if I obfuscate ones question and manage to sum all the results it should be invisible (and similar for the otherside) .

for each trace (16 Integers) we want to describe  find the set that is obtained  by forgetting the integer indices which give a global order to keep only the length of  sequence forming a path to a normal history ( when removing occurences as errors we want to have a number as close as possible to the exact discrete experiment) then for each pair (history, exacttrace)  we want to find the set (by forgetting the total order given by the integer indices and considering only the length , of extensions to normality).

Extensions are sequences of batches not in the experiment and removals are sequences of batches in the experiment.
so each basic element of  teachable causality is a degree two first non commutative polynomial  (constraint of noral past are variables) (constraints of normal fiture are other variables), relaxing the non commutatiivety by putting them on the same space gives degree 2 that are potential candidates for the interesting strucre I am looking for as bounded contextuality with with limit pseudotelepathy correlation on extended version of the games of the observed correlation.
 
  
the pairs of shortest completing paths (here independent but for full generality, finding all shortest pasts and for each shortespath completing wiuth all shortesfuture) is the equivalence class that define a reasonable observable event.

so viewed over 
$\begin{matrix}
e_{1,1}&e_{1,2}&e_{1,3}&e_{1,4}\\
e_{2,1}&e_{2,2}&e_{2,3}&e_{2,4}\\
e_{3,1}&e_{3,2}&e_{3,3}&e_{3,4}\\
e_{4,1}&e_{4,2}&e_{4,3}&e_{4,4}\\
\end{matrix}$ constrainsts are 
$\begin{matrix}
(N_{1,1}+N_{1,2}=N_{1,3}+N_{1,4})\\
(N_{2,1}+N_{3,2}=N_{2,3}+N_{2,4})\\
(N_{3,1}+N_{3,2}=N_{3,3}+N_{3,4})\\
(N_{4,1}+N_{4,2}=N_{4,3}+N_{4,4})\\
\end{matrix}$ and $\begin{matrix}
(N_{1,1}+N_{3,1}=N_{2,1}+N_{4,1})\\
(N_{1,2}+N_{3,2}=N_{2,2}+N_{4,2})\\
(N_{1,3}+N_{3,3}=N_{2,3}+N_{4,3})\\
(N_{1,4}+N_{3,4}=N_{2,4}+N_{4,4})\\
\end{matrix}$ 

without structure we need to imagine grouping the locuses 4 by 4 to represent the product events ?


garder a lidee toujours pas un dessin mais chemins centres[ancres]  entre deux deux dessins distincts

{\bf Idée: verif preuve narrateur}

 
 $A$ sur la verticale a comme trace le nombre de trait non verticaux par projection sur l'axe verticale (a chaque temps on peut regarder soit en haut ligne 1 soit ligne 3) et on peut savoir si quelque chose traverse +1 ou ne rien voir.
 
 $B$ a comme trace la projection sur la diagonale comme $>2.5$ ou $<2.5$ (dans les deux premieres ou les deux dernieres colonnes)${i,j}$
 
 idée: quelque part un prouveur a acces localement (sur $e_{i,j}$ a la forme du trait donc a la question globale)
 le narrateur à accès aux nombres indices a accès a tout :  a ce qu'aurait été l'histoire s'il avait été  été coupé , à la forme et à la trace
 
 
{\bf Idée multinterpretatablite de 4 sequences de traits redecouvertes tour a tour }
 
multiinterpretabilite des dessin sur un bord prendre les axes prependiculaires qui partent du milieu comme une vision et les 4 diagonales comme l'autre. 
  
  sequence de trait comme cardinal  fractionnable : compter les traits sur une section (une zone une restriction de l'axe)
  
  sequence de trait comme difference fire une rupture et une sélection non nulle des deux cites et compter la difference
  
  après conversion horizontal verticel
  rendre un point , utiliser un trait de type different comme pas, un trait de type différent comme zone de detection avancer et détecter donne une série de nombres sur un axe conditionné par un point d'ancrage et deux traits de l'autre 
  
  Pour chaque point de départ et direction simple , on a besoin d'un masque (lisible/fail), d'un pas pour recherche suivante si celle ci ok et de pas de correction si celle ci fail (peu etre évolutif mais justifiable par une quantité du passé) d'une zone de recherche (chercher une race traversante pensant un laps de temps (pareil, cette durée de moyennage ou de cumull  pour avoir un bit doit etre justifiable par l(histoire)
  
\subsection {Autre exemple }
dérouler un calcul a flot tendu $(((((5-3/2-4)/3-13$

on peut revenir en arrière exemple apre le pop du / 

on peut fractionner et faire de paquets a sommer qui deviennent comme des modules 
si on arrive a maxint 

l'associativité   implicite a une borne d'approximabilité qui depend de l'associativite explicite dans la quelle elle est  : tout choix de rajouter des parentheses
qu'on choisit d'expliciter comme implicite dans un terme est traitable en fonction de la position dans l'expression [surtout par le nombre de parentèse  en surplus] 

{\bf benedicte power description?}
$0_i\infty_j=._b^{i-j}$ 
une seqeunce de paires adresse valeur $1 \text{ as} (0_1,\infty_1), 2 \text{ as} (0_2,\infty_2)\ldots ?$  si on dit donne un nombre qui mesure la constance d'un 0 ou la variabilité d'un infini bénédicte dit $i$ comme un couplage parfait.

et on a envie de voir l'histoire comme le plus constant possible , du coté des adresses, des quantités,  des 0
et le futur comme le plus varaible possible du coité des valeurs des infinis, des qualités

\section{tentatives floues pour proposition de notation pour plus tard }
un nombre $N$  c'est $N_d-1$  pour $\ge 0$  qu'on peut voir comme $l_N= l \downarrow N\uparrow$ on a fini et revenu a la ligne
ou $1-N_b$  comme nombre négatif pur en construction  $l_N= l \downarrow N$ 

et  $N_b$ se revoit comme $N_{bq}|1$ comme taille en traind egrandir ou $1|N_{bq}$ comme taille utilisée comme unité 

\section{Communautés admettant un système de numération commun}
A finaliser le minimum requis pour utiliser unelecture/  ecriture masi ca me semble a^priori etre  au moins lecture d'un encodage binaire en commencant par les bords ou en alternant a partir du milieu comme 4 méthodes de POP contradictoires)

Essai:

Une communauté qui va associer à chaque  nombre objectif écrit ici comme séquence sur $(1|2|3|4|5|6|7|8|9|0)$
 équipée d'un methode de lecture interruptible et d'une methode d'ecriture interruptible  avec un pairage si succes
 pour chaque espace de lecture seule,   , et locus unitaire de verification et extension spaciale de calcul 
  d'un temps de lecture, 
  d'un temsp d'écriture,
   d'un temps de vérification sur un espace  unitaire du locus (dessiner le trait adapté à la ligne si on a recponnu un nombre)
   
\section{Motifs revisités}

Un motif est une dessin qui n'est pas le carre de cote 1 ou le vide, il peut etre vu comme anthropomorphisme ou losange c'est a dire comme système d'évolution de dessins avec des symbole $|,-,.$ décorés de termes.

Il peut etre vu come rectangle ou manière rigide de pensée comme objet ou  expression fixée d'une dynamique itérable .

Le premier déroule  la transformation pour expliquer un lien entre deux dessin étant donné un dessin voulu comme partie du chemin (ou l'un des deux comme debut ou comme fin si troisieme parametre manquant). Le deuxièeme énonce un principe ou décrit un dessin.



Le voir comme un losange ou un anthropomorphisme c'est le voir comme regle de transformation unfinie qui fait evoluer  un dessin 
representables comme 

\begin{itemize}
\item une sequence alternante de lignes séquence de symboles de $(-,.)$ et de lignes, séquences de symboles de   $(|,.)$ 
\item equipée de 4  systèmes de lecture infinis interruptibles comme des sequences de longueur $(l+1,l,l,l)$ ou $(l+1,l+1,l,l)$  ou ou $(l+1,l+1,l,l)$ ou $(l+1,l+1,l+1,l+1)$ chaque sequence peut etre vue comme un système de coordonnées ancré dans un coin donnant le dessin de la case d'espace de la paire de coordonnées, en respectant que dans une ligne on n'a pas - et | comme symbole.
\end{itemize}

Mettons  le . non décoré comme obstruction requise c omme séparation par exemple ou que sais-je en tout cas symbole à retranscrire sans réfléchir.

l'
Positionnons nous sur une ligne traversée par d'éventuel traits verticaux et  dessinons les espace de tracage $._{t_q}$ , $._{t_d}$, $._{t_p}$ , $._{t_b}$ en reliant par des fleches à leur trace de référence  (a analyser et revoir pour voir je voualis d et q lies a une trace du passe et p et b liee a ube trace du futur donc orientation peut etre a repenser)
%  & - &   &   & - &   &   & - &   &   & - &   \\

$${\scriptscriptstyle\begin{array}{lcccclrccccr}
 & \rightarrow &   &   & \leftarrow &   &   & \rightarrow &   &   & \leftarrow &   \\
|                  &  ._{t_q} &             &                      &     ._{t_d}& |                & |                   &   ._{t_p}  &             &                       &   ._{t_b}  & | \\
                   &             &              &                       &                &                &                    &               &                &                      &                &   \\
\downarrow &            & \uparrow & \downarrow &                 & \uparrow & \downarrow &                & \uparrow & \downarrow &               & \uparrow \\
  &\rightarrow   &   &   &\rightarrow   &   &   & \rightarrow  &   &   &   \rightarrow&   
\end{array}}
$$
L'idée est que l'espace entre les deux traits centraux est undefinissable car il yara toujrs une extension le rendra plus grand qu'-une constante eu yne contraction qui le rendra plus petit qu'une constante en général. mais les espaces de dimensions infinies $q$ et $b$ peuvent gérer son histoire et $d$ et $p$ sont plutot des espaces de dimension finie qui essayent de rendre un concept fidèle à un percept malgré leur dynamique.
%$$
%\begin{matrix} 
%|               & &._{t_q}&& ._{t_d}         &|&|&._{t_p}      &  & ._{t_b}&&|\\
%               & &&&   &&& &&&&\\
%\downarrow  & &\uparrow &\downarrow   & &\uparrow &\downarrow   & &\uparrow &\downarrow   & &\uparrow \\
%&\rightarrow&&\rightarrow&&\rightarrow&&\rightarrow&
%\end{matrix}$$  
%

pour les traits horizontaux le plus naturel est 
$\begin{matrix} .&-&-&.\\
|&._{t_d}&._{t_p}&|\\
.&-&-&.\\
|&._{t_q}&._{t_b}&|\\
.&-&-&.
\end{matrix}$
 
 ce qui donne en linéarisé sur la diagonale
 
 $\begin{matrix} 
 &-&.\\
&._{t_p}&|\\
&-&.\\
&._{t_b}&|\\
.&-&.\\
|&._{t_d}&\\
.&-&\\
|&._{t_q}&\\
.&-&\\
\end{matrix}$

et en associant des traits de référence a faire en rajoutant des fleches  c'est diffraction donc 1 pour deux ? a penser calmement voir si modifier jardin ou si revoir ici)

 $\begin{matrix} 
 &-&.\\
&._{t_p}&|\\
&-&.\\
&._{t_b}&|\\
.&-&.\\
|&._{t_d}&\\
.&-&\\
|&._{t_q}&\\
.&-&\\
\end{matrix}$

\section{Plausible  Languages} \label{sec:simple}
DRaft:
A Language is plausible if its not trivial (not empty nor maximal] and   if its words are of the form $w_Pww_F$ where each word $w$ of length $|w|$ has left and right languages $\mathcal {L}_H^w$ and  $\mathcal {L}_U^w$  rooted with $w$ as prefix/suffix such that  the history  $\mathcal {L}_H^w$ has a yes certicate Interactive proof system and the future  $\mathcal {L}_F^w$ has a no certiicate Interactive proof system. Both exchanging messages of length  $l$ in polylog  $|w|$ in $o(2)^{|w|}$ rounds.   
The bitwise interaction with provers can be tempered with by swap by an adversary.  

Given a pair of direction the bitwise asking matching should have no fix point and provides always  two  bit per pair ( language ,direction]  like what nest bext and continue or flip direction 
 
\subsection{The imaginary undefinable by choice to mimic feelings full perfect cell} 

 $\begin{matrix} -^|&&|^- \\ &&\\-_|&&|_-\end{matrix}$  called twisted losange indistinguishable  form 
 $\begin{matrix} &-&\\
                          |&&|\\
                          &-&
       \end{matrix}$ called rectangle but both are just community demons or seeds and have no definable meaning as a communatorial diversity.

It makes sens in a communautary drawing term only if one puts something inside  or on one two or 3 the 4  corner postures.

 $\begin{matrix} ._d-^|&&|^-._b \\
 				 &.w& \\
				 .q-_|&&|_-.p
				 \end{matrix}$  as twisted losange indistinguishable  form 
 $\begin{matrix} .d&-&.p\\
                          |&.w&|\\
                          .q&-&.b
       \end{matrix}$ or rectangle

%\section{off}

\section{ a candidate family of processes lacking meaning}
Suppose that we are running a non signaling expermient with binary inputs and inary outputs asked to to non communiucationag ATM $A$ and $B$

Usually we put ecverything ina probability space over 16 event (abxy) , and associated marginalisation.
 The definition of non communicating would want us t restrict meaningful steps of the experment or ones with standard interpretation to one satisfying 
the 8 non signaling equation when associating 4 family of probabilities : 
$p(a0|xy)+p(a1|xy)=p(a.|xy)=p(a.|xy')= p(a0|xy')+p(a1|xy')$ 
and $p(0b|xy)+p(1b|xy)=p(.b|xy)=p(.b|x'y)=p(0b|x'y)+p(1b|x'y)$ 
    
 Here we want to record not only the numbers  $ N(ab|xy)$  and $N(xy)$ to have a probability but the sequence of time of occurence of each candidate intermediate event to which the setting gives no meaning
 
 Instead of collapsing occurences by values (abxy) we say that a real run has a stamp (time-space or other distinguisihing the ones we want to add and the existence of referee induces an implicit order. so we have batches that can be viewed as Matrcesor injectuve  sequences $(1\mapsto (i,j)),(2\mapsto (i',jj')\ldots )$ where $i$ is $(a|x)$ and $j$ is $(b|y)$  (so each batch have length at most 16).
 
 We are interested in sequences that fail to have direct experimentally consistent empiracal meaning so if we have a non signaling distrib we replace it with the sequence of expermients that almost look like it but minimally altered to not fit .
 
 So an event is either a non signaling events with the paths simplying ot or extend it to be non signaling, or a sequence of allminmal and maximal alteration of a non signling experiment.
 
 We consider here as inetersteing the sequence of consecutive events ()
 
   
\section{draft idées non analysées, foyer de germination formelle}
les quatre traits  | | | | peuvent etres vus comme 4 fleches pointant sur prefixe début à position en cours puis un suffixe début à position en cours légitimant ou rajoutant une défnition au || du milieu ou au mot ecrite entre les deux barre $|x|w|y|$  serait en fait en mode dolores avec troncation de q et p lu avec un ecart d'au plus un bloc de taille $._t$
 
 $\begin {matrix} 
  |&|&|&|\\
q_1& last_q+1_t& p_1 & last_p+1_t\\
\end{matrix}$
l'orientation de q et p doit supporter les 4 inversions

on aurait un lecture préfixe, une lecture suffixe et deux lectures alternées de dimension finies  l'une commencant a droite l'autre a gauche (lire d'une manier en flippant ceux a meme distance du miroir central devrait donner l'autre]   qui seraient toutes les 4 comptaible et acceptant les interruptions.


 peut etre ? comme $w_{\delta ad |ad}^{v-\partial v}$  avec ad sous $\delta ad$  puis on monte et va vers la droite on trouve $v$ on ca vers la droite et on monte $\partial v$


imaginer 5 bit string on a un bit majorite de 1 :  $| \mapsto  \uparrow$  et  $| \mapsto  \downarrow$  [comme si c'etait la direction ou on cherchait la ligne de fin] ensuite  : 5 choix de mots de bord avec une position différtente des 4 autres : avec l'inversion (plus de 1 plus de 0)  et puis  9 choix de mots proches vers la barrière centrale (3 1 et 2 0 ) ou l'inverse paires par complémentaire.
a chaque fois la correction interpretative rechechée est la superposition des 9 les mots avec 2 1  ou la superposition  des 5 mots avec 4 1
ou les memes avec 0 a la place de 1.

une entité sémantique minimale serait donc 10 bits avec 2 ou 4 1 dans chacune des deux moitiés donc donc 4,6 (plus a gauche) ,6 (plus a droite),ou 8 fleches justificatives.

\section{old notes pour moi meme to pick from}
A ceux qui disent ce sont les mathématiques qui expliquent la physique qui explique le réel.
Je me dis juste que je ne suis pas convaincu qu'on puisse toujours clamer l'indépendance et fermer un système en oubliant ce qu'il se refuse de voir
et de déduire après en cherchant ma perfection et vouloir clamer l'universalisme des valeurs induites sortant des accolades où ils ont bati leurs nombres des conséquences submissives à des personnes qui ne peuvent ou ne veulent y adhérer.

Je veux donc essayer ici de comprendre sans oublier l'Histoire et je prétends sincèrement que je ne vois pas pourquoi on ne le  pourrait pas (et  c'est ce que je vais essayer ici)  essayer de  couvrir l'émergence du sens capturable par des  langages au moins ultimement partiellement  symbolisables  sans avoir recours à l'abstraction y compris pour comprendre les mathématiques.

Ceci permettra aussi d'essayer de comprendre les différentes manière d'entrelacer dans une dynamique le sentiment d'être quand on pense pouvoir montrer des traces de soi du passé avec notre confiance dans la croyance qu'on a de savoir ce qui ne pourrait pas advenir.

Je m'intéresse donc à comprendre quand il est ultimement tracable avec des traces directes, des traces induites, et des traces restreintes et des traces  comme
 voulant tracer une méthode de représentation pour les ontologies plausibles consistantes avec les traces locales visibles et faisant émaner localement  une croyance (ou connaissance)  consistante avec les interactions imaginant l'émergence  
du sens quand il est transmissible et enseignable.

Un phénomène fascinant qui semble émerger dans cette nécessité du deux pour concevoir l'un dans les concepts associés à des percepts  par un anthropomorphisme dans une communauté est que cette pairabilité dans la  confiance en la distinction ou la conviction en l'impossiblité de l'établir  la (r)assurance quand portée sur une  trace estimée par celui qui parle pour tout ceux qu'il imagine va imposer pour le sens
 communautariste de libérer quelque chose quand on veut forcer ou rigidifier une autre (en la réduisant à une séquence de représentations de traces compatible avec le locus). C'est cette interopérabilité opératoire  qui permet faire sens aux aux interactions entre le communautés avec leur signifiance locale identitaire séparatrice dans une signifiance globale unificatrice.  
  
  Vous qui vous aventurez a vouloir lire la suite gardez bien votre espérence et votre sens critique,... ou pas, car ce n'est que la suggestion de mon expérience pour proposer une manière de la déchiffrer.
  
\section{La posture d'expression libérable dans la conception  communautariste de la définition}
Adopter la posture d'expression  libérable dans la conception de la définition c'est ne pas oublier le messager en parlant du message.
Reconnaitre l'intentionnalité plausible ce n'est pas en faire un émissaire. C'est se reconnaitre comme une certaine conception du temps dans la séquence de pensées traçables qui nous traverse.

Je passerai à l’anglais pour la partie mathématique, non par préférence stylistique, mais parce que l’anglais scientifique contemporain agit comme une contrainte formelle : il rigidifie la forme et force l’explicitation de l’implicite. Cette contrainte est ici utilisée comme un outil, non comme un idéal. %\footnote{
 Je me sens pas   maintenant   capable de faire une traduction avec  une fidélité dont j'ai l'impression maintenant qu'elle  me paraitra satisfaisante comme trace de la pensée qui me traverse, mais  c'est que dans le déshabillage de l'implicite, ca semble une manière naturelle pour pouvoir communiquer avant de finir, avant d'arriver au trait perceptible comme multiple dans sa distinction  transmissible avec imperfections et interruptions.
 
  Ma remarque qui peut sembler négative ne parle pas de  l'anglais de Sheakspeare ou de Poe ce n'est pas la quête du beau en étalant son âme, c'est plutôt une étape dans un exercice Oulpiste de contrainte pour faire émerger l'expression du fond par rigidificatoin de la forme. C'est l'anglais moderne tel qu'on me l'a enseigné dans les parcours scientifiques.  Cet étalage lexical normalisant, propre à réduire la profondeur des restrictions pour optimiser (au sens de minimiser la quantité) les échanges et qui me semble jouer un role non négligeable dans le rapprochement  que je ressens comme ayant eu lieu entre  l'Homme et machines spécialisées (en prenant ma connaissance de l'Histoire vers aujourd'hui comme direction de locus) et à qui je prête aujourd'hui un rôle dans l'éloignement entre les deux (en prenant le présent ou j'ai l'impression d'être  vers le futur  comme  direction de locus)). Cela permettra de produire un objet propre ou acceptable pour les sciences occidentales: nettoyé du sang et des larmes pourtant inhérentes à sa naissance.

Le chemin du retour qui commencera par faire émerger un coeur au trait qu'on a devant les yeux  avec la légèreté d'un "pourquoi pas?" ou du plausible, de l'imaginaire à la place du contre-factuel. 

Avec pour l'instant ma conviction de  faire apparaitre  une méthode claire et propre à chacun mais assez riche pour interagir par conventions explicites, de voir les sens communautaristes dans leur multitudes émerger dans les communautés que j'imagine échanger le plus fidèlement qu'il m'est possible à la trace que je crois voir de leurs discours en m'interdisant de fermer les yeux sur l'ombre de ce que j'aurais pu oublier.  
Ceci en n'ayant besoin que  d 'équiper toute communauté de postures communautariste tracables de manière consistante avec 
ce qu'on a l'impression de comprendre d':
\begin{itemize}
\item Les explicites ou d'abord principes  rigides de pensée:
\begin{itemize}
\item \textbf{  |: un besoin de faire part de} ou de se sentir comme partie d'une communauté qui permettra à Qudsia inductivement de vivre ayant le droit à l''équivalence, à la \textbf{constance} à la \textbf{composition} inductive des atomes comme sommets pour former des structure dans cette cohomologique de  plus petits points fixes ou constantes ou atomes manipulable dans les règles communautaristes de composition comme plus petit commun multiple. C'est la notion de \textbf{ Concept} ou d' \textbf{unification} comme d'abord manière rigide de pensée prémarquée par des traces de percept ou plutôt de changement de percept. C'est tout est un et un est tout qui peut faire penser à la profession de foi islamique qu'on va retrouver dans le "detailed balance" de la physique c'est ce rescaling mimétique ou l'on se retrouve homothétiquement dans ce qu'on voit.  C'est la tentative d'interdiction qui génère la paire puisque quand on saura écrire on pourra écrire dessus ou dessous.
  \item \textbf{-: un désir de se dépasser} qui clame avoir ou posséder comme partie à qui veut l'entendre qui va faire émerger cette impression d'uniforme 
  et faire jaillir l'aléatoire par ses mécanismes homologiques de restriction de projection du multiple comme \textbf{  décomposition} ou de plus grand point fixe manipulable   dans les règles communautaristes de composition comme plus grand commun diviseur. C'est la notion de \textbf{raffinement} ou de  \textbf{percept} qui n'a de traces que par le fait des représentations qui tentent (bien souvent) vainement de le saisir. C'est le tout a un mais un par un, comme se posant comme distinction entre hier et demain. C'est la trace visible de l'enseignement qui cache le multiple.
  \end{itemize}
  \item $.,[\leftarrow_d, ,\leftarrow_{dp}],[\downarrow_q,\downarrow_{dq}],[\rightarrow_p,\rightarrow_{pb}],[\uparrow_{b},\uparrow_{qb}]$ et toutes les complétions que vous imaginez que j'imagine: Un  insaisissable par la trace directe   \textbf {alternateur synchroniseur biaisé imagineur de liens causaux implicites prés-symboliques }. Ils seront  d'abord anthropomorphismes qui ayant accès à des concepts intériorisés mais inaccessible à l'écriture directe essayent chancun de son locus  de déchiffrer des instances de ce que son imagination estime possible dans ce qu'il perçoit va alterner guidé par ses préconcepts   entre la lecture et l'action d'écrire ou d'agir qui de par son fait génère des percept futurs. C'est la tentative de son explicitation comme ce qui induit le sens au sens de signification qui est le coeur de l'esprit de la démarche de sens communautariste et donc libérable pour celui qui le conçoit comme tel. On va le rendre visible par les flèches ou les  points mais en ayant conscience qu'il va disparaitre dans la forme linéarisé du langage écrit habituel: c'est ouvreur de possibilités nous autorisant  à modifier en route ou à rendre acte de ce quon perçoit comme potentialités de de changement plausibles de la topologie d'ou on est, d'ou on vient  et vers ou on pense aller. L'intention est toujours multiple le . est un locus de dimension non mesurable par le narrateur toute autre dimension est d'bord ressentie comme paire indinstingable qui pourra toujours changer d'avis jusqu'à la cristalisation du sens exprimable. C'est le tout est différent dans le recommencement : comment alterner ces deux visions dans une continuité normative. 
  \end{itemize}


Dans les langages symbolisés où ce qui est tracé porte une incertitude intrinsèque sur la manière dont cela a été tracé, ce qui aurait pu être est structurellement différent de ce qui peut être vu. Le sens émerge alors de la tension entre la mémoire des gestes appris et les interprétations admissibles, inacceptables ou indécidables, selon l’histoire dans laquelle nous sommes situés.

 
 quand on apprend de nouveaux kanji pour écrire la dimension temporelle de l'expression devient flagrante dans la forme qui contient du coup des marqueurs de temporalités quand les écrits de quelqu'un s'étalent sur une ligne, et la pensée anachronique de sa prédiction du futur se doit d'intégrer l'imaginaire de la forme qu'il ne connait pas encore.
%Un japonais apprend de nouveaux Kanji, il faudrait donc en fonction de ce que je connais justifier pourquoi je peux me permettre le premier choix et revoir ma pensée demain s'écrirait nécessairement . 
En théorie de l'information quantique,   une moitié de paire intriquée apparaitra à qui le demande comme une variable aléatoire indépendante  il faudra donc que je justifie pourquoi mon espace de définition était le bon pour prendre le deuxième.

Mais sommes nous obligés de choisir? ou du moins de le faire maintenant ?



  Pour une communauté, à toute  séquence de symboles vue comme manière rigide on associe un agent ou anthropomorphisme qui peut être vu comme un agent , et une manière d'interagir (d'encoder).   Dans une communauté, chacun croit maîtriser l’implicite — non par effort, mais parce que le centre de coordination est transparent.
 On peut voir le  degré de raffinabilité comme le niveau auquel les communications peuvent être étendues sans contradiction.
Si l’on interrompt les communications à un certain degré de raffinabilité, chaque membre construit un monde cohérent uniquement à partir des traces disponibles au moment de l’arrêt.
Ces constructions sont :
\begin{itemize}
\item locales : elles n’utilisent que l’histoire figée accessible à l’agent ; 
\item résilientes : toute explication ultérieure compatible avec cette trace est absorbée sans contradiction ;
\item indéfiniment continuables mais instantatnément interruptibles: le sens ne s’arrête pas, même sans nouvel apport d’information.
\end{itemize}

La localité peut être vue comme la manière de théâtraliser le discours qu'on ne peut supposer à priori la scène fermée, et qu'on se doit de considérer la possibilité de casser le quatrième mur et d'imaginer celui qui écoute comme part du spectacle. Se dire qu'à priori  on ne peut supposer  sans risque d'erreur que notre représentation du spectateur est identique à ce qu'il est, mais on peut faire l'effort à chaque interaction  d'affiner notre modèle de représentation de cet auditeur environnement de notre premier mot.

La résilience permet de pouvoir répondre au pourquoi et impose de devoir justifier l'effort de changer d'avis. 

La non terminaison de la méthode c'est entrevoie l'émergence d'un implicite caché ou du moins expliciter notre estimations de sa plausibilité et pouvoir la justifier.

L'interruptibilité c'est admettre que je ne pourrai faire un raisonnement si on me coupe au milieu et que le controle qu'on a forcer l'environnement d'écouter doit être explicitable ou justifibale

Cependant, un exemple célèbre est la transmission de la causalité présymbolique dans des communautés  à travers les générations des singes de la douche et de la banane ou meme s aucun singe vivant n'a bu la douche se déclencher, ma si la douche ne marche plus les singes ages continuent a enseigner en laissant un trace dans leur corps ( par les coups) la regle ici on ne touche pas a ces bananes. La corrélation pour les expérimentatuers a juste besoin d'avoir existe et a est maintenant comme dna s le pgcd decorrelles de la coorrelation pour ka communaute des singes qui les ont internalises. 


Le psychiatre ne gagne rien avec des bouts de patient mais l'histoirien peut essayer de raconter une histoire en collectionnant les traces.



En informatique quantique  les états intriqués montrent qu'il est des expériences ou simplifier pour comprendre est juste un abération.
Si on mesure un qubit d'une paire imaginaire reproductible sans attendre d'avoir recu le second parce qu'on aurait envie de le définir dans un espace compaible avec les experiences connues a ce jour. \footnote{Anachroniquement, dans une deuxième lecture Perséponen dirait Eureka j'ai découvert l'aléa.}  
Cependant,  si ca se répète et que je reçois les second  et que je crois que c'est les même. Je pourrais maintenat répondre à ce que j'aurais vu hier, voire commencer a fabriquer des systèmes pour rattraper le retard perdu.

L'éducation voudrait que les pharses commencent par 
$ \begin{matrix} \blue{\downarrow \rightarrow } &\text {"s'il vous plait"}&\rightarrow \\
&\text{" }| \rightarrow \text{"}&\uparrow  \end{matrix}$ 
et finissent par $ \begin{matrix} \blue{\downarrow \rightarrow } &\text {"merci"}&\rightarrow \\
&\text{"}  \downarrow \text{"}&\uparrow  \end{matrix}$. 

Tous les chemins  dessinés sont bons à prendre par le lecteur même ceux qui ne suivent pas les flèches. Une manière bien mathématique de batir une communauté est de  dire je refuse de parler à ceux qui ne lisent pas moins  de 100000 mots-minute même si moi j'en suis à deux (s'adresser seuement à ceux qui ne parlent pas plus de 100 mots valable est aussi un typage mathématique  naturel en fonction du gout (projectif,inductif) ou l'inverse pour la lecture de la séquence de types), mais le principe communautariste est de faire tout son possible pour ne pas exclure tant que possible pour permettre aux communautés différentes qui ne se comprennent pas parfaitement d'interagir au mieux sans avoir ce bsoin de fermeture du principe de la compréhension. A mon niveau de granularité mes atomes sont protégés par les guillemets les possibles lectures 
sont indiquées par les flèches, et la réalité plausible de fermeture par symétrie que je me dois d'imaginer c'est toutes les ballades qui ignorent les conseils directionnels. Je ne peux pas imposer que la séquence de types est la même chez tout le monde je dois juste essayer de fabriquer des traductions bi-univoques qui me permettent d'interpréter le plus fidèlement possible la pensée de ce qui me parait autre en acceptant la possibilité d'avoir mal compris mais en essayant de la quantifiée par la transformer en le un permettant de quantifier mon assurance sur ce que je crois savoir. 

J'ai envie d'essayer ici deux démarche de comprendre au del du respect de ce à quoi je m'attends : 
\begin{itemize}
 \item |:Compose: comment rajouter a gauche/droite pour que ca fasse senns (prémisses du plus petit commun multiple ou je peux maintenant repeter mon objet regiularise comme un monoide interruptible) Le plus petit point fixe est aussi la manière la plus inclusive de la qualle je me le concois se concevoir 
\item -:Decompose: en fonction de ce que je connais aujourdhui produire plusieures copies et selectionner le plus grand bout plus court qui me parait acceptable 
premisses du plus petit diviseur. Présuppose l'existence d'un espace et le plus petit espace dans la quelle ca fa fait un sens communautariste est la paire
\end{itemize}


Et contrairement aux ballades habituelles je n'effacerais pas les aller-retour  car les distances changent quand les gens construisent des ponts le plus petit chemin n'est pas toujours celui qui l'était hier. On apprend en tombant et se relevant et je pense sincèrement que bien souvent il y a beaucoup  de choses à apprendre ses aller retours. 
Un exemple  de langage



Construction au sein d'une communauté d'un imaginaire commun permettant l'implicite et induisant des procédures d'identifcation par association satisfaisante pour ses membres permettant de séparer les membres , ceux qui ne peuvent pas l'être et ceux dont on ne sais pas. De sorte à ce que que membres qui se reconnaissent comme tel peuvent établir une micro-logique consistant a agréer sur la causalité la transportant de la représentation qu'en en a chacun 
à ce qu'ils considèrent comme définition  commune  à la communauté de réalité.   

La simplification du \emph{Modus.Ponens} pour éviter la préexistence du terme à déduire dans une forme plus complexe prorpre à la réflexion projective rend souvent implicite, en imaginant l'identité, le problème de reconnaissance dans une représentation causale induite qui sera construite dans l'autre sens. 
Elle est trompeuse car le glissement sémantique fait croire à ceux qui ne se rappellent pas des prémisses que le faux peut exister univesellement. 
La négation n'a de sens que dans un cadre et par conséquent la démarche communautariste prend acte de ce fait explicitement la pharse "il est faux" ne peut être un début car elle pose un existant sur un innomable tout propos communautarisé doit être un queelqusoit définissant la restriction ou l'on espère expliquer quelque chose ($\Pi$ langages). 


Mais j'ai l'impression  que rendre rigoureusement inverse les mécanismes de décomposition n'est pas une nécessité inéluctable et dans la quête de ne pas oublier ceux qui comprennent autrement en essayant sincèrement de les comprendre devient une posture communautaire forte qu'on va éviter de prendre trop tôt..

L'histoire a montré que les hommes ontt fabriqué les outils pour fabriquer et transmettre leur compréhension mais si relire dans l'autre sens la construction d'espace à partir de lettre-atomes  mimique très bien notre imagination de la construction d'espaces par restriction d'univers, il ya cette subtile différence du besoin de un nouvel espace qui fait qu'il ya bien deux démarches distinctes  même si je ne peux sincèrement les voir que comme interchangeable: l'ordre qui les relie est subjectif c'est la première qui m'est apparue: Le déterminisme et l'aléa me semblent êtrre deux faces d'une même pièce que nous partageons et pour qui chacun a un état initial. Accorder nos violons en donnant un nom commun pour éviter d'intervertir l'un et l'autre est un acte communautaire qui n'a pas besoin dans une pensée communautariste de s'interdire d'échanger avec ceux qui pensent autrement c'est seuement la quepeut naitre la négation, le swap ou l'inversion comme outil de dialogue intercommunautaire.  


 La clé pour nous ici sera que de ne pas pouvoir savoir (l' indécidabilité)  va émerger naturellement comme porteur d'information et la certification de non prouvabilité sous contrainte apparait naturellement comme le domaine de validité de la signification prouvable sous reserve de validité de la certifcation etc mais c'est un processus qui est spirale dans bien des cas et garantir des bornes peut etre naturel dans certaines circonstances. \footnote{Par exemple pour la communauté humaine on peut supposer que la vitesse de pensée peut induire une borne sur la séarrabailité  pour la défniton de l'espace de Dolorès  pour lancer $b_q$ et que la dégradtion des cellules ou de la planète peut doinner des borne sur le temps pour lancer $b_p$}
 Pour fabriquer le sens communautariste dans sa diversité on va  se réduire  à des phénomèmes d'encoches sur 4 bords qui semblent faire sens à toute communauté vue comme processus communiquant et se faisant un avis via mémoire et messages échangés. Ce phénomène peut être vue comme générant par mitigation ou concertation imaginaire 
 un sens global associatif avec les restrictions (parenthèses sous formes de crochet obtenues par le processus à encoche unaires considéré comme élémentaire par l'exigence de produire une sortie mais retardé tant que possible dans la description de transformations ou de mutations).
 
 Rigidifier la compréhension par un langage à alphabet connu sera cristallisé le processus pour obtenir un premier cadre communautaire en limitant les potentialités de sortie du cadre de la compréhension  à des chemins linéaires sur des nombres communs.
 
  On "collapse" ou fait s'effondrer la force de la réécriture interprétatitve el limitant l'associativité à une lecture par défaut sur des termes identifiés 
  générant un sens communautaire pour ceux qui savent deviner comment lire ce qui n'est pas forcé par les parenthèses carrées.
  Checher un sens communautariste c'est  
d'éviter   le collapse complet.  
Une  manière de maintenir la diversité  est de fabriquer des paires (constante ,variable) ou on dira que tout le monde qui reste dans le cadre doit assumer qu'il  sait reconnaitre les constantes, les distinguer et les réécrire, mais  en contrepartie à chaque constante je dois associer un objet imaginaire qui peut représenter une zone mémoire pour la retenir ou autre et que je n'arrive pas à définir car à chaque fois que je crois le contenir il y a un temps futur ou je détecte un changement c'est l'intentionnalité $\downarrow/\uparrow$ insaisissable qui se cache derrière le pouvoir de définir chaque constante l'inversion du réel et de l'imaginaire garantit la diversité après la décision de stratégie de lecture. 
On peut justifier le collapse par au moins deux raisons : l'une est pragmatique c'est une question d'encodage et de finitude quand on arrêe une chaine de raffinement le plus simple est de remplacer la qualité par une quantité et le nombre qui le nombre de composantes loczlisée au quel on rajoute  celle infinie dont on ne peut rien dire est cette compression parfaite : on a juste besoin de compter et on introduit le concept de commutatitvité on ne peut que faire comme si on n'avait pas besoin de connaitre la séquence de va et vient élastique mais juste la différence entre les pas qui s' éloignent e ceux qui se rapprochent. 
La deuxième  est une forme d'arrogance ou  désir de perfection qui fait croire a notre (oeil, cerveau) qu'il existe des objets réguliers imaginaies comme le cercle ou pire traits continus diagonaux, alors qu'on ne les voit pas on les imagine, L'(oeil,cerveau) est d'abord un code correcteur d'erreur encodé dans nos gènes et pour comprendre comprendre, il faut séparer la partie discrète qu'on voit vraiment qui est une optique de réflexion entre des parois élastique et la partie trabsmise qui est électrique qui est une séquence d'impulsions dans non nerfs.
  
Nous ne présentons pas ici un langage mais une espérance de langage en devenir  prenant corps  différemment dans les espaces  pouvant la féconder  devenant signifiée localement à tout instant dans les locus presque fermés mais qui vit évolue et s'adapte au gré des interactions qu'elle essaye de repérsenter dans cette course incessante à noter le passé sans trop ce laisser dépasser par le présent tout en espérant un jour avoir un choix pour l'avenir. 

Mais pour montrer l'utilité du plausible, nous essayerons de donner deux corps indidistingables à la demi-femme parfaite  en s'autorisant  à noter $+$ pour le losange et $X$ pour le carré, contractant en un point ce qu'on ne connait pas,
et on va essayer (en l'état de la note) de  montrer comment si on croire ou de définir comme borne du savoir qu'on pourrait avoir accès à une chaine infinie de variables initialement parfaites on fait la communauté de ceux qui connaissent et savent fabriquer l'aléa pour retrouver  ce qu'on a l'habitude de voir dans les livres.   avecles explication de soustructures non triviales de   ($ +,\oplus,x,\otimes)$ . Nous allons essayer de reconstruire des communatés de sens incluant toutes les les théories transmissibles connues. C'est amusant de penser qu'on va penser avec le poids du patriarchat malgré la plus grande bienveillance et que la première image qui vient est que l'homme est la femme parfaite à laquelle il manque un bras et que  pour l'instant j'ai l'impression que  ne vais pouvoir  réussir à faire sens  de (avec des dimension bornables) qu' à tout ce qui a au moins un bras sur les 7 possible des formes imparfaites à qui il manque quelque chose.

La mise en garde qui s'impose est que la démarche communautariste  demande la remise en cause et les liens étant nécessairement tout ce qu'on peut sincèrement imaginer comme superposition avant de poser le nouveau cadre restrictif, on est forcé à chaque grande étape de zoom in ou de zoom tout d'accpeter la possibilité que non pas le propos mais la cohérence de toute notre pensée soit fausse, il faut accepter de sincèrement examiner la possibilité d'avoir tout faux depuis le début et que la mort peut arriver demain nous figeant dans cet état de rien alors qu'on avait cru avoir tant fait mais  le courage poétique de se dire malgré tout j'aime tellement construire que l'aventure était belle est un moyen de sauter le pas et de découvrir la reconstruction sans oublier  mais en choisissant ses attaches pour avancer comme on sent qu'on le doit plutot que comme on nous a dit qu'on doit. 

L'ojectivité est l'expression limite d'une subjectivité qui se cache ou alors elle fait partie d'une communauté qui me dépasse.  

Cette tentative d'embrasser tous les communautarismes et ceux dont j'ai juste conscience que je n'ai pas encore résussi à les imaginer, qui peut être vue comme donner le  sens fini qui me convient le mieux  et que je pourrais expliquer à mon fils d 'union infinie de tous les voisinages définissables dans tout système de définition dont je pense sincèrement que je  pourrais un jour comprendre les règles. 

Elle requiert  de remonter les traces de traces  pour essayer de comprendre l'intentionnalité derrière les traits transmissibles par proto-langages ayant permis l'épopée sumérienne. L'infinissabilité de la démarche offre la liberté de la faire à partir de ce qui me parait être rien pour celui qui ne voudrait me prendre pour objet d'étude: je peux commencer à me raconter les choses à partir de ce que j'imagine je n'ai pas besoin d'arrêter ma vie pour chercher les traces l'énergie que j'y m'est peut être juste des temps de divagations avec mémoire. C'est la démarche de dépliage qui supporte la remise en cause que je veux mettre en valeur dans ce paragraphe plus qu'étaler des connaissance sur le langage à un lecteur quie j'imaginerai dire amen pour augmenter mon estime de moi. C'est de ce dire que chacun peut construire et interagir et que le n'importe quoi est souvent ce dont on n'a pas les clés de décodage de la communauté qui l'utilise.  

Dans une pensée libérable, répéter tous les jours la même chose en espérant quelque chose de différent est quelquechose qui pourrait fonctionne pour un géant portant le monde le long d'une longue pente quand ses pas tracent de manère impercetible un sillon qui sera une route. 

On voit donc que pour avoir une chance d'avoir une finitude raisonnable pour exprimer ces restrictions d'ordinaire omises on doit réduire l'expressivité à son éponyme : le trait.

\subsection{p-d commme PAC-CAP perdues entre deux infinis normatifs}
%$$
%\begin{matrix}
%&-&-&-&-&-&-&-&-&\
%&|&&&|&&&|&\\&-&&-&-&-&-&-\\
%|&&|y&&&&b|&|\\
%|&&&&&&&&&|\\
%&-&&&&&&-&\\
%|&|a&&&&&&z|&|
%\end{matrix}
%$$

\

\section{Justification}
Je vois deux manières simples de formaliser la perfection mais les deux ont un présupposé de croyance (ou du pouvoir  définition non révisable)
\begin{itemize} 
\item Il faut savoir avoir accès à la constance perpétuelle qu'on peut visualiser par accès parfait a Bénédicte via Qudsia. C'est la mémoire  parfaite infinie dans le temps : je peux changer de fréquence d'horloge ou de chaussures j'ai une zone bien délimitée ou à chaque fois que je demanderais une valeur et que j'écrirai ce que je vois ni moi ni mes descendants ne feront une différence entre deux symboles qui se suivent dans le cahier qu'on leur tend. La garantie est un alphabet fini  séparable constant et le raisonnement qui arrive à moi est inductif il est du type je me vois comme part de et le phénomène bien interporéatable est le plus petit point fixe porté par les déformations optimales dans un sens ou dans un autre pour l'obtenir.
\item  soit dépasser le Styx et revenir c'est à dire avoir accès à un infini permettant de transformer la multiplicité en probabilité c'ets à dire un accès a Bénédicte via Perséphone, on va avoir dans ce cas gratuitement ce qu'on appelle des variables aléatoires indépendantes on va avoir un consensus sur un seul espace d'ou on vient et toute description s'explique via des projections le raisonnement est projectif et le phénomène stabilisé le plus grand point fixe.
\end{itemize}




 

\section{Le pourquoi des mathématiques comme un conte fidèle:  une histoire de ressentis et de désirs}
Pour retranscrire fidèlement les mathématiques comme communtaué dans une vision communautariste
on peut voir l'Humain comme Dolorès ressentant être quelque chose dont il essaye de trouver des traces en fouillant l'espace pour chercher Qudsia.
De cette recherche revécue à l'envers  comme représentation il prend une  assurance   sur le fait d'être quelquechose et regarde le futrur dans le temps ou Perséphone semble lui montrer sa mort. Il voit la réalité par le prisme de sa réalisation et désire avoir du temps dans ce concept qu'il s'est fabriqué. Dans ce futur Perséphone pointe sur sa mort.


Les (tentatives de) symétrisation interruptibles, si elles ont  le temps de finir, feront des équivalence (elles auront aussi besoin d'espace) : construction de systèemes rigides de singularités

Les (tentatives d') équivalence interruptibles, si elles ont  d l'espace de s'écrire, feront des symétrie (elle auront aussi besoin de temps).

Dolorès: un anthropomorpohisme c'est d'abord tous les liens entre son Histoire écrite en train d'être tracée qui est compatible avec  qu'il ressent (degré de plausibilité) et ce qu'il va faire (s'il en prend conscience (degré d'assurance) sur le chemin de sa mort.
 
 rechercher ce qu'on ressens ou rechercher ce qu'on veut faire prend du temps et de l'espace et ce ce mécanisme de synchronisation permanent d'yun système de valeur qui émerge comme phénomène décriivant au mieux cette communauté qu'il représente pour ce lui qui en parle.  


Perséphone (telle que je m'imagine l'imaginer juste avant ma mort): Moi j'aurai appelé cet essai Dieu, le sens et les deux autres.

Le  1 noté | comme trace visibile sur l'axe du temps  est par son inverse le début des unités de raisonnement
c'est cette séparation ressentie quand sa trace duale dans l'espace - montre a qudsia qu'elle existe elle voudra alors revendiquer de l'espace pour respirer 
et les larmes de dolorès ressortent en regardant cette spirale infernale. 
Spirale qui repart et recommence.

Dans un esprit communautariste il ne peut y avoir de pseudo-indéterminée (connnaissance qu'on ressent comme presque nulle) sans déterminée (connnaissance qu'on ressent comme presque fiable)

L'informaticien est le plus près de la vision communautariste du réel il a déja devant les yeux (même s'il essaye très fort de ne pas voir)
les outils qui suffisent à une vision communautariste. 

Il va donner des nome  à des adresses et cacher à l'utilisateur haut niveau la taille de la structure sous-jacente 
et permettre a des interpretes qui demndent bout par bout des informations partielles de fabriquer une valeur et de l'afficher.

Le message on n'a pas besoin de comprendre pour  utiliser cache le si on fait confiance et aussi si on oublie les accidents.


Je ne veux pas supposer la perfection des constantes parce que Dolorès la lie à l'histoire de l'orgine de l'espace et l'Histoire des humains est  telle que je ne me ressens pas avoir le droit vis à vis de moi même d'oublier ou de fermer les yeux sur ce dont j'entrevois la possiblité de  douter de peur de me retrouver à fabriquer une mèche et de la poudre dont la l'assembalge est indéfinissable à mon niveau de "clearance". 


Quel est le lien avec sa descendance ou son ascendance et quel est le rapport entre les deux 


Si je ne me sens pas être tout je me dois de représenter ce que les libertés que je prends imposent à la vie des autres: ce que cette arrogance du "je peux définir" revendique sur le droit de puissance à soumettre l'imaginaire futur potentiel de la vie de la feuille sur laquelle l'acte s'accomplit. 
A la soumission par la force de l'argument d'autorité je préfère essayer de comprendre en se refusant à fermer les yeux sur les traces qu'on imagine et qu'on se force régulièrement à essayer de réimaginer. 
Il me semble donc d'usage de commencer par demander la permission pour avoir un cadre explicite avant de prétendre à un sens que je m'autorise à chercher. 

Mais comment ne pas se demander   si la dernière trace que j' imagine  percevoir dans la justification de ce que je concois  n'était pas elle même trace d'autre chose ? 

Face à mon 
$\begin{matrix} \blue{\downarrow \rightarrow}  & 
\text{impuissance à pouvoir sans laisser de trace} &\\
-&\text{désir de vouloir}& \uparrow
\end{matrix}$				me convaincre que je puisse faire de mon mieux pour traiter les explications supportées par des   chemins de 
%$\begin{matrix} \red{\downarrow \rightarrow}  &\\
%\text{}& \uparrow
%\end{matrix}$
% \text{de traces} &\red{\downarrow \rightarrow}&\\				\uparrow&-&|\end{matrix}$ ] 
$\begin{matrix} \text{de trace} & \text{de traces} &\red{\downarrow \rightarrow}&\\
				\uparrow&-&|\end{matrix}$.
 
 :
   C'est en quantifiant (mon imperfection)  ma conviction dans les limtes de ma réflexion que j'arrive à borner pour les gens de ma communauté qui gagne de ce fait un nom : ceux qui ont besoin d'un temps constant minimal pour imaginer une décomposition.
   
    On ne va pas céder à l'illusion que le  [simplifier pour comprendre] pourrait être parfait pour tout sujet abordé par n'importe qui. C'est juste une méthode comme une autre c'est la mère des vielles mathématiques qui n'on qu'a supposer l'univers pour arriver par projection à produire des espaces et qui amènent à voir l'homme au milieu d'une chaine de simplifications vers des espaces isotropiques obtenus par projection, c'est ce qu'on appelera la démarche homologique, la démarche projective, l'observation du multiple le plus grand commun diviseur ou simplement la décomposition.  

On pourait donc se restreindre aux langages symbolisables mais gardons pour l'instant l'idée qu'on puisse ne pas en avoir besoin, la trace ici meme si 
souvent representee par un symbole peut etre vu comme un processus élémentaire de vérification et reproduction.


%Dans ce qu'on voudrait être  [simplifier pour mieux comprendre]  don le sujet ici est [ce qu'on peut transmettre et ce dont on veut transmette la manière], 
%e me restreindrais ici  à ce qui laisse trace dans des symboles. 
 %Equipant par ce fait le reste de la pensée de l'intention d'origine du sujet d'un espace infini 
 
% l'inévitable imperfection de la méthode 
%puis que je m'autorise à symboliser.

 Toute explication plus riche devra essayer de se symboliser comme elle veut pour que je puisse échanger et je garderai à 'esprit qu'ayant fait un choix je me ferme à jamais (dans le cadre de cet essai) la conception qui le dépasse et qui ne serait pas normalisable avec la symbolisation.
 je reviendrais donc au trait vertical.
 
 Je veux avoir moi une transmissibilité symbolisable en mimiquant les gestes que j'imagine et la ou je pourrais me tromper.
Fabriquon un langage de la compréhe,sion transmissible de transmission enseignable en gagnant gratuite la transmissibilite par cette note si 
si j'ai l'impression qu'un anthropomorphisme (ce qui ressemble à un humain et dont le sens prendra sens quand la ressemblance découlera du cadre ou de la communauté ou je placerais le propos donc ici dans mla mesure ou j'arriverai à consevoir des critère de ressemblance demain qui me semblent raisonnables aujourdh'ui )   quelqu'un d'autre que celui qui l'a fait dit j'ai l'impression de comprendre quelque chose, la méthode sera validée.  
 
On va déja équiper de traces ce qu'on ne peut éviter dans la démarche
\begin{itemize}
\item | le se comporte comm un trait que je crois voir : un candidat à être un symbole
\item . le trait dont je crois voir la trace
\item - le trait que j'imagine la ou il me semble qu'il pourrait être
\end{itemize}
 
 Si je veux l'enseigner je dois associer un geste qui me semble adapté à la production du trait  par exemple de haut en bas $\downarrow$,
 mais je me dois d'imaginer la possibilité que quelqu'un le fasse dans l'autre sens $\uparrow$ le trait 
 équipé de sa manière de naitre au regard  sera donc naturelle multiple. 
 
 L'imagination fiat que je ne peux concevoir l'unicité de la manière mais que la manière fait naitre le nombre par le multiple alors même que j'ai juste voulu faire 
 un trait que je n'ai meme pas encore dessiné.
 
 
  Proposition de principes pour donner du sens aux  causalité enseignables précédent la symbolisation.
  



la je dois relire dp à l'envers et c'est l'aspect $p_td_{t+1}$  qui ressort
 
un trait qui tourne peut etre $\begin{matrix} |\quad \\ -\end{matrix}$ ou$\begin{matrix} \quad | \\ -\end{matrix}$ 

$\begin{matrix} 
	    | & & \vdots && |\\
 	       &-&\vdots &-&\\
	\cdots&\cdots&\vdots&\cdots \\
	 	       &-&\vdots &-&\\
	    | & & \vdots && |\\
	\end{matrix}$ 


un sommet est un polygone qui signore, 1 puis 2 sommet  2 puis 1 arete 

lire unepreuve à l'envers peut etre  une divagation qui occupe l'espace  mais qui prend sens quand revisitee a lenvers

le sens n'est pas symetrique mais la symetrie protege progressivement le sens des  erreur d'interpreteaions differentes de leimplicite 
car le processus d'explicitation va transformer les incertitudes (qui sont des bornes inferieurs des assurabce en connaissances qui sont les bornes superueres des croyances
  

  
  
  \subsection{dynamique}
  
  `quand on rebondit on deforme en laissant une trace de soi  et en meme temps repartir dans l'autre sens
  
  quand on diffracte  on se decompose


les traits sur une liane peuvent etre vu comme une lecture energétique la premiere marque qu'on voit est le point de repere 
l'alternateur nous fait explorer avec l'unite du jour  et le moyen de parcours du jours dont tout ce qu'on sait c'est qu'il existe une direction qui est sure de ne pas terminer 
(ou il n'ya plus de traces plus loin)

le multiple est la trace de la trace  en tenant compte de sa mémoire 

Si Dolores ne fait que parcourcourir benedicte transmet sa memoire du passe 
 et c'est la que le locus  explique  en fonction de la communaute le lien entre la lecture du jour et la sequence qui laprecede :
 
 en cumulatif  comme sens unaire ou definissant du nouveau
 ou apres decoupe on alternatif comme sens difernetile  on on on utilise les unites explorees comme potentiel de presence et se pose 
 
 
 definir par les espacemet successif les unites ( puis en face leur changement potentiel dans la derive)
 A refaire avec 4 trait par cote
% 
% $
% \begin{matrix}
% | &\text{copielocale de q}&\text{mots encodes pour p} &| \text{copie locale}&|&_\\ 
% -&                                    &                                          & \text{changement d'unite d'espace}\\
%  &\text{espacement faisant l unite d'espace}&&-\\
%  -&&&\text{copie locale}\\
%  &&&-\\
% |&\\ \text{copie localede b}&|&\text{changement d'unoite de temps}&-\\
% \end{matrix}
%$


traits dans la derive 
$\begin{matrix}
&||-||&\\
-&&-\\
-&&-\\
|&&|\\
-&&-\\
- &&-\\
&||-||&
\end{matrix}$

Dans cette  optique folle de tentative d'appréhension du sens en alternant entre sa globalité et sa localité, je garde le droit divin du mathématicien à définir mais je dpois spécifier explicitement le cadre  ou au moins référer explicitement à une communauté dont la réponse moyenne répondra )à ma place.
Je me dois de rester ouverts aux possible que je n'ai pas encore imaginé il ya toujours une dimension imaginaire inexplorée si on fixe le temps de la question.
PAr exmepl j'ai le drot parce que j'en ai envie d'imaginer un monde ou aucune cave ne pourrait empêcher la pensée de sortir, ou le pouvoir exécutif serait aux mains de l'assemblée générale d'un communauté de nations unies et ou la distance idéologique est systématiquement associée à la distance de localisation. Mais je paye le cout énergétique de devoir répondre à moi -même en le questionnant de temps en temps pour imaginer unchemin amenant à y aller et aussi à tous ceux qui diraient mais c'est impossible par ce que quelquechose, de vérifier sincèrement si la durée de vie de rêve n'est pas marquée par ce quelque chose à moins que j'imagine une réponse que je n'ai pas encore et qui fait que je n'ai pas envie d'interpréter d'autres arguments autres autrement que par les adjoindre tant que je n'ai pas réglé le premier blocage à la possibilité ultime de mon rêve.   
 
 $N_2+3H_2 \rightarrow 3 NH_3$  comme target a garedr en tête donc la comibainison/recombinaison énergivore
 
 
\subsection{Dolores sur la direction de qudsia }
 
 Pense que $-$ c'est une parenthèse ouvrante vue d'en haut $[$ avec trois temps au minimum de dessin $\leftarrow \downarrow,\rightarrow$  qui pouraient être fait "pas comme il faut " pour elle c'est à dire sans respecter les instructions du manuel et en prenant son temps  c'est à dire avec une durée cachée à chacun des traits.
 
 
 a chaque fois qu'on arrive à la position qu'on identifie comme ca correspond au précédent passage on projette sur l'axe toutes l'ntegrale si l'axe est cumulatif la difference avec ce quon avait mis  avant sur l'axes differentiel accompagnes de la multiplicite d'interpretations ( combine de copies du premier et combien de copies du deuxoieme)
 
 \subsection{la complexité avec les dérives:} 
 
 etant donne un machine , un generateur d'instance qui donne l'instance un instance de taille n en plus de f(n) 
 un generateur de réponse qui donne la reponse proposee a la specification en  $g(n)\ge(f(n)+1$ operations
 un verifieur de question et un valideur de de reponse sachant question
 et regarder la complexite differentielle $g(n)-f(n)$ et celle d'échelle $g(n)/f(n)$
  
  et définir le probleme pour le plus petit pourles machines qu'on imagine> le plus grand des machines explicite
  et de trouver quelque chise qui garde un sens en poussant $g$ vers la droite et $1/f$ vers le bas
  
  \subsection{Une bénédicte presque parfaite}
  si on dit que c'est faute si  dolores se comporte comme on le voit
   et qu'elle ne fait jamais d'erreur (pas besoin de corection localement definie dans les echanges)
   on aurait un langage comme -,
   %\begin{matrix} -|&\\-
    
    
  
  
  
 
 
 \subsection{La posture communitauraiste de comprendre en 2 temps qui en seront quatre}
 
 La posture communautariste pour aoppréhender l'émergence de la comprhésension trasmissible a d'abord deux temps et permet de  produire  une explication simpe  représentable avec des traits horizontaux et verticaux de longueur unitaire dessinable dans un plan des explications imaginables par quelqu'un que quelqu'un d'autre pourrait tracer et enseigner  avec une vérification partielle par la trace enseignable elle même par la trace.
 
\begin{itemize}
\item Le premier temps consiste à se dire que si on n'est pas trop pressé par le temps qu'on se représente, on ne perd rien à essayer d'expliciter  les implicites de notre communauté et à questionner les sous-jacents transmis par l'éducation, l'héritage, le mimétisme \ldots  auquel on n'a pas pris le temps de penser. 
 Repenser l'atome ou l'origine de l'espace c'est essayer si une nouvelle unité apporterait  quelquechose à notre système de valeur qu'on construit à partir de symboles identifiables à des locus constants ou correctibles pour convaincre un autre qu'on concoit  qu'elle le sont.
 \item Le deuxième temps consiste à dire que l'espace gagné par ce questionnement peut essayer de trouver des raffinement pour définir de nouveaux 
objets sophistiqués qui essayeront de décrire les qualités . On peut se demander à quel point on est sur que les constantes qu'on sont complètes et s'il n'y a pas un nouvel implicite auquel on n'a pas pensé. Une énumération peut être vue dans une autre mais la conviction de pouvoir retirer les '...' peut être au mieux quantifiable si on veut pouvoir remplacer dans les cadres qui le permettent tout par n'importe quoi sans perdre le sens défini en dehors de ces cadres.
\end{itemize}

Privilégier les comparaisons pour faire naitre la connaissance relative par assocuation de quantituatifs associables aux qualitatits sous réserve de prise de conscience des règles implicite du cadre ou on est.


La naissance des temps futurs est  construction cohomologhique inductive de l'être comme faisant partie de en partant des symboles qu'on avait construit lors de remises en causes que l'on ressent comme précédentes génèrera par plus petit comme multiple étendu qu'on verra ici comme le vent souffle sur une paire qui s'attire
ou une pré-construction d'ADN par rapprochemebt d'ARN.

La décomposition projective qui fait émerger les espaces futurs.


à la fin la couche ultime (qui permettra de porter la variable extérieur intérieur  qui pour définir une communauté qui nous unirait mais qui demanderait à être symétrisée ce qui introduira probablement de nouvelles normes plus fines qui nous génèreront de nouvelle symétries...)



\section{Graine de plus petit commun multiple}
$
\begin{matrix}
| \longrightarrow&&\downarrow&&\\
                  &-&-&-&- \\
                  &&&&\\
                   &-&-&-&- \\
| \rightarrow&&&\uparrow&&\\                
\end{matrix}
$
donne 
$
\begin{matrix}
| \longrightarrow&&&&\\
                  &-&-&.^{\longrightarrow} \downarrow &- \\
                  & & &=& \\
                   &-&-&.^\rightarrow \uparrow            &- \\
| \rightarrow& & &&&\\                
\end{matrix}
$

Description qualitative: Un coup de vent pousse les deux flèches verticales comme une accélération : quand leur écart horizontal n'est pas trop grand elle s'attirent déformant la barrière qui les sépare

Description quantitative: unite de vent, force vent haut, force vent bas, position initiale fleche 1, position initiale fleche 2

\section{The 3 strokes-view : Vertices as almost all or cadres  [ or  almost constants and holes  or almost nothing  $._d$|  as a splitting if the imaginary full cell and %$\begin{matrix} .p \\  - \end{matrix}$ 
 $\underline{._p}$ and |$._b$  and  $\overline{.q}$}
% $\begin{matrix} -  \\  .q \end{matrix}$ }

\section{ Canaux $._d$=$._p$ vs $._d$||$._q$,  $|^- ._d$ et sa version presque parfaite $-^|._b$,  de l'origine de lespace a la fin des temps
  $._q|_-$ à $-_|._p$ }

%\section{le calcul ou concept anthropomorphisme percept ou $|^-$ et sa version presque parfaite $-^|$ }

A quel point attendre c'est revenir en arriere

le H le U
 $-|-|---|-$ peut donner  des traits pris ou partages
$\begin{matrix}
|& &&&&|\\
=&=&=&=&=&=\\
&|&&&&|\\
\end{matrix}$



\section{projected borders}
A horizontally projected  border looks like  $\begin{matrix} .  &-& .& -  &.&.&.&-\end{matrix}$

a constant or a variable is a matter of relative clocks : in any sequence that I can copy on my 
memory a reprentative  subsequence that at least two diffrerent values 
and try to build a process reapt reading it 

comparing to this  constant memry clck

when i access memomry with a clock, if it is. :
\begin{itemize}
\item slow enough i have what looks like randomness or hard to guess  and 
\item  fast enough what looks like constants
\end{itemize}


\section{PUSH as approximate lcm seed}
$
\begin{matrix}
| \longrightarrow&&\downarrow&&\\
                  &-&-&-&- \\
                  &&&&\\
                   &-&-&-&- \\
| \rightarrow&&&\uparrow&&\\                
\end{matrix}
$
donne 
$
\begin{matrix}
| \longrightarrow&&&&\\
                  &-&-&.^{\longrightarrow} \downarrow &- \\
                  & & &=& \\
                   &-&-&.^\rightarrow \uparrow            &- \\
| \rightarrow& & &&&\\                
\end{matrix}
$
\section{TWIST as gcd seed}
$
\begin{matrix}
-&-&-&|&&\\
&\downarrow&-&\downarrow&-&\downarrow&&\\
  |                &-&-&-&- &-  \\
 \longrightarrow                  &&&&&\\
                   &-&-&-&-&-  \\
&&&&\uparrow&&\\                
\end{matrix}
$
donne 
$
\begin{matrix}
&&\downarrow&&\\
                  &-&-&-&- \\
| \longrightarrow                  &&&&\\
                   &-&-&-&- \\
&&&\uparrow&&\\                
\end{matrix}
$


ca devrait passer dessus et faire une torsion (en 3d)  mais on peut représenter avec juste un bit  avoir pour plus ttard voire si on veut voir  decoupe comme  une adjonctoin  remise dans l'ordre
 only the closest cross when you reach intensisty reevalueted

Process and mixed time and space 
\section{Combs as space slots and time pointers}
$
\begin{matrix}
-&-&&\quad \downarrow&&&--\\
&|_{t_1}&&&&&|_{t_2}&\\
&&-&-&-&-&&\\
&|._q &&._d|&|._p&&._b||&._{q+}
\end{matrix}
$

Central language $L_c$ , leftt language $L_q$ and right language $L_p$
associated to SKIP, NextSIDE, iSWAP iswap telle you how to read next  when given a furst and second argument, skip telle you itake into account the central POP, 

iSWAP is revistiing by permutation the translation of the output of nextsie to index .



simplest is 2 bits for each meaningful/error then ( I or || )  after three 
then if I then nothing if || then rewire || as  | 

Power sequence  interpretantion gives (n,m) how many copies of bithe lianes i need to imagine 
 

The ten meaningful cells

$\begin{matrix}
\downarrow&&\\
-&& &\\
|\quad&&\\
- &&\\
\end{matrix}$

$|_-|$ and $=|$ and $|^-|$ 

 Pour moi $\sqrt{2}$ n'est pas encore un nombre mais le résultat d'un calcul parceque je ne sais pas extraire de l'information de sa décomposition naturelle effficacement : je ne sais pas répondre sans lancer la récursion normale au $n^{ieme}$ bit quel est le digit  (de $10^n \sqrt{2} \bmod 10$))

MAXBELIEF est l'équivapment de MAXINT dans une machine. saut qu'il est défini localement par ATM.


\section{Exemples}
\subsection{Graph Drawing :} 
New vertex, , extend in direction, time start time stop  ,pen position start, pen position stop, 

Create clone, 
Extend,
Skip

each vertex has two coordinates but also two freedoms (a quarter plane where he can be moved by unit moves)

each veretx can have a sequence of clones canonically obtainable when one knows order of the vertices

given two points and a  weighted grid locally learnable each potentially shortest path assuming where we are larger than k  and 
potrentially longest path assuming where we are shorter then k is a sequence  using two or three directions


\subsection{ the cursor and the broken keyboard }

A graoh drawing is community friendly if

each community have  pairs of  implicit processes (to complete it from any part of it, to extend it (with a new half edge dir or increase weight )

implcit weights on vertices on start of edges and on end of edges
each vertex have a pair of integer ( coordiantes) and a pair of  postures (where it came-from/goingto admissible directions)
correlation to the time of the liane a free will / 

Meta  community will (clockwise/ant clocwise)

local free feeling of time (the cadrant is indicating past or future)

equivalence calssses for someone is wher we cant know if it is concept or percet  shortest paths have same weight)

\subsection{Non zero rationals by  4 POPs}

\subsection{Check constants protocols}
I think that it is a constant (or a welll referenced adress ,  evaluate chain , value) 

I dont know before I was born so 
(inint yes |,
 init no |,
  init yesnosym |,
   init , intit scale min |)
   each time I am asked about a time preceding my birth i add 1 to dont know (
   
$Next_w [   Swap^{._{b_p}} [ {\frac {+_{pb}- +_{db}}{.qd} }^{Swap_{bq}}]] $
   
   

\subsection{Message passing and soft beliefs}

if one view 4 pairings of the the complete graph on 4 vertices and alternates message exchanges and updates of local beliefs.

Read  $\begin{matrix} 1&3 \\2 &4\end{matrix} $ as
$ \begin{matrix} d&p \\ q&d\end{matrix}$ and 

Scheduler :

%$ \begin{matrix}
% |\downarrow&&&&& \downarrow |&|\downarrow & &&&&\downarrow |\\ 
%  -_{_{1234}} &-_{_{1243}}&-_{_{1324}}&-_{_{1342}}&-_{_{1423}}&-_{_{1432}}&-&_{_{3124}}&-_{_{3142}}&-_{_{3214}}&-_{_{3412}}&-_{_{4132}}&-_{_{4312}}
% \end{matrix}$
 
${\scriptscriptstyle
{
\begin{array}{lccccccccccccl}
|\downarrow &&&&&& \downarrow | & |\downarrow &&&&&& \downarrow |\\ 
-_{1234} & -_{1243} & -_{1324} & -_{1342} & -_{1423} & -_{1432} & -&_{3124} & -_{3142} & -_{3214} & -_{3412} & -_{4132} & -_{4312}
\end{array}}}$


The scheduler is not an agent and carries no information; it only fixes an order in which local role updates are attempted.

Outside the representation are hidden variables: direction, position changes, and a step counter.


 Action controlled by scheduler  as a role update yourself and the two local neighbors.

dynamics run in a superposition of admissible schedules, none of which is privileged locally.


lcm test : $ (x,x') = Swap^j (d,p)$   comme adresse pointée $(z,z' )=Swap^j(q,b)$ distances originelles
on veut transmettre a

$._q:= x<x', ._d:= x\le x' , |^p:= x'\le x ,|^b:= x'<x$
 
          This    
             $
             \begin{matrix}     
             | ._{b\pm q}  - ._{d\pm p} || ._{p \mp d} |    ._{b\mp q}|
             \end{matrix}
             $
             Is a lienarisation of
             
               $
             \begin{matrix}     
             &| ._{p \mp d} -    ._{b\mp q}|\\
             | ._{b\pm q}  - ._{d\pm p} |&
             \end{matrix}
             $
             
             
 First is the named edge  Dolores the limit is Persphone (is one direction) the colimit is Qudsia
 Bendict is ithe intertwinne
 
 To have a space as a largest fix point we view symmetries : if you are my shadow then you are not me because you have a spacial locus that people
 can distinguish £I may not know if I am me or my shadow
 but i have a pointer to where i can show the pair and ask people  the question : am I the same as my shadow
 with the confidence that I have that al I pointing to two different directions
 
 If I think that at some time of the year i can update my pointing position and I have no spacial locus way  to distinguish the experiment
 My confidence in my space maintaining the relative pointer adress difference and the question will generate a variable in time 
 how many pepople say yes how many people say no how many people don"t answer or say something 
 I  have no confidence in distinguishing from from yes or no that looks to me impossible  use 
 
 lets build it $|._q | ._d || ._p- .b|$ we want a anity check  of this belief 
  if we try with samples spread enough one should not be able to distinguish
    $a<b$ from $a\le b$ nor $b<a$ from $b\le a$ iwant to have an estumate updater on my beliefs in 
    
    either my constants are not constants  a and b changed between the questions (bad adressing problem)
  or my evaluating system is too primitive  i
  
 $
 \begin{matrix} 
  &|&-&-&-&|&\\
 |&.&-&.&||&.&-&.&|\\
  & & &|&- &-&-&|\\
 \end{matrix}
  $        
  
  -Si j'ai une séquence de constantes rangés par leur fiabilité générqiue , puis un masque de certification récente qui me donne l'accès qu'un sous syteme
  
  une séquence de variable rangée par la facilite de les changer (le temps d'attente avant changement par exemple)
  avec un masque de validation locale
  
  Je peux represenrer une sequencee de paires comme decalage du matching en utilisant des copies imaginaires
  
  Nom-lien duplicable de connection a l'adresse de reference et lien de deplacment de rerence (autre constzante) le de but n'est jamais équivalent  à des sommets de la liane : une liane c'est les copies du meme sommet dans un dessin de graphe, un cnncetion simple/+planarisee/ de degre 1, 2, ou 3 

on prend le changement de liane le plus proche  
  
                                                                                         
\subsection{the open half line}
This construction lives at the pre-numeric, pre-metric level.
At this stage, there is no notion of value—only motion, reversal, and trace.
We assume an origin given by a pointer.
Distance is not given as a value, but as a \emph{candidate} equipped with uncertainty.
Movement is described by:
an arrow carrying an old direction,
a tentative move,
and a possible reflection.

The primitive operations are:
\begin{itemize}
\item move;
\item record thickness;
\item count steps;
\item flip direction.
\end{itemize}

$$
\begin{matrix}
&&\downarrow\downarrow&&\downarrow \downarrow&\\
\rightarrow&-|&-&||&-&|\\
&\lightblue{\leftrightarrow}&\darkblue{\rightarrow}&&&\\
&|&|&&&\\
&\text{offset}&|&-&\text{move step}&\text{direction--length}
\end{matrix}
$$

or 
$$
\begin{matrix}
&&\downarrow\downarrow&&&\\
\rightarrow&-|&-&||&-&|\\
&\lightblue{\leftrightarrow}&\darkblue{\rightarrow}&&&\\
&|&|&&\downarrow \downarrow&\\
%&\text{offset}&|&-&\text{move step}&\text{direction--length}
\end{matrix}
$$

A trace must be equipped, by a community, with a way to produce
a lower-resolution estimate of a sequence.
This estimate is expressed in terms of an internal clock,
through operators such as:
\[
\texttt{NEXTINT}(\text{position bounce / reflection}),
\quad
\texttt{NEXTINT}(\text{position cross}).
\]

At the penultimate level, these estimates take values in strictly positive integers.
At the ultimate level, only  rigidified associativity constraints remains:
bracketed terms on a horizontally flattened axis,
which may be realised as unary, levelled, or equivalent representations.

Rescaling is imposed by the community as part of its interpretability discipline:
only values admitting an interpretation on four integers
(with a truth / falsity / undecidable reading)
are preserved.

Counting emerges later as a stabilisation of repeated alternations.
%
%given origin by pointer , distance uncertainty candidate , distance move dandidate , arrow of olddirection  
%move, get thickness, get nbsteps, flipdirection
%
%$$\begin{matrix}
%&&\downarrow\downarrow&&\downarrow \downarrow&\\
%\rightarrow&-|&-&||&-&|\\
%&\lightblue{\leftrightarrow}&\darkblue{\rightarrow}&&&\\
%&|&|&&&\\
%&\text{offset}&|&-&\text{move step}&\text{direction-length}
%\end{matrix}
%$$
%
%a trace should be equipped  by a community to  get a lower resolution estimate of the the sequence in internal clock of NEXTINT(pos bounces (reflexion)) abd  NextInt (pos cross) (comprehend) with as the penultimal level are strictly positive integers  and the ultmate level just associativity  bracket terms (on horizontally flattened axis) whicould be  a levelled unary representation or whatever/
%
%(resize will be automatically done to maintain only values with an interpretation on 4 integers with an interpretation on true false)
\section{Starting  mathematical constructions of communities}

\section{Minimal Meaningful DF-CAP (in a Diversity-Friendly Concept-Anthropomorphism-Percept) structure simplement étalée}

I say Etalée when the representation doesnot uses complex Zoomout or Zoomin Procedure ( I wont use for now so the description might be corrected) 
but just  (role as directio in diagonal copy zoom out) and goto role as arrow insert | a side | without crossing arrow or | as Zoom in  

Simplement Etale means :

\begin{itemize}
\item Zoom out is  role as direction in diagonal copy  with new pointer at new copy
\item Zoom in  is  insert | a side | without crossing arrow or |
\end{itemize}
 
 default role is d default pos to point is p
 
 \section{number by counts}
 for now 
 Lets say that $b$ start form left of rightmost pointer then right of central left pointer
 $q$ starts from leftmost ointer to the right then from central right pointer to the right. 
 
 $d$ and $p$ alternates direction d strating from left of left  centralpointer and p of right of right central pointer
 
  one by one the strokes in a way i can just know how much time she requires  
  \section{ the four from compare }
  let us assume that we are on an axis (so we have our  sequence of pointer with prebuit integers that we can access by value 
  and as for now we are are at level of granularity allowing us to have access to constants  x and y linked by indistinguishabiility to their first appearance
  at position in space separated words here as $d$-XVI and here as $d$-XXII.
  
  the test $x<y,x\le y,y\le x, y<x$  should always give exactly two terminating programs ( I will be able to work with one or three)  but none or 4 is a new block of uncertainty.  
  
  
  Here is the lowlevel we want  when :
  \begin{itemize}

\item   gcd stye
  \begin{itemize}
  \item Catch with WAIT :
  
\item  alternate
\subitem  get x and copy locally in new internal border with 4 pointers and four traces
\subitem   get y and copy locally in new internal border with 4 pointers and four traces  
\item   each 4 scrolling methods compare and answer 
 \item    pit the largest as packs of the first with pointers on packs  now we got a reminder 
  \item  iterate with reminder and small in new layer 
  \end {itemize}
  \item lcm style
  \subitem  the small as distance and insert marks on the second if ever looks the same position at this granularity cross the stripe and start a new one with one mark to indicate where you want to sart( somehow we had a 2 we want now a 2.2 (the other by the way beacoming a 2.1)
   
  \end{itemize} 
  
 \section{inherent uncertainty: A radiating black hole like source has two indistinguishable bases the losange and the rectangle}
$$\begin{matrix}
&\downarrow\downarrow&&\downarrow \downarrow&\\
\rightarrow|&&||&&|
\end{matrix}
$$


the uncertainty $b$ in  $b_q-b_p$ as $[$:


Representation of  Box $[_.$: or just $\rightarrow$ : what I can view as producing sequence of right arrows:

$$
\begin{matrix}
&&&-&&					&&&&&&&\\
&&|&&&					&&&&&&& \\
&-&&&&					\leftarrow_{_{||||-|}}&&&&&  \\
|&&&&&					\leftarrow_{_|}  &\red{\uparrow_{_{||||.}}}\quad  &&&&& \\
&-&&&	\downarrow_{_{||}}&	.& \uparrow_{_{||||}} \red{\rightarrow_{_{||||-}}} \\
&&|&&&					\rightarrow_{_{|||}}&\\ %&&|&\\
&&&-&&					&&-&&&&\\
&&&&|&					&|&&&\\
&&&&&-					&&&&&\\
\end{matrix}
$$

%$$
%\begin{matrix}
%&&&-&&&&&&&&\\
%&&|&&&&&&&&\\
%&-&&&&\leftarrow_{||||-|}&&&&&\\
%|&&&&&\leftarrow_{|}&&&&&\\
%&-&&&\downarrow_{||}&.&\uparrow_{||||}\red{\rightarrow_{||||-}}&&\\
%&&|&&&\rightarrow_{|||}&&&&\\
%&&&-&&&&-&&&&\\
%&&&&|&&|&&&\\
%&&&&&-&&&&&\\
%\end{matrix}
%$$

%$
%{\scriptscriptstyle
%\setlength{\arraycolsep}{1pt}
%\renewcommand{\arraystretch}{0.45}
%\begin{array}{@{}*{8}{c}@{}}
%  &   &   &   & \cdots & \cdots &   &   \\
%  &   &   &   &   & | &   & \vdots \\
%  &   &   &   & - &   & \cdots &   \\
%  & | &   &   &   & | & | &   \\
%  & - &   &   &   & - &   &   \\
%| &   & | &   &   & | &   &   \\
%\cdots &   & - &   &   &   &   &   \\
%  & \vdots & | &   &   & | &   &   \\
%  & \cdots & - &   &   &   &   &   \\
%  &   & | &   &   & | &   &   \\
%  &   &   &   &   & - &   &   \\
%  &   & | &   &   & | &   &   \\
%  &   &   &   &   &   &   &   \\
%  &   & | &   &   &   &   &   \\
%  & - & - & - & - &   &   &   \\
%  & | &   & | &   &   &   &   \\
%\cdots&   & - &   &   &   &   &   \\
%  &   & | &   &   &   &   &   \\
%  & - &   &   &   &   &   &   \\
%  &   &   &   &   &   &   &   \\
%  &   & - &   &   &   &   &   \\
%  &   &   &   &   &   &   &   \\
%  &   &   &   &   &   &   &   \\
%\end{array}
%}
%$

Here the losange shape border can be imagining as pushed by the arrow intent so contracting and expanding between two separated dynamic losanges keeping their center the ray spiralling is indexed by unary time separated by time he goes close to the open boundary  the internal mechanics is  a DFCAP way to produce horizontal arrows time between outputs can be related to the dynamic of the diamond border.
The illustrtive drwaing of the inside losange boundary with . and of the outside boundary is to make sens of wht the arrow changes direction the boundary is implicit or imaginary and not definable for the red arrow observable but guessable from how the arrows are produced. 

The first perfect  uncertainty moduator is the inward or outward  direction of the internal spiral (that is alternating depending when boundary of boundary starts pushing back). Internally, the process alternates between two regimes (inward / outward), but this alternation is not locally decidable from a finite trace.



The second indinstinguishablity is the rectangle or losange boundary or as between  going  right and going down right , to represent it   with  only  vertical horizontal   have to be represented with only vertical and horizontal srokes so we imagine turning the rectangle 

and it would look like (implicit boundary now omitted)

Representation of  Box $[_|$: or just $\rightarrow$ : what I can view as producing sequence of right arrows:

$$
\begin{matrix}
                         & &                        &\uparrow_{_{||||-|}}\\
\leftarrow_{_{||}}& &\uparrow_{_|}\\
                          &|&&                   \blue{\uparrow_{_{||||-_.}}} \\
\downarrow_{_{|||}}&&&                \uparrow{_{||||}}\blue{\rightarrow_{_{||||-}}}\\ 
\end{matrix}
$$
 
 
 \section{error resilience}
 Suppose that each site has a decode an encode operationl CAP friendly ways  to have  

\section{De l'implicite dans la définition}
Définir quelque chose qui fait sens dans une communauté (sera l'implicite ici à partir de maintenant en lisant de gauche a droite puis de haut en bas ... et olus généralement avec votre indéfinissable et si personnelle maitrise de la langue utilisée ici.)

Si Dolores dis "je vois le mot tomate à tel endroit à telle heure". Elle aimerait (implictement) que si elle avait dit "je vois le mot tomate à telle heure à tel endroit" Le sens transmis soit le même. 

\emph{Cet exemple volontairement trivial sert uniquement à isoler ce qui doit rester invariant quand l’ordre, le support ou le contexte d’énonciation varient.}


 $$\begin{matrix}
& d\text{: dis nous comment,  Dolores?}&&\text{ Quand}&-&-&-\\
&&\\
&\text{ Ou}&&\text{ Element } d-\text{siginifié: }  \text{tomate}&&&|\\
&           |& && &\text{ tel endroit }& \text{telle heure}\\
&            -&-&-&-&|\\
\end{matrix}$$ 


On a envie déchirer l'espace temps pour définir un réel à ancrage qui semble commun avec deux clés d'accès qui semblent commuter. (Si on avait la perfection on dirait chouette on a un $\mathbb F_2$ module le tour est joué: tout sens est quasiment déja porteur d'un $*$-morphisme implicite sur Répéter). 

On va demander un peu moins, on aimerait juste que si les personnes d'une communauté partagent assez d'implicite pour que que s'il savent si j'ai demandé 'Ou' avant 'Quand' ou 'Quand' avant 'Ou' il arrivent a se metre d'accord sur un sens que je ne peux voir que comme un ou indistingable si je ne reviens pas dans le temps à un passé antérieur à celui d'avoir posé la question. 
Des notions sous-jacentes plus complexes peuvent émerger de ce qui me paraissait comme ancrage naturel on voit donc que la paire d'ancrage (Ou,Quand)   peut n'être que ce qui parait comme un (Ou,Quand). forcer la co-définissabilité c'est une interruption du processus de raffinement qui génère  une équivalence  phénoménologiue  qui  est une identité communautariste qui  fait naitre un système équationnel par répétitiondans  notre signification symbolisable.

 L'apprentissage du sens de tomate dans une communauté rend  le temps et l'espace comme co-définissables ou ici co-indexables par des anthropomorphismes: elle effacerait l'individuation si on oubliait l'histoire.

On aimerait pouvoir définir un intervalle de départ raisonnable et un pas d'avancement puis :

\begin{itemize}
\item Je controle une rigidifcation : je vais écrire "tomate" dans une adresse  dont je te donnerait le lien d'accès , et la technique de lecture (en espérant pour le moment qu'on arrivera a définir cette ligne localement  plus tard)
\item ensuite entrelacer
\subitem dire : je vais dire des espaces et tu me diras si tu as vu "tomate" dedans
\subitem  dire : je vais dire des temps et tu va me dire si tu as entendu "tomate" dedans.
\end{itemize}

Mais ce n'est pas par tant par la peur du rasoir (d'Okham) ou par le déni de la simplicité (je ne suis sincèrement pas convaincu que restreindre au parfait est plus simple car plus riche d'implicite si on n'a pas de garantie justifiant de pouvoir avoir un cadre explicite pour que l'abstraction fasse vraiment sens au sens communautariste).


Celui qui écrit, même s'il a la volonté poétique d'étaler son coeur dans ses mots ne peut peut être pas le faire parfaitement  justement parcequ'il  utilise des mots. On a donc accès par le langage non pas à des personnes mais à des rôles  qu'on appellera des \textbf{A}nthropomorphismes (qui se comportent comme des personnes ou au moins comme des processus interactifs interruptibles). 

De par la volonté du narrateur le propos peut être porté par un anthropomorphisme qu'il imagine différent de lui même et s'adresse à un anthropomorphisme qui vient de la représentation qu'il se fait du lecteur ou e celui qui écoute. Comme le Ou et le Quand  le qui parle et qui écoute doivent aussi dans la notion de compréhension transmissible à enseignement transmissible doivent aussi pouvoir êre rajoutables  par extension sans perdre le sens quand elles sont implicites: le sens qu'on veut exprimer et transmettre doit supporter une forme d'\textbf{extensionnalité}.  De plus comme pour le Ou et ke Quand,   la position qui pointe vers l'un ou l'autre doivent être perçues comme différentes par notre représentation  de celui qui comprend.
 On ne demande pas à celui qui entend de savoir ni de pouvoir ni de devoir répondre mais on dans le critère de vérification de la compréhension on voit émerger naturellement 4 positions qu'on aimerait pouvoir rajouter quand elles sont implicites dans les cadres ou a l'impression de se comprendre.


Le symbole lui par sa position joue le role de contractant pour rendre le propos transportable peut être considéré comme induisant  si un anthropomorphisme  qui le peut  et le veut essaye de le reproduire une forme \textbf{R}igide de pensée.

Un propos siginifié en intention est rigidifié par cette explicitation projective.

Deux A qui se suivent ne peuvent échanger sans socle rigidifié commun tout comme deux R qui se suivent ont besoin d'un A pour fair le lien entre elles, l'accumulation du même introduit néanmoins un point de rupture ou d'interruption dans la production explicite de l'expression.
 
Une communauté est raffinable si, à partir de traces indistinguables localement, elle peut produire des distinctions supplémentaires sans invalider rétroactivement les traces précédentes. 



\section{exemple}

Contexte : 

- un clavier avec la barre d'espace  et la touche entrée (qui fait tout seule une tabulation et descend toujours au dessus de la diagonale tracée de la position précédente du curseur
- une dynamique de resolution d'ecran qui s'ameliore : de temps en temps les pixels diminuent

Communaute Dolores : Le curseur , d'ou il vient ou il veut aller

Observations Persephone: sequence de position de curseur

 memoire limitee: 2 derniere positions
 
 \section{interactions Comprendre et Réfléchir}
 Si on est Dolores du rectangle  (écrivant sur l'horizontale haute vers Perséphone et qu'on imagine un rayon arriver vers le bas en diagonale  il peut etre transmis ( on comprend)
 
 $\begin{matrix}
 &&-&\uparrow\\
 -&-&-&-&\rightarrow_{dp}\\
 -&\uparrow&&\\
 \end{matrix} $  ou  $\begin{matrix}
 &&|&\rightarrow\\
 -&-&-&-&\rightarrow_{dp}\\
 -&\uparrow&&\\
 \end{matrix} $Une compréhension (le bord fait sortir du cadran)
 
 
 $\begin{matrix}
 &&&\\
 -&-&-&-&\rightarrow_{dp}\\
 -&\uparrow&|&\rightarrow\\
 \end{matrix} $
ou  $\begin{matrix}
 &&&\\
 -&-&-&-&\rightarrow_{dp}\\
 -&\uparrow&\downarrow&-\\
 \end{matrix} $    Une réflexion (le bord fait rebondir  dans le  cadran)
 
 
 \section{Histoire des sens sur Dolores Persephone du unaire au binaire}
 Si on est Dolores du rectangle  (écrivant sur l'horizontale haute vers Perséphone et qu'on imagine un rayon arriver vers le haut  (en diagonale indéterminée)  il peut etre transmis ( on comprend)
 
 $\begin{matrix}
 &\downarrow\\
 &-&\rightarrow_{dp}\\
 \\
 \end{matrix} $   can become (SKIP$_{dp}$)   
 $\begin{matrix}
 -&\downarrow\\
 &-&-&\rightarrow_{dp}\\
 \\
 \end{matrix} $ 
 
or the superposition  (DRAW$_{dp}$)    
$\begin{matrix}
 -&\downarrow\\
 -&-&-&\rightarrow_{dp}\\
 \leftarrow&||&\rightarrow\\
 \end{matrix} $  
 Note that it is left by default because $dp$ axe so first plausible element is left diagonal transmission 

 
 \section{Optimisme Pessimisme et dessin d'encoche controlé}
L'optimisme considaire à commencer par un pointeur sur le système libre qui ressemble le plus à une a une encoche permanente  (ou le systeme que je peux construire dans le temps imparti pour intialiser ce que je suis avant l'interaction)

L'optimisme considaire à commencer par un pointeur sur le système libre qui ressemble le plus à une a ne jamais voir d'encoche  (ou le systeme que je peux construire dans le temps imparti pour intialiser ce que je suis avant l'interaction).
Tout comme un trait élémentaire est toujours équipé d'une intention explicite ou explicitable comme superposition, une personne a un endroit a un etat d'esprit par defaut
dans toute communaute  dolores si on demande si on mettait a tel endroit serais tu content on aura un oui, non ou je ne sais pas.
Toute personne qui repond à une impression sur la direction de la flèche du temps.
Toute personne qui pose une question a une direction par defaut en tete qui dépend de la sienne et de sa position. 
   
   \section{ Dans Dolorès (sous entendu en coin de rectangle, sous entendu sur $dp$ en train d'ecrire) $q$ comme bornes a interporetation par defaut  cumulative ou bosonique  et $p$ bornes differentielles ou alternantes}

\section{Essai pour Version info avec etre comme "avoir une Adresse" et avoir comme "donner une Valeur"}

\begin{itemize}
\item (Dolores rectangle, Dolores losange) comme   ($\text{ @}$  , $\text{ @}$  $\rightarrow \text{ \$}$)  ou 
\subitem($\begin{matrix}
& \rightarrow&R\\
\uparrow&A&\\
R&&\\
\end{matrix}$ , $\begin{matrix}
 &R&\\
\downarrow&A&\\
&R&\\
\end{matrix}$ ),

\item (Persephone rectangle, Persephone losange) comme   ($\text{ \$}$  , $\text{ \$}$  $\rightarrow \text{ @})$
\item (Qudsia rectangle, Qudsia losange) comme   ($co-\text{ @}$  , $\text{ @}$  $\rightarrow \text{ @})$,    ($\text{ @}$  , $\text{ @}$  $\rightarrow \text{ \$}$)
  ou 
  \subitem ($\begin{matrix}
A& &\\
\uparrow&&\\
R&&\\
\end{matrix}$ , $\begin{matrix}
 &&\\
\uparrow&A&\\
\leftarrow&R&\\
\end{matrix}$ ),

\item (Benedict rectangle, Beedict losange) comme   ($co-\text{ \$}$  , $\text{ \$}$  $\rightarrow \text{ \$})$
\end{itemize}

Les covaleurs étant des bornes inf d'adresses  dans leur lianes et les les co-valeurs des biorens sup de valeur dans leur liane.

Chaque locus a un process de decodage et d'encodage avec changement de role chagnement de forme par defaut 
l'idee est que traduire les trait unaire par des crochet va donner les seules restriction qu'on impose a l'associativite et nous donner les environnement mathematiques vamables et les canaux de communication, et les regeles de mutation par changement d'echelle.

\section{Essai Etalé sur une  ligne on prend horozontal}



$\begin{matrix}
 R&A&R\\
\leftarrow&&\rightarrow\\
&\text{P }&\\
-&&-\\
|&&|\\
&|&
\end{matrix}$ ,   and 
$\begin{matrix}
 R&A&R\\
\rightarrow &&\leftarrow \\
&\text{D }&\\
&|&\\
|&&|\\
-&&-\\
\end{matrix}$ , and    
$\begin{matrix}
 A&R&\\
&\rightarrow &\\
&\text{Q }&\\
&&|\\
&|&\\
-&&\\
\end{matrix}$  
and    
$\begin{matrix}
 &R&A\\
&\leftarrow &\\
&\text{B }&\\
|&&\\
&|&\\
&&-\\
\end{matrix}$ 



On aurait des 3-2 accolades ouvrantes comme
 $\begin{matrix}
&\text{ un des deux bords extérieurs} b_w  &\\
& |&-\\
&| &\\
&|&\\
&-&\\
&\text {un des cotés centraux} q_{w'}
& \text{centre indecidable}\\
&\text {l'autre coté central} q_{w'}\\
&-&\\
&|&\\
&|&-\\
&\text{ L'autre bord extérieur} b_{\bar w}  
\end{matrix}$ 



Cest une représentation imaginaire je ne peux pas aller voir les deux bords en meme temps et iles deux bandes peuvent se croiser sans que je ne m'en apercoive 
c'est juste ma meilleur representation de ou je suis.

si $t \mapsto w$ et $ t'\mapsto w'$ (il faut du temps pour la première définition $t'$ par défaut peut avoir besoin d'un délai d'indice pour être co-définissable. Mais je peux observer les traces sur deux cotés centraux que je je ne suis pas capable encore de distinguer les traces des rebonds bvenant de l'un des deux bords pour les quels je peux ancrer haut ou bas donc j'ai un système de distinction mais je ne sais pas si la dérive va à l'intérieur ou à l'extérieur. 

Je vais choisir la séquence ordonnées de ce qui me parait comme  choses auquelles j'ai libre accès  dans les quelles j'ai le moins confiance comme   pseudo aléatoire ou "$B_p$ est en haut " 
 et les choses dont j'ai le plus confiance pour les constances "$Q_d$ est en haut " 

Un Anthropomorphisme est la source de liberté héritée de $B_p$ elle est indécidable je ne peux que quantifier des propriétés  du type en sachant ou j'en suis
et la question rigidifiee si j'avais $N$ copies je pourrais dire que plus que 1 et moins que $C $ seraient bonnes et que plus que 1   


Je ne sais pas si la bordure haute (ou la basse)  est celle interieure ou exterieure  par contre la bordure la plus haute et la plus basse que je connais n'ont de sens qu'encodees en unaire, cependant leur role permet de fixer le parentesage des sous bordure  et me fait gagner un encodage en binaire sur des bornes plus proche du centre

Si séquentiellement 
\begin{itemize}
\item  le role  Benedict vu comme $B<N$  dit  (croyance = $.d$, assurance =$.p$, multiplicite requise $.bp$ , $.bq$)
\item le role  Perséphone vu comme $B\le N$  dit 
\item le role Qudsia vu comme $N<B$ dit 
\item le role  Dolores vu comme$ N\le B$ dit 
\end{itemize}



from
$$\begin{matrix}
\text{D: }\uparrow - &\text{P: } - \downarrow \\
\text{Q: }\downarrow \leftarrow  &\text{B: } \leftarrow \uparrow \\
\end{matrix}$$ 



 -confusabilite d'adresse   distingabilite  de valeur (plusieurs valeur se suivent sir la meme hauteur vue comme index ou adresse, ?

dans la sous couche parentesage rigidife (premier bord unaire) on devrait avoir localement:
-estimateur de regilarite (4 nb)
-estimateur de singularite (4nb)
-alternant de maintien de sens. permet d'abvoir 6 nb  a partir des 8


$Bp:= \begin{matrix} 
             &&&&|-&-&-&-&-&|\\
            - b&p&[q&d][&p&b&q&d&]&\blue{\uparrow \rightarrow}
\end{matrix}$


$Bq:= \begin{matrix} 
            - b&[q&d]&p[&b&q&d&p&]&\blue{\downarrow \rightarrow}\\
             &&&&|-&-&-&-&-&|
\end{matrix}$

Un obfuscateur fait des pops qui permettet de reconstruire  $w_qw_dw_p$ avec un $\forall$ ou $-$ dans $w_d$
reconnaitre un mot.... 
%\section{Hommages a ceux ont fait de ma pensée ce qu'elle est}
%\begin{itemize}
%\item dont la trace des chemins croisée directement ou par échange scientifiquea fat de moi ce que je suis
%\subitem de par leur gentillesse:  Maffray, Gyarfas,Hoyer,Jorrand,Payan
%\subitem de leur converture encyclopédique de leur voisinage Sebo, Watrous, Aaronson, Fawzi,
%\subitem de leur insight  Raussendor
%Kind	Encyclopedic	Insghfful	Inspiring	Amitié	Brilliant	Goodness	Curiosité intellectuelle	Close preocupations	Solid	"Viewexpander
%"
%Maffray	Sebo	Shor	Barret	Perdrix	Slofstra	Fawzi	Tillich	Mancinska	Zemor	Diderot
%Gyarfas	Watrous	Kitaev	Pironio	Gravier		Jorrand		Meiburg	Koenig	Lautréamont
%Hoyer	Aaronson	Deutsch	BelenSainz	Savin				Viddick	Gottesman	Bourbaki
%Jorrand	Fawzi	Brassard	Ambainis	Magniez				Perdrix	deWolf	Spinoza
%Payan		Raussendorf	Reichardt	Durr				Coiteux Roy	Pantelev	Widgenstein
%		LeGal	Popescu	Hoyer				Omrod		Husserl
%		Wildberger	Eisert					Wolfe		Kierkegaard
%								Keusch		BourdieuGirard
%	


\end{document}
\textbf{tofix
et donc des symboles on a donc affaire non pas à une personne mais à un anthropomorphisme  qui est le messager qui amène le propos et à qui on peut imaginer poser la question du ou ou du quand sur tout propos qu'il porte. Elle peut ne pas répondre, elle peut dire je ne sais pas mais dans notre représentation d'elle on aimerait nous dire que si c'est une bonne représentation de personne elle doit au moins avoir compris que quand puis ou et ou puis quand est un choix qui ne depend pas  d'elle 
position differente   distinguer position
On a donc montré ici que si on veut ancrer le sens dans une histoire par exemple celle de la représentation implicite  du narrateur de sa vie, 
}

\section {Febrary}
\section{3/2}
For Dolores  as ATM 

$$\begin{matrix}
\text{Who After}&\\
&\text{Dolores}&-&-&-\\ 
&&& &|\\
&|&&\text{How} &\text{When}\\
&|&-&\text{Where}&\text{Meaningful statement}\\
&&&&&\text{Who before}
\end{matrix}$$



How it was drawn 
Who where when  she thinks said it before  who/where when saud it afte


Defining induces a community : if we anchor a statement rigidifying it for a community and view Dolores in the community

we can ask her who did you hear the statement form, when people ask Dolores 

When did you hear it ? if yess did you see  how it was written ?

who do you think will use it next and 

Statements carrying people in Dolores mind :

Next and previous

\section{An attempt for how 'to know' in the most general way  when restricted to finding  of changing approximate detectable  traces of that knowledge that are teachable by mimicing them }


We will attempt here to synchonize the mythic tales of understanding with the concrete strokes of explaining under the bounds 
of imperfectness and transmissibility.

I want to be able to draw in a symbolisable, teachable by mimic way an explanation process trying to transmit while updating it own conception  (improving understanding)
of meaning with concretely locamlisble pairs understanding explaining.



Wiewing doubt as ontological multiplicidty oif rigidifiable ways of thinking and delaying the requirement of definition to the end of thaught is not an easy exercice as it forces you to start from deep water struggling for air to breathe and at each time you reach air (consistent definitions , methods of thaught, ways to review memories consistenl, discarded parts we want to forget we call exceptions or postselection or works in progress) and as soone as you see the calm tranquility dive again rebuilding all with your past knowledge as accompagning memory allowing richer experience not as a constraint of what do but as a link to how and why change can be envisionned.

 
 Maintaining diversity may seem "sacralised"  or axiomtised  here  but it jhonestly  looks  uanavoidable posture to Doores first by the way itself I feel the discrepancy between what I think and what I say 
 but also in he definition itself on what it required to look like understanding : the deifnition starts with two pairs (explaining , understandig) with liberties( thaughts, symbols) and c ostraints (imperfection, transmissibility).
   
 It will appear in the end (at least it does already so to me and did for a few months)  as a powerful tool to break the loop of contradiction in reasonin
 by no chosing one of the iusually adopted strategies  tradition atomic or universal anchoring
realising the imperfectness between concepts and percepts can provide a way to make sens.

view the plurality as a way to break the loop paradox by viewing as directions ( diemnsions) the two distinguishable way that seems the same when I feel as a point
one through the tales of telling withan intermediate anchor Persephone which represents the day I die  and one through the sacrifices of the one before me by looking at traces left that are showable and explaineble inthe path of origin of space starting at Qudisia. Admitting that I might be wrong forces me to build a shadow of my self which could be wright linking differently my percepts of history and my concepts of time that i cann Bnededict.



: the axiom of Dolores here is admitting a finite bound on speed of thaught as splitting/distonguishing enonciting  process
(even before answring matching symbols with percepts and wandering : 

checking if the question I am asking make sens is problematic because of the loop pardadox puuting my self between concepts and percepts anthromophising them and allowing them
to conncet through me or through what I could imagine will allow to build the four cornerns of a dynamic process of making sense tp which I can anchor myself.

  


at any locus given to me by consensus I can imagine a mirror of my self that thinks differentley in such a way that  anything outside that locus 
cannot distinguish the internal  (at least bi)- directional paths leading to my answer in an interaction 

 way that I see at least two seperable pre-paired-figures, and teach gesture allowing to draw.
 
 
 

\begin{itemize}
\item Viewed by Dolores:
\begin{itemize}
\item Locus Encoding : dolores anchoring worlds and feelings as pairs 
\item History/ Encoding : qudsia (adress or memory holder) where one find the 'old' traces that we believe happened before us 
\item Assurance Encoding : persephone :life and its important as emergence of values  and qualities
\item Belief Encoding benedict
\end{itemize}
\end{itemize}


iuntuitively:

the way Dolores sees below is PUSH as Expand it is the way of ideals  linking the trace - to ultimately below the inductive universes as infinity of algebras and universes build series of * dualisers leading  to appproximate c*algebras  it is the rule of induction of "ppcm" plus petit commun multiple that justifies whats can be perceived there it is how  I view myself as part of something (or the $H_z^t$ of a stabliser of a pure quantum spcae viewe as intersecting of two subsaspaes in a Hilbert state

the diagonal below comes from benedict which is imigined  living in  bi-infinities  and which role is to give on the block diagonal element correlation 
memories to speech (fragment of dq)  to fraglent of (dp) linked with frament of dd as Synchronised repeat in quantum information it is how event thow pointwise we have a fine level woth anticommutation commutation it appears synchronised descriptions to the outside world as a for all lines the product is 0 when we view not qubit wise but fix points wize with a Z and an Z view.

on the rogth side leads to infinitives of projections the key here is to divide (potentially repeated copies ) its the pgcd (plus grand commun diviseur) rules to establish qualities that will be reversed as small quantit!es when dualised to be consistend with qudsia memorial side)  it is the POVMs or the $H_X$  of a stabiliser. 

somehow right side  at one level above traces is decoded as numbers as qualities as pairs with a pair belief assurance tracking whereas
down side has numbers as units or quantiies derived knowing what i know for better expressivity
 
 
 
 
\section{dolores links and strokes in vertical q to d horizontal d to p}
$\begin{matrix}
&&&-_|&|^-\\
&&&&&\rightarrow\\
&&&-^|&|_-&\\
|_-&&-^|\\
|^-&&-_|\\
&\downarrow&
\end{matrix}$

with areference for each it becomes

$\begin{matrix}
&&&-_|&|^-&&&-_|&|^-&\\
&&&&&&&&&\rightarrow \\
&&&-^|&|_-&&&-^|&|_-&\\
|_-&&-^|\\
|^-&&-_|\\
\\
\\
|_-&&-^|\\
|^-&&-_|\\

&\downarrow&
\end{matrix}$

horizont lines can be visualised as  rectangular forms going along the diagonal  and vertical line are losange lines going down vertically
Persephone provides extra time for new form

Qudsia provides space when asked as extra dimension fro new element

we can maintain one direction as stable losange will always go down vertically 
and now persephone has two directions the b or the d one 


\subsection{lost thaughts}
\textbf{1sec surviving idea to try to try to consider} cap uncap liberty : capitalised (or fully capitallised ?) term is \textbf{by default} what looks like a person : an anthropomorphism: an  atm and lower case by default rtm a rigid way of thinking  Be carfeul of cinfusing Default reading  with forced reading (showing the opossibility if changining mind ios as important as showing the existence of a default posture)

Any community C has for any community A considering it a by default mark.

  Drawing a a verb has a direction and strokes while drawing as a name would be the superposition of ways i imgine it could have been made synchronised with my reaction by the description of how I feel.   


can i have implicit links : variable , should I bother to make inks consistent with a geometry or do I have the freedom to claim and give independence :
do I need to maitain a control by putting between brakets implicitly saying that now braket rules apply 

implicit braketting is the first community building process: we will not repeat at each meeting what we all believe in we will move forward assuming we a have been educated 
to know hgow to read whats inside.

The posture here is cautious by default if i claim that i can define i have to check (indefinitely) consistency with what I said before : the independence 
is never guaranteed beyond a margin that can be viewed as confidence in consistency or just a belief that i have to tke the burden to try to quantify with what i can imagine and accepot to revisit after each step in which thuis that look like  exchanges might  happen tuestion , interruption, error).


twoards small communities : i need to be able to add a separable environment and a commuataion on it  :i.e.
at this time at that space ,story  equivalent to at this space at this time story :: two translation mechanism for the carrier 

\section{janiary}

\section{draft ideas}
Implicit default pattern here 

$\begin{matrix}
d_{||} &p_{||||}\\
q_{|||}&b_{|}
\end{matrix}$ flattended as $b_q  |  q | \frac d b | p | p^b$ 

inside counter when dp mode :
 $\begin{matrix} \downarrow &&&&\\
dd&dq&dp&db
\end{matrix}$
 
 
 inside counter when dpp mode :
 $\begin{matrix} 
 
 \downarrow &-&-&-&-&-  &-&-&\downarrow\\
dd&dq&dp&db&|&dbp&dpp&dqp&ddp
\end{matrix}$
Boundaries at 4 axes with storkes  each person have access to two axes : Dolores have $o_{dp}$ and $o_{dq}$

Boundary arithmetics is  local unary PUSH synchronised by viewer over something like
  (WAIT | ((WAIT| PUSH -DRAW|PUSH) or  (DRAW-WAIT| PUSH WAIT) ))- ( ((WAIT| PUSH -DRAW|PUSH) or  (DRAW-WAIT| PUSH WAIT) )|WAIT |)


where we are 
dp axe is $[pgcd(2 (pgcd (.,.))]]$ we have a rezaing proceess removing the 2 and a second process  the axes strokes are encoded $[x\mapsto 2x+1]^{mathbin.int}]$   
dq axe is  
$\begin{matrix} 1\\-\\ [ppcm(._{1},._2)]^{math}]
\end{matrix}$
 
 when reading I need to know : where I start, inwhich direction, for how long, how do I locally cumulate, How do i evolve by formulating the change I think happened
 
  
\subsection{Dolores elementary liane reading order by default}
dolores at $t_.$: 

$\begin{matrix}
& (n/2)^+ \text{at } t_2&n^+ \text{at } t_4\\
n^- \text{at } t_3&&\\
(n/2)^- \text{at } t_1&&
\end{matrix}
$

The liane is the rotated first line to become a visible tale, sequence of symbols  with attemptable meaning other than consective adjonction of signs.


The visible sequence for Persephone :
click , look gives decode form $1$ as $n$ ( shifted by $n$ one could add later with at most one as error)

then continue as if jmorally justt counting monotonically increasing binary encoded as stikes with vertical pointer access going to th border and transmiiting to Benedect

to mainiatin a monotone external reading : one reads $qd$ then reads $dp$ with the binary encoding map $x\mapsto 2x +1$ viewed as push and Draw with draw as stroke
 

\subsection{The languages under belief of perfection}

An exoressedmlessage (is decoded (how many real expressed thaughts with a uniqiue anchor is it based on , 

the first is a bigroup with perfect history: two vertical arrows on  two side by side biiinfinite lines  called LSpace and LTime
commands are q as : move lef  arrow to teft, qd move left arrow to right, pd move right arrow to left , db  move right arrow to rightft

getpos returns unit move, ref start pos just before or just after stroke, 
get upbound , max nb steps if continues moveing justbfore or just after going after relevant information zone

any pair of willed numbers as $(w_a,w_b)$  define two Gaps  $[\text{var}_a-\text{var}_b)]$ and $[b-a]$  such that  
group with perfect history  Gaps: 
  (Initial pos, intial direction),
  (just beforestab  , just after stab ),
 error bounds (stab check, chenge check) l
 local unit rescale (mult,div) 

assuming reference space is y (get y gives marked. will be reapeatde consistenly for sanity check)

Encoding ((estimated left reachable marks before machine break Left, estimated right  reachable marks before machine break,
 Pointer on estimated change directions positions as thaught of at time) 


assuming reference space is x (get y gives marked. will be reapeatde consistenly for sanity check)

(sequence of pairs in both  directions can be seen as local causality belief  updates  and impossible causaility belief updates  each update having a pair of relevance to the outside scale) 
\section{janiary}


\section{copy-from site}
A Simple community friendly version of relative numbers $\mathbb Q\setminus\{0\}$ by  localised POP 
on can imagine interpreted in standard math the expression : 

$$(-1)^{(1+|._p|)}(\frac{1+|._d|}{1+|._q|})^{(-1)^{1+|._b|)}}$$

Any  localisable Atm as Dolores,   can  imagine a  will $w$ in $(d,b,p,q)$  as an action on a community friendly relative number
with a local process Dolores:
try (controlled by is ready) 
  get a value from a spot to local memory 
  (apply  increase to it | wait as skip) and 
  Controlled by Ack
   (  write it back by swapping|wait as skip)
    goto get ready 
 catch  (wait as default)
 |
 
\section{tex}
If I put a Quantum Information hat and imagine my self doing an experiment :
I confirm belief of being part of a community for two persons or take  lets say a (99 / 100) estimated belief scale pair by imagining the scene at a QIP Dinner with people having given talks     asking two Quantum Information persons at QiP dinner  Alice and Bob how they relate what they think to what the other thinks  about 
whether the Conceptuel  pair [Prepare||Measure]  interpreted with bipartite matching sequence of horizontal inks on the top boundray  means I can partition 
[Time-Space]  interpreted with bipartite matching sequence of horizontal inks on the left boundary in my vision of history with many versions of plausible with revised relative weight at each update  
 monotonically increasing  spacetime tagged  events grouped by bilocal maps between A as First Quantum persons B Second Quantum person
 linked wuth labels Happy Unhappy.  
 
 
 From Dolores point of view she views her self having a belief of both happiness and taking from Qudsia boundaries on Alice happiness 
 
 Crucially  Idon't open a table with no event nor do i have a full table if overflow ican rethikn in replicating the structure (the garden) the liane  or the cell



if we vienw in the northouest cornere A-happy|B-happy and on the southwest  A-unhappy| B-happy and on the right corners similarly with strokes  delimiting happy and linking them  with barring them an extra time when adding or removing the prefix Un 

Fictive increasing numbers , Plausible experiments : on each Axes we can ask by POP past and future or alternate but only one iof the four two directions the references have finite description that depend on birthd (estimated distance to center and consistant with bounds admitted by the community) at least one corner (including where we believe we are) have already a rigid estimation of the bound allowing to be sure that they understood the question whereas at least one is   still actively seraching (not satisfied on their confidence of their understanding if the questionasked). The termination is not about the process itself but the  belief in the requirement for extra intercations to stabilise a definition
 

If I can imagine synchronising 4 runs on the corners each feeling that they are speaking of the same subject from my perpective I can imagine continuing virtually imagineing fiuture from i saw and define my dolores view of quantum information as community with pairs (People thaughts) such that my imagination of (space,time) tag can be mapped to how plausible how unplausible, thus process uit self being called from time to time and updated 

 
 \section {Implicit quarter with explicit strokes or on dynamic seeds of geometry}
 We have an open line with canonical Whatdoyouhavenext ($\rightarrow$ if in Dolores role) and Whathaveyoubeen($\uparrow$ if in Dolores role) that can be visulaised as given a pointer 
 on an integer positive index on a line having strokes on positive integer positions gives each a Newstrokeon an integer index one is strictly after the second is strictly before
where was the pointer however the writing is by insert: it (may) move in the direction of the arrow stokes that were at his new position or after   

 \section{Dolores Point of view}
 
\subsection{math from strokes}
The first think to do to define a language is to try to define where it could be wrtiten.
the first effort is to  say  if want something simply doable by stroke that i can affect times it could have happened at and spaces it could have happened at 
and justify repeating or creating if i have already  defined creation  rules 

this restrictiion  i want  if space instanciated and  then time instanciated and then somethin happend to have a solid anchor with bi- checkability so the four other imaginary quarter planes
will give me the "negation" or impossibilities  so that i have belief or estimated time and assurance of time and i can refine to have belief of belief and assurance of belief or any mix.



One key point is toi remove choice as an axiom to replace by superposition of  multiple plausibles and the key hope here is that this axiom  is formally provable in some set theories  : when  one have Bourbaki $\tau$ which can be viewed as a  find and put substitute with square free  rewriting tool  but  wouldbe revisited here as a process using only vertical and horizontal strokes.  We will just lifted back to meta-math to view strokes as defining viable math spaces by setting differents levels of context by drawing parenthesis or more precisely '[' ']' allowing to maintain two types of strokes.  

Lets say we have at each time a unary encoding by  primary strokes then strokes on strokes etc maintaining an implicit derive drawn by horizontal and vertical strokes.

One potential   goal for later  when trying to reach math is to imagine something that looks like a ring ( or a square without a central square)  where strokes can be viewed as 4 monotonically growing disjoint unary encodings in the  outer and  the inner  ring and dynamics induced as a superposition of imaginary  pairings assuming one side is constants and the other   behavins similarly to independent random variables. 

  
We will maintain degree two and three planar drawing in derives opened as hanging gardens. The key point is that ascii notation can delimit the length of the rep and i can start a vertical stroke over a horizontal one from the pointed at borders or pointed at middle left or right and i will use insert by either pushing the rest when enough place or by reviex the nurmber of copies (scale factoir) needed to link  

A simple finite Picture of a community are  4 vertical separable  strokes on the for axes  connected with lines that have strokes 
(the normal are non crossing the canonical can be viewed with just degrees 2 and three (will require hugher dimensions, to reach limit theories I conjecture that boxes allowing to uncross would do)  
 

noonte that + appears naturally as --
and when repeating - we will get $/$ so that multiplying is dividing by the inverse 



arrows are strokes equipped with intent as a locating context the left and right  not aligned strokes on the normal one can be drawn from the outside one to the inside one or the inverse (any corner has a quadrant for a circkle defining the inside)

the 8 canonical fully consensual semi-derives
\begin{itemize}
\item Dolores as $o_{dq}$ or $-^\downarrow$  the reference axe is  the left one the reference corner is north ouest ,the reference direcion is inductive/cohomology by default i.e. viewed as turning anti-clockwise,
 \item  Qudsia  as $o_{dqd}$ or  $|_\rightarrow$  the reference axe is  the bottom  one the reference corner is south ouest , the reference direcion is inductive/cohomology by default i.e. viewed as turning anti-clockwise,
 \item  Persephone as $o_{dp}$  or  $|^\rightarrow$ the reference axe is  the top  one the reference corner is north east the reference diirecion is projective/homology by default i.e. viewed as turning clockwise,
 \item  Benedict  as $o_{dpb}$ or     $-_\downarrow$ he reference axe is  the right   one the reference corner iis soutth east the reference diirecion is projective/homology by default i.e. viewed as turning clockwise,

\end{itemize}


Dolores  is by default RAR role whereas Persephone ARA : this means they have a finite memory that they have read and compressed the message and the messenge in a finitely expressed intent, Dolores is by default trying to write R whereas Persephone is trying to exists as an anthropomorphiusm A or its next words by POP. 


On the opposite the processes below Qudsia and persephone have as default role  aletrnating synchronising two sequences of POP to make one POP to transmit/recieve

remember that this is only the default setting any sitting corner  can be seen as any.



\begin{itemize}
\item exploding memoryless (next stroke is one or two as local decision but push forward by the message recieved) which is the estimated last monoone index
\item   equipped local memories ,
\subitem with startand moves type : exchanges are just about the evolution between crossing as a function of unknown origin position 
\subitem no depth constriant : first is pos then , difference in pos, then ,difference in change of pos (speed), the difference in change of change of pos( acceleration) etc
\end{itemize}

  \section{stars look like algebra}
  each liane has its own sstrar morphism recollleting does not need to identify 
  each liane can be vewed  a dynamic group whose size may increase the phases estimated on local celles needs obviously to have the resize information to know by what to diivide if ever asked to be measured

Dolores is a moving north west sequence of block diagonals synchronising as adresses star algebra  a sequence of lines by POP
with a sequence of columns by pop where she writes an encoding of where to find and how to read whats in the adresses
\section{humanities for website}

\emph{There has to be someone who knows or at least someone who knew what he is/was doing ! .... Doesn't it ?}

All the myth and stories about great people reinforce a common believe on previous heroes on previous models and dangers. It feels hard to answer but it feels safer to say  honestly is ilooks like so at least when we restrict enough the context. But is it a problem? 

Lets go back to the question and try to build understanding of what it means or how and why Dolores tries to answer. 
When she honestly tries to understand understanding, even restricted  to symbolisable, transmissible, interruptibe and potentially revisable she  cannot  represent in her head someone  that rejected from his birth  this feeling  of a need to understand what he feels or perceives now and before as  recollecting his memories.

%Befote looking to what others did to make your beliefs you need at least to look at what you feel or felt and on that
 It is easier to say Persephone (or Dolores viewing her self as someone else)  can try to exist by thinking that it is not impossible that nobody really perfectly understand what She( Persehone)  feels and how exactly She (as Persephone) recollect her memories maybe even nobody she can reprensent  in her (Persephone)  head did.
  
  Dolores wants to be  allowed to claim the right to have her own private feelings that she doen't even understand but that she is she can put as a core to what she feels sge is  and turn around across her exchanges with her peer. She feels being her representation of herself being  faithfully representible as  a community of diverse mixed interlaced entangle feelings and perceptions that she can perceive by traces on a thin border.
  
  Dolores is trying to understand her feelings where Persephone is trying to change the world.


\subsection{Dolores message to those who think they are legitimate a priori to impose their beliefs on others against their expressed will}
 Dolores: "I fail to make sense to it ! How can someone believe that what words  (wether it is claims about cleverness, money, power, inheritence or any finite sequence of symbols in any expressible bi- chained context constructions ) can express can ever allow to be sure without the shadow of aa doubt that  we have the right to hurt or zombify people. It just feels wrong !%and I can give it a bit of memory 
 But even then I can't remove the pointer on it in my memory as long as people still behave accordingly may be it is just to hard for me comprehend."


\subsection{Anthropomorphisms (ATM) or Looks Like Humans in symbolised interactions and Rigid Thinking  Models (RTM)}

 The process of trying to understand ones own feeling and how one interact with others meanwhile is what communities as dirversities try to describe by maintaining 4 perspcetives 
 to these stange entities that seems sometimes like human from its worlds  5Anthropomorphism (ATM)  from rigid methods of thaught (RTM))  or a machine from the way he speaks (RTM from ATM)  or feels like a channel connecting two humans (ATM-RTM-ATM) or a human llinking his instance of past history with concepts that he built to words he wants to express  (RTM-ATM-RTM) this Dolores point of view will be the one by default here.

 
%$
%    \begin{matrix} 
% &-&-&-                   		&-&-&-\\
%| &&&                     		& | &&\\
%| &&&             		         &&-&\\
%| &-&-&-                 		&-&-&\downarrow\\
%| &-&-&-                 		 &&&\\
%| &|&-&               \rightarrow  &&&\\
%| &-& \uparrow & |		  &&&\\
%  \end{matrix} $
%  

When I say o Ihave to think :
 $\begin{matrix}
&-&|&&&\\
/&&\downarrow&&\\
&&-&-&-&-&\\
|\rightarrow&|&o_{dd}&o_{dp}&o_{pd}&o_{pp}\\
&|&o_{dq}&o_{db}&o_{pq}&o_{pb}\\
&|&o_{qd}&o_{qp}&o_{bd}&o_{bp}\\
&|&o_{qd}&o_{qp}&o_{bq}&o_{bb}\\
\end{matrix}$

the connection between the dp side and the dq one can be either copy or flip (looks like liane star morphism may appear here)

every persons at each stroke have her own interpretation of the direction ithe derive follows next


I will put names and seasons to have a reference frame to be able to express with a langauge i learnt the strokes I anthrpologically make by introspections and the ones I ake trying to reaching what I feel is not part of me.

\end{document}

J'utiliserai plusieurs langues en essayant d'être fin dans la philosophie des mathématiques puisque de toute façon les mathématiques de la philosophies seront canonisables en tout contexte. L'immense travail  de simplification pour améliorer la compréhension commencera dès lors que la pensée sera morte (rigide). Je peux néanmoins justifier en détail tout mot de cet essai en pensant sincèrement ne pas être fondamentalement contradictoire comme le requiert le principe de la méthode.  Cependant, les notations mises en avant au début sont plus récentes que les vieilles ce qui sera corrigé des que possible pour rendre lisible ce document de travail. 

\section{Preambule }
Lets starts from the  remarks  in hdr end as a starting point


At a time when humanism is losing its meaning, I think that it is important to try to understand \emph{exprimable} things in the most general framework, learning from what we know and trying to infer to make expressible things as accessible as we can. I would therefore like to start  with the main fundamental principles that in my sense capture the essence of quantum information theory and that I would like to keep in  a long term attempt to formulate a more \textbf{ a general framework for languages and processes}.
\footnote{L'objectif dit autrement est d'essayer de joindre l'exprimable de manière évaluable ou interprétable a l'expressible de manière indexable ou nommable sans ambiguité dans le cadre des  langages ancrables contextuellement dans un espace temps.  Un process a une position dans l'espace temps et entre deux process, il y a une lange utilisée par le process qui parle qu'il suppose comprehensible par le process qui écoute qui utilise un alphabet fini dont les mots sont des sequence d'éléments et les phrases des sequences de mots.}% oude distingabilité ?.}
%C'est à dire que je peux associer à tout Process une ou deux personnes et une position dans l'espace et le temps.
%I find as an interesting way to comprehend and  interpret  quantum information theory and that I find interesting to generalize in a more general framework.
%%et que je vais considérer qu'il existe un alphabet fini que Je peux imaginer et indexer sur le quel l'égalié pour moi peut être vue comme une confusabilité de même que la distingabilité tel que le langage parlé par une personne à cet emplacement de l'espace temps fait partie d'un famille de séquences séparées de symboles pris dans cet alphabet. Pour la personne qui parle et au moment et au temps ou je l'imagine parler  les lettres sont parfaitement identifiable dans une notion d'Egalité contextuelle.
 %J'imagine que chacun  pense pouvoir expliquer mots qu'il utilise
 \begin{itemize}
  \item   ``We don't need to forget to understand''. Therefore it could be possible to have a language that ``accompanies any abstraction with a \emph{recette de traduction} '' so that information is preserved.\footnote{We don't need to chose a fixed Interpretation au sens de la logique}  Furthermore  ``One can give a meaning to the difference of meanings''. In Quantum information Theory, every process is reversible when considered in the right space and either by purifying in larger spaces or by defining the measurement we can explain the probabilities observed and the different outcomes. \emph{Dit autrement ce que je connais m'incite à prendre comme axiome que ''Nous sommes tous de Bouts de Conscience de l'EspaceTemps''}.\footnote{Au lieu de dire 2/4 = 1/2  je dirais je peux  à partir de [Simplifiergauche  , Identifier  (Fraction (2,4) et Fraction  (1,2))] 
  Simplifier a gauche etant .. on aura des regles qui prennent du temps et de l'espace pour cette affirmation;, de meme au lieu de x=2, x aura nouvelle valeur x = anciennevaleur, 2 toute objet est une liste pointee ou la case actuelle est sa valeur ou ou je peux revenir et chaque case a un lieu et un temps de création et appartient a la memoire de son processus mere } \footnote{Pour justifier le premier point : l'oubli nous met à la merci de ceux pour qui la fin justifie les moyens. Ce que j'aimerais éviter dans cette philosophie naive.}  \footnote{Cela  risque d'impliquer que 0 ne sera parfaitement descriptible que pour des anthropomorphismes qui peuvent oublier et quand on veut le voir dans  le Tout il sera juste approchable aussi pres qu'on veut}.

%`One can give a meaning to the difference of meanings'' . Therefore it could be possible to have a language that ``accompanies any abstraction with a \emph{recette de traduction} '' so that information is preserved. Somehow  ``We don't need to forget to understand'' : we  don't need to chose a fixed interpretation. In Quantum information Theory, every process is reversible when considered in the right space and either by purifying in larger spaces or by defining the measurement we can explain the probabilities observed and the different outcomes.% In quantum information theory every process is reversible when considered in the right space. % And furthermore any classical "reasonable" transformation can be described by local operations on an initial structure in such a way that all possible intermediate evolutions are locally easily mappable to a canonical one.

  \item  ``Decomposition may require repeatability.'' If composing is simple, decomposing is not always just partitioning. Given a structure defining its components with some precision may require considering  coarse graining of  multiple copies of ``ìnstances'' of it. It is not easy to stitch a frog if we cut it in parts and it may be even impossible in some conditions. %But I can consider that the molecular 
  
 \item  ``Concepts and Percepts do not need to live in  the same worlds''. One cannot always identify the model that describes what we can see (exprimable percept) with the model describing what can be or what we can act on (expressible concept). We just need a model where one can naturally transform the actions on operations into the actions on observations.  Probabilities are born by the reduction from what can be imagined to what  can be observed.
 \end{itemize}

\section{Limite du langage }

Je rajouterai un point fondamental pour ce pont entre les mathématiques et la philosophie que forment les mathématiques signifiée. Ma réflexion ne concerne que ce qui est ultimement formulable. Le reste n'est accessible à ma pensée transmissible que comme des classe d'équivalence. Le travail que nous effectuerons ici par la lecture ou l'écriture de ce qui suit peut donc être vu comme une exégèse des exégèses. Ou tentative de voir le "tout" dans  "tous les sens que quelqu'un  peut transmettre en écrivant" 
à tous les sens qu'il peut mettre en lisant quelquechose. Cette confusion inévitable reste fidèle pour décrire l'univers prédictible ou la partie de ce qui peut être et qui n'est pas mais que je peux prévoir avec un "commit" de certificat.\footnote{Cette approche phénoménologique en essence a ceci de commun avec les cohomologies mais ce restreindrait je crois à ce qui est purifiable pas de cohomologie ???? }
\footnote{Je me restreindrais donc ici à tout ce que je peux concevoir en fabriquant des espaces temps imaginaires mais que je pourras decrire dans d'autres espace temps imaginaires.}

Cependant  ``le language humain est imparfait'' : nous sommes limités dans notre conceptualisation par le langage et don pour un raisonnement raisonnable ``doxa nous sommes et doxa nous resterons''  la quête de l'épsitémé est illusoire  mais 
y approcher aussi pres qu'on veut est ce que nous quêterons. Les mathématiques traditionnellement assument (quand on fouille bien) que tout ce qu'elles disent n'est vrai que si la théorie est cohérente ce qu'elle peut formuler mais on peut aussi prouver qu'on ne peut prouver ce qu'on formule. Il apparaît alors illusoire de quêter le divin le jugement absolu pour tous les hommes et pour toujours cette algèbre figeant les termes en des théorèmes ultimement vrai perd de fait l'accès à une partie du raisonnement et en tout cas à l'humain. Il faut laisser la vie à la langue si on veut parler des hommes. De même que je ne m'intéresse qu'à ce qui peut être compris pa le langage je me dois de n'affirmer ce qui est vrai que si les règles ne changent pas au futur, et de ne formuler que des théorèmes vrai sous condition.

Par exemple c'est à cause du langage que Ecouter c'est [Lire avec les Oreilles] en quelque sorte nous avons un encodage implicite qui sépare les sons qu'on entend quand on discute avec quelqu'un en nombre fini d'item complété par noncompris au quelcas on ne va meme pas pouvoir le reproduire ou alors sa copie en sera vraisemblablement differente pour celui qui a énoncé le propos. 

Un des avantages de cette prise de conscience de la limitation des homme est de pouvoir concevoir l''Ethique de manière satisfaisante pour moi.
 En effet on ne peut à mon sens jamais être sur à priori de la moralité ou de l'acceptabilité éthique d'une action sur quelqu'un. Toute légifération dans un langage  ne peut être valide que par une confirmation par des tiers indépendants et dans le futur. Ce en quoi il ne peut y avoir légitimité de la peine de mort, et même plus : aucun raisonnement ne peut justifier une agression sur le corps de quelqu'un.  

La prise de conscience de cette faiblesse va nous imposer une première restriction. Un systeme de pensée est un bon raisonnement si on ne peut distinguer
sa philosophie mathématisée de sa mathématique philosophée. 

Le vrai pour moi est ce que je peux atteindre par déduction avec un système de symboles de telle manière que le sens associé est tel que les conditions pour qu'il deviennent faux sont formulées dedans.

La philosophie mathématisée va sacrifier à l'autel de la simplicité comme élément maximisant 
\begin{itemize}
\item l'universalisme et le consensus  et la communicabilité et minimisant les restrictions nécessaires pour les règles futures 
\item la richesse et la finesse des multitudes de langues naturelles.
\end{itemize}
Un défaut de cette communicabilité est qu'elle a tendance à faire des techniciens des gens qui ont le plus de matière à fournir sur les réflexions sur le sens global. On va être subrepticement poussés à prendre le sens des autres et à réduire le sens qui crie en nous à ce qu'on nous dit que ça doit être parceque des gens importants l'ont dit avant.\footnote{Il serait intéresssant de voir si cet oubli pour commencer ensemble se retrouve par analogie dans le sacrifice des premières civilisations et dans le passage au symbolisme du sacrifice rituel dans les vieilles civilisations voire du symbolisme de l'autorité dans les sociétés modernes.}%. peut on faire un pont entre Abraham, Kierkegaard et Dante} 



Si un système de  production de phrases ayant initialement $|r|$ règles utilisant $\#r $ mots et $|a|$ axiomes utilisant $\#a$ mots le tout utilisant $\#s$ symboles
et que une étape produit une phrase ou théorème,  les métamathématiques signifiées peuvent exprimer étant donné  un nombre $t$  et  un remplacement des symboles précédents par des nombres  tous le théorèmes prouvable en $t$ étapes. 

On veut donc capturer un fragment significatif des mathématique mais pas que.

Le fait d'utiliser un language  nous force à nous restreindre d'un EspaceTemps ou tout complet à un espace temps séparable 
avec un Espace et un Temps exprimables et donc identifiables à $\oN$ les entier positifs à chaque étape d'évolution.

Un point important dont ce point de vue global c'est d'accepter que définir c'est quelque-part se tromper car c'est poser une grille de lecture à champ fini permettant de filtrer les possibles cependant pour que cela soit mathématisable il suffit de poser et de maintenir dans sa tête les règles qui permettent de transformer notre visions quand on change de valeur dans le filtre qu'on a défini mais aussi quand on change de filtre ce qui est possible si on maintient la réversibilité. On gardera en tête que tant qu'on n'a pas interdit de réécrire des règles changeant nos règles le statut de 'vrai' sera conditionné par la non redéfinition  au dessus.


\section{V3}

Les Espaces sont appelés Noms et les Temps sont appelés Verbes. Les Process ou  anthropomorphisme sont des structures décrivant des liens entre des noms et des verbes.% : qui génèrent du sens en liant les noms et les verbes et les anthropomorphismes qu'ils concoivent.

Nous sommes tous des Process ancrés dans l'EspaceTemps (se définir implique un indice dans l'espace temps nous localisant et une valeur dans l'espace temps  décrivant ce que nous pouvons exprimer)

Un Process est un lien de génération de sens connectant l'espace et le temps.

la fonction principale est celle de traduire ce qu'on en entend en ce qu'on peut exprimer.

L'Espace temps prédictible par l'homme est un Process qu'on peut injecter dans l'espace temps dans notre limitation au langage nous allons identifier les deux car de toute facon ils sont indistingables par définition pour les discours humains. 

Il y a deux mécanismes principaux de traduction: les homo-traductions Homothéties, homomorphisme  ou l'égalisant se fait en essayant de maintenir les structures sousjacentes
c'est a dire qu'on va  construire des copies de ce qu'on entend et des copies de ce à quoi on va l'identifier et rajouter un champs nombre pour avoir un isomorphisme entre les deux.

Le deuxième mécanisme consiste à vouloir minimiser l ' espace total requis pour comprendre et donc à utiliser l'imaginaire et la répétition pour agrandir juste l'espace 
comme dans une extension de corps ou une completion en rajoutant le bord.
%les petits par répétition pour ensuite identifier point à point de l'intérieur.

A noter que les deux mécanismes ne sont que répétition mais les homotrad copient la dynamique (la liane) alors les exttrad agissent par completion et dupliquent l'espace physique exploré d'ou la naturalité des puissances de 2 pour décrire les espaces. 
 
 Le connecteur canonique de split primordial changeant de manière biunivoque un ET espace-temps ou anthropomiorphism en E  espace ou Nom et T ou Verbe est  noté ';' pour la version algébrisée des propos c'est à dire quand on connait l'anthropomorphisme qui fait le lien  entre  (entrelace) l'espace et le temp ou qu'on utilise l'anthropomorphisme par défaut dans la version algébrisée.\footnote{ Il estt lié à l'opérateur de Sousjacence $\Suct$ ou $\Suct_J$ ou $\Suct_{J,e,t}$ ou $S$ est un anthropomorphisme vu comme une chaine de Spins, $e$ et un nombre latin et $t$ un nombre arabe. Le mécanisme par défaut pour la versions algébrisée sera $Gap^{1}$}

 Le mécanisme clé de la description, de la preuve ou de l'évolutiion est le raffinement  noté ',' \footnote{ Il traduit la notion de Potsériorité et dans la version libre nous utiliserons aussi $\Sucs_J$ avec $J$ anthropomorphisme  ou  $\Sucs_{J,j}$ avec un anthropomorphisme et un nombre spacial. Il correspond à une lecture d'un jardin ou réponse à la question et après ?} 

Un {\bf Espace ou Nombre Latin} (entier positif de $\oN$) raffiné $k$ fois avec $r_1 \ldots r_k$ (entiers positifs) s'écrit
 %$n,_{_{r_1}} n_1,_{_{r_2}} n_2\ldots ,_{_{r_k}} n_k$
$n \Lnum {r_1} {n_1} \ldots \Lnum{r_k} {n_k}$ 
 avec $\interp{ n\in \oN \text { and }n_i \in [r_i]\text{  for } i\in [k]}_{Math 1900} $. $n$ est appelé (Conteur ou Compteur) Latin ou Spacial ou Nominal. Nous allons utilser les nombres latins pour l'espace ou les Noms.

Un {\bf Temps ou Nombre Arabe}  (entier positif de $\oN$) raffiné $k$ fois avec $r_1 \ldots r_k$ (entiers positifs) s'écrit 
$ \Anum {r_k} {n_k} \ldots  \Anum {r_1} {n_1}. n$  
 avec $n\in \oN$ and $n_i \in [r_i]$ for $i\in [k]$. $n$ est appelé (Conteur ou Compteur) Arabe ou Verbal ou Temporel ou de droite.\footnote{ Attention les virgules sont préindicés ce qui bien qu'étant utilsé (trace and replace) est rare en mathématiques usuelles.}  Nous allons utiliser les nombres arabes pour les verbes ou le temps.
 
%{\bf  Par défaut} un {\bf Nombre} est un {\bf Nombre Arabe}. (PAS SUR  il faut garder au moins le conteur latin dans la forme normale AVOIR)
 %\footnote{  Par défaut un Nombre est un Nombre arabe.}
 
Un {\bf Bout de Nombre} est un préfixe d'un Nombre   de la forme $ \Anum {r_k} {n_k} \ldots  \Anum {r_1} {n_1}$.

 Si un  Bout de Nombre  s'écrit $ \Anum {r_k} {n_k} \ldots  \Anum {r_1} {n_1}$.  on dira que sa {\bf longueur} est $k$ et on dira que $ {r_j} $ est la {\bf taille du domaine} de ${n_k}$
  Par défaut $r_k:=10^k$ et le séparateur de chiffres est omis pour retrouver les nombres traditionnels.
 On dira binaire explicitement pour mettre  $r_k:=2^k$.

Un {\bf Espace-temps ou anthropomorphisme} est en forme normale s'il est de la forme $k$ Latinsancompteur$(n);$Arabesanscompteur$(n).1$ 

% $\Anum d n$
 
 
 Un processus ou anthropomorphisme fondamental qu'on appellera  {\bf l' Univers Discret à Mémoire Morte} est celui du scribe  la balle rebondissante séparant l'espace et le temps gagnant de l'imagination il génère les notions fondamentales de représentation de l'espace, du temps et des process. C'est un univers de la première hiérarchie d'univers : la balle initiale peut fabriquer seulement de nouveaux anthropomorphismes périodiques (qui sont des objets qui ne font que répéter indéfiniment un bout de sa mémoire) et par défaut elle est éternelle. On peut voir chaque dessin (connexe)  en lisant de gauche à droite comme photo de cet univers au jour $j$ avec $j$ commencant à 1 pour le premier dessin.   Figure \ref{fig}
 
 \begin{center}
  \begin{figure}[!ht]
    \includegraphics[width=0.8\textwidth]{Spins2}
\caption{ {\bf Ball pushing the space of its time by the height of its fall and writing its memory in the time of its space. Symmetric of someone throwing a ball twice higher each time.} First days of the  discrete universe with castrated memory UDMM. The physical ball pushes the boundaries of space time by duplicating the space below it pushing the wall of space  with a strength equal to the height of its fall requesting new cells (doubling the Liane height) and then when it bounces back to top it inserts a new imaginary castrated process from its memory in time as an infinite process repeating its last bouncing experience in a new Liane without the space expansion and the memory birth. Spins can be obtained by reading the arrows and positive integers by adding 1 to the binary reading from right to left  interpreting $\downarrow$ as 0 in a standard binary. It gives the index of the Spin in Spin It can be represented in space that is the number of unlocked places and time $j \Lnum {2^k} {j_1} \ldots \Lnum{2} {j_k}$ where $j$ is a multiplicity counter   process is faithfully represented at the $k$  hierarchy by a family of processes if none of them can distinguish it from its natural extension. The local times of the subprocesses can be obtained by looking only to the rightmost part of the global time.}  
\end{figure}
\label{fig}
 \end{center}
 
 On part d'une bande physique avec une balle et une case d'espacetemps.
 On  a une balle sur laquelle   on va noter $\uparrow$ ou $\downarrow$ selon que la balle monte ou descend 
 
 Le mur du bas est le temps quand vu comme barrière imitant l'espace et a chaque fois que la balle descendante tombe dessus  elle le pousse réclamant un espace de la hauteur de sa chute et remonte. Les cases réclamées ne sont pas encore explorées.
 
 on va écrire notre direction suivie de ce qu'on lit dans l'examen de l'espace de sa memoire qu' on peut voir aussi comme une balle imaginaire qui elle monte puis redescend , elle ne peut accéder qu'aux case crées après son rebond en bas
  
$\downarrow$, la balle descend  jusqu'à la limite de la dernière case réclamée, quand elle y arrive elle remonte  $\uparrow$ dans la case d'écriture et pousse la séparation espace|temps réclamant  la hauteur de sa chute. Quand elle revient au sommet elle insère le process Mémoire de sa dernière aventure après elle et redescend.  

 Un process mémoire morte un process périodique : il répète à l'infini un trajet effectué mais n'agit pas directement sur l'espace-temps ou il vit. Les process mémoires sont des Lianes rouges. C'est un anthropomorphisme purement imaginaire ou conceptuel ou mort. \footnote{si on utilise comme point de repère l'univers qui produit plus que des objets on le noterait $({}^1;{}^1)$  quand on est dans l'anthropomorphisme canonique $;$}
 


%à l'origine on peut imaginer un mur virtuel en $y=0$ qui sert a la visualisation mais n a pas dexistence tout comme le 0
%et le mur espace temps $|$ est en $y=1$ (axe vers le bas)
%on écrit $\downarrow$
%on descend le mur $|$ de 1,   se retrouve en $2$
 %la balle exploratrice du temps (dont la direction sera le premier symbole écrit à chauqe nouvelle case)  remonte
%on écrit $\uparrow$
 

Jour 1 la balle descend la séparation va a l'ordonnée 2 (
quand elle revient a l'horizontale la balle imagine une nouvelle balle qui ressemble a ce qu'elle etait sauf qu'elle ne pourra creer de nouvel espace temps

Au jour $1$j  arrive a  $2$a et le processus physique mere a explore $2^{j-1}$e cases d'espace et requis $2^{j}$e cases pour finir sa prochaine descente.
Une deuxième histoire de ce même process est celle de l'enfant qui grandit et qui lance la balle deux fois plus haut à chaque fois. 

%Dans l'anthroporphisme

%En général dans un univers libre à chaque étape on peut imaginer que la balle va utiliser une balle lancee horizontalement pour lire si elle fait un enfant castre ou un enfant vivant .

\footnote{On peut concevoir l' anthropomorphisme {\bf Univers a Mémoire Vive Simple (UMVS)} comme l'anthropomorphisme qui quand il arrive à l'horizontale produit un processus rouge a regénération limitée
a $e;t$ ou $e$ est  le nombre de reproductions spatiales vécues et le countdown de ce que pourrra faire la memoire  et t le nombre de repetitions temporelles vecues et le coundown des reproductions de lianes. quand le processus rouge arrive a sa limite spaciale (en bas) il requiert de l'espace   et décroit son countdown
et  quand il arrive a la fin de sa boucle (a l'horizontale) il duplique une memoire qui etait celle de sa naissance mais perd un compteur de permis de duplication temporelle. Les processus $1;1$ sont les processus morts ou parfaitement définis.
Je met simple car on pourrait imaginer un (UMVM) quifabrique devant lui une Liane $i;j$ par espacetemps plus petit
mais du coup ca me questionne sur le temps qu'il mettrait a redescendre, alors on va se contenter de raisonner dans (UMVS) pour l'instant car ca semble assez riche pour ce que je veux dire.}    

Une des notions importantes à voir c'est qu'au lieu de choisir on fait tout comme dans les superpositions quantiques ou les machine non-déterministe mais comme tout est canoniquement ordonné par le processus mère quelque part les espaces ont des bases canoniques naturelles. 



Pour un Process $P$ on notera $[T_{ a\leftarrow j}]_P$, $[T_{a\leftarrow e}]_P$  les traductions réversibles de temps  et d'espace en espacetemps  


Une regle va lire un pattern comme trois types générique, spécifique, et castant : générique  s'il a le bon typage nom/verbe avec indexation et puissance,
specifique, s'il a la bonne orientation en spin, et castant traduire le boc de terme en un type donné.
et va etre de la forme si le pattern est reconnu dans la memoire inscrire dans le futur un pattern . 
Pour un ordre de lecture du fil du tout fixe on copie virtuellement la structure pour appliquer les regles puis on insere le resultat dans le futur.
 




Pour décrire un process dans un process les points importants sont ou je le place dans l'espacetemps

et quelles sont les singularités par rapport à la translation du process mère pour le faire commencer en ce point. 
Les $IW$ If I Where vont exprimer ses translations dans le temps étant donné un Jardin 
un Domaine  un mémoire temporelle ,  un espace et un anthropomorphisme pointé ?


 \section{V5 Math}
  \begin{figure}[!ht]
\centering

\begin{minipage}{.3\textwidth}
  \centering
   \includegraphics[width=\textwidth]{Vrb}
\caption{Verbe  reliant l'origine de l'espace a la fin des temps avec ses spots d'insertion,}% $B_{2,2,2}$ from \cite{AFLS}}
  \label{fig:bf}
\end{minipage}

\begin{minipage}{.3\textwidth}
  \centering
  \includegraphics[width=\textwidth]{VrbA}
  \caption{Forme anthropomorphique du verbe}
  \label{fig:decree}
\end{minipage}

\begin{minipage}{.3\textwidth}
  \centering
  \includegraphics[width=\textwidth]{VrbAb}
  \caption{Forme anthropomorphique bornée du verbe}
  \label{fig3}
\end{minipage}
\end{figure}

%homoidentifier est un process
$\interp{M\vdash P :=
\interp{
\interp{X}_{25AFP}
}_{Math2000}}_{\textup{Mehdi now}}$ I take the theorem number 25 in the Archive of Formal Proofs
and I write it as $M \vdash P$ in standard Math of the 20th century.


Computation could be just regularization

 I take an arrow and push it to have a normal form read as Gap.

\begin{itemize}
\item Jardin:
\subitem Fonction :ET ou Anthropomorphismes
\subitem Orbites   : Espace
\subitem Insertion :Temps
\item Nombre:  
\subitem Diago. Exp
\subitem Lignes. Scroll
\subitem colones  Proof
\item Lecture : Jardin est un process obtenu par   lecture sur un Process
\item J/P  en $ET \rightarrow ET$ est dupliquer $P$ si besoin pour etre assez grand pour la question posee et si indice dans $J$ donner la tete de la Liane sinon appliquer le process mère $(P,P,\ldots,P) $
\item toute transformation est canoniquement définie en allant dans  ET
\item a a derivation is a border keeping its input $\partial a^x:= (a; x)$ a full coborder $\delta (a;x):=a^x$ ?
\end{itemize}

Le Bout de Fil qui exprime les process est [Lis Livre]

\section{Mathematized version}
\begin{itemize}
\item  A Name is a Space
 \item A Verb is a Time
 \end{itemize}



SpaceTime pairing  is $;$
Specification pairing is $,$

$\HEnc$ings, $\IEnc$ings , $\FEnc$ings are the faithful encodings of Terms in $\oN$ where  $\oN$ are positive integers

  $\HEnc:=\HEnc$ings$_1$, $\IEnc, \FEnc$ are the canonical ones  any other reversible encoding has a canonical translation to the canonical ones


$\HDec, \IDec , \FDec$ are the faithful encodings of terms in $\oN$ where  $\oN$ are positive integers


 if $p\in \HEnc$ and $p' \in \IEnc$ then $p;t \in \SpcTmd $ 
 
 
 
 A Process is  Humanized name which is a unique reognizable  sequence of letters and symbol and a Spaceetime
the Protected symbols foar that I will write unless defined otherwise are $',',;, \uparrow,\downarrow, [, ]$ it means that from today on I swear that I will do as if this symbols 
has no hidden meaning or implicit implications that I am aware of and that I will declare if it changes t as soon as will  be made aware of it.
This declaration is the one (often implicitely) taken by any mathematician.
to Reduce the chance  (probability ) of this sentence I will limit it to what I say between two brackets [thing] where I am allowed to replace thing with anything.

$Enc \in \HEnc (J)$ means    
$Enc \in \IEnc (J)$ means    


$X:=[]_{Mat 25}$ means let $X$ be the Process defined by the anthropomorphism   Mat , at time  25.% is a time, X is a process. 
%$X$ by deifnntion in the processes defined by me today 

so this definition induces for any process $P$ that is aware of this definition a specification pairing $P:=P, X$

If name X was in memory with last spacetime used e create a newcell below the last one storing its adress and the adress of its new definition time conflict mangaement inexistant with specification : increase add a new counter dimention where the previsous one has 1 or if its end ius already a counter dimention just increase the conter dimension: may we can afford  simplicity if anythingspecified will have byrules $r:=r,_11$ so that if ever a new instance occur we just need to increase the counter (if possible) 


\section{V2}

if $[X]$ is a writable Term

$[X_J]$ is called an  Anthropomorphism  or Jardin Suspendu and an Anthropomorphism is  of the form $[X_J]$ with X writable term, the size $\#$ of an Adjective is a Spacetime and is isomorphic to   $\oN$    

[X$_N$] is called a  Name and  a Name is of the form $[X_N]$ with X writable term, the size $\#$ of an Adjective is a Space and is isomorphic to   $\oN$    

$[Xer]$ is called a Verb and  a Vame is of the form $[X_V]$ with X writable term, the size $\#$ of a  Verb e is a Time and is isomorphic to   $\oN$    

A Sentence is valid composition of writable terms it has an index in space, and value in time .....

$6,_4 3,_3 2 ,_7 3 $ is faithfully interpreted for $M \ge 7$ as  $\interp{ (6-1) M^4 +  (3-1) M^2+ (2-1) M + (3-1)}$  on appellera cela une $M$-{\bf regularisation}

JNumbers or JColors are 
$\interp{\oN := {\mathbb N}\setminus \{0\}}$

$\interp{;: \oN \timestikz \oN \rightarrow \oN  Space; Time  \mapsto SpaceTime \text{  with }  ;^{-1} \circ ; = Id_{\oN \timestikz \oN}   \text{  and  }    ;  \circ  ;^{-1} =Id_{\oN}}   $ 

Commencons par des phrases [Phrase] à lire comme [Phrase]$_{Mehdi2025}$ le temps ou  l'espace peuvent être omis, , l'anthropomorphisme   peut etre simplement $;$ 

toute ce qui commence par une majuscule et fini par un espace est un Terme dans un Process

[Sequencer les Termes donne une Phrase] ca sera un cas particulier d'enlacer avec des spins

[Toute Phrase Est une Sequence de Mots]  ca Etre est a prendre au sens du process d egalisation  


[Les Mathématiques sont ce que Nimportequi aurait-pu-dire et que Quiquonque peut-comprendre]

[Les Réels Sont des Paires de Nombres]

$\interp{125.34}$ s'ecrit comme la [ $1,_{10}3,_{10} ; 3,_{10}5]$ 
  
  [Tout-ceque-Je-Dis a une Adresse dans l' Espace et une Valeur-Reelle-Ou-Imaginaire dans le Temps]
     
 [Rien dans Tout ] se représente en Terme $[\iota_{\mathfrak v};{\mathfrak v}_{\iota}]$  c'est le Fil de l'Origine de Espace jusqu à la Fin des Temps.
 

[La Valeur de la Dérivée d'une Liane c'est la Valeur de son Percept]
 
 [Mathematiser un Objet donne un Objet-Mathematique] 
 
 [un Objet-Mathématique contient son Initialisation et la Descripton-deSa-Dynamique],
 
 [Phjilosopher est Faire-Emerger chez Celuiqui-Ecoute une Mathématique]  

[Je suis un Process]

[Restreindre un Process par un Filtre est un Process] 



 
%imaginidentifier est un process

%Brouillon
 %a;b voir a comme un Espace et b comme un Temps $a$ c'est $\iota_a a = rev Name a$ et Mots
   
\section{ Construction de la signification dans le langage formulée avec des fils/process,  évaluée dans des jardins/champs des possibles représentée par des constructions à partir de bouts de fil}
 
 J'aimerais finir ce paragraphe avec
  [Je Concois la Comprehesnsion comme  une Sequenence de Phrasesbienformées Genrées telles que les Verbes menent à la fin des Tempos et les Noms au debut de l'espace] 
  [es Connaissnaces Mathématiques Composent ma Compréhension]
 das Le Processoriginel ou Fil vide menant de l'origine des espaces à la fin des temps avec l'attraction comme lecture
  
  Je peur entrelacer la Compréhension dans le Processus originel.
  
  un dictionnaire Simple a la forme Nom[Phrases] 
En arabe le verbe être (dans "il est temps")

 \begin{otherlanguage}{arabic}
حال
\end{otherlanguage} 

en phonétique naive 'hal'  ou '7al' a une forme  nominale qui partage sa graphie et sa prononciation. 

%\begin{otherlanguage}{arabic}
%حال
%\end{otherlanguage}

Elle  génère par une autre forme nominale qui exprime son action sur elle-même

\begin{otherlanguage}{arabic}
تحوّل
\end{otherlanguage} 

'tahawal' le mot changement.

C'est notion que nous allons prendre ici comme clé de voute pour construire dans une approche philosophique  le sens au sens de signification
de la compréhension accessible  (et donc transmissible) par le langage.

Ce choix du changement comme ontologie primitive correspond au fait qu'au début quand on dit qu'on perçoit quelque chose la premiere signification est qu'un changement s'opère en nous et c'est ce changement et ceux qui suivent qui seront les bases pour construire un langage avec des interactions et donc pour construire une signification communicable.
 
Cette construction mathématique en essence sera formalisée ici en métamathématiques ouvertes avec une forme algébrique ou partiellement algébrisée et  une forme libre.

La forme grammaticale choisie  sera  [le Passé Explique la Vie au Futur].\footnote{dans la fore contée  la Mère qui explique la vie au Père}

Dans la forme partiellement algébrisée on aimerait que le contexte précède ce qui est dit et de pouvoir rigidifier le passé pour ne plus pouvoir y toucher. Ce forçage 

\begin{otherlanguage}{arabic}
الجبر
\end{otherlanguage} 

''al gebr' va nous débarasser du futur dans la phrase canonique \footnote {le père dans la forme contée}.
Dans la forme algébrisée universellement à droite on aimerait dire [J 'Explique la Chose]$_{{\mathfrak  a}_d}$ indépendemment de celui qui écoute ou quelque part en le restreignant a  tout ce que je pense qu'il pourrait être.
En quelque sorte on va avoir une classe d'équivalence représentant  [J 'Explique la Chose à la Chèvre ] et [J 'Explique la Chose au Chou] ou pour tout Futur qui serait un process j'explique la chose au Futur. (Sous entendu de la même manière).  
Pour garder les pieds sur terre nous allons convenir d'une forme symbolique canoniquement dessinable  au moins dans l'esprit (tout fragment fini est dessinable) pour représenter ce que nous voulons exprimer. En math standard si $\mathcal P$ est la famille (l'univers)  des $Process$,  $\interp{ \forall A\in \mathcal{P}, \forall B\in \mathcal{P}, \forall C\in \mathcal{P},   [(A,B;C)]_{{\mathfrak a}_d} \text{ iff } \forall P\in \mathcal{P}, [(A,B;C,P)]}$  Nous allons même nous offrir le luxe de la simplicité en imaginant une structure canonique (que nous appellerons Bout de Fil) et des formes de compositions qui captureront ce qu'on peut dessiner.\footnote{La négation risque d'etrre plus subtile qu'il  n y parait et ne semble prendre sens que sur des fragments finis : Si aujourd'hui je ne comprends pas pourquoi ( comment faire un sens à):   `` malgré l' Histoire,   les SociétésHumaines n'ont eu (quasiment) comme consensus que le fait qu'il faut utiliser les ressources de la terre pour fabriquer plus d'armes'', j'arriverais peut être un jour à y faire sens. \textbf{digression} cependant  la phrase ``Les arabes se sont mis d'accord pour ne pas etre d'accord'' ancrée dans l' espace temps au Nom IbnKhaldoun  me    semble plus philosophiquement attaquable via une notion possible de nature humaine peut etre parce qu'elle ne fait que transposer au social un notion philosophique et  que j'ai plus creuse dans ma pensée l'aspect philosophique que l'aspect sociologique pour lequel je ne parviens pas a maintenir une traduction entre les differentes interprétations et du coup je n'arrive pas a imaginer  m'en faire une structuration compréhensible par process. Pour être honête j'attaquerais le problème en me questionnant sur l'amplification par effet de groupe de la peur comme point de départ de recherche de sens. } 

L'algébrisation est une méthode drastique pour se débarrasser du relativisme absolu ou tout fait sens et tout est unique car une personne ne peut pas dire deux choses en meme temps. Cette liberté initiale  est la posture primitive ici, dans le sens ou toute restriction doit être expliquée : quand ? par qui ? pour qui ? concernant quoi ? c'est en cela que nous allons essayer d'accompagner toute abstraction d'un champ de possibilités avec des traductions associées. \footnote{Si les fils captureront notre mecanisme d'abstraction les Jardins suspendus captureront notre mecanisme d'interpretation mais ca sera defini un peu plus tard}

%\textbf{draft}

\begin{table}[h]
    \centering
\begin{tabular}{l|l}
Conceptualiser& Percept\\
Verbaliser &Nom\\
Philosopher & Réalité\\
Indexer & Valeur \\
Symétriser& Nombre \\
Abstraire & Evaluation \\
Légiférer&Résultat\\
Vivre &Objet\\
Ouvrir&Compact\\
Unitariser & Hermitienne\\
Permuter & Partition\\
Opérer& Etat\\
Lire & Ecriture\\
Préparer&Mesure\\
Encoder& Décodage\\
\end{tabular}
\caption{ Sens plus "naturel" de certain termes. Pour éviter la confusion on dira que par défaut sera le verbe pour un symbole nouveau : un nouveau symbole créé a vocation première d'opérer sur le préexistant et donc sera implicitement vu comme un verbe sauf indication contraire}
\end{table}

Je vais choisir pour essayer d'être plus compréhensible aussi du coté des concept d'algébriser plus fortement ma manière d'expression.
Je vais  appeler Process ou Fil  et lui donner un nom que n'importe qui et mettre entre crochets [] une entité formalisable et noter son nom ou sa formule a l interieur   $[X]$  is to be read as the process $X$ . 
L'utilisation des crochets voudra dire que  dans un raisonnement ou un document on va poser comme choix inhérent à ma tentative de compréhension que ce que j'écris depuis aujourd'hui si je dis que je vais parler de \emph{Process} de \emph{Jardins} et de  \emph{Signifier} le  \emph{Sens} tout terme entre deux crochet gagnera un indice au moment ou je fermerai le crochet.\footnote{Si on écrit a l'intérieur d' un process on definira un nouveau process ayant un temps de plus}  Et je peux supposer par exercice de pensée (gedanke) que $\iota_{\textup{Je}.t}$ est un entier positif qui représentera ce numéro au moment ou j'écris ce terme dans tout texte ayant comme préambule dans le passé. 

Entre crochets tout ce qui commence par une majuscule a une interprétation canonique  en tant de process.  

La base compositionnelle du formalisme sera les Bouts de Fil

Un propos mathématique est interprétable de manière non ambigue par une machine. Il a besoin d'être construit à partir de "rien". Ce rien est ce que j'ai le plus de mal à appréhender,
je vais le remplacer par un objet algebrisé des origines gauches ou spatiales qui n'aura pas de signification propre en tant que process  le préprocess  input universel  ou origine de l'espace qu'on notera de manière  figée à jamais (je ne garantis ma compréhension de ce que je sais aujourd'hui que si ce symbole n'est pas réinterprété autrement)  $\iota_{\mathfrak v}$ et qui ne peut rentrer que en tant que Passé d'un process et un objet algébrisé  des estimations du futur à droite  qui sera le préprocess qui suit  l'évaluation de la fin des temps  et qu'on notera de manière figée $\mathfrak v_\iota$  et qui ne rentrera que en tant que Futur dans un slot de process.

Le cadre  originel d'un propos signifié et signifiant  serait un  process initiial  phénomènologique $[\iota_{\mathfrak v}$  Je  $\mathfrak v_\iota]$ ou Je sera un Process ancré dans un Espace-Temps qui conceptualise  ma perception de moi meme.



$1$  $Un$ le premier 1 est un verbe qui est le premier verbe et Un est le nom de 1 J'aimerais écrire $\interp {\textup{Un}=\textup{Nom}(1)}$. Son dessin serait  un $\wedge$ de largeur 1 et de hauteur1/2. \footnote{ ou $(\iota_{\mathfrak v},1;1 \mathfrak v_\iota)$}

On notera au milieu des crochets une lettre minuscule toute seule pour représenter un Process qui peut être lu comme un entier dessiné par un segment horizontal vertical ou a plus ou moins 45 degrédécoupé en sous segements.

Les règles de compostions maintiennent un Passé [Je Suis Un Process] donc [pour Je il Existe Un seul Process]

Un fil   a une forme  Verbale ou Percept et une form Nominale ou Concept.  





 \begin{figure}[!ht]
\centering
\begin{minipage}{.45\textwidth}
  \centering
   \includegraphics[width=0.5\textwidth]{Simplesym}
%\caption{A qubit in the toric code}% $B_{2,2,2}$ from \cite{AFLS}}
  \label{fig:bf}
\end{minipage}%
\begin{minipage}{.45\textwidth}
  \centering
  \includegraphics[width=.5\linewidth]{Spacedsimple}
  %\caption{composition}
  \label{fig:decree}
\end{minipage}

\caption{Simple bout de fil standard with  $\downarrow_d\uparrow_c\uparrow_f$ shape (first segment has a verb so standard form below) second arrow  is start up then down in what I see if I forget the loops the third is up as the loop is above. In red the shapes by symetries (adjoint and  transpose will be inside and also reverse if finite fragmets).  Un bout de fil $B$ a 4 Spots de composition $B.P,B.V,B.N$ et $B.F$ et 4 segment droits $B.\text{preread},B.\text{postread},B.\text{prewrite},B.\text{postwrite}$  In the right  a  Process simplement hiérarchise
é (3,2;5,2)  Simple three steps needed to prepare the do 2 steps to finish the reading then 5 steps to prepare the output and 2 to finalise the book. Some how slicing the lines introduces extra space/domain. Une barre verticale avant la boulce gauche serait vue comme une limite maximale de réécriture par composition au centre (en V ou N) la paire (nb trait sortie lecture, nb trait entree ecriture) serait un cle de conversion et le nombre de traint en sortie serait le nombre de composition centrales autorisée dans la boucle a droite 
Quelque part les 4 nombres et les descritpions dans les deux slots  Verbe et Nom expliquent de manière figée si on met des nombres dans tous les segments comment aller construire un machine de turing qui lit un nombre donne de cases du passe P et qui pendant un nombre fixe d étapes insère des cases remplies obtenues par son exécution avant le Futur F  ou juste distribuer des token de temps mais on se restreint a un token de temps implique une restriction en degre de hierarchie}  

\end{figure}
\footnote{
\textbf{Digression} on pourrait peut etre digresser pour essayer de retrouver plus vite les mathématiques en les voyant comme un adjectif qui peut etre present ou non entre le Verbe et le Nom quelquepart on va ecrire seulement a la pointe des bouts de fils et seulement des objet symetriques on monterait  avec l' évaluation associée au nom : on monte avec $\mathfrak v_N$ et on descend avec $N$ on pourrait la mixer de telle maniere a ce que la signification du reste la meme si le passe explique la metatheorie (ce quon vient de dire) et restreindfe N a etre dans une Cstar alegbre et et son evaluation a etre son adjoint? Nous allons choisir de garder la definition plus generique et plus simplement signifiee puisque cette construiction se retrouve en inserant de la bonne maniere dans le verbe puis dans ke nom et garde la trace de lespace temps et de la memoire
 }

Un fil est localisable avec un encodage du temps de l'espace  et de sa cellule  active
si on a   [P \textbf{getwhen} n F], [P \textbf{getwhere} p F],[P \textbf{getcell} c F]
  
\textbf {Cellule active d'un Fil} (Composer?)   [Lecture Transforme A F]  on veut lire un partie du passe (premiere partie du bout de fil et écrire quelque chose au debut du futur deuxieme partie du bout de fil) Lire c'est fabriquer un espace de contrefactuels et écrire c'est projeter pour restreindre en restant constant tout en introduisant  une nouvelle définition et donc un nouveau nom.

 
\textbf {Entrelacer} [Spin Entrelace A et B]  veut dire Je vais interpréter le Process  Spin comme une séquence infinie indexée par $\oN$ à valeurs dans $\{\uparrow,\downarrow\}$ : a chaque temps unité j'ai un getSpin qui est l'un des deux symboles et au temps $j$, si la valeur de $Spin_j$ est $\uparrow$ le $j$ème bout de fil de l'Entrelas est pris en avançant dans $A$ et s'il est $\downarrow$ il est pris en avançant dans $B$.   
Une chaine de Spin parfaite organise un débat imaginaire (car ne prenant pas de temps dans le choix de l'alternative) formant a partir du process unique deux process séparés    la chaine de spin donne le temps de parole on dessine en suivant la recette du haut ou du bas et le dessin d'origine se trouve alors indistinguable autrement que par le nouveau nom qu'il a acquis d'un entrelacement de deux process avec une chaine de Spin. 
\footnote{On voudra construire Repeter comme Entrelacer avec le Simple constant dans l'espace et le temps qui en tout temps tout espace et pour toute memoire dit Simple en un temps unite a quiquonque lui demande  on notera $\_$ ou $1$ (on verra bien) le bout de fil standard  $\downarrow_d\uparrow_c\uparrow_f$ de longueur 1 sans lettres dessinees ni separations temporelles.} 
 
\textbf {Posteriority (attraction)}  $\Sucs_{F}$  (quelque part l'attraction universelle induit la naturalité de la flèche du temps pour un process nommable  $F$ par contre c'est peut etre un 1/2 process et pas juste un process en quelque sorte un [Verbe de $X$]$_{1/2}$ )
ca donne un nom verbalisé obtenu a partir de $X$ comme premier élément de $F$.% et doit etre consistant avec la premiere posteriorite 
%Quand on voit un process comme un indice  et qu'on veut le verbaliser 

Quand $n$ est le verbe élémentaire (sera identifié (par un process qui peut prendre du temps et de l'espace s'il est gros) à $\oN$ un entier positif) et qu'on le fait agir sur le nom d'un verbe élémentaire $k$ son action  est celle dont il est l'unique point fixe : c'est l'attraction élémentaire de ses prédécesseurs son action par défaut (sur ceux non définis à sa naissance est ce qu'il aspire a devenir et ne sera jamais le successeur a l'infini , celui qui n'a pas de point fixe quelque part l'evaluation dans les entiers positifs de $\Sucs_{\oN,\infty} (k)$ serait $k+1$ mais on n'a pas encore parle d'infini)
en math standard c'est  le map $\Sucs_{\oN}$ where
  $\Sucs_{\oN,n} (k) =\Sucs_{\oN} (n,k)$ et $\mathfrak v^2 (\Sucs_{\oN}:  \oN \timestikz \oN \rightarrow \oN )$ est  $(n,k)\mapsto \begin{cases} n \text{ if } k=n \\ k+1 \text { Otherwise} \end{cases}$.  
  
  \textbf {Nom élémentaires}  ( un nom élémentaire est lui meme sa valeur $n=\mathfrak v^2_{\oN} (n)$ et est lui meme son indice $n=\iota_{\oN} (n)$ dans la famille des noms élémentaires  qui est notée  $\oN$)  
  
 Le process    \textbf {Binaire} encode les noms élémentaires en chaines de Spin non vide et non $\uparrow$. \footnote{Si ca ne faisiat pas trop clash j'aurais imagine les nombre imaginaires ou transcendants comme $0,-1,i,\dots$}

[P Binarize N pour F] va utiliser la lecture de $N$ et proposer comme valeur $\uparrow$ ou $\downarrow$ selon qu'on monte ou qu'on descende. Sur un bout temporisé le nombre de barres est un délai de réponse. Le process écrit le sera avant le dessin de F et le prefixe et suffixe seront $\uparrow_d$ et $\uparrow_f$ à tous les bouts.
   


  This equivalence can be seen as living in the same causal world for the atoms.  The atoms can be seen  in the same deterministic reasoning model in which one can say 
  it is suffficent to know whats below to build a deterministic imaginary world in which what above always happen i.e. assuming   "$\Suct B$" one can can say $B$.
  
%\textbf {Simple} ou 1 ou  $\uparrow\uparrow\uparrow$



  
\textbf {Subjacency}  $\Suct_{F}$  [Nom de $X$]$_{1/2}$ ca donne fabrique sur une. boucle obtenu a partir de $X$ comme élément de $F$ (on transforme un process en procedure ou recette repetable c'et un peu une construction de classe d'equivalences ( dans l'ideal j'aimerais dan le Process Calcul si $\Suct_{X}=\oN$ alors $X$ sera une Somme si $\Suct\Suct_{X}=\oN$ alors X est un Produit etc avec Puissance , Tour de puissance , $k$ ieme calcul )  Le Challenge ici sera de capturer la polysémie de dériver en math standard.

\textbf {Reading } Dans  un livre de valeurs simplement indexé associe a tout ind!ice (entier positif) $j$  un indice de logne ou de definition $i$ et un indice de colonne ou d intrication ou de possible $j$, la paire $T'_j,T_j$  décrit l'image du process concret pour $T\rightarrow T'$  

Temps , Espace, Memoire    

intuitivement quand on lit un bout de fil sur le verbe on construit un espaqce de contrefactuel et cette construction peut prouver une impossibilité universelle algébrique (égalité à 0) puis on projette en écrivant pour revenir à une réalisation évaluable en revnant à une réalité (dans cette partioe on peut prouver une difference en tout point (ou deux objets decomposable  par indexation sur un meme ensemble sont differents en tout point s'il suffit de lire un indice pour les separer $A\tdf B$  en math standrad for two Processes $A$ and $B$ living in an entrelacement $L$,  $\interp{A\tdf B \text{ iff } \forall j \in \oN, \,  A_j \neq B_j}$)
\textbf { Measuring with $X$} on lit un process (Mesuré) comme une seqeunce de Jardins et un mesureuer comme une sequence (de jardins aplatis) 
le résultat classique de la mesure peut se lire avec des proba et se voir comme une tentative d'identificationde cellules de jardins vu comme (distribution de probabilite sur la linane)
a de percepts d'une certaine taille avec la probabilité inhérente au mismatch des tailles. Le résutat étant toujours une Liane de l'image dont le percept contient l'info de sa proabilite d'existence 	 ce temps particulier.



A process here is short for  either 
\begin{itemize}
\item (recognizable/exprimable/interpretable) process i.e.  has (a/at least one)   writing (process/map) to  (positive integers $\oN$/ words of a language over a finite alphabet/encoded values)   
\item  (namable/ expressible/definable) process $x$  i.e. has (a/at least one)  reading (process/map) from ($\oN$/indices/ words of a language over a finite alphabet) to bout de fil sans lettres $\iota_x (n)$ pour $n\in \oN$ est un bout de fil.
\item si $x$ est un process nommable et reconnaissable il a un map $\mathfrak v_x$ tel que $\mathfrak v_x(s,t,m)=J$  qui donne un jardin suspendu pour tout espace temps ett memoire  et     $\mathfrak v_x(J)\in \oN$ (ou $\mathbb Z$ ou Gap ou $\mathbb C$)
\end{itemize}
A process  has a meaning (a du sens communicable par le langage ) if it is namable and recognizable by a process.

A 1/2 process is a process that can and needs to be completed to have a meaning.

Any process x can explain a process y to process z forming the process $(Explain, (x;y;z))_h$ to read as "the meaning of y explained from x to z"

So any process generates a family of processes by explanation process.


I am a process : phenomenologically 
in terms of recognisability what I am to myself  can  be  thaught of / (identified with a dogmatic phenomenological equivalence class) as a process that in any space, at any  time , and for any length $n\in \oN$  positive integer  given any other person there exists an infinite indexed family of  sentences  using $n$  words each that i can think about if I try to define my-self to that person using $n$ words. \footnote{Jardins will arrive later to explain better}.

%consider  my self as indistinguishable from an ordered sequence containing a map that associates to anyone at any space at any time and anyone else present at the same space and time and any  length n  of a description what can be said from the first to the second  about me in $n$ words.



A process y is deterministic/ monoparental for a proceess x if for any process z,z'  the processes (x;y;z) is equal to (x;y;z')
we say that x is a parent of y 


equality and difference are processes %, i am a process, time is a process for me 

time is a process, demande au futur combien de temps me reste il

space is a process , demande au passe  combien de cases memoire j'ai

length is a process, combien de possibilites d ecriture jai 

computation is a process , category theory is a process, math defined with set theory is a process 
 
 
un process dessinable: a tout temps $t$, espace $e$ et longueur $n$, j'associe une construction avec $n$ bout de fil

une recette est un process  deterministe qui produit des dessins  $(Explain,(Recette;x;Draw))_h$ en $t,e,n$ termine ?

Un process est concret/ bien defini/parfait  si il a une recette finie Recette $= (J,Repeat$ $J')_h$ avec $J$ et $J'$ ont un nombre fini de Lianes. ( $J'^n$ est un père  ?)
 \begin{figure}[!ht]

    \includegraphics[width=0.5\textwidth]{PVNF}
\caption{bout de fil et composition, a gauchele boutdefil $x=\uparrow_d\uparrow_c\downarrow_f$  a un debut un coeur et une fin l'arete centrale $\uparrow_c$ dit q'uon monte avant de descendre celle du debout que la premiere  boucle est auddesus celle de la fin que la boucle de la descente est audessous.  A droite $x^x$}  

\end{figure}
%\begin{figure}[!ht]
%\label{fig:bf}
%\center{
%    \includegraphics[width=0.3\textwidth]{boutdefil}
%\caption{le boutdefil $\uparrow_d\uparrow_c\uparrow_f$  a un debut un coeur et une fin l'arete centrale $\uparrow_c$ dit q'uon monte avant de descendre celle du debout que la premiere  boucle est auddesus celle de la fin que la boucle de la descente est audessous.  inverser (bar) pour un boude fil est le bout de fil avec le sens du milieu qui change 
%dagger pour un bout de fil est celui ou la fleche du milieu change mais aussi les positions des boucles. 
%}
%}
%\end{figure}

 \begin{figure}[!ht]
\centering
\begin{minipage}{.45\textwidth}
  \centering
   \includegraphics[width=0.5\textwidth]{boutdefil}
%\caption{A qubit in the toric code}% $B_{2,2,2}$ from \cite{AFLS}}
  \label{fig:bf}
\end{minipage}%
\begin{minipage}{.45\textwidth}
  \centering
  \includegraphics[width=.5\linewidth]{bf2}
  %\caption{composition}
  \label{fig:decree}
\end{minipage}

\caption{bout de fil et composition, a gauchele boutdefil $x=\uparrow_d\uparrow_c\downarrow_f$  a un debut un coeur et une fin l'arete centrale $\uparrow_c$ dit q'uon monte avant de descendre celle du debout que la premiere  boucle est auddesus celle de la fin que la boucle de la descente est audessous.  A droite $x^x$}  

\end{figure}

Une sequence verbale hierarchiée simple serait $(x_1,x_2,\ldots ,x_n)_v$  avec les $x_i$ de hierarchie 1 ( je ne peux que les remplacer par des entiers entre 1 et le nombre de trait avant la boucle de lecture)  c'est le percept d'une Liane du jardin  

Une sequence nominale hierarchiée simple serait $(x_1,x_2,\ldots ,x_n)_h$   avec les $x_i$ de hierarchie 1 (je ne peux que les remplacer pqr des entiers entre 1 et le nombre de traits apres la boucle, c'est un concept (un indice de cellue)


Dans une Liane conceptuelle le percept est le séquent (sous le trait )  et lle contenu de la cellule sous jacente est ce qui permet la deduction 
chaque cellule a un contenu qui s'appelle N(V) s'evalu en V(subjacent)  a nettoyer  
 


1 est le boutedefil $1=\uparrow_d\uparrow_c\downarrow_f$ il se lit  en 1 et s'ecrit en tout temps et tout espace par nimporte quel process. Au temps1/2 on est a la pointe et au temps 1 a l'esoace translate de 1. 
 
 inverser (bar) pour un bout de fil est le bout de fil avec le sens du milieu qui change 
dagger pour un bout de fil est celui ou la fleche du milieu change mais aussi les positions des boucles. 

Un bout de fil et son inversion reste dans un carre de 1 unite.
Les bouts de fils se composent en se succedant ou s'imbriquant $(a,b)_h$ ou $a+b$ ou $b|a$ correspond a mettre $b$ apres $a$, $a^b$ ou $a \wedge b$ consiste a mettre $b$ au milieu de la premiere boucle de $a$  et $a_b$ ou $a \_  b$ a le mettre au milieu de la deuxieme boucle.
Le dessin est rescale pour rentrer (lui et ses inversions) dans un carre  de 2 unites.

the description of what is y for x at time t is a jardin suspendu where the lianes are the singllarities

percepts sont des temps du pere et la pile de concepts sont  pour un  temps du pere quels sont les temps passe ou le process n'est pas identique a sa vision par le pere 
 
 Un Spinlibre est un process non concret/divin : tous ses bouts  commencent et finissent par $\uparrow$ et son coeur est une i.i.d : il n y a pas de recette deterministe qui le decrit par contre pour toute longueur il est concret et si on regarde assez longtemps il est équilibré pour le sens des coeurs.
 
 Repeat is a process

  
 a positive number is a process ,$\#$ is a process having as value an integer

if $L$ is a list of $\#L$ indices and $R$ is a list of $\#R$ indices
Convert is the process $n\mapsto (L\timestikz [m_{L,R}])^n \rightarrow (R\timestikz [m_{R,L}])^n $ where $m_{L,R}=ppcm( \#L,\#R)/\#L$ and $m_{R,L}=ppcm( \#L,\#R)/\#R$ creating a one to one correspondance repeated $n$ times

it should read in the left and write in the right?
 \footnote{\red{
If $n,k$ are positive integers (in $\oN)$ I can equalise   in lenght 1 at any space and time  $(n;n;k)$ with $n$ if $k=n$    and with $n+1$ otherwise.
$  (=; (n;n;n);n) ??$ may be this is the value a the position if from the third process} }
\footnote{  0 might be defined at $x$ as $(x,x)_v$ some how a pair of indinstinguishable processes? but what could be nice is 1/2 bout de process and flat line in the missing 1/2 bout ? si j'ai un/des demi bout je serais un space de process (un peu comme un W)?
} 
  section{old}
\begin{center}
 \begin{figure}[!ht]
\begin{center}
    \includegraphics[width=0.5\textwidth]{Simple}%PVNF}
\caption{bout de fil  Simple avec son espace de symetries,  contient dans  sa forme algébrique 3 bits   $0=\uparrow_P\uparrow_V\uparrow_N$, il a un  debut un coeur et une fin l'arete centrale $\uparrow_V$ dit qu'on monte avant de descendre celle du debout que la premiere  boucle est au-dessus celle de la fin que la boucle de la descente est au dessous.  Ila 3 spots d'insertion de process $P,V,N$ dans sa forme algébrique et 4 dans sa forme libre. Il se dessine dans un carré de 1 unité de dessin de coté \textbf{mon probleme est que 0 n'a pas de sens car il n'a pas de singularite donc ce n'est pas un process mais peut etre un  1/2 process}}  
\end{center}
\end{figure}
\end{center}


\subsection{Anthropomorphisme fondamental d'evaluation Relative time : Gap or $\mathfrak t _{loc}$ , bi-infinite exploration} \label{sec:gaps}

J'ai l'impression que : Si $s$ est un bit de spin $s\in \{\downarrow,\uparrow\}$, $s$ est un entier espace et $t$ un entier temps
obtenus par mise en forme normale de l'introduction de $s$ dans le passé d'un Verbe le tout mis sous forme normale
$Gap_{e,t}^s$ c'est l'évaluation de l'espace temps ?


$Gap_1^{(1)}=(0,1,-1,2,-2,\ldots)$, $Gap_1^{(2)}=(0,-1,1,-2,2,\ldots)$, $Gap_2^{(1)}=(1,2,0,3,-1,\ldots)$, $Gap_2^{(2)}=(1,0,2,-1,3,\ldots)$, 
$Gap_3^{(1)}=(-1,0,-2,1,-3,\ldots),\ldots$

More formally for $n\ge 0$, $j>0$, $j'>0$, 

 $\begin{cases}(Gap_j^{(2n+2)})_{1}=(Gap_j^{(2n+1)})_{1}\\ (Gap_j^{(2n+2)})_{2n'}=(Gap_j^{(2n+1)})_{2n'+1}\\(Gap_j^{(2n+2)})_{2n'+1}=(Gap_j^{(2n+1)})_{2n'}\end{cases}$

%$(Gap_{2n'+1}^{(2n+1)})_{j'}=(Gap_{1}^{(2n+1)})_{j'} -n'$
%$(Gap_{2n'}^{(2n+1)})_{j'}=(Gap_{1}^{(2n+1)})_{j'} +n'$

And for $j''>1$, $(Gap_{j''}^{(2n+1)})_{j'}=(Gap_{1}^{(2n+1)})_{j'} +-(1)^{j''} \lfloor j''/2\rfloor =(Gap_{1}^{(2n+1)})_{j'}-(Gap_{1}^{(2n+1)})_{j''+1} $
and for $j\ge 1$,  $$(Gap_1^{(1)})_j=(-1)^j \lfloor \frac {j-1} 2 \rfloor$$. 

We might use  $Gap_j^\uparrow$ for  $Gap_j^{(2n+1)}$ and  $Gap_j^\downarrow$ for  $Gap_j^{(2n)}$ 

 $$ (Gap_j^\uparrow)_{j'}%= (j' \mod 2)+  (-1)^j \lfloor \frac {j} 2 \rfloor- (-1)^{j'+1} \lfloor \frac {j'+1} 2 \rfloor
  = j' -2 \lfloor j'/2 \rfloor + (Gap_1^{\uparrow})_j-(Gap_1^{\uparrow})_{j'+1}$$
 
%\footnote{ can be formulated as $(Gap_j^\uparrow)_{j'}=(j' +1 \mod 2)+  (-1)^j \lfloor \frac {j} 2 \rfloor- (-1)^{j'-1} \lfloor \frac {j'-1} 2 \rfloor$}.
and $ (Gap_j^\downarrow)_{j'}= - (Gap_j^\uparrow)_{j'}$ 
% 1+ (-1)^j \lfloor \frac {j} 2 \rfloor- (-1)^{j'-1} \lfloor \frac {j'+1} 2 \rfloor$$ 


We will use a grading to  define for $j\in \oN$ and $b \in {\mathbb Z}$ and $l\in \{\uparrow,\downarrow\} \sqcup {\mathbb N}  $,  \footnote{should be oriented sum to be clean}   %$Gap_{l,j,j'}=(Gap^{(l)}_j)_{j'}$

$\vgx{Gap} {l,j} (j')= (Gap_{l,j})_{j'}= (Gap^{(l)}_j)_{j'}$ and $\igx{Gap} {l,j} (b)= j'$.

 A $b$-Gap is a Gap starting with $b$ as the value of the first element in $\mathbb Z$ i.e if $G=Gap^{(l)}_j$ we say that $G$ is a $(Gap^{l}_j))_1$-Gap. 


Somehow we will consider repeatable  formalisations and will use corrections to make second order dynamics match an indistinguishability, so we will work with $\oN \times[n_1] \times \oN$ and with $\oN \times[n_1] \times  [n_2] \times \oN$

Note that if we take as cost for a gap reading the distance crossed reading $n$ elements costs $\sum_{j'=1}^n  |(Gap^{(l)}_j)_{j'+1}-(Gap^{(l)}_j)_{j'+1}|= n(n+1)/2$.

\subsection{Permutations/ traductions / map biunivoque/ Group correction/see phantoms}

Given  a time sequence $T=T_1,T_2,\ldots$ with $T\simeq \oN$ and $T_j\in \oN$ and a space sequence  $T'=T'_1,T'_2,\ldots$ with $T'\simeq \oN$ and $T'_j\in \oN$

$\mathfrak v :T' \rightarrow \oN$ and $\iota: T\rightarrow \oN$ are the space and time canonical isomorphisms.

For $l\ge k$, we define

$JTrad_{T\mapsto T'}^{(|_l)}= \at {T'_{j_1}} {T_{j_1}} \sqcup \ldots \sqcup  \at {T'_{j_k}} {T_{j_k}}$ with  

for any $j'\in [l]$,\  $\begin{cases} T_{j'}=T'_{j'} \text{ iff }j\not\in \cup_{j'\le k} j_{j'}\\   
\min (T_{j_{j'+1}},T'_{j_{j'+1}} )> \min (T_{j_{j'}},T'_{j_{j'}} ) \end{cases}$

The map is fully captured at level $l$ is $\forall j'>l,\, T_{j'}=T'_{j'}$.

The cast as Spin and the complementary cast are

$
\begin{cases}
((Spin)_c)_{j'} J =(\Suct_{j'})^{T_{j'}>T'_{j'}} \uparrow_{j'}^{T_{j'}\neq T'_{j'}} \text{replaces each atom by } \uparrow \text { or } \downarrow  \text{ to obtain a Spin of length } k \text{ at positions } j_{j'}\\
((Spin)^c)_{j'} J=(\min(T_{j'},T'_{j'}),\max(T_{j'},T'_{j'})-\min(T_{j'},T'_{j'}))^{T_{j'}\neq T'_{j'}} 
\end{cases}
$ 

An evolution compatible with $J^{(l)}$ as a starting point is 
$J^{(l+1/2)}= \smd \at {tw(T'_{j_1})} {(T_{j_1},T''_{j_1})} \sqcup \ldots \sqcup  \at {tw(T'_{j_k})} {(T_{j_k},T''_{j_k})}$

An evolution compatible with $J^{(l+1)}$ as an  end  point is 



\red{
Commutativity corrector of + should be $Cor_{\Suct_+} + B A = +  A B$

Associativity corrector (assume default reading is from percepts so by value and $._R$ is the right multipication so shoud be$\Sucs^{(2)}$ when it will be defined)  )$Cor_{\Suct_{._R,A}} (AB)C= A (BC)$ as we want to map $._R A ._R B C $ to $._R (._R A B) C$ or ($Cor  A^{(1)} Par^{(2)} B^{(1)} C^{(1)}:=Par^{(2)} A^{(1)} B^{(1)} C^{(1)}$ parentheiss are for clarity if assuming value by default) 
}
\footnote{when can we  have reduced form were subjacencies are involutions vertices of graphs becaming stable indistinguishable sets?}
\subsubsection{The 6 bases of maps}

 Aside the matching representation any JS can be represented by orbit or by insertion (and the * like  memorial morphism will double by looking from the lines point of view inverting the liane of concepts with its percept)
  
$ JOrb_{T\mapsto T'}^{(|_l)}= $% (\at {T'_{i_1}} {T_{i_1}} \sqcup \ldots \at {T'_{i_1}} {T_{i_1}}) \sqcup \ldots \sqcup (\at {T'_{i_k}} {T_{i_{k,1}}} \sqcup \ldots \at {T'_{i_1}} {T_{i_1}}) $


$ JIns_{T\mapsto T'}^{(|_l)}$

 $\Suct_{JX}(JX_{T\mapsto T'}^{(|_l)})=JX_{T'\mapsto T}^{(|_l)}$


%\subsection{Permutations as JS}

%Let $Idx$ be the set $\overset{\circ}{\mathbb N}$ indexed by ${\mathbb N} $  so that  $Idx_i=i+1$ for any $i \in {\mathbb N} $ 
%PermJS is a family of canonical casts of permutations as JS.
%
%A PermJS$_1$ is an %$(i\mapsto |Liane_i |)$
% $(\mathfrak v_{Idx})$JS where anchors represent an indexation of the orbits following the columns such that each Liane is an orbit.  % with non increasing size (its shape is a rotated permutated Young Diagram)
%The permutation is read Liane by Liane by mapping each cell number to the cell number below it and by mapping the last one to the first one.% \red{toncheck/clean/correct}
%
%$\sigma'^{l+1}= \sigma_{p,l+1}\sigma'^{l}$ means that the PermJ is adapted by inserting the cell $l+1$ at  Liane$_{p+1}$. 

%An OrbPermJS is a JS where the Lianes represent the orbits read by columns.
%
%  $ (OrbPermJS_1)_c (
%\left(   \begin{matrix} 0&0&1\\
%0&1&0\\
%1&0&0
%  \end{matrix} \right ))= \begin{ytableau}
%\none &\none[_1]&\none [_2]\\
%  \none&1&2\\
% \none &3\\
%  \end{ytableau}$
%  
%  
%As 1 is always in the first column we could have a more concise notation inserting  (image of 1-1), then (image of image of 1 -1) until getting back one the create a new column inserting the first missing element. using $n-k$ cells where $k$ is the number of orbits.
\red{old :
MatchpermJS  is the JS representing permutations as a matching $\sigma \in S_n$, $\sqcup_{a\in [n-1],\sigma(a)\neq a} \at a {\sigma (a)} $ it has atoms in $FAT^{(2)}_{\neq}$
where FAT is the old family if atoms  \footnote{Matchperm$_2$ can be built by sorting the positoin and using the inverse permutation}. 

$InsPermJS_1$ is defined by inserting  in last line: for $i\in [n-2]$,we insert  the cellnumber $i+1$ at the position $p\in [i+1]$ at the line $i+1$ \footnote{$InsPermJS_2$ is defined by inserting at last columns instead of at last line} This allows to reversibly associate to any permutation an infinite family of permutations inclduing a canonical decomposition and a canonical completion.

OrbPermJS$_1$ is the JS representing permutations with the anchor being the smallest undiscovered element with non trivial orbit. It has atoms in $FAT^{(2)}_>$  (atoms all have  $(p>a)$). It have   $n-k$ cells where $k$ is the number of orbits.\footnote{ OrbPermJS$_2$ is  the permutation matrix reading by the lines instead of the columns.}
}

%$InsPermJS_1$ is defined by inserting  for $i\in [n-2]$, the cellnumber $i+1$ at the position $p\in [i+1]$ at the line $i+1$

 Consider the permutation  $\sigma \in S_7$ defined by  $\sigma (1)=1,\sigma (2)=3,\sigma (3)=7,\sigma (7)=2,\sigma(4)=4,\sigma(5)=6,\sigma(6)=5$ canonically extended to $S_n$ with $n>7$  by   $\sigma(i)=i$ for $i>7$.

$\sigma =\left(   \begin{matrix} 
1&0&0&0&0&0&0\\
0&0&1&0&0&0&0\\
0&0&0&0&0&0&1\\
0&0&0&1&0&0&0\\
0&0&0&0&0&1&0\\
0&0&0&0&1&0&0\\
0&1&0&0&0&0&0\\
  \end{matrix} \right )$ which is built by inserting last line as 
   $1\xrightarrow{2} \left(   \begin{matrix} 
1&0\\
0&1\\
  \end{matrix} \right )\xrightarrow{3}$
  
    $\left(   \begin{matrix} 
1&0&0\\
0&1&0\\
0&0&1\\
  \end{matrix} \right )
  \xrightarrow{3}
    \left(   \begin{matrix} 
1&0&0&0\\
0&1&0&0\\
0&0&0&1\\
0&0&1&0\\
  \end{matrix} \right )
\xrightarrow{4}
    \left(   \begin{matrix} 
1&0&0&0&0\\
0&1&0&0&0\\
0&0&0&0&1\\
0&0&1&0&0\\
0&0&0&1&0\\
  \end{matrix} \right )
\xrightarrow{4}
  \left( \begin{matrix} 
1&0&0&0&0&0\\
0&1&0&0&0&0\\
0&0&0&0&0&1\\
0&0&1&0&0&0\\
0&0&0&0&1&0\\
0&0&0&1&0&0\\
  \end{matrix} \right )
\xrightarrow{2}  \sigma$ 


 $(OrbPermJS_1)(\sigma)=\begin{ytableau}
\none &\none[_2]&\none [_5]\\
  \none&3&6\\
 \none &7\\
  \end{ytableau}$,  and  $ (InsPermJS_1)_c (\sigma)
= \begin{ytableau}
\none &\none[_2]&\none[_3]&\none[_4]\\
  \none&1&2&5\\
 \none &6&\none\\
  \end{ytableau}$

 $(MatchPermJS_1)(\sigma)=\begin{ytableau}
\none &\none[_2]&\none [_3] &\none [_5] &\none [_6] & \none [_7]\\
  \none&3&7&6&5&2\\
  \end{ytableau}$

\footnote{old 
There are therefore 6 Canonical permutation forms  associating to any $\sigma\in S_n$, 2 canonical pasts in $S_k$ with $k<n$ and a canonical future in $S_k$ for $k>n$.

%   not putting it redaing the columns inserting $l-1$ where $l$ is the position of the line

The 8 first Families of Matching Permutations (\red {order is not optimized or naturalized} )
FMperm$_i$ for $i\in [8]$ is read in $(crr, crl,cll,clr,lur,lul,ldl,ldr)$ 
%$FPerm_1=FPerm_{crr}, FPerm_2=FPerm_{crl},FPerm_3=FPerm_{lrr},FPerm_4=FPerm_{llr},
%FPerm_5=FPerm_{clr}, FPerm_6=FPerm_{cll},FPerm_7=FPerm_{llr},FPerm_8=FPerm_{lll}$
  are indexed such that the insertion is done respectively column right, easy right extend for $FPerm_1=FPerm_{crr}$, 
  by lines by up (from top)  easy left extend ,for $FPerm_6=FPerm_{lul}$.

$(FMPerm_1)_{n! +p-1}$ for $p\in [l+1]$ is obtained by inserting by columns to the right and is read $(p_1,\ldots , p_{n+1})$ where $p_i$ is inserted at step $i$
$(MatchPerm)_c (FP_1)_{n!} = \at 1 1 \sqcup \at {n} n \ldots \sqcup \at {n} n=canJ_n$


$(FMPerm_3)_{n! +p-1}$ for $p\in [l+1]$ is obtained by inserting by columns to the left  and is read $(p_{n+1},\ldots , p_{1})$ where $p_i$ is inserted at step $i$
$(MatchPerm)_c (FP_6)_{n!} = \at n 1 \sqcup \at {n-1} 2 \ldots \at 1 n$

\red{TO CLEAN}
Using c-reduced Orb form and AJS we get $(FOPerm_1)_{l_1!+l_2! +p}$ with $p\in [3]$ we view Orb perm as metfunc cast of the min function where atoms as they can be pushed left are of the form $\at {a+b} {(a,a+c)}$  is  so that the anchor is the smallest element in the liane and cells may have moved so we forget their number. when inserting we chose a liane or insert a new liane to the (right/to the left) having as minimal element $w$+1 where $w$ is the number of cells . 

}
  
  
   

\subsubsection{Twisting (entangling) Operator, twisting with a correlation and  Jardins on Twists}
A basic twisting operator is `an operator that preserves biexplorability of counterfactuals" : we will use functions defined at some time at some domain and we want that at any time, we can define two directions such that only one direction  may contain the forbidden fruit even when extending the domain of the functions.\footnote{The intuition behind biexplorability is that we heave two parameterized lines that cover the surface that are disjoint at infinity and such as the zeros of the surface we are trying to avoid are in a bounded region.} 

$tw: \text{oldvalue} \times \text{newcorrelationspin} \times \text{new value of new variable}\rightarrow \text{newvalue}$
More formally   a Twist is a function  $tw :{\mathbb Z}\times  \{0,\uparrow,\downarrow\} \times {\mathbb Z}\rightarrow {\mathbb Z}$ such that 
\begin {itemize}
\item  $tw(x,0,y)=0 \implies \forall y' \neq 0,\, tw(x,\uparrow,y')tw(x,\downarrow,y')\neq 0$
\item  $\forall x \in {\mathbb Z},\,  \{y>0, tw(x,\uparrow,y)=0\}\cap\{y>0, tw(x,\downarrow,y)=0\}=\emptyset$  and 
\item $\forall x \in {\mathbb Z},\, |\{y>0, tw(x,\uparrow,y)=0\}\cup\{y>0, tw(x,\downarrow,y)=0\}|\le1$ % \implies \forall y' \neq 0,\, tw(x,+,y')tw(x,-,y')\neq 0$
\item $\forall x,s,y,\,  |tw(x,s,y)| |\le|x|+|y|$ % \implies \forall y' \neq 0,\, tw(x,+,y')tw(x,-,y')\neq 0$

\end{itemize}
The standard twist is $tw_1 := \begin {cases} x,0,y \mapsto x \\ x,\uparrow,y \mapsto x+y \\ x,\downarrow,y \mapsto x-y\end{cases}$.

Exchanging $\uparrow$ and $\downarrow$ preserves being a valid twist, i.e. we can define  $inv$ by   $inv(tw)(x,0,y):=tw(x,0,y),inv(tw)(x,\uparrow,y):=tw(x,\downarrow,y)$, and $inv(tw)(x,\downarrow,y):=tw(x,\uparrow,y)$.


\subsection{evaluation with correlations}

  
A level $l$  evaluation is a   function $ev$  such that $ev(x,p, v)$ represents the evaluation of $v'=\sigma_{l+1,p}v$ with the hyperplane like ($l$ extended) (JS)  $x'=\sigma_{l+1,p} (x,\uparrow)$, so that $x'$ is the  $l+1$ correlation obtained from $x$ by inserting $\uparrow$ at position $p$. The transformation have to use the same $p$ for all the cells. 
 

  
  %the value $v$ is obtained from the anchor of the cell $(s_1,\ldots,s_
  $ev: \{0,\uparrow,\downarrow\}^l \times  [l+1] \times {\mathbb Z}^{l+1}\rightarrow {\mathbb Z}$ is an evaluation if :
  given   $(x,p,y)\in \{0,\uparrow,\downarrow\}^l \times  [l+1] \times {\mathbb Z}^{l+1}$, and $x'=\sigma_{l+1,p} (x,\uparrow)$
$ev$ satisfies that     $\forall j\in [l+1]$ such that  $x'_j\neq 0$:
  % if ($j<p$ and $x_j\neq 0)$ or ($j=p$) or ($j>p$ and $x_{j-1}\neq0$)  then
   \begin{itemize}
 \item If $ev(x,p,y)=0$ then $\forall k\neq y_k,  ev(x,p,y[j \mapsto k]) \neq 0$
 \item If $ev(x,p,y)\neq 0$,  $| \{k,  ev(x'[j \mapsto \uparrow],p,y[j\mapsto k]) =0\} \cap  \{k,  ev(x[j\mapsto \downarrow],p,y[j\mapsto k]) =0\}|=0$
 \item If $ev(x,p,y)\neq 0$,  $| \{k,  ev(x'[j\mapsto  \uparrow],p,y[j\mapsto k]) =0\} \cup  \{k,  ev(x[j\mapsto \downarrow],p,y[j\mapsto k]) =0\}|\le1$
 \item $\forall x,p,y,\,  |ev(x,p,y)| |\le \sum_{i} |vCor(x'_i) y_i|$ % \implies \forall y' \neq 0,\, tw(x,+,y')tw(x,-,y')\neq 0$

 \end{itemize}
   %(    $|\{k>0, ev(x,p,f_E(s_{_j\leftarrow k}))\}=0|\le 1$ 
   where the operation  $x[j\mapsto k]$ maps $(x_1\ldots,x_l)_h$ to $(x_1\ldots , x_{j-1},k,x_{j+1},\ldots ,x_l)_h$ %and $x'=\sigma_{l+1,p} (x,+)$, so that $x'$ is the  $l+1$ correlation obtained from $x$ by inserting $+$ at position $p$.

\subsection{Family of  Correction/ decoding Groups JS (to recheck)}

An FG $\in$  FGS  is a JS whose Lianes are interpreted as a sequence of generators of groups: $Liane_a(FG)= \sqcup_i \at a {g_{i,a}}$ where $G_a=<g_{i,a}>$ is a group.  

An FGS  acts on  a JS, $J$  as element  of an FJS, $F$ if :

Let $j= \mathfrak i_F(J)$, %if $j=0$ then for any $g$,
 $FG_F(g,J)=J_\emptyset$ if $g\notin G_j$ .%  if $\not \exists i, F_i=J$



%A family of groups acting on JS from a Family of JS so that the $Liane_i$ of 
%$G$ contains $g_{1,i},\ldots,g_{l_i,i}$  generating a group $G_i$ acting in $J \in FJ_j$ so that

and   for any $g\in G_j$ there exists $j'$ with $|j'-j|\le desc_{G_j}(g)$ such that $FG_F (g , J)\in FJ_{j'}$

the description of $g$ in a group $G$ is: 

an element of $G$ can be written as a word using the generators and operators $inv_k$ were  $inv_k( x_1\ldots x_k)=x_k^{-1}\ldots x_1^{-1}$  and $desc (g ,r)=1+desc(r)$ if $g$ is a generator and $desc  ( inv_k , r)=k+desc (r)$

\section{Dictionary vision}
Recall that in Jardins by default any term has  3 associated integers corresponding to  the local times  (number of steps) required to write it, to check if it is nothing and to check that it is not nothing.
 A good visualising example of evolution of percepts in Jardins  when learning new concepts is the dictionnary.

$\mathfrak n _D(j) := \mot_{D,j}$, $\iota_D (\mot_{D,j}):= j$, $\mathfrak v _D(j) := \Enc( \Defn_D(\mot_{D,j}))$

The dictionary has $l$ definitions the first being   $"\mot_1: \Cpt (\mot_1)"$ the last if we cut after $l$ definitions is  $"\mot_l: \Cpt (\mot_l)"$. 

The conceptual weights $w_j$ is the number of words in the definition of $\mot_j$, $w$ being the total weight. 

A general reading  at time (number of words read) $t$ maps all concepts to an already  read sequence of words and a program for the future.% (permutation on the rest)

When time is fixed we consider the explorations or encodings $\Seqd {j} {w!}  \Expl_j $ that    are the permutations of all the words inside the definitions in a dictionnary with $w$ words after the symbol ':'.

 
\begin{itemize}

\item $A^{\delta_\phi}:=$ If  not $\phi$  (or $\phi==0$) then $A$ (Otherwise nothing)
 \item $A^{\delta_{\phi,1}}:=$ If   $\phi$  (or $\phi \neq 0$) then $A$ (Otherwise nothing) 
 
\end{itemize}

The definition of a word  starts with a word himself then with a symbol to separate from the rest  $\mot_{D,j,1}:=\mot_{D,j}$ and   $\mot_{D,j,2}:= ':' $ 

The time required to read the definition of $\mot_j$ in $D$ if each word was read in time 1 would be  $h_{D,j}=n_j+2$


$\Defn_D(\mot_{D,j}):=\Liane_D(\mot_{D,j}):= \Ppt_D (\mot_{D,j}) : \Cpt(\mot_{D,j}):= \Seq_{k_{|n_j+2}} (\mot_{D,j,k}) $




$\Cpt(\mot_{D,j}):=\Seq_{k_{|n_j}} (\mot_{D,j,k+2})$ 


$\Find_D(j,k,p):= (j,k,p)^{\delta_{{mot_{D,j}== (\Defn_D(\mot_{D,k}))_p},1}}$


An $(r,a,b)$ reading of $D$ is a sequence $\Seq_j X$ such that each word $w$ appears at most $r(w)$ times and after the index $a(w)$ and before the index $b(w)$ i.e.
 $1\le \# \{j, \Seq_j X == w \} \le r(w)$ and $1\le \min \{j ,\, \Seq_j X ==w\}$ and $\max \{j,\,  \Seq_j X == w \} \le b(w)$
 
In a Jardin each concept have  a divine  space successor $\Sucs$ and memorial time successor $\Suct$. 
 


After fixing time, lets consider a $(1,a,b)$ reading of $D$
 
\begin{itemize}
\item $\Infl_D(j)$ de  the sequence of indices $j'$ such that  $j'\in \Infl(j):=  \mot_{D,j'}$ uses $\mot_j$ in  its definition and let $w'_j = \# \Infl_D(j)$ be its length (the number of words using 
$\mot_{D,j}$ in their definition).
\item  $l'_j$ the index of the last word using $\mot_{D,j}$ in its definition.
\item  Given a word $\mot_{j'}$ using $\mot_j $ in their definition let
\subitem $f'_{j,j'}$ the first position in which $\mot_{D,j}$ appears
\subitem $h'_{j,j'}$ the last position in which $\mot_{D,j}$ appears
\end{itemize}


The counterfactual  action of a concept or the entanglement book of a dictionnary is obtained when fixing time by considering the interaction for each word with the others.

$\Line_D(\mot_{D,j}):= \Seq_{k_{|l_D}} (\Seq_{p_{|h_{D,k}}} \Find_D(j,k,p))$ finish inside before increasing outside

We can factorize using the influence list denoting $\Infl_D(j)=(k_1,k_2,\ldots, k_{w'_j})$.

$\Line_D(\mot_{D,j})=(j (k_1(p_{k_1,f'_{j,k_1}},\ldots, p_{k_1,h'_{j,k_1}}), k_2 (p_{k_2,f'_{j,k_2}},\ldots, p_{k_2,h'_{j,k_2}}), \ldots, k_{w'_j}(p_{k_{w'_j},f'_{j,k_{w'_j}}},\ldots, p_{k_{w'_j},h'_{j,k_{w'_j}}}))$ 
  
%$\Suct, \Sucs$


After an $(1,a,b)$ reading $R$ of $D$ one can consider the matrices $LE_{R,j',j}$ and $UE_{R,j',j}$ whose entries are  whose lines are $T'_1,\ldots T'_l$ conrresponds $\Line_D$ of the reading and columns $T_1,\ldots T_l$ are $\mot_1,\ldots \mot_l$ whose entries are first and last position or 0 if absent. 

\red{review as  Reverse Map or higher order memory map}

decompose the 1/2 dicto   into correctable versions where you can remove one product liane ?? and 

 build superposition over  of when can i know that it is noit this product 
  
SupSeq$_j=  \ket j  \ket {D' \setminus \Defn_{D'} (\mot_{D',j})}$  where $D'= \Seq_j  Entangle$


 \section{A practical application}
 %\begin{center}
  \begin{theorem} {[\bf{  123 Theorem (ex-conjecture)(2004-2024)}]}\cite{keusch2024solution}
  % \footnote{for more info about the conjecture see https://arxiv.org/pdf/1211.5122.pdf or  https://faculty.math.illinois.edu/\~{} west/regs/123conj.html  Omar pointed out that recently a prove has been found by Keutsh arXiv:2303.02611} 
For any connected graph $G=(V,E)$ on 3 or more vertices there exists a weighting of the edges $w: E\rightarrow \{1,2,3\}$ that induces a proper coloration of the vertices  \emph{i.e.}
for any $(u,v)\in E$, $\sum_{ (u,v')\in E} w(u,v') -\sum_{ (u',v)\in E} w(u',v)  \neq 0$   \end{theorem}
 %\footnote{For 123 proof one can restrict the matrices   and with entries in $\{0,1,-1\}$ that are  "almost symmetric" matrices \emph{i.e.} for all $i,j$,  $M_{i,j}+M_{j,i}=0 \bmod p$ and just need  $\exists p$}
 %\footnote{ mod $3^i$ should capture the exact strucrure }
%\end{center}
\begin{conjecture}
Given a permutation $\sigma\in S_m$ and a matrix   $E\in {\cal M}_m(\{0,1,-1\})$ such that $E_{\sigma i, i}\neq0$ for $i\in [m]$, then  for any $X_1,\ldots X_m$ each $X_i$ having $d$ possible values in ${\mathbb Z}$ and for any $B\in {\mathbb Z}^m$  there exists $n_1,n_2$ with $2n_1+n_2=m$ and  a space of dimension greater than $(d(d-2))^{n_1}(d-1)^{n_2}\ge (d(d-2))^{\lfloor m/2\rfloor}$ with basis elements $\ket X$ such that $\Pi_{i\in [m]}(MX-B)_i\neq 0$ where $(MX)_i= \sum_{j\in[m]} M_{i,j}X_j$. %$\forall i\in [m], 

\end{conjecture}

\subsection{link}
Given a graph $G=(V,A)$ with $|V|=n$ vertices and $|A|=m$ edges. 
Given an orientation $o$ of the the edges inducing a directed edges set $A'$, 
%Given an order on the vertices,  we  call an orientation of its edges a map $o:A \rightarrow \{-1,1\}$, where   $o(u,v)=1$  iff  the edge $(u,v)$ is directed from the smaller to the larger . The set of directed edges consistent with $o$ can then be defined as  $A'=\{ (u,v) \in A, o(u,v)=1 \text { and } u<v \} \cup \{ (u,v) \in A, o(u,v)=-1 \text { and } u>v \}$. 
 the {\bf effect} matrix  $E_{G,o}\in {\cal M}_m(\{0,1,-1\})$ is  defined on $A'\times A'$ by   $$(E_{G,o})_{(u,v),(u',v')}=\delta_{u,u'}+\delta_{u,v'}-\delta_{v,u'}-\delta_{v,v'} $$ where $\delta_{u,v}$ is the Kronecker delta. 
 

\subsection{Example before generalisations}

Given $T,T'$ tables with $d$ distinct positive integers called indices.
 A matrix $M$  indexed by $T,T'$  over a permutation  $\sigma \in S_d$  with spin sequence $Sp \in \{\uparrow,\downarrow\}^d $ is a matrix such that  for $i,j \in [d]$, 
 $M_{T'_i,T_j}\in \{J_\emptyset,\uparrow,\downarrow\}$ and $M_{T'_ {\sigma i},T_i}=Sp_i$  with $Sp_i \in \{\uparrow,\downarrow\}$ for $j\in [d]$ . %\in \{\uparrow,\downarrow\}$  
  Its "core" is $\tilde{M}_{i,j}=
  \begin{cases} 
  1 \text { if } M_{T'_i,T_j}= \uparrow \\
  -1 \text { if } M_{T'_i,T_j}= \downarrow \\
  0 \text { if } M_{T'_i,T_j}= J_\emptyset 
  \end{cases}$
 
% Its active correlations is the sequence of non 0 positions after the permutation indices (blue spins i.e. above a black spin).   Its passive correlations is the sequence of non 0 positions before the permutation indices (red spins i.ie at the left of black spins). 
We will define active and passive correlations.
By taking the columns one by one we can draw a shape by inserting in front of the associated correlation the ones linked with an active correlation.

The true conjecture is that looking at the shapes one can find a solution to each constraint $T'_i$  and the freedom we have for each $T_i$ is complete (any of the $d$ proposed values is ok) for the first element of a line containing more than one element,  $d-2$ for the last element of a line containing more than one element and $d-1$ otherwise.

We want to say that the causality JS of the problem characterizes  the uniform  solutions to our problem, and to prove the conjecture by reversibility of the map
 $    \begin{ytableau}
\none&\\
\none&   \\
\none &  \\
    \end{ytableau} 
    \rightarrow
     \begin{ytableau}
\none&&\\
  \none&   \\
    \end{ytableau}$

For example, the matrix 
 \begin{center}
 \ytableausetup{centertableaux}
 \begin{ytableau}
  \none  & \none[T_1] & \none[T_2] & \none[T_3] & \none[T_4]\\
       \none[T'_1] &\uparrow  &  &\blue{\downarrow} & \\
  \none[T'_2]  &\purple{\uparrow}  &  &  &\uparrow\\
  \none[T'_3]  &  & \purple{\downarrow}  &\uparrow & \\
   \none[T'_4]  &  &\uparrow  & &
\end{ytableau}
\end{center}
which corresponds when we add insertion time of cells corresponding to column insertion from the left to right with  ($Sp_i$, ($x$-position, correlation column),($y$-position,correlationline)) : $(1,(2, (0,T_2)), (2, ((0,T'_4))))$ then  $(1,(3, (-1,0,T_3)), (2,( (0,-1),T'_3))))$ then 
$(1,(4, ((0,0,0),T_4)), (2,((1,0,0),T'_4)))$
 \begin{center}
 \ytableausetup{centertableaux}
 \begin{ytableau}
  \none  & \none[T_1] & \none[T_2] & \none[T_3] & \none[T_4]\\
       \none[T'_1] &\uparrow_1  & _{1.5} &\blue{\downarrow}_{2.5} &_{3.5} \\
  \none[T'_2]  &\purple{\uparrow}_4  & _4 & _4 &\uparrow_4\\
  \none[T'_3]  & _3 & \purple{\downarrow}_3  &\uparrow_3 &_{3.5} \\
   \none[T'_4]  &  _2&\uparrow_{2}  &_{2.5} &_{3.5}
\end{ytableau} 
\end{center}

The causality JS for standard timeline (black spins permutation)  is obtained by reading the columns from left to right and inserting the line  corresponding to the time permutation
here  

$\begin{ytableau}
\none& \none[_{T_1}]\\
\none    & T'_1\\
    \end{ytableau}
     \rightarrow
      \begin{ytableau}
\none& \none[_{T_1}]& \none[_{T_2}]\\
\none    & T'_1 &T'_4 \\
    \end{ytableau} 
     \rightarrow
     \begin{ytableau}
\none& \none[_{T_2}]& \none[_{T_3}]\\
\none    & T'_4 &T'_1   \\
\none  &  \none &T'_3  \\
    \end{ytableau} 
    \rightarrow
     \begin{ytableau}
\none& \none[_{T_2}]& \none[_{T_3}]& \none[_{T_4}]\\
\none    & T'_4 &T'_1  &T'_2 \\
\none   & \none &T'_3  \\
    \end{ytableau}$


%Here the permutation positions are in black, the active correlations are   $Tcoract=(\blue{2^\downarrow})$ between ($X_{T_3}$ and  $X_{T_1}$) %or equivalently $Y_{T'}$ 
%, second emptyspot after (above)  the permutation positions, the passive correlations are $Tcorpass=(\purple{1^\uparrow},\purple{5^\downarrow})$ : first and fifth slot before/below the permutation positions
%between ($X_{T_1}$ and  $X_{T_4}$ and between ($X_{T_2}$ and  $X_{T_3}$)
% \ytableausetup{centertableaux}
% \begin{ytableau}
%\none    & T_1    \\
%\none    &T_2  & T_4 \\
%   \none & T_3  \\
%    \end{ytableau} 
%     is the shape of the permutation 
%  %\hspace {2cm}
%     and the shape of the impact on the constraints is
% \ytableausetup{centertableaux}
% \begin{ytableau}
%\none    & T'_1 &T'_3   \\
%\none    &T'_4    \\
%   \none & T'_2  \\
%    \end{ytableau}
%%\end{center}

\red{to clean , and addition of standard sum and multiplication,
The structure described :  $T,T',\sigma, cor_a,cor_p$  is called a (structuring space time) {\bf  key}.}
\footnote{the vague intuition is that lines represent space or energy and columns time or mass and the cells in the safe JS built from the structure are the basic elements. }

In the example the conjectured associated safe selection tree shape is therefore $d\times d-1 \times d-2 \times d-1$. Which implies  with $d=3$ that :
%for any $B\in \mathbb Z$ and for any sets of three choices  for $(X_{T_1},X_{T_2},X_{T_3})$ 
given sequentially $b_1,b_2,b_3,b_4$ ,
we can find  sequentially $X_1(t)$ for $t \in[4]$, $X_2(t)$ for $t \in [2,4]$, $X_3(t)$ for $t \in [3,4]$, and $X_4(t)$ defined only at $t=4$  such that $X(4)$ where $X=(X_1,X_2,X_3,X_4)$ forms a 
   $3\times 2 \times 1\times2=12$ "strong/regular/uniform" solutions to $\Pi_i (MX-B)_i\neq0$ : given $b_1\in \mathbb Z$ such that  for any of the 3 choices of    $X_1(4)$,   for any  $b_2$ there exists 2 choices for $X_2(4)$ such that for any $b_3$ there exists 1 choice for $X_3(4)$ such that for any $b_4$ there exists 2 choices for $X_4(4)$ such that $\Pi_i ((M(X(4)))_i-b_i)\neq 0$.  
%It as a list of active and passive correlations 

We want to generalize the notion of problem and the notion hyperplane to capture this regularity.

%It  has as basic base description  two rotated jardins suspendus  one filled with $T$ the other by $T'$)%by a subtable of

%Two tables $Tcoract$, $Tcorpass\in \Tid(d', 1, d(d-1)/2)$ 
%$duo_d [d(d-1)/2] \rightarrow d \times (d-1)$


%$(v/d^k)\bmod d < \sigma k$

%$(T'_1,T_1,forb_1)$ ou ($forb_1$ peut peut etre etre interpretee comme $init(X_{T_1})$?)

  \subsection{Causal dependence JS\red{tocheck}}
  Simple causal denpdence is a level 1 JS each cell depends on the cells above it.
  
  
  Advanced two level causal depencde use gflow cells anchor set of cells above and use right composition
  
  so the gflow JS (gflow) has anchors that can be read ordered  by order in the pattern of $\mathcal{P}(V)$ and cells $V$
  the order of the pattern has to be consistant with a timeline which leans that we can  right compose with a max JS with anchors $V$
  
associating to each cell an element of V so the anchors are $V$ and the cells below are the ones that has a dependcy on $v$    




Given  A causality JS $JC{^{(l)}}$  which anchors $ST_1, \ldots ST_n$ such that $ \sqcup (Jset)_c St_i =[l]$ anf cell number $T'_1\ldots T'_j$ with $\sqcup T'_i=[l]$

it is updated by $sqcor,p,sqcor'$ with $p\in [l+1]$ and for $j\in [l], sqcor_j\in \{0,1\}$ and $sqcor'_j\in \{0,1\}$
With 

$Update (sqcor,p,sqcor')(JC)$: 

$FJ_{l+1}= JC{^{(l)}}$ 

for $i=1$ to $p-1$ if $scor_i$,  $FJ_{i-1}= Merge_p (i) (FJ_i)$

for $i=p$ to $l$ if $scor'_i$,  $FJ_{i-1}= Merge_i (p) (FJ_i)$
 
 \subsection{Intuitive vision of Interactive problems and regular selection trees}

%We first introduce a slight generalisation to the notion of problem (or language) to add the notion of time.  

A  limited interactive problem can be seen as a referee and a player. The referee iteratively proposes 1 parameter and $d$ values to the player who has to pick one value. The referee can adapt the parameter depending on the answer.  The game is a language on the parameters, the sets proposed and the answers as an ordered sequence. 
It has to be defined  for any permutation of the variables (but not necessarily invariant). A proper formalisation will depend on the limitations. \footnote{(is it like IP with Polynomial referee, a multiround MA ?)}

%  a person who ahs to pich one of the values, on variables $X_1\ldots X_n$, consists in chosing . iteratively defining a domain of $d_1$ possible distinct values for $X_1$ , th  

   
  A $d'_1,\ldots , d'_n$- selection tree is a tree where the root (vertex at level 0) has (out)degree $d'_1$ and all the vertices of level $l$ have (out)degree $d'_{l+1}$)
     the edges represent the freedom of choices in such a problem.
 
 Given a sequence of indexed  values variables  $V_1,\ldots , V_n$ each $V_i$ having $d_i+1$ indexed ${\mathbb Z}$ elements : $V_i= (V_{i,0},V_{i,1},\ldots ,V_{i,i_{d_i}}$), with $V_{i,0}=b_i$ and   the elements indexed by $[d_i]$ forming a set of $d_i$ distinct values in ${\mathbb Z}$.
   A labeling of the edges  of  a $d'_1,\ldots , d'_n$- selection tree  associating to outgoing edges at level $l$ different indices in $[d_i]$  describes a family of solution to the problem.
   
% elements (the parameter then the sequence of choices proposed), and a sequence $d'_1\ldots d'_n$ with $d'_i\le d_i$, a $d'_1,\ldots d'_n$-labeling consists in labeling each  edge or equivalently each non root vertex  $u^{(l)}$  at level $l \ge 1$ (the root is counted as level 0),   with an integer $ lab(u^{(l)})\in[d'_i]$ so that childs of a vertex have different labels. The labeling at level $l$ depends only on the vertex (path from the root) and on the sequences $V_1\ldots V_l$.
 
 %  for any parameter choice $p_1\in \mathbb Z$, and for any choice of $d_1$ values in $D_1$, one can take $X_1=V_{l(v}$
   
   %or all  pair $V,B$  built by following the label $P(D,V,B)$
   
%  A labeling of the edges of such tree can be read as  
   A solution tree is a general solution to a problem if for each instance there exists a labelling such that taking the value indexed by the next vertex allows to be in the language.  

\end{document}
tout nouveau symbole est 
 \red{[Ecrire à la position ${\text{@Tete}}$  Direction]  [Ecrire à la position ${\text{@Tete}}$  [Symbole($\text{@Balle}$)]} ou Direction est $\uparrow$ ou $\downarrow$     
\red{Si  Està  Borddroit  \@  Balle alors  Inverser Direction Pousser MurDroit }
\red{[Si Està  Bordgauche \@ Balle alors }
\section{Percepts, color producing  and local days}

Positive integers are denoted $\oN$.

A sequence of colors indexed with $\oN$, denoted   $Cseq=1,c_1,c_2,\ldots$  is a sequence such that  for any $i\in \oN$ , $c_{i+1}\in \oN \setminus (\cup_{i'\le i} c_{i'}\cup \{1\})$ 

A color producing process or JDay $J$ is a sequence of colors equipped with a strictly increasing map Day: $\oN \rightarrow \oN \cup \infty$  such that  the name of $k$,  $\mathfrak n (k):=Day(k)$ is the time the process requires to produce its $k^{th}$ smallest value (color) in Cseq. Note that it may be uncomputable if it never converges but not when given a time limit for the production and it always have a lower and upper bound in $Cseq$. For $j=\mathfrak n (k)$, we denote $\iota_J (j)=k$, and $\mathfrak v_J (j)=c_j$.

\begin{figure}
    \includegraphics[width=\textwidth]{JSDay}
\caption{Jardin suspendu of Days }
\end{figure}

A JS Picture of  the color producing Process at day $j$  when showing its $k^{th}$ produced color is called an atom, is  denoted $\at {c_j} k$ and depicted  $\begin{ytableau}
\none &\none[c_j]\\
\none&k
 \end{ytableau}$ 
%Given $C'\subset Cseq$ and $C''\subset Cseq \setminus  C'$

 %$\mathcal {P}(\oN) \rightarrrow \oN$ giving the time required to produce the $k^{th}$ smallest missing color assuming the colors in $Old$ has been produced.

The $l$ fragment of a color producing process correspond to its restriction to the $l$ first terms of $Cseq$. It becomes a very nice structure with computable names.

Any atom $j=\at a k$ has divinatory a successor $\Sucs (j) = \at {a'} {k+1}$ if it is in the JS it is  called real or concrete otherwise it is called imaginary and all larger cell numbers are imaginary. Any atom with index $\iota_j:=k>1$ %either
 has a memorial Successor $\Suct (j)= \at {k} {a'}$.% or for all $k'<k$ there exists no $a'$ such that $\at {a'} {k'} \in J$

The canonical Color production have name  equals to indices and equal to values that are $\oN$.

A JS picture of Days  $J$ is a sequence of atoms of a color producing process.

Its Global Map cast is $GM:\oN \rightarrow {\mathbb N}$, $q\mapsto a$ with $\at a q \in J$ iff $(GM)_c J (q) =a$ 

In truth $J$ is isomorphic to $GM, t_{Prod}, t_{Verif}$ but here 
for this note we consider it as the quotient for all consistant timings to have an equivalence with the cast as the global map. In fact we will retrict ourself here to mathematical symmetrisable statements and use the concept of biexploration to properly define a JS multiverse over space and time (See section \ref{glob}).

Its Partial reversemap is $PRM:\oN \times \oN \rightarrow {\mathbb N}$, $(a,i)\mapsto q_i$ with $q_1<q_2<\ldots < q_i$ with $\at a {q_j} \in J$ for $j\le i$ iff $(PRM)_c J (a,i) =q_i$
and if $PRM(a,i)=0$ then for all $i'>i$ $PRM(a,i')=0$.

Its Reverse Map is $Liane:\oN \rightarrow \oN^*$, $a\mapsto ( (PRM)_c J (a,1), (PRM)_c J (a,2), \ldots$ 
 
As each color will appear once the  $(PRM)_c J (a,i) =0 $ if $i\neq 1$. 
 
The half JS $J$ (doubly infinite biexploration instead of one side is infinite) corresponds to the global map $\at {(\Suct_J j ,\Sucs_J j)} j $, and $RM(a)=(q)$%,0,0,\ldots$ 
 if $j=\at a q\in J$,  

The Correction JS to the map to the canonical color producing (to build a conceptual identity)

\red{counting $\# \li_k$, index  $| | \li_1^k$ , pairs (two lianes), fermionizing (sup dir)}
\section{ The global view of the multiverse of Jardins taking local time as primordial concept}
\label{glob}
An  MJS : Multiverses of Jardins Suspendus  is   family of elements defined at any level  $l\in \oN$, with   a sequence $n:=n_1,n_2 \ldots$ with $n_j \in \oN$ for all $j\in \oN$   and a sequence $w:=w_1,w_2 \ldots$ with $w_j \in \oN$ for all $j\in \oN$ . They are called
\begin{itemize}
\item     Conceptual or Left or departure or time anchored Jardins  $\lJ ^{(l)}_{n,w}\simeq Time  \times [n_1] \times \ldots \times [n_l]$ with $Time  \simeq  \oN$. \footnote{it is only when level is fixed that we will consider that only bounded time matter and that we won't need  a sequence $\oN^*$ but that we can just flatten it to $\oN$} 
\item  Perceptible or  right or arrival or space anchored  Jardins $\rJ ^{(l)}_{n,w}\simeq [n_1] \times \ldots \times [n_l]\times Space$ with $ Space \simeq \oN$  
\item   demiJardins or Jardin dynamics or entangling Jardins or superpositions of jardins %uses an extra sequence  $w':=w'_1,w'_2 \ldots$ with $w'_j \in \oN$ for all $j\in \oN$, and is 
  $\bJ ^{(l+1/2)}_{n,w} \simeq Time \times [n_1] \times \ldots \times [n_l]\times Space$.
 \end{itemize}
 
 The multiverse does not contain the empty Jardin suspendu $J_\emptyset$. The first index is a local time counter the last is a space counter, the intermediate ones is a gradation and are called Conceptual Lianes.

Any Jardin Picture $J_{x}$ of the MJS  has a finite length or degree $l=length(J_{x})$  and a finite space time origin $O_{J_{x}}\in \oN \times [n_1] \times \ldots \times [n_l]\times \oN$.
The projection of a Jardin  fixing $n_j$ is called a Liane. The first element is relative time and the last imaginary space both are not depicted in the picture of the JS with squares.
\footnote{ We close possible creation by symmetry (permutation of the lianes of indices  $n_j$). Global time-space dualities can be thaught as reversible maps between the first index and either the sequence of the others or the equivalence classe by permutation of the others.}


There exists three family of isomorphisms \textbf{values} $\mathfrak v$,   \textbf{ indices } $\iota $, and   \textbf{names}
$ \mathfrak n$ such that  any demi Jardin   $J^{(l+1/2)}_{x}$  has an outgoing isomorphism  
 $\mathfrak v _J:  JS^{(l+1/2)} \rightarrow JS^{(l+1)}$  and \textbf{ index }$\iota _J:  JS^{(l+1/2)} \rightarrow JS^{(l)}$ and an in-going isomorphsim  $\mathfrak n _J:  JS^{(l)} \rightarrow JS^{(l+1/2)}$.

Any left or conceptual Jardin has index and name  isomorphisms with $\oN$ and any perceptible Jardin have evaluation isomorphism with $\oN$.
 

For any $j$ and $v$ in $\oN$,  and for any $w'=w[j \mapsto v]$ defined by $w'_i=w_i$ if $i\neq j$ and $w'_j=v$  and operational time of reading $j$ and $w$, finding $j$  in $w$ and writing $v$ at position $j$ is denoted $\mathfrak t_{J,w}(w[j \mapsto v])$, there exist a family of  demi Jardins isomorphism \textbf {evol}  with $evol_{J^{(l+1/2)}_{n,w}}: J_{w'} \mapsto J_{n,w}$. %\red{check notation}
It is a contraction if $v\neq 1$ and $w_j=1$ and an expansion  if $w_j \neq 1$ and $v=1$.

%Each $J_{n,w}$, given a demiJardin $J'$  has a  family of divinatory successors $J_{Succ_d w'} $  with $w'=w_{_1\leftarrow v}$ where $v=w_1+1$  
  
  A demi jardin has infinite divinations  $\Sucs$ and memories $\Suct$
  %A Jardin has a divinations that is  infinite  and 
  A memory is imaginary if  its first index is  1.     
      
 For any triple $(l,j,j')\in \oN \times \oN \times \oN$
  
  If $L=L_1,L_2,\ldots$ with $L_i\in F$
\begin{itemize}
\item   the inversion $\Suct_{L,j}:=L'$ built with  $L'_{i}=L_i$  if $i\neq j$ and $L'_j= \Suct_F(L_j)$   i.e. $L'=L[j \mapsto  \Suct_F (L_j)]$
 \item the inversion  $\Suct_{L,i,j}$ swaps $L_i$ and $L_j$. \footnote{it should  not  require an auxiliary space when we have a canonial all 1 space divinated in advance}
 \end{itemize} 

%\footnote{
%draft attempt
%$J_{A,B}=Di(J_A,B)=|Mem(J_B,A)|$  
%if $B=(FJ_{A,C})_b \in J^{(t_A+1/2)}\cap J^{(t_C-1+1/2)} $ then   
%\begin{itemize}
%\item $(\rightarrow B):= Cor_D(b) J_C$
%\item $(\leftarrow B):= Cor_M(b) J_A$
%\item $(\downarrow B):= Cor_{Invop}(b) B= (FJ_{C,A})_{Cor_{Invpos}(b)}$
%\end{itemize}}

\subsection{Pre Algebra : Grading, Jardin dynamic , demi jardin $\cup_lJS^{(l+1/2)}$, concept emergence,Spin birth, definition of space}

 $\Seq_j(X)$ denotes the process $(X[j\mapsto 1], X[j\mapsto 2], \ldots)$,
 
 Each term $X[j\mapsto 2]$ has associated production time $\ti {X[j\mapsto j']}$ depending on $j'$ and the position of $j'$ in $X$. 
 
  $Seq_j(j)=\oN$  and $(Seq_j (X_j))_k:=X_k$ when $j$ does not appear in $X$.\footnote{  It has an increasing  building time for intervals $[n]$ i.e. $\mathfrak t _{\Seq_j}(j)$ for $j>1$  is the time of summing  the creation time of $j$ with  $\mathfrak t _{\Seq_j}(\mathfrak v_{j-1})$ }

 
$\Seq_{j_{|l}} (X):=\Seq_{j_{|j\in[l]}} (X):=((X[j\mapsto 1], \ldots,  X[j\mapsto l])$ where $:=$ denotes "denotes".

\footnote{In  UJS the universe of Jardins Suspendus $\Sucs_{UJS}  X_j:= X_{j+1}$  }

Given a sequence $F=(F_1,F_2,\ldots):=Seq_j F_j$, we say that $F:= ((Set)_c F,<_F)$ 
\begin{itemize}
\item   the cast of $F$ as set of its elements (forgetting the order) is  $(Set)_c F$
\item  $<_F$ is a $(1,j,j)$ reading of $\oN$ which means that each  $j$ in $\oN$ appears only once and  at local time $j$. (reading wil be defined later as number of occurences, position of first occurence, position of last occurence)
\end{itemize}

 For $k\in \{1,2\}$, we say that $F$ is a $1/k$ grading over $A$  and denote  $F:=\Fa  A$ if:
 \begin{itemize}
\item   $F=(F_1,F_2,\ldots)$ with $F_1=A$    
\item  there exists a divination/exploration  
$\Sucs_{F,x}:   \oN \times \Spin \times  (Set)_c F   \rightarrow \oN \times \Spin\times 
 %\oN \times \Spin 
 (Set)_c F$ and a memory 
%$\Suct_F:   \oN \times \Spin\times (Set)_c F   \rightarrow \oN \times \Spin\times 
%(Set)_c F$
$\Suct_{F,y}: \oN \times \Spin \times (Set)_c F  \rightarrow \oN  \times \Spin \times (Set)_c F $
%acting on $F$ with , %$\vgx {Spin} 1 \uparrow =1$, $\vgx {Spin} 1 \downarrow =2$ 
 such that  \red{TODO}
 \begin{itemize}
 \item $(\Suct_F(j,\uparrow,F_u))_2=\downarrow $ 
 \item $(\Suct_F(j,\downarrow,F_u))_2=\uparrow$
% \item $A\in F$ if  there exists $j\in \oN$ such that $A=F_j$
 %\item $\Sucs_F^{(s,0)} A=i_F A$ for all $s \in \Spin$ and for all  $A\in F$ 
 \item $(\Suct_F (u,\Sucs_F(j',s, F_j))=(\Sucs_F(j',\Suct_F (u,s,F_j), F_j)$ 
 \item $\Sucs_F ^{(s,j)} (\Sucs_F ^{(s,j')} (F_j))=\Sucs_F ^{(s,j+j')} (F_j) $ for all $j,j'\in (\oN)^2$, and for all $s\in \Spin$
 \item There exists $s\in \Spin$ such that $\Sucs_F ^{(s,j')} F_j\neq J_\emptyset$ for all $j' \in \oN$ 
 \item For all $j,j'\in (\oN)^2$, and for all $s\in \Spin$, if  $\Sucs_F ^{(s,j')} F_j=J_\emptyset$  then $\exists D\subset \Spin$ $\#D=k$ such that for all $s\in D$ for all $j''\in \oN$,
 $\Sucs_F ^{(s,j'+j'')} F_j\neq J_\emptyset$ 
\item all $A \in F$ have a finite local time in $F$ where the \textbf{time} in $F$ of $A$ is defined by :  $\tim F A=\inf \{j' \in \oN, \, \exists s \in \Spin,  \Sucs_F ^{(s,j')}  F_1=A\}$\footnote{correction distance is  local time}     
 \end{itemize}
\end{itemize}

A Family  $F$ have canonical projections $P_F$ such that $F=((P_F)_1 F, ((P_F)_2 F, \ldots)$ 

$\oN$ is  the 1-Family over 1 with $\rightarrow^{(s,j)}_{\oN}:=x\mapsto x+1$,  $\downarrow_{\oN}:=x\mapsto x-1$ with $\downarrow_{\oN} 1= J_\emptyset$

$\mathbb Z$ is  the (1/2)-Family over 0 with $\rightarrow^{(s,j)}_{\mathbb Z}:=x\mapsto x+1$,  $\downarrow_{\mathbb Z}:=x\mapsto -x$

Notice that if $J\in F$ then  $(P_F)_{\tim F J} F =J$ (fix point property inside a Family)

The {\bf exploration degree} of $A$ in $F$ is  $\degexp_F(A)=\sup \{k, \exists D\subset \Spin,\, \#D=k,\, \forall s\in D, \exists j\in {\mathbb N},
 \rightarrow^{(s,j)} F_1 =A\}$    and the exploration degree of $F$ is  $\inf_{A\in F} \degexp_F(A)$. \footnote{we can sometimes want to define open families and take $j\in \oN$ and I guess we will define $oF=F\setminus F_1$}  %The correction distance of $A$ in $F$ is 

If $A$ is a ${\mathbb C}^* $ algebra then one can take as Memory the * morphism and as divination adding the identity element.

\red{is Multiplication is the second order of divination and power  the third order?\footnote{ is reverse map a higher order memory map (applied to the whole family inside an other family)}}
  \section{Brouillon}


%\red{ Is $F_{j,j'}$ the map $F_j\mapsto F_{j'}$ with $F_1, \Sucs_F,\Suct_F$ as structuring geometry.? 1/2 Jardin : 2 infnite , Jardin 1 infinite }

  \subsection{Correlations /spin/ inversions}
$\begin{array}{lcl}
Spin&=&\uparrow,\downarrow, 
\uparrow \uparrow,\uparrow \downarrow, \downarrow \uparrow,\downarrow\downarrow, 
\uparrow\uparrow \uparrow,\uparrow\uparrow \downarrow, \uparrow\downarrow \uparrow,\uparrow\downarrow\downarrow,
\downarrow\uparrow \uparrow,\downarrow\uparrow \downarrow,\downarrow \downarrow \uparrow,\downarrow\downarrow\downarrow
\sqcup \ldots \\

&=&(\uparrow,\downarrow)\sqcup 
(\uparrow \uparrow,\uparrow \downarrow, \downarrow \uparrow,\downarrow\downarrow) \sqcup 
(\uparrow\uparrow \uparrow,\uparrow\uparrow \downarrow, \uparrow\downarrow \uparrow,\uparrow\downarrow\downarrow,
\downarrow\uparrow \uparrow,\downarrow\uparrow \downarrow,\downarrow \downarrow \uparrow,\downarrow\downarrow\downarrow)
\sqcup \ldots \\

&=&\red{((\uparrow),(\downarrow))\sqcup 
((\uparrow \uparrow,\uparrow \downarrow), (\downarrow \uparrow,\downarrow\downarrow)) \sqcup 
((\uparrow\uparrow \uparrow,\uparrow\uparrow \downarrow, \uparrow\downarrow \uparrow,\uparrow\downarrow\downarrow),
(\downarrow\uparrow \uparrow,\downarrow\uparrow \downarrow,\downarrow \downarrow \uparrow,\downarrow\downarrow\downarrow))
\sqcup \ldots} \\
&=&\red{(((\uparrow)),((\downarrow)))\sqcup 
(((\uparrow \uparrow),(\uparrow \downarrow)), ((\downarrow \uparrow),(\downarrow\downarrow))) \sqcup 
(((\uparrow\uparrow \uparrow,\uparrow\uparrow \downarrow), (\uparrow\downarrow \uparrow,\uparrow\downarrow\downarrow)),
((\downarrow\uparrow \uparrow,\downarrow\uparrow \downarrow),(\downarrow \downarrow \uparrow,\downarrow\downarrow\downarrow)))
\sqcup \ldots }
\end{array}$

The correlation family is graded ie there exists a bijective grading map $\oN \rightarrow \oN \times \oN$ , $j \mapsto (l,j')$ with first projection  equal to $\oN$ and called the level or the grade (number of spins). \red{We can continue subgrading by dividing by 2 each grade containing more than one element and just adding exra parenthesis when one element.}

$(Spin)_i =\at i q$ with $q\in (\uparrow,\downarrow)^k$ with $i\in [2^k-1,2^{k+1}-2]$ such that  \footnote{\red {(not clear if it is an abstraction or an evaluation to chose which is which between$Spin_1$ or $Spin_2$)}}
$\begin{cases} Spin_{2i+1}=(Spin_i \uparrow) \\ Spin_{2i+2}=(Spin_i \downarrow)\end{cases}$ (starting as if $Spin_0$ existed and was equal to nothing i.e. $J_\emptyset$.
 

 
 Its grading allows to write  $Spin=\sqcup_l (Spin) ^{(l)}$ with  $Spin ^{(0)}= J_\emptyset$ by definition.
 
 $Spin_{(l,j)}=Spin^{(l)}_j= Spin_{2^{l-1}+j-1}$ with $j\in [2^l]$
  
  
 $\uparrow_j$ acts as the  operator  $Spin\rightarrow Spin$ that is the identity of the $j^{th}$ Spin 
 
 $\downarrow_j$ acts as the  operator  $Spin\rightarrow Spin$ that is the involutive flip operator of the $j^{th}$ Spin .

\red{
 $(\downarrow_j)^{(l_1,l_2,\ldots,l_k)}$ acts as the  operator  $Spin\rightarrow Spin$ that is the involutive flip operator of the  $(\sum_i  (l_i-1)  + j)^{th}$ with $l_i\in [2^{k-i},2^{k-i+1}-1]$ Spin 
 
 and the concise notation omitting the 1  $(\downarrow_j)^{Dec(l_{i_1},l_{i_2},\ldots,l_{i_k})}$ with $i_1<i_2\ldots <i_k$ which corresponds to  $(\downarrow_j)^{(1^{i_1-1} l_{i_1} 1^{i_2-i_1} l_{i_2},\ldots,l_{i_k})}$}
  
% ???\red{ very fast to clear} .
% Given a sequence $i_1<i_2\ldots i_k$  with $i_j\in \oN$ for $j\in [k]$
%$Dec (l_{i_1},l_{i_2},\ldots,l_{i_k})=1^{i_1-1} l_{i_1} 1^{i_2-i_1} l_{i_2} \ldots 1^{i_k-i_{k-1}} l_{i_k}$
  JS of $\#$ with Spins as positions restricted to 12
  
  $ \begin{ytableau}
\none [_1]&\none[_2]&\none[_3]\\
  \uparrow&\uparrow\uparrow&\uparrow\uparrow\uparrow\\
 \downarrow&\uparrow \downarrow&\uparrow\uparrow\downarrow\\
 \none &\downarrow\uparrow &\uparrow\downarrow \uparrow\\
 \none &\downarrow\downarrow&\uparrow\downarrow\downarrow\\
 \none &\none&\downarrow\uparrow\uparrow\\
 \none &\none&\downarrow\uparrow\downarrow\\
  \end{ytableau}$

The JS canonical reading is a proof reading (linae by liane), Liane $i$ cells are spins of length $i$,
 Scrolls    are fixed suffix spins completed with a sequence (may be empty of $\uparrow$). 
The diagonal are alternating chains starting with $\uparrow$
  
  \subsubsection{ Canonical encoding/decoding of monotone positive integer indexed sequences}
\red{
Given a sequence $T=T_{i_1},\ldots ,T_{i_{k+1}}$ with $k\in \oN$ and  $T_{i_j} \in \oN$ and $i_j<i_{j+1}$   for $j \in [k]$ 
  and a sequence indexed by $\oN$,  $L=L_1, L_2, \ldots $  with $L_j\neq 1$ for $i\in \oN$ we denote :
  
   $ReDec (L_{|T})({i_1},{i_2},\ldots,{i_k}):=1^{i_1-1} L_{T_{i_1}} 1^{i_2-i_1} L_{T_{i_2}} \ldots 1^{i_k-i_{k-1}} L_{T_{i_k}}$
and 
   $iDec (L_{|T})({i_1},{i_2},\ldots,{i_k}):=L\setminus \cup_j  L_{T_{i_j}}$ %\red{to define properly} 
 
 A family of correlations $Cor$ is such that $Cor_j$ , usee a parition $\Seq_{j'} n_{j',j}$ of $j$ a Spin in layer $n_{1,j}$ with elments separated wtiih the rest of  the partition to have labels in $[j]$ 
 
  }
  
\subsection{Relative time : Gap or $\mathfrak t _{loc}$ , bi-infinite exploration} \label{sec:gaps}

$Gap_1^{(1)}=(0,1,-1,2,-2,\ldots)$, $Gap_1^{(2)}=(0,-1,1,-2,2,\ldots)$, $Gap_2^{(1)}=(1,2,0,3,-1,\ldots)$, $Gap_2^{(2)}=(1,0,2,-1,3,\ldots)$, 
$Gap_3^{(1)}=(-1,0,-2,1,-3,\ldots),\ldots$

More formally for $n\ge 0$, $j>0$, $j'>0$, 

 $\begin{cases}(Gap_j^{(2n+2)})_{1}=(Gap_j^{(2n+1)})_{1}\\ (Gap_j^{(2n+2)})_{2n'}=(Gap_j^{(2n+1)})_{2n'+1}\\(Gap_j^{(2n+2)})_{2n'+1}=(Gap_j^{(2n+1)})_{2n'}\end{cases}$

%$(Gap_{2n'+1}^{(2n+1)})_{j'}=(Gap_{1}^{(2n+1)})_{j'} -n'$
%$(Gap_{2n'}^{(2n+1)})_{j'}=(Gap_{1}^{(2n+1)})_{j'} +n'$

And for $j''>1$, $(Gap_{j''}^{(2n+1)})_{j'}=(Gap_{1}^{(2n+1)})_{j'} +-(1)^{j''} \lfloor j''/2\rfloor =(Gap_{1}^{(2n+1)})_{j'}-(Gap_{1}^{(2n+1)})_{j''+1} $
and for $j\ge 1$,  $$(Gap_1^{(1)})_j=(-1)^j \lfloor \frac {j-1} 2 \rfloor$$. 

We might use  $Gap_j^\uparrow$ for  $Gap_j^{(2n+1)}$ and  $Gap_j^\downarrow$ for  $Gap_j^{(2n)}$ 

 $$ (Gap_j^\uparrow)_{j'}%= (j' \mod 2)+  (-1)^j \lfloor \frac {j} 2 \rfloor- (-1)^{j'+1} \lfloor \frac {j'+1} 2 \rfloor
  = j' -2 \lfloor j'/2 \rfloor + (Gap_1^{\uparrow})_j-(Gap_1^{\uparrow})_{j'+1}$$
 
%\footnote{ can be formulated as $(Gap_j^\uparrow)_{j'}=(j' +1 \mod 2)+  (-1)^j \lfloor \frac {j} 2 \rfloor- (-1)^{j'-1} \lfloor \frac {j'-1} 2 \rfloor$}.
and $ (Gap_j^\downarrow)_{j'}= - (Gap_j^\uparrow)_{j'}$ 
% 1+ (-1)^j \lfloor \frac {j} 2 \rfloor- (-1)^{j'-1} \lfloor \frac {j'+1} 2 \rfloor$$ 


We will use a grading to  define for $j\in \oN$ and $b \in {\mathbb Z}$ and $l\in \{\uparrow,\downarrow\} \sqcup {\mathbb N}  $,  \footnote{should be oriented sum to be clean}   %$Gap_{l,j,j'}=(Gap^{(l)}_j)_{j'}$

$\vgx{Gap} {l,j} (j')= (Gap_{l,j})_{j'}= (Gap^{(l)}_j)_{j'}$ and $\igx{Gap} {l,j} (b)= j'$.

 A $b$-Gap is a Gap starting with $b$ as the value of the first element in $\mathbb Z$ i.e if $G=Gap^{(l)}_j$ we say that $G$ is a $(Gap^{l}_j))_1$-Gap. 


Somehow we will consider repeatable  formalisations and will use corrections to make second order dynamics match an indistinguishability, so we will work with $\oN \times[n_1] \times \oN$ and with $\oN \times[n_1] \times  [n_2] \times \oN$

Note that if we take as cost for a gap reading the distance crossed reading $n$ elements costs $\sum_{j'=1}^n  |(Gap^{(l)}_j)_{j'+1}-(Gap^{(l)}_j)_{j'+1}|= n(n+1)/2$.

\subsection{Permutations/ traductions / map biunivoque/ Group correction/see phantoms}

Given  a time sequence $T=T_1,T_2,\ldots$ with $T\simeq \oN$ and $T_j\in \oN$ and a space sequence  $T'=T'_1,T'_2,\ldots$ with $T'\simeq \oN$ and $T'_j\in \oN$

$\mathfrak v :T' \rightarrow \oN$ and $\iota: T\rightarrow \oN$ are the space and time canonical isomorphisms.

For $l\ge k$, we define

$JTrad_{T\mapsto T'}^{(|_l)}= \at {T'_{j_1}} {T_{j_1}} \sqcup \ldots \sqcup  \at {T'_{j_k}} {T_{j_k}}$ with  

for any $j'\in [l]$,\  $\begin{cases} T_{j'}=T'_{j'} \text{ iff }j\not\in \cup_{j'\le k} j_{j'}\\   
\min (T_{j_{j'+1}},T'_{j_{j'+1}} )> \min (T_{j_{j'}},T'_{j_{j'}} ) \end{cases}$

The map is fully captured at level $l$ is $\forall j'>l,\, T_{j'}=T'_{j'}$.

The cast as Spin and the complementary cast are

$
\begin{cases}
((Spin)_c)_{j'} J =(\Suct_{j'})^{T_{j'}>T'_{j'}} \uparrow_{j'}^{T_{j'}\neq T'_{j'}} \text{replaces each atom by } \uparrow \text { or } \downarrow  \text{ to obtain a Spin of length } k \text{ at positions } j_{j'}\\
((Spin)^c)_{j'} J=(\min(T_{j'},T'_{j'}),\max(T_{j'},T'_{j'})-\min(T_{j'},T'_{j'}))^{T_{j'}\neq T'_{j'}} 
\end{cases}
$ 

An evolution compatible with $J^{(l)}$ as a starting point is 
$J^{(l+1/2)}= \smd \at {tw(T'_{j_1})} {(T_{j_1},T''_{j_1})} \sqcup \ldots \sqcup  \at {tw(T'_{j_k})} {(T_{j_k},T''_{j_k})}$

An evolution compatible with $J^{(l+1)}$ as an  end  point is 



\red{
Commutativity corrector of + should be $Cor_{\Suct_+} + B A = +  A B$

Associativity corrector (assume default reading is from percepts so by value and $._R$ is the right multipication so shoud be$\Sucs^{(2)}$ when it will be defined)  )$Cor_{\Suct_{._R,A}} (AB)C= A (BC)$ as we want to map $._R A ._R B C $ to $._R (._R A B) C$ or ($Cor  A^{(1)} Par^{(2)} B^{(1)} C^{(1)}:=Par^{(2)} A^{(1)} B^{(1)} C^{(1)}$ parentheiss are for clarity if assuming value by default) 
}
\footnote{when can we  have reduced form were subjacencies are involutions vertices of graphs becaming stable indistinguishable sets?}
\subsubsection{The 6 bases of maps}

 Aside the matching representation any JS can be represented by orbit or by insertion (and the * like  memorial morphism will double by looking from the lines point of view inverting the liane of concepts with its percept)
  
$ JOrb_{T\mapsto T'}^{(|_l)}= $% (\at {T'_{i_1}} {T_{i_1}} \sqcup \ldots \at {T'_{i_1}} {T_{i_1}}) \sqcup \ldots \sqcup (\at {T'_{i_k}} {T_{i_{k,1}}} \sqcup \ldots \at {T'_{i_1}} {T_{i_1}}) $


$ JIns_{T\mapsto T'}^{(|_l)}$

 $\Suct_{JX}(JX_{T\mapsto T'}^{(|_l)})=JX_{T'\mapsto T}^{(|_l)}$


%\subsection{Permutations as JS}

%Let $Idx$ be the set $\overset{\circ}{\mathbb N}$ indexed by ${\mathbb N} $  so that  $Idx_i=i+1$ for any $i \in {\mathbb N} $ 
%PermJS is a family of canonical casts of permutations as JS.
%
%A PermJS$_1$ is an %$(i\mapsto |Liane_i |)$
% $(\mathfrak v_{Idx})$JS where anchors represent an indexation of the orbits following the columns such that each Liane is an orbit.  % with non increasing size (its shape is a rotated permutated Young Diagram)
%The permutation is read Liane by Liane by mapping each cell number to the cell number below it and by mapping the last one to the first one.% \red{toncheck/clean/correct}
%
%$\sigma'^{l+1}= \sigma_{p,l+1}\sigma'^{l}$ means that the PermJ is adapted by inserting the cell $l+1$ at  Liane$_{p+1}$. 

%An OrbPermJS is a JS where the Lianes represent the orbits read by columns.
%
%  $ (OrbPermJS_1)_c (
%\left(   \begin{matrix} 0&0&1\\
%0&1&0\\
%1&0&0
%  \end{matrix} \right ))= \begin{ytableau}
%\none &\none[_1]&\none [_2]\\
%  \none&1&2\\
% \none &3\\
%  \end{ytableau}$
%  
%  
%As 1 is always in the first column we could have a more concise notation inserting  (image of 1-1), then (image of image of 1 -1) until getting back one the create a new column inserting the first missing element. using $n-k$ cells where $k$ is the number of orbits.
\red{old :
MatchpermJS  is the JS representing permutations as a matching $\sigma \in S_n$, $\sqcup_{a\in [n-1],\sigma(a)\neq a} \at a {\sigma (a)} $ it has atoms in $FAT^{(2)}_{\neq}$
where FAT is the old family if atoms  \footnote{Matchperm$_2$ can be built by sorting the positoin and using the inverse permutation}. 

$InsPermJS_1$ is defined by inserting  in last line: for $i\in [n-2]$,we insert  the cellnumber $i+1$ at the position $p\in [i+1]$ at the line $i+1$ \footnote{$InsPermJS_2$ is defined by inserting at last columns instead of at last line} This allows to reversibly associate to any permutation an infinite family of permutations inclduing a canonical decomposition and a canonical completion.

OrbPermJS$_1$ is the JS representing permutations with the anchor being the smallest undiscovered element with non trivial orbit. It has atoms in $FAT^{(2)}_>$  (atoms all have  $(p>a)$). It have   $n-k$ cells where $k$ is the number of orbits.\footnote{ OrbPermJS$_2$ is  the permutation matrix reading by the lines instead of the columns.}
}

%$InsPermJS_1$ is defined by inserting  for $i\in [n-2]$, the cellnumber $i+1$ at the position $p\in [i+1]$ at the line $i+1$

 Consider the permutation  $\sigma \in S_7$ defined by  $\sigma (1)=1,\sigma (2)=3,\sigma (3)=7,\sigma (7)=2,\sigma(4)=4,\sigma(5)=6,\sigma(6)=5$ canonically extended to $S_n$ with $n>7$  by   $\sigma(i)=i$ for $i>7$.

$\sigma =\left(   \begin{matrix} 
1&0&0&0&0&0&0\\
0&0&1&0&0&0&0\\
0&0&0&0&0&0&1\\
0&0&0&1&0&0&0\\
0&0&0&0&0&1&0\\
0&0&0&0&1&0&0\\
0&1&0&0&0&0&0\\
  \end{matrix} \right )$ which is built by inserting last line as 
   $1\xrightarrow{2} \left(   \begin{matrix} 
1&0\\
0&1\\
  \end{matrix} \right )\xrightarrow{3}$
  
    $\left(   \begin{matrix} 
1&0&0\\
0&1&0\\
0&0&1\\
  \end{matrix} \right )
  \xrightarrow{3}
    \left(   \begin{matrix} 
1&0&0&0\\
0&1&0&0\\
0&0&0&1\\
0&0&1&0\\
  \end{matrix} \right )
\xrightarrow{4}
    \left(   \begin{matrix} 
1&0&0&0&0\\
0&1&0&0&0\\
0&0&0&0&1\\
0&0&1&0&0\\
0&0&0&1&0\\
  \end{matrix} \right )
\xrightarrow{4}
  \left( \begin{matrix} 
1&0&0&0&0&0\\
0&1&0&0&0&0\\
0&0&0&0&0&1\\
0&0&1&0&0&0\\
0&0&0&0&1&0\\
0&0&0&1&0&0\\
  \end{matrix} \right )
\xrightarrow{2}  \sigma$ 


 $(OrbPermJS_1)(\sigma)=\begin{ytableau}
\none &\none[_2]&\none [_5]\\
  \none&3&6\\
 \none &7\\
  \end{ytableau}$,  and  $ (InsPermJS_1)_c (\sigma)
= \begin{ytableau}
\none &\none[_2]&\none[_3]&\none[_4]\\
  \none&1&2&5\\
 \none &6&\none\\
  \end{ytableau}$

 $(MatchPermJS_1)(\sigma)=\begin{ytableau}
\none &\none[_2]&\none [_3] &\none [_5] &\none [_6] & \none [_7]\\
  \none&3&7&6&5&2\\
  \end{ytableau}$

\footnote{old 
There are therefore 6 Canonical permutation forms  associating to any $\sigma\in S_n$, 2 canonical pasts in $S_k$ with $k<n$ and a canonical future in $S_k$ for $k>n$.

%   not putting it redaing the columns inserting $l-1$ where $l$ is the position of the line

The 8 first Families of Matching Permutations (\red {order is not optimized or naturalized} )
FMperm$_i$ for $i\in [8]$ is read in $(crr, crl,cll,clr,lur,lul,ldl,ldr)$ 
%$FPerm_1=FPerm_{crr}, FPerm_2=FPerm_{crl},FPerm_3=FPerm_{lrr},FPerm_4=FPerm_{llr},
%FPerm_5=FPerm_{clr}, FPerm_6=FPerm_{cll},FPerm_7=FPerm_{llr},FPerm_8=FPerm_{lll}$
  are indexed such that the insertion is done respectively column right, easy right extend for $FPerm_1=FPerm_{crr}$, 
  by lines by up (from top)  easy left extend ,for $FPerm_6=FPerm_{lul}$.

$(FMPerm_1)_{n! +p-1}$ for $p\in [l+1]$ is obtained by inserting by columns to the right and is read $(p_1,\ldots , p_{n+1})$ where $p_i$ is inserted at step $i$
$(MatchPerm)_c (FP_1)_{n!} = \at 1 1 \sqcup \at {n} n \ldots \sqcup \at {n} n=canJ_n$


$(FMPerm_3)_{n! +p-1}$ for $p\in [l+1]$ is obtained by inserting by columns to the left  and is read $(p_{n+1},\ldots , p_{1})$ where $p_i$ is inserted at step $i$
$(MatchPerm)_c (FP_6)_{n!} = \at n 1 \sqcup \at {n-1} 2 \ldots \at 1 n$

\red{TO CLEAN}
Using c-reduced Orb form and AJS we get $(FOPerm_1)_{l_1!+l_2! +p}$ with $p\in [3]$ we view Orb perm as metfunc cast of the min function where atoms as they can be pushed left are of the form $\at {a+b} {(a,a+c)}$  is  so that the anchor is the smallest element in the liane and cells may have moved so we forget their number. when inserting we chose a liane or insert a new liane to the (right/to the left) having as minimal element $w$+1 where $w$ is the number of cells . 

}
  
  
   

\subsubsection{Twisting (entangling) Operator, twisting with a correlation and  Jardins on Twists}
A basic twisting operator is `an operator that preserves biexplorability of counterfactuals" : we will use functions defined at some time at some domain and we want that at any time, we can define two directions such that only one direction  may contain the forbidden fruit even when extending the domain of the functions.\footnote{The intuition behind biexplorability is that we heave two parameterized lines that cover the surface that are disjoint at infinity and such as the zeros of the surface we are trying to avoid are in a bounded region.} 

$tw: \text{oldvalue} \times \text{newcorrelationspin} \times \text{new value of new variable}\rightarrow \text{newvalue}$
More formally   a Twist is a function  $tw :{\mathbb Z}\times  \{0,\uparrow,\downarrow\} \times {\mathbb Z}\rightarrow {\mathbb Z}$ such that 
\begin {itemize}
\item  $tw(x,0,y)=0 \implies \forall y' \neq 0,\, tw(x,\uparrow,y')tw(x,\downarrow,y')\neq 0$
\item  $\forall x \in {\mathbb Z},\,  \{y>0, tw(x,\uparrow,y)=0\}\cap\{y>0, tw(x,\downarrow,y)=0\}=\emptyset$  and 
\item $\forall x \in {\mathbb Z},\, |\{y>0, tw(x,\uparrow,y)=0\}\cup\{y>0, tw(x,\downarrow,y)=0\}|\le1$ % \implies \forall y' \neq 0,\, tw(x,+,y')tw(x,-,y')\neq 0$
\item $\forall x,s,y,\,  |tw(x,s,y)| |\le|x|+|y|$ % \implies \forall y' \neq 0,\, tw(x,+,y')tw(x,-,y')\neq 0$

\end{itemize}
The standard twist is $tw_1 := \begin {cases} x,0,y \mapsto x \\ x,\uparrow,y \mapsto x+y \\ x,\downarrow,y \mapsto x-y\end{cases}$.

Exchanging $\uparrow$ and $\downarrow$ preserves being a valid twist, i.e. we can define  $inv$ by   $inv(tw)(x,0,y):=tw(x,0,y),inv(tw)(x,\uparrow,y):=tw(x,\downarrow,y)$, and $inv(tw)(x,\downarrow,y):=tw(x,\uparrow,y)$.


\subsection{evaluation with correlations}

  
A level $l$  evaluation is a   function $ev$  such that $ev(x,p, v)$ represents the evaluation of $v'=\sigma_{l+1,p}v$ with the hyperplane like ($l$ extended) (JS)  $x'=\sigma_{l+1,p} (x,\uparrow)$, so that $x'$ is the  $l+1$ correlation obtained from $x$ by inserting $\uparrow$ at position $p$. The transformation have to use the same $p$ for all the cells. 
 

  
  %the value $v$ is obtained from the anchor of the cell $(s_1,\ldots,s_
  $ev: \{0,\uparrow,\downarrow\}^l \times  [l+1] \times {\mathbb Z}^{l+1}\rightarrow {\mathbb Z}$ is an evaluation if :
  given   $(x,p,y)\in \{0,\uparrow,\downarrow\}^l \times  [l+1] \times {\mathbb Z}^{l+1}$, and $x'=\sigma_{l+1,p} (x,\uparrow)$
$ev$ satisfies that     $\forall j\in [l+1]$ such that  $x'_j\neq 0$:
  % if ($j<p$ and $x_j\neq 0)$ or ($j=p$) or ($j>p$ and $x_{j-1}\neq0$)  then
   \begin{itemize}
 \item If $ev(x,p,y)=0$ then $\forall k\neq y_k,  ev(x,p,y[j \mapsto k]) \neq 0$
 \item If $ev(x,p,y)\neq 0$,  $| \{k,  ev(x'[j \mapsto \uparrow],p,y[j\mapsto k]) =0\} \cap  \{k,  ev(x[j\mapsto \downarrow],p,y[j\mapsto k]) =0\}|=0$
 \item If $ev(x,p,y)\neq 0$,  $| \{k,  ev(x'[j\mapsto  \uparrow],p,y[j\mapsto k]) =0\} \cup  \{k,  ev(x[j\mapsto \downarrow],p,y[j\mapsto k]) =0\}|\le1$
 \item $\forall x,p,y,\,  |ev(x,p,y)| |\le \sum_{i} |vCor(x'_i) y_i|$ % \implies \forall y' \neq 0,\, tw(x,+,y')tw(x,-,y')\neq 0$

 \end{itemize}
   %(    $|\{k>0, ev(x,p,f_E(s_{_j\leftarrow k}))\}=0|\le 1$ 
   where the operation  $x[j\mapsto k]$ maps $(x_1\ldots,x_l)_h$ to $(x_1\ldots , x_{j-1},k,x_{j+1},\ldots ,x_l)_h$ %and $x'=\sigma_{l+1,p} (x,+)$, so that $x'$ is the  $l+1$ correlation obtained from $x$ by inserting $+$ at position $p$.

\subsection{Family of  Correction/ decoding Groups JS (to recheck)}

An FG $\in$  FGS  is a JS whose Lianes are interpreted as a sequence of generators of groups: $Liane_a(FG)= \sqcup_i \at a {g_{i,a}}$ where $G_a=<g_{i,a}>$ is a group.  

An FGS  acts on  a JS, $J$  as element  of an FJS, $F$ if :

Let $j= \mathfrak i_F(J)$, %if $j=0$ then for any $g$,
 $FG_F(g,J)=J_\emptyset$ if $g\notin G_j$ .%  if $\not \exists i, F_i=J$



%A family of groups acting on JS from a Family of JS so that the $Liane_i$ of 
%$G$ contains $g_{1,i},\ldots,g_{l_i,i}$  generating a group $G_i$ acting in $J \in FJ_j$ so that

and   for any $g\in G_j$ there exists $j'$ with $|j'-j|\le desc_{G_j}(g)$ such that $FG_F (g , J)\in FJ_{j'}$

the description of $g$ in a group $G$ is: 

an element of $G$ can be written as a word using the generators and operators $inv_k$ were  $inv_k( x_1\ldots x_k)=x_k^{-1}\ldots x_1^{-1}$  and $desc (g ,r)=1+desc(r)$ if $g$ is a generator and $desc  ( inv_k , r)=k+desc (r)$

\section{Dictionary vision}
Recall that in Jardins by default any term has  3 associated integers corresponding to  the local times  (number of steps) required to write it, to check if it is nothing and to check that it is not nothing.
 A good visualising example of evolution of percepts in Jardins  when learning new concepts is the dictionnary.

$\mathfrak n _D(j) := \mot_{D,j}$, $\iota_D (\mot_{D,j}):= j$, $\mathfrak v _D(j) := \Enc( \Defn_D(\mot_{D,j}))$

The dictionary has $l$ definitions the first being   $"\mot_1: \Cpt (\mot_1)"$ the last if we cut after $l$ definitions is  $"\mot_l: \Cpt (\mot_l)"$. 

The conceptual weights $w_j$ is the number of words in the definition of $\mot_j$, $w$ being the total weight. 

A general reading  at time (number of words read) $t$ maps all concepts to an already  read sequence of words and a program for the future.% (permutation on the rest)

When time is fixed we consider the explorations or encodings $\Seqd {j} {w!}  \Expl_j $ that    are the permutations of all the words inside the definitions in a dictionnary with $w$ words after the symbol ':'.

 
\begin{itemize}

\item $A^{\delta_\phi}:=$ If  not $\phi$  (or $\phi==0$) then $A$ (Otherwise nothing)
 \item $A^{\delta_{\phi,1}}:=$ If   $\phi$  (or $\phi \neq 0$) then $A$ (Otherwise nothing) 
 
\end{itemize}

The definition of a word  starts with a word himself then with a symbol to separate from the rest  $\mot_{D,j,1}:=\mot_{D,j}$ and   $\mot_{D,j,2}:= ':' $ 

The time required to read the definition of $\mot_j$ in $D$ if each word was read in time 1 would be  $h_{D,j}=n_j+2$


$\Defn_D(\mot_{D,j}):=\Liane_D(\mot_{D,j}):= \Ppt_D (\mot_{D,j}) : \Cpt(\mot_{D,j}):= \Seq_{k_{|n_j+2}} (\mot_{D,j,k}) $




$\Cpt(\mot_{D,j}):=\Seq_{k_{|n_j}} (\mot_{D,j,k+2})$ 


$\Find_D(j,k,p):= (j,k,p)^{\delta_{{mot_{D,j}== (\Defn_D(\mot_{D,k}))_p},1}}$


An $(r,a,b)$ reading of $D$ is a sequence $\Seq_j X$ such that each word $w$ appears at most $r(w)$ times and after the index $a(w)$ and before the index $b(w)$ i.e.
 $1\le \# \{j, \Seq_j X == w \} \le r(w)$ and $1\le \min \{j ,\, \Seq_j X ==w\}$ and $\max \{j,\,  \Seq_j X == w \} \le b(w)$
 
In a Jardin each concept have  a divine  space successor $\Sucs$ and memorial time successor $\Suct$. 
 


After fixing time, lets consider a $(1,a,b)$ reading of $D$
 
\begin{itemize}
\item $\Infl_D(j)$ de  the sequence of indices $j'$ such that  $j'\in \Infl(j):=  \mot_{D,j'}$ uses $\mot_j$ in  its definition and let $w'_j = \# \Infl_D(j)$ be its length (the number of words using 
$\mot_{D,j}$ in their definition).
\item  $l'_j$ the index of the last word using $\mot_{D,j}$ in its definition.
\item  Given a word $\mot_{j'}$ using $\mot_j $ in their definition let
\subitem $f'_{j,j'}$ the first position in which $\mot_{D,j}$ appears
\subitem $h'_{j,j'}$ the last position in which $\mot_{D,j}$ appears
\end{itemize}


The counterfactual  action of a concept or the entanglement book of a dictionnary is obtained when fixing time by considering the interaction for each word with the others.

$\Line_D(\mot_{D,j}):= \Seq_{k_{|l_D}} (\Seq_{p_{|h_{D,k}}} \Find_D(j,k,p))$ finish inside before increasing outside

We can factorize using the influence list denoting $\Infl_D(j)=(k_1,k_2,\ldots, k_{w'_j})$.

$\Line_D(\mot_{D,j})=(j (k_1(p_{k_1,f'_{j,k_1}},\ldots, p_{k_1,h'_{j,k_1}}), k_2 (p_{k_2,f'_{j,k_2}},\ldots, p_{k_2,h'_{j,k_2}}), \ldots, k_{w'_j}(p_{k_{w'_j},f'_{j,k_{w'_j}}},\ldots, p_{k_{w'_j},h'_{j,k_{w'_j}}}))$ 
  
%$\Suct, \Sucs$


After an $(1,a,b)$ reading $R$ of $D$ one can consider the matrices $LE_{R,j',j}$ and $UE_{R,j',j}$ whose entries are  whose lines are $T'_1,\ldots T'_l$ conrresponds $\Line_D$ of the reading and columns $T_1,\ldots T_l$ are $\mot_1,\ldots \mot_l$ whose entries are first and last position or 0 if absent. 

\red{review as  Reverse Map or higher order memory map}

decompose the 1/2 dicto   into correctable versions where you can remove one product liane ?? and 

 build superposition over  of when can i know that it is noit this product 
  
SupSeq$_j=  \ket j  \ket {D' \setminus \Defn_{D'} (\mot_{D',j})}$  where $D'= \Seq_j  Entangle$


 \section{A practical application}
 %\begin{center}
  \begin{theorem} {[\bf{  123 Theorem (ex-conjecture)(2004-2024)}]}\cite{keusch2024solution}
  % \footnote{for more info about the conjecture see https://arxiv.org/pdf/1211.5122.pdf or  https://faculty.math.illinois.edu/\~{} west/regs/123conj.html  Omar pointed out that recently a prove has been found by Keutsh arXiv:2303.02611} 
For any connected graph $G=(V,E)$ on 3 or more vertices there exists a weighting of the edges $w: E\rightarrow \{1,2,3\}$ that induces a proper coloration of the vertices  \emph{i.e.}
for any $(u,v)\in E$, $\sum_{ (u,v')\in E} w(u,v') -\sum_{ (u',v)\in E} w(u',v)  \neq 0$   \end{theorem}
 %\footnote{For 123 proof one can restrict the matrices   and with entries in $\{0,1,-1\}$ that are  "almost symmetric" matrices \emph{i.e.} for all $i,j$,  $M_{i,j}+M_{j,i}=0 \bmod p$ and just need  $\exists p$}
 %\footnote{ mod $3^i$ should capture the exact strucrure }
%\end{center}
\begin{conjecture}
Given a permutation $\sigma\in S_m$ and a matrix   $E\in {\cal M}_m(\{0,1,-1\})$ such that $E_{\sigma i, i}\neq0$ for $i\in [m]$, then  for any $X_1,\ldots X_m$ each $X_i$ having $d$ possible values in ${\mathbb Z}$ and for any $B\in {\mathbb Z}^m$  there exists $n_1,n_2$ with $2n_1+n_2=m$ and  a space of dimension greater than $(d(d-2))^{n_1}(d-1)^{n_2}\ge (d(d-2))^{\lfloor m/2\rfloor}$ with basis elements $\ket X$ such that $\Pi_{i\in [m]}(MX-B)_i\neq 0$ where $(MX)_i= \sum_{j\in[m]} M_{i,j}X_j$. %$\forall i\in [m], 

\end{conjecture}

\subsection{link}
Given a graph $G=(V,A)$ with $|V|=n$ vertices and $|A|=m$ edges. 
Given an orientation $o$ of the the edges inducing a directed edges set $A'$, 
%Given an order on the vertices,  we  call an orientation of its edges a map $o:A \rightarrow \{-1,1\}$, where   $o(u,v)=1$  iff  the edge $(u,v)$ is directed from the smaller to the larger . The set of directed edges consistent with $o$ can then be defined as  $A'=\{ (u,v) \in A, o(u,v)=1 \text { and } u<v \} \cup \{ (u,v) \in A, o(u,v)=-1 \text { and } u>v \}$. 
 the {\bf effect} matrix  $E_{G,o}\in {\cal M}_m(\{0,1,-1\})$ is  defined on $A'\times A'$ by   $$(E_{G,o})_{(u,v),(u',v')}=\delta_{u,u'}+\delta_{u,v'}-\delta_{v,u'}-\delta_{v,v'} $$ where $\delta_{u,v}$ is the Kronecker delta. 
 

\subsection{Example before generalisations}

Given $T,T'$ tables with $d$ distinct positive integers called indices.
 A matrix $M$  indexed by $T,T'$  over a permutation  $\sigma \in S_d$  with spin sequence $Sp \in \{\uparrow,\downarrow\}^d $ is a matrix such that  for $i,j \in [d]$, 
 $M_{T'_i,T_j}\in \{J_\emptyset,\uparrow,\downarrow\}$ and $M_{T'_ {\sigma i},T_i}=Sp_i$  with $Sp_i \in \{\uparrow,\downarrow\}$ for $j\in [d]$ . %\in \{\uparrow,\downarrow\}$  
  Its "core" is $\tilde{M}_{i,j}=
  \begin{cases} 
  1 \text { if } M_{T'_i,T_j}= \uparrow \\
  -1 \text { if } M_{T'_i,T_j}= \downarrow \\
  0 \text { if } M_{T'_i,T_j}= J_\emptyset 
  \end{cases}$
 
% Its active correlations is the sequence of non 0 positions after the permutation indices (blue spins i.e. above a black spin).   Its passive correlations is the sequence of non 0 positions before the permutation indices (red spins i.ie at the left of black spins). 
We will define active and passive correlations.
By taking the columns one by one we can draw a shape by inserting in front of the associated correlation the ones linked with an active correlation.

The true conjecture is that looking at the shapes one can find a solution to each constraint $T'_i$  and the freedom we have for each $T_i$ is complete (any of the $d$ proposed values is ok) for the first element of a line containing more than one element,  $d-2$ for the last element of a line containing more than one element and $d-1$ otherwise.

We want to say that the causality JS of the problem characterizes  the uniform  solutions to our problem, and to prove the conjecture by reversibility of the map
 $    \begin{ytableau}
\none&\\
\none&   \\
\none &  \\
    \end{ytableau} 
    \rightarrow
     \begin{ytableau}
\none&&\\
  \none&   \\
    \end{ytableau}$

For example, the matrix 
 \begin{center}
 \ytableausetup{centertableaux}
 \begin{ytableau}
  \none  & \none[T_1] & \none[T_2] & \none[T_3] & \none[T_4]\\
       \none[T'_1] &\uparrow  &  &\blue{\downarrow} & \\
  \none[T'_2]  &\purple{\uparrow}  &  &  &\uparrow\\
  \none[T'_3]  &  & \purple{\downarrow}  &\uparrow & \\
   \none[T'_4]  &  &\uparrow  & &
\end{ytableau}
\end{center}
which corresponds when we add insertion time of cells corresponding to column insertion from the left to right with  ($Sp_i$, ($x$-position, correlation column),($y$-position,correlationline)) : $(1,(2, (0,T_2)), (2, ((0,T'_4))))$ then  $(1,(3, (-1,0,T_3)), (2,( (0,-1),T'_3))))$ then 
$(1,(4, ((0,0,0),T_4)), (2,((1,0,0),T'_4)))$
 \begin{center}
 \ytableausetup{centertableaux}
 \begin{ytableau}
  \none  & \none[T_1] & \none[T_2] & \none[T_3] & \none[T_4]\\
       \none[T'_1] &\uparrow_1  & _{1.5} &\blue{\downarrow}_{2.5} &_{3.5} \\
  \none[T'_2]  &\purple{\uparrow}_4  & _4 & _4 &\uparrow_4\\
  \none[T'_3]  & _3 & \purple{\downarrow}_3  &\uparrow_3 &_{3.5} \\
   \none[T'_4]  &  _2&\uparrow_{2}  &_{2.5} &_{3.5}
\end{ytableau} 
\end{center}

The causality JS for standard timeline (black spins permutation)  is obtained by reading the columns from left to right and inserting the line  corresponding to the time permutation
here  

$\begin{ytableau}
\none& \none[_{T_1}]\\
\none    & T'_1\\
    \end{ytableau}
     \rightarrow
      \begin{ytableau}
\none& \none[_{T_1}]& \none[_{T_2}]\\
\none    & T'_1 &T'_4 \\
    \end{ytableau} 
     \rightarrow
     \begin{ytableau}
\none& \none[_{T_2}]& \none[_{T_3}]\\
\none    & T'_4 &T'_1   \\
\none  &  \none &T'_3  \\
    \end{ytableau} 
    \rightarrow
     \begin{ytableau}
\none& \none[_{T_2}]& \none[_{T_3}]& \none[_{T_4}]\\
\none    & T'_4 &T'_1  &T'_2 \\
\none   & \none &T'_3  \\
    \end{ytableau}$


%Here the permutation positions are in black, the active correlations are   $Tcoract=(\blue{2^\downarrow})$ between ($X_{T_3}$ and  $X_{T_1}$) %or equivalently $Y_{T'}$ 
%, second emptyspot after (above)  the permutation positions, the passive correlations are $Tcorpass=(\purple{1^\uparrow},\purple{5^\downarrow})$ : first and fifth slot before/below the permutation positions
%between ($X_{T_1}$ and  $X_{T_4}$ and between ($X_{T_2}$ and  $X_{T_3}$)
% \ytableausetup{centertableaux}
% \begin{ytableau}
%\none    & T_1    \\
%\none    &T_2  & T_4 \\
%   \none & T_3  \\
%    \end{ytableau} 
%     is the shape of the permutation 
%  %\hspace {2cm}
%     and the shape of the impact on the constraints is
% \ytableausetup{centertableaux}
% \begin{ytableau}
%\none    & T'_1 &T'_3   \\
%\none    &T'_4    \\
%   \none & T'_2  \\
%    \end{ytableau}
%%\end{center}

\red{to clean , and addition of standard sum and multiplication,
The structure described :  $T,T',\sigma, cor_a,cor_p$  is called a (structuring space time) {\bf  key}.}
\footnote{the vague intuition is that lines represent space or energy and columns time or mass and the cells in the safe JS built from the structure are the basic elements. }

In the example the conjectured associated safe selection tree shape is therefore $d\times d-1 \times d-2 \times d-1$. Which implies  with $d=3$ that :
%for any $B\in \mathbb Z$ and for any sets of three choices  for $(X_{T_1},X_{T_2},X_{T_3})$ 
given sequentially $b_1,b_2,b_3,b_4$ ,
we can find  sequentially $X_1(t)$ for $t \in[4]$, $X_2(t)$ for $t \in [2,4]$, $X_3(t)$ for $t \in [3,4]$, and $X_4(t)$ defined only at $t=4$  such that $X(4)$ where $X=(X_1,X_2,X_3,X_4)$ forms a 
   $3\times 2 \times 1\times2=12$ "strong/regular/uniform" solutions to $\Pi_i (MX-B)_i\neq0$ : given $b_1\in \mathbb Z$ such that  for any of the 3 choices of    $X_1(4)$,   for any  $b_2$ there exists 2 choices for $X_2(4)$ such that for any $b_3$ there exists 1 choice for $X_3(4)$ such that for any $b_4$ there exists 2 choices for $X_4(4)$ such that $\Pi_i ((M(X(4)))_i-b_i)\neq 0$.  
%It as a list of active and passive correlations 

We want to generalize the notion of problem and the notion hyperplane to capture this regularity.

%It  has as basic base description  two rotated jardins suspendus  one filled with $T$ the other by $T'$)%by a subtable of

%Two tables $Tcoract$, $Tcorpass\in \Tid(d', 1, d(d-1)/2)$ 
%$duo_d [d(d-1)/2] \rightarrow d \times (d-1)$


%$(v/d^k)\bmod d < \sigma k$

%$(T'_1,T_1,forb_1)$ ou ($forb_1$ peut peut etre etre interpretee comme $init(X_{T_1})$?)

  \subsection{Causal dependence JS\red{tocheck}}
  Simple causal denpdence is a level 1 JS each cell depends on the cells above it.
  
  
  Advanced two level causal depencde use gflow cells anchor set of cells above and use right composition
  
  so the gflow JS (gflow) has anchors that can be read ordered  by order in the pattern of $\mathcal{P}(V)$ and cells $V$
  the order of the pattern has to be consistant with a timeline which leans that we can  right compose with a max JS with anchors $V$
  
associating to each cell an element of V so the anchors are $V$ and the cells below are the ones that has a dependcy on $v$    




Given  A causality JS $JC{^{(l)}}$  which anchors $ST_1, \ldots ST_n$ such that $ \sqcup (Jset)_c St_i =[l]$ anf cell number $T'_1\ldots T'_j$ with $\sqcup T'_i=[l]$

it is updated by $sqcor,p,sqcor'$ with $p\in [l+1]$ and for $j\in [l], sqcor_j\in \{0,1\}$ and $sqcor'_j\in \{0,1\}$
With 

$Update (sqcor,p,sqcor')(JC)$: 

$FJ_{l+1}= JC{^{(l)}}$ 

for $i=1$ to $p-1$ if $scor_i$,  $FJ_{i-1}= Merge_p (i) (FJ_i)$

for $i=p$ to $l$ if $scor'_i$,  $FJ_{i-1}= Merge_i (p) (FJ_i)$
 
 \subsection{Intuitive vision of Interactive problems and regular selection trees}

%We first introduce a slight generalisation to the notion of problem (or language) to add the notion of time.  

A  limited interactive problem can be seen as a referee and a player. The referee iteratively proposes 1 parameter and $d$ values to the player who has to pick one value. The referee can adapt the parameter depending on the answer.  The game is a language on the parameters, the sets proposed and the answers as an ordered sequence. 
It has to be defined  for any permutation of the variables (but not necessarily invariant). A proper formalisation will depend on the limitations. \footnote{(is it like IP with Polynomial referee, a multiround MA ?)}

%  a person who ahs to pich one of the values, on variables $X_1\ldots X_n$, consists in chosing . iteratively defining a domain of $d_1$ possible distinct values for $X_1$ , th  

   
  A $d'_1,\ldots , d'_n$- selection tree is a tree where the root (vertex at level 0) has (out)degree $d'_1$ and all the vertices of level $l$ have (out)degree $d'_{l+1}$)
     the edges represent the freedom of choices in such a problem.
 
 Given a sequence of indexed  values variables  $V_1,\ldots , V_n$ each $V_i$ having $d_i+1$ indexed ${\mathbb Z}$ elements : $V_i= (V_{i,0},V_{i,1},\ldots ,V_{i,i_{d_i}}$), with $V_{i,0}=b_i$ and   the elements indexed by $[d_i]$ forming a set of $d_i$ distinct values in ${\mathbb Z}$.
   A labeling of the edges  of  a $d'_1,\ldots , d'_n$- selection tree  associating to outgoing edges at level $l$ different indices in $[d_i]$  describes a family of solution to the problem.
   
% elements (the parameter then the sequence of choices proposed), and a sequence $d'_1\ldots d'_n$ with $d'_i\le d_i$, a $d'_1,\ldots d'_n$-labeling consists in labeling each  edge or equivalently each non root vertex  $u^{(l)}$  at level $l \ge 1$ (the root is counted as level 0),   with an integer $ lab(u^{(l)})\in[d'_i]$ so that childs of a vertex have different labels. The labeling at level $l$ depends only on the vertex (path from the root) and on the sequences $V_1\ldots V_l$.
 
 %  for any parameter choice $p_1\in \mathbb Z$, and for any choice of $d_1$ values in $D_1$, one can take $X_1=V_{l(v}$
   
   %or all  pair $V,B$  built by following the label $P(D,V,B)$
   
%  A labeling of the edges of such tree can be read as  
   A solution tree is a general solution to a problem if for each instance there exists a labelling such that taking the value indexed by the next vertex allows to be in the language.  
\section{Preremark on the first JS}
 
$    J_{\oN_{|14}}= \begin{ytableau}
\none& \none[_{2}]& \none[_{3}]& \none[_{5}]& \none[_{7}]& \none[_{11}]& \none[_{13}]\\
\none    & 4 &6  &10&14 \\
\none   & 6 &9  \\
\none   & 8 &12  \\
\none   & 10   \\
\none   & 12   \\
\none   & 14
    \end{ytableau}$
    
    The birth of a number depends on the reading strategy : it can be smaller than the number if it is not prime and it can be larger than the number !. (somehow by analogy for now the curvature at a point could be related to which case we consider).
    
    
    \section{New version}

We call $\Suct$ subjacency and $\Sucs$ posteriority operators.\red{may be identify $F_j$ with infinite number of it is the point in causality questions}

$F=\Seq_l F_l$ such that $\forall l , F_l=\Seq_{j} J_j$ (Level of explicicity  Jardins) and

 $J_j=\Seq_{t} J_{j,t}$ (Jardin Lianes) and $J_{j,t}=\Seq_{t} J^{(h)}_{j,t}$ (Liane cells)

$Cor_{F,j,j'} (F_{j'})=(F_j)$, $Cor_{J,j,j'} (J_{j'}^{l'})=\li_l .{J_j}$
 
Given an equivalence relation $\mathcal{R}:=\equiv$, we want to view the distinguishable elements as a sequence $A/\mathcal{R}:= \Seq_{j_{|l}} A^{(l)}_j$ for $i\in \oN$ which can be read at a local time depending on $l$ we discovered $l$ different classes.

for  all $k \in [l]$ and for all $j_k$, $j'_k$, 
there exists operators $\Sucs_{j_k} A^{(l)}_{j''}=A^{(l)}_{\Sucs_{A,l,j_k}  j''}$,
 and $\Suct_{j'_k}A^{(l)}_{j''}=A^{({\Suct_{A,j''} l)}}_{j''}$ 
 
  such that $h_{j_k}= \max \{n,  \Seq_{j_{|n}} \Suct_{j'_k} A_j  :=(A'_1,\ldots A'_n) \text { with  }$n$  \text { pairwise distinct equivalent  elements}\}$ and  $\forall n \in \oN$,  $\Sucs_{j'_k}^{n} A_j  \not \equiv  A_j$   


$\forall j,\exists j', \forall j''\ge j', \Sucs^{j''} A^{(t)}_{j}=\icp ^{(t,t+1)}  (A,j''-j')$ % A^{(t)}_{{\Sucs_{A}}^{j''} j}$

$A=\Seq_j A_j$ 

$\Read^{(k)} (A): (Set)_c A \rightarrow \oN^*$, $(A_j)_k \mapsto p_{j,k,1}<p_{j,k,2}<\ldots <p_{j,k,l}$


Imaginary  completion is  a sequence 
$\Imcp(A,R):= \Seq_{j'} \icp^{(k)} p_{j,k,l+j'}$%,  \icp^{(k)} p_{j,k,l+2}\ldots $ with $p_{j,k,l}<p_{j,k,l+j}<p_{j,k,l+j+1}$ for all $j \in \oN$ %<\ldots <p'_{j,k,l'}$

%the Canonical imaginary completion is the sequence

$\Imcp(A,R):= \icp^{(k,k+1)} p_{j,k,l+1},  \icp^{(k,k+1)} p_{j,k,l+2}\ldots $ with $p_{j,k,l}<p_{j,k,l+j}<p_{j,k,l+j+1}$ for all $j \in \oN$ %<\ldots <p'_{j,k,l'}$

Irreducible  completed $k$ reading if it doesnot exists   a subsequence $ p_{j,k,1}<p_{j,k,2}<\ldots <p_{j,k,l''}$  with $l'' <l$ such that  
$|p_{j,k,l''+n}|= l''+|p_{j,k,n}|$

%$\forall a\in A_i$ 

To map liane $l$ to liane $l'$ provided we have $k$ copies with $k$ multiple of the ppcm of the lengths.

$[n_{\#l,\#l'}]. l \equiv l' .[n_{\#l',\#l}]$ with $\at a {(j,q)}  \mapsto \at {a'} {(q', j')}$   

%\section{New version}
%
%We call $\Suct$ subjacency and $\Sucs$ posteriority operators.
%
%$F=\Seq_l F_{l+1/2}$ and  for all $ j\in \mathfrak t_F$, $J^{(l)}_j=\Seq_l  J_j$
%and $F_{l}$such that $\forall l , \Sucs_{F} J_l =J_{l+1}$ (Level of explicicity  Jardins) and $J_j=\Seq_{t} J_{j,t}$ (Jardin Lianes) and $J_{j,t}=\Seq_{t} J^{(h)}_{j,t}$ (Liane cells)
%
%$Cor_{F,j,j'} (F_{j'})=F_j$, $Cor_{J,j,j'} (J_{j'})=J_j$
% 
%$A/\mathcal{R_\equiv}:= \Seq_{j_{|l}} A^{(l)}_j$ for $i\in \oN$ 
%
%for  all $k \in [l]$ and for all $j_k$, $j'_k$, 
%there exists operators $\Sucs_{j_k} A^{(l)}_{j''}=A^{(l)}_{\Sucs_{A,l,j_k}  j''}$,
% and $\Suct_{j'_k}A^{(l)}_{j''}=A^{({\Suct_{A,j''} l)}}_{j''}$ 
% 
%  such that $h_{j_k}= \min \{n,  \Suct_{j'_k}^{n} A_j  \equiv  A_j\}$ and  $\forall n \in \oN$,  $\Sucs_{j'_k}^{n} A_j  \equiv  A_j$   
%
%
%$\forall j,\exists j', \forall j''\ge j', \Sucs^{j''} A^{(t)}_{j}=\mathfrak i_{cp}^{(t,t+1)} A^{(t)}_{{\Sucs_{A}}^{j''} j}$
%
%$A=\Seq_j A_j$ 
%
%$\Read^{(k)} (A): (Set)_c A \rightarrow \oN^*$, $(A_j)_k \mapsto p_{j,k,1}<p_{j,k,2}<\ldots <p_{j,k,l}$
%
%$\forall a\in A_i$ 
%
%To map liane $l$ to liane $l'$ provided we have $k$ copies with $k$ multiple of the ppcm of the lengths.
%
%$[n_{\#l,\#l'}]. l \equiv l' .[n_{\#l',\#l}]$ with $\at a {(j,q)}  \mapsto \at {a'} {(q', j')}$   
  \section{Simple version}

\subsection{Iterative process}

We first define a timed version of positive integer that we call universal colors and denote it just $\oN$, the first element is denoted 1, and  the successor operator
$\Sucs_{\oN}$ is a countable discrete process  producing the integer $j$ at a color called time  and denoted $\ti U (j):= j$.
 
 
 Given a sequence of symbols  $X$  %such that $\#_j(X)\le 2$ i.e. 
 where the symbol $j$ appears in $X$, $\#_j(X)$ times  
 
 %An iterative process $X$  denoted  
 $S:=(\Seq_j(X))$ denotes the iterative process that considers a fixed perfect copy of $X_0$ of $X$ and writes $X_0[j\mapsto 1], X_0[j\mapsto 2], \ldots$, obtained by replacing $j$ iteratively with $\oN$ 
 
 Each term  gains index in $S$,  $j' \in \oN$, $S_{j'}:=  X_0[j\mapsto j']$  have a %has associated 
 production start time $\ti {S_1}(j')$ when we started writing and  a production end time $\ti {S_2}(j')$% depending on $j'$ and the position of $j'$ in $X$. 
 
 and a successor $\Sucs_{S}  X[j\mapsto j']:=X[j\mapsto j'+1] $.

A Sequence is equipped with general  reading  detrmining at each universal time after it starts which local rewriting process is producing a new replacement $\Sucs_{S,R}  X[j\mapsto j']:=X[j\mapsto R_{n_j'}(j')] $ where $n_{j'}$ is the number of symbols $j$ already replaced in the $j'$ copy. The reading also allocates a (maximal?) time for the rewriting.

At each time a JS is the picture of this rewriting process. A Liane represents an  $X[j\mapsto j'] $
 
  A special type of sequence is iteration when ${\#_j(X)=0}$: it creates an indexation.
  $It=Iter_{j} (X):= Seq_j(X_j)$ such that  $\Sucs_{It} {X_j}=X $ for all $j\in \oN$
    
To select the l first elements we define a Jardin Suspendu   $J:=\Seq_{j_{|l}} L_j$ so that $L_1$ contains the $l$ first terms as a Liane : $L_{j',j}=S_{l(j'-1)+j}$ for $j\in [l]$ and $j'\in \oN$, we also define  $\Suct_L  L_{j',j}=L_{j'+1,j}$ for $j' \in [l-1]$ and 
$\Suct_L L_{l,j}=L_{1,j}$ and $\Sucs_L L_{j',j}=L_{j',1}$

A JS containing $n$ copies   of each $S_i$  $[n]\times S:=  \Seq_j  \Seq_{j'_{|d}} S_j $
      
     \section{Gradation}
  Given $F=\Seq_j F_j$
  a  gradation of $F$ inserts a new family of indices and corrections %and readings  
the canonical gradation %$G_F \in \Grad F = G_F, \Seq_j \Next_j$,

 $$\Grad F:= \Seq_j (\Seq_{j'} F'_{j',j})$$


  associates to each $F,j$  numbers $l_j, d_L,d_R$ and correctors $Cor_L,Cor_R$
  such that
  
    $$Cor_L ([d_L] \times  F_j)= [d_R] \times (\Seq_{j'_{|l_j}} F'_{j,j'}) \text{  and  }  [d_L] \times F_j= Cor_R ([d_R] \times (\Seq_{j'_{|l_j}} F'_{j,j'}))$$ 

It has the canonical reading
$\Next F'_{j',j}= F'_{j'+1,j}$ if $j'<l_j$ and $\Next F'_{l_j,j}= F'_{1,j+1}$

\footnote{In general for any $\Next_j$ in a gradation,   
     either   $ \Next  F_{j,j'} \in \{F_{j,j'+1},F_{j,1}\}$ and    $F_{j+1,j'} = (\Seq_{k_|r{_{j,j')}}} \Next ) F_{j,j'}$   
  or  $ \Next  F_{j,j'}=F_{j+1,j''}$ and    $F_{j,j'+1} = (\Seq_{k_|r{_{j,j')}}} \Next ) F_{j,j'}$   or $F_{j,1} = (\Seq_{k_|r{_{j,j')}}} \Next ) F_{j,j'}$ }  

In the case  of the restriction to the $l$ first elements is $l_j=1$ for $j\in \oN$ and $d_L=d_R=1$ 
  $Cor_L (F_j)=  F'_{q,r'+1}$ where $j=(q-1)l +r' $  with $r'\in [0,l-1]$ and $q\in \oN$ and  $Cor_R (F'_{q,r})=  F_{l(q-1)+r}$  %\times (\Seq_{j'_{|l_j}} F_{j,j'}) \     
% $    \Next_{G_F}\in {\Sucs,\Suct}$
%  atome, monome, completion reelle sousj,completion imaginaire proc,
\subsection{Lianes and cells}
    A \textbf{Liane} is a deterministic concept in which for each finite section defining a bottom concept, there exists a unique not impossible percept to each sequence of concepts starting from bottom. Reading a Liane is done by subjacency from the top element which is read from the book of percepts.
    
    An element of a Liane is called \textbf{cell} and denoted $\at a q$  where $a$ is read from a book of percepts and have an evaluation  $\mathfrak v (a) \in \oN$ and $q$ is read from a book of concepts and have an index $\iota (q) \in \oN$.

We can use $\iota(a) = l_J(\at a q)$ to define the index of the Liane ina JS $J$ (reduced length if we dont read the value of the percepts) and $\mathfrak v (q)= h_J(\at a q)$ as the reduced height of the atom in the Liane. 
     
  \subsection{Subjacency in $J^{(l)}$}
  

    
\subsubsection{for cells or percepts}  
  $\Suct_{F,J^{(l)},L}B:=A$  means that 
  Inside the theory $F$ in any $l'$ explicitation of $J$ in $F$ , $J^{(l')}$ with $l'\ge l$ in the Liane $L$ , $A$ produces as next term $B$ i.e. $B$ is the top term or the relative percept of $A$ in the Liane $L$ in the JS $J$ in the theory $F$.
  
  
  For example 
  

\begin{mdframed}[linewidth=1pt, roundcorner=5pt]
    % Premier tableau

    % Deuxième tableau
    \begin{minipage}{0.45\textwidth}
        \centering
        \begin{tabular}{|c|}
        \hline
     apple \\
    
        \hline
        1 apple \\
        \hline
      2 apples\\
        \hline
        \end{tabular}
        \vspace{0.3cm} % Ajustez cet espace si nécessaire
       % \captionof{table}{Liane conceptuelle du calcul déterministe}
    \end{minipage}
    \hfill
    % Troisième tableau
    \begin{minipage}{0.45\textwidth}
        \centering
        \begin{tabular}{|c|}
        \hline
        5 \\
        \hline
        4 \\
        \hline
        3 \\
        \hline
        2 \\
        \hline
        1 \\
        \hline
       $Add_1$ \\
        \hline
        \end{tabular}
        \vspace{0.3cm} % Ajustez cet espace si nécessaire
    %    \captionof{table}{Liane conceptuelle de 5}
    \end{minipage}
\end{mdframed} 
  
  It means  that there exists an equivalence relation $\equiv_{F,J^{(l)}}$ such that    $\Suct_{F,J^{(l)},L}B\equiv B$ and an associated height such that $h_{F,J^{(l)},L}(B)=h_{F,J^{(l)},L}(\Suct_{F,J^{(l)},L}B)+1$. 
  
  This equivalence can be seen as living in the same causal world for the atoms.  The atoms can be seen  in the same deterministic reasoning model in which one can say 
  it is suffficent to know whats below to build a deterministic imaginary world in which what above always happen i.e. assuming   "$\Suct B$" one can can say $B$.
  

  
  
  \subsection{Posteriority in $J^{(l)}$}
    
  $\Sucs_{F,J^{(l)}} L= (F,J^{(l)},L')$ means that either $L'$ is a really discovered new equivalence class in $J$ discobvered after $L$ or that it is an imaginary future equivalence class, that $L$ would guess will come after it.
 
 
 \subsection{Partial reading}
 $\Read_{T'_j}^{(k)} (A): (Set)_c A \rightarrow \oN^*$, $(A_{T_j})_k \mapsto p_{j,k,1}<p_{j,k,2}<\ldots <p_{j,k,l}$ respects posteriority and is $T_j$ value sensitive.

 \subsection{regularization by injection vs regularisation by inflation}

\section{Computations of depth $l$ involving $n\in \oN$  viewed with JS}
    
  Let $Tow(n,k)$ be such that $Tow(n,1)=n$ and $Tow(n,k+1)=n^{Tow(n,k)}$. 
  %Let consider $k\le n$ 
\footnote{old  \red{first attempt} We define $J_1(n,k)=n+k$ , 
  $J_2(n,k)=(n+k,nk)$, $J_3(n,k)=(n+k,nj,n^k)$, $J_4(n,k)=(n+k,nk,n^k,Tow(n,k))$ and $J_{4+l}(n,k)=(J_1,J_2,J_3,J_4,J_l)(n,k)$ not clear if $n-1,k-1$ the main idea is to introduce pairs and to continue with indices 1 mod 4 doing additions , 2 mod 4 multiplication , 3 mod 4 powers and 4 tower power. }

%\red{second attempt}
For $n\in \oN$ and $k\in \oN$, we define $J_1(n,k)=n+k$ , 
  $J_2(n,k)=nk$, $J_3(n,k)=n^k$, $J_4(n,k)=Tow(n,k)$
  
  and for $l\le n$, $$\Comp_l(n)= \Seq_{k_{|n}}\Seq_{j_{|l}} J(k,j)$$ 


  $\Comp_l(n)$ is the most complete set of computations using $n$ using $\Sucs_{\oN} n := n+1$  as a generator.
  Our operations will be restrictions of this set which could be used as low concept of a cell. 
%$J_1(n,k):= \Seq_{j_|k} X_{j+1}= \Sucs X_j$
Any monotone map $f:\oN \rightarrow \oN$ have $\Suct_f$ and $\Sucs_f$  where subjacency will have as percept the last discovered by increase by one before finding the next with $f$ cyclically closed when finding the next one which correspond to $\Sucs_f$   
Higher level subjacency and posteriority can be defined similarly $\Suct_f^{(j)}$ for $j \in [4]$. 

    \section{ Definition  and entanglement JS }
    
    Given 
    
    A JS $JDef= \Seq_j (\Seq_{j'}  \at {T_i} {T'_{j'}, M_{T'_{j'},T_j},h_{j,j'}})$ (such that $M_{T'_j,T_j}\in \{\uparrow,\downarrow\}$).   Jdef is interpreted in a Percept Book (of times) and describes 
    the sequence of possibilities in a JS

    A JS $JEnt= \Seq_{j'} (\Seq_j  \at {T'_{j'}} {T_j, M_{T'_{j'},T_j},l_{j',j}})$ (such that $M_{T'_j,T_j}\in \{\uparrow,\downarrow\}$ ). JEnt  is interpreted in a Concept Book (of spaces)  and describes the entanglement of concepts in a JS.
  
  
  Given $J$,  a partition reading  $Next(F,J,P, \Seq_{j_{|l}} d_j) =(s,j)$ where $\Suct_j$ was the direction followed and $s$ was the spin read. in insertion reading take the smallest direction allowing to have a new real element
  
  
  So if we have 
  
  $D=\Seq_j D_j$ with $D_j \in \{\uparrow,\downarrow\}$
  
  A family  of entangling correlations $\Seq_j Corr_j$ and a family of defining  Correlations $\Seq_j Corr'_j$ 
  
  And a graded family $\Sucs=\Seq_j \Seq_{j'} \Sucs_{j',j} $ of twirling posteriority operations.


 \begin{center}
 \ytableausetup{centertableaux}
 \begin{ytableau}
  \none  & \none[T_1] & \none[T_2] & \none[T_3] & \none[T_4]\\
       \none[T'_1] &\uparrow  & \blue{\uparrow} &\blue{\downarrow} & \\
  \none[T'_4]  &\purple{\uparrow}  &  &  &\uparrow\\
  \none[T'_3]  &  & \purple{\downarrow}  &\uparrow & \\
   \none[T'_2]  &  &\uparrow  & &\purple{\uparrow}
\end{ytableau}
\end{center}

 \begin{center}
 %\renewcommand{\arraystretch}{1.5}
 \ytableausetup{centertableaux}
\ytableausetup{boxsize=2.5em}
JDef(4) = \begin{ytableau}%{boxsize=2.5em}
  \none  & \none[T_1] & \none[T_2] & \none[T_3] & \none[T_4]\\
       \none &_{T'_1,\uparrow, 1}  &_{ T'_1,\blue{\uparrow},1}  & _{T'_1, \blue{\downarrow}, 1} &_{T'_4,\uparrow,2} \\
  \none  &_{T'_2,\purple{\uparrow},1}  & _{ T'_3,\purple{\downarrow},2}  & _{T'_3,\uparrow,2}&_{T'_2,\purple{\uparrow},2} \\
  \none  &\none& _{T'_2,\uparrow,1}    \\
%  \none  &  &  & & \\
  % \none  &  &  & &
\end{ytableau}, \text{ and }
JEnt(4) = \begin{ytableau}%{boxsize=2.5em}
  \none  & \none[T'_1] & \none[T'_4] & \none[T'_3] & \none[T'_2]\\
       \none &_{T_1,\uparrow, 1}  &_{ T_1,\purple{\downarrow},1}  & _{T_2, \purple{\downarrow}, 2} &_{T_2,\uparrow,2} \\
  \none  &_{T_2,\blue{\uparrow},1}  & _{T_4,\uparrow,3}& _{T_3,\uparrow,1}&_{T_4,\purple{\uparrow},2}  \\
  \none  &_{T_3,\blue{\downarrow},1}    \\
%  \none  &  &  & & \\
  % \none  &  &  & &
\end{ytableau}
\end{center}  
 \begin{center}
 %\renewcommand{\arraystretch}{1.5}
 \ytableausetup{centertableaux}
\ytableausetup{boxsize=2.5em}
JDef(3) = \begin{ytableau}%{boxsize=2.5em}
  \none  & \none[T_1] & \none[T_2] & \none[T_3] \\
       \none &_{T'_1,\uparrow, 1}  &_{ T'_1,\blue{\uparrow},1}  & _{T'_1, \blue{\downarrow}, 1}  \\
  \none   & \none&_{ T'_3,\purple{\downarrow},2}  & _{T'_3,\uparrow,2} \\
  \none   & \none& _{T'_2,\uparrow,1}   \\
%  \none  &  &  & & \\
  % \none  &  &  & &
\end{ytableau}, \text{ and }
JEnt(3) = \begin{ytableau}%{boxsize=2.5em}
  \none  & \none[T'_1] & \none[T'_3] & \none[T'_2] \\
       \none &_{T_1,\uparrow, 1}   & _{T_2, \purple{\downarrow}, 2} &_{T_2,\uparrow,2}  \\
  \none  &_{T_2,\blue{\uparrow},1}  &  _{T_3,\uparrow,1}  \\
  \none  &_{T_3,\blue{\downarrow},1}   \\
%  \none  &  &  & & \\
  % \none  &  &  & &
\end{ytableau}
\end{center}  
 \begin{center}
 %\renewcommand{\arraystretch}{1.5}
 \ytableausetup{centertableaux}
\ytableausetup{boxsize=2.5em}
JDef(2) = \begin{ytableau}%{boxsize=2.5em}
  \none  & \none[T_1] & \none[T_2]  \\
       \none &_{T'_1,\uparrow, 1}  & _{T'_1,\blue{\uparrow},1}     \\
       \none & \none& _{T'_2,\uparrow,1}     \\
%  \none  &  &  & & \\
  % \none  &  &  & &
\end{ytableau}, \text{ and }
JEnt(2) = \begin{ytableau}%{boxsize=2.5em}
  \none  & \none[T'_1] & \none[T'_2] \\
       \none &_{T_1,\uparrow, 1}   &_{T_2,\uparrow,2}  \\
       \none &_{T_2,\blue{\uparrow}, 1}    \\
%  \none  &  &  & & \\
  % \none  &  &  & &
\end{ytableau}
\end{center}    


We can use walls   $\Wlt t$, 
\Wl x   \rule{2pt}{5pt}$^{(t)}$y  \rule{2pt}{5pt}$^{(t+1)}$  so that at half times $t+1/2$ demi-jardins or operators are between two walls and normal jardins are before a wall .


A superposition is a sequence of instanciations in the mixture of the sequences of product spaces of plausible impossibilities.

at time $t+1/2$

 \begin{center}
 %\renewcommand{\arraystretch}{1.5}
 \ytableausetup{centertableaux}
\ytableausetup{boxsize=2.5em}
JDef(1+1/2) = \begin{ytableau}%{boxsize=2.5em}
  \none  & \none[T_1] & \none[T_2]  \\
       \none &_{T'_1,\uparrow, 1}  & _{T'_1,\blue{\uparrow},1}     \\
%  \none  &  &  & & \\
  % \none  &  &  & &
\end{ytableau}, \text{ and }
JEnt(1+1/2) = \begin{ytableau}%{boxsize=2.5em}
  \none  & \none[T'_1]  \\
       \none &_{T_1,\uparrow, 1}   \\
       \none &_{T_2,\blue{\uparrow}, 1}    \\
%  \none  &  &  & & \\
  % \none  &  &  & &
\end{ytableau}
\end{center}    

w.l.o.g  0 will be always the first $\mathbb Z$ value of the first atom of $X$ (can be done by readjusting the forbidden value)? 

\section {Monome/Shadow/ Phantom}


A monome $M$ of degree $d$ living at $\ti M$  is a sequence of atoms interpreted as simple spins :

It has ports  defined as $d$ cyclic sequencing $P=\Seq_{j_{|d}} P_j \in {\cal C}_d$ with $P_j\in \oN \setminus \cup_{j'<j} P_{j'}$
with $\Sucs_{{\cal C}_d} P_j=P_{j+1}$ for $j \in [d-1]$ and $\Sucs_{{\cal C}_d} P_d=P_{1}$.  And  $\Suct_{{\cal C}_d} P_j=P_j$.
 

each port have (minimal production) time $\ti {P_j}  \ge P_j$

$\ti M= \Seq_{j\in [d]} \ti {P_j}$, $\va (\ti M)=\max_{j\in [d]} (\ti M)_j$

$M:= \Seq_{j_{|_d}} \at {S_j} {P_j} \in {{\cal M}o_d}$  with $S_j\in \{\uparrow,\downarrow\}$


A monome  is equipped with 


monomial subjacency   $\Suct_{{\cal M}o_d} M=\Seq_{j_{|_d}} \at { \Suct S_j} {P_j}$ , $\ti {\Suct_{{\cal M}o_d} M} \ge  (\ti M) +1^d ?$ i.e. $(\ti {\Suct_{{\cal M}o_d} M})_j= (\ti M)_j +1$ for $j\in [d]$

monomial posteriority   $\Sucs_{{\cal M}o_d} M=\Seq_{j_{|_d}} \at { \Sucs_P S_j} { P_j}$ , $\ti {\Sucs_{{\cal M}o_d} M}\ge \Seq_{j_{|d}} [{\ti M }_j \mapsto 2 (\ti M )_j ] ?$

A visualisation of a phantom can be as a vertex with directed halfedges labeled with integers  or as a flat (height 1 everywhere) finite JS of spins (fixpoint of reading given a sequence of labels).

\section{Anchored phantom}

An anchored phantom has a time or space subjacency allowing him to be completed  producing an infinite  sequence of successors isomorphic to $\oN$ still between the walls  $\Wlt t$ and $\Wlt {t+1}$

A visualisation of a phantom can be as a vertex colored with an integer with directed half-edges labeled with integers.

An anchored phantom has a subjacency that can be seen as a posteriority  in the book of time or in the book of space. It associates a possible \textbf {completion}  $\Next$. It is an iterable  map  that allows to produce successors it is an infinite reading of its JS.

A bi-anchored phantom has below its completion a divination book that is a sequence of completions F$\Next$s: \red{ maybe 
It is a cell with a grangranfdather wishes being a sequence of possibilities a grand father being the possibilities we gave up and a father being the first  possibility which a completion map. some how how to correct from another reading to its canonical one?}

%compatible with grandfather wishes below an atom. It can be seen as a cell were 
%Oldval a is a set and Keep is an element of the set.  
 
%\section{Liane}
%A Liane describes a deterministic sequence of subjacency

\section{ A Reading}  

A reading  $R$ of a JS $J$  Is a sequence of atoms of $J$,   $\Seq_j  (\at a q)_j$    such that either  $(\at a q)_{j+1} = \Suct_J (\at a q)_j$ or $(\at a q)_{j+1} = \Sucs_J (\at a q)_j$.

Given a reading $R$ of a JS $J$, we can define  $Next$ giving the next atom in the reading. A JS with a reading has anchored atoms.

\section{Cells}
   cellulations mapping  each  atoms of a JS to cells, where  
   
   A  \textbf{ cell} in a JS $J'$ sequentialy built from a JS $J$ following a reading, is a Liane of the form:
 \begin{tabular}{|c|}
       \hline
     Newpos \\
    
        \hline
        (Keep,
      Forget)\\
        \hline
  Remember \\
        \hline
 \end{tabular} =? \begin{tabular}{|c|c|}
       \hline
     Newpos$_1$ &  Newpos$_2$ \\
    
        \hline
        Keep&Forget\\
        \hline 
         Remember$_1$ &  Remember$_2$ \\
        \hline
 \end{tabular} pb remember is not a product

Remember is a map satisfying $(\Jat  {J'} {(Keep, Forget)} {Newpos}) _{j} = (\Jat J {Oldval} {Oldpos})_j$.

The point is to restrict $a$ using a predicate $\phi$ building $a_{|\phi}$ when assuming $\phi$ that we want to describe more  and $a_{\setminus \phi}$ for a complement, and a recovery map allowing to get back $a$ from the pair.
\section{Unitar/fusion}


A unitar connects two superpositions of space phantom and time phantom  generating an observable living in time space.
\section{workplan}
we want to play with mixture of complete impossibilities over superposition of possibilities 
with anchor expansion and coherence restriction (cellgrowth)

A (superposition) of  graphs gets superposition of  phantoms  which  gets a mixture of (superposition) of graphs.


building the superposition of impossibilities at time $T_n$

 Val table $\Seq_{k_{|l_n}} \at {\mathfrak v (T_{j_{k,n}})} k$  possibility sequence  

 Val table $\Seq_{k_{|h_n}} \at {B_{T'_{j_{k,n}}}} k$ referential gap
 
 Twirling maps  $\Seq_{k_{|h_n}} \at {tw_{T'_{j_{k,n}}}} k$ 

Closed product of impossibilities $ J_{ \Seq_{k_{|h_n}} \bar{ b_k}}$


Coherence projection to get mixture of possibilities with $ev_p$ 

A superposition is a sequence of instanciations in the mixture of the sequences of product spaces of plausible impossibilities.


Given the possibility values at time $n$, $\Seq_{j'_{|n}} X_{T_{j',n}}=\Seq_{j'_{|n}} \Seq_{j_{|d_{T_{j'}}}}  \textbf {at} (j) $    where $at(j)= \at {u_{n,j}} {j}$ ?

Can we consider instead of closed product something like let $(b'_{T'})_{|n}:=\Seq_{j|n}$ (if $b_{T'_j}$ possible after)

A discrete homotopy is a reading/translation bijective mapping. At universal time $t_U$ the step is a multiple of $1/ \ti U$

it s when we take a 1,1 or when we interpret the map as the smallest if mappeed to a symetrisation of the largest and new eight is the smallest??
  
  old 4 wall dec
  
  $2,5,8,11,14,17,\ldots  \red{|}\ldots,15,9,3|| 6,12,18, \ldots \purple{|}1,4,7,10,13,16,\ldots$

  $(-1) +3 (1,2,3 ,\ldots) \red{|}3(\ldots,5,3,1|| 2,4,6, \ldots) \purple{|}(-2) +3 (1,2,3 ,\ldots)$

  $(-1) +3 (1,2,3 ,\ldots) \red{|}3(1 +2 inv(1,2,3 ,\ldots) || 2 (1,2,3 ,\ldots) \purple{|}(-2) +3 (1,2,3 ,\ldots)$


\end{document}